\documentclass[12pt]{article}
\title{Digest of The New York Times}
\author{The New York Times}

% % \definecolor{dark-green}{rgb}{0,0.25,0}
% \usepackage[utf8]{inputenc}

\usepackage{fontspec,xunicode,xltxtra}
\setmainfont[Ligatures=Required,Mapping=tex-text]{Georgia}
% \setmainfont[Ligatures=Required,Mapping=tex-text]{Palatino Linotype}
\setsansfont[Ligatures=Required,Mapping=tex-text]{Myriad Pro}
\setmonofont[Ligatures=Required,Mapping=tex-text]{Courier New}
\usepackage{xeCJK}
\setCJKmainfont{宋体}
\setCJKmonofont{黑体}

% \usepackage{indentfirst}


\usepackage[dvips,dvipsnames,svgnames]{xcolor}
\definecolor{light-gray}{gray}{0.95}

%生成PDF的链接
\usepackage[dvipdfmx]{hyperref}
\hypersetup{
    bookmarks=true,         % show bookmarks bar?
    bookmarksopen=true,
    pdfpagemode=UseNone,    % options: UseNode, UseThumbs, UseOutlines, FullScreen
    pdfstartview=FitH,
    pdfborder=1,
    pdfhighlight=/P,
    pdfauthor={wuxch},
    unicode=true,           % non-Latin characters in Acrobat’s bookmarks
    colorlinks,             % false: boxed links; true: colored links
    linkcolor=blue,         % color of internal links
    citecolor=blue,        % color of links to bibliography
    filecolor=magenta,      % color of file links
    urlcolor=cyan           % color of external links
}
% 使用目录
\usepackage{makeidx}
% 使用图形
\usepackage{graphicx}


\usepackage{titletoc}
\dottedcontents{section}[1.5em]{}{2.3em}{1pc}

% % 页眉页脚的设定
% \usepackage{fancyhdr}
% \pagestyle{fancy}
% \lhead{\includegraphics[height=0.2in]{zte_logo}}

% \usepackage{footmisc}
\usepackage{pfnote}
\usepackage{fnpos}

% % 嵌入的代码显示
% \usepackage{listings}
% \lstset{language=C++, breaklines, extendedchars=false}
% \lstset{keywordstyle=\color{blue},
%         commentstyle=\color{SeaGreen},
%         stringstyle=\ttfamily,
%         showstringspaces=false,
%         backgroundcolor=\color{light-gray}}
% \lstset{basicstyle=\small}

\usepackage{booktabs}

\usepackage{tipa}
\usepackage{phonetic}
\usepackage[gen]{eurosym}

% \usepackage{CJK}
\usepackage{makeidx}
% \usepackage{glossary}

\usepackage{fancyhdr}

\usepackage{lastpage}
\cfoot{\textsf{\thepage\ of \pageref*{LastPage}}}

\pagestyle{fancy}
\fancyhead{}
\renewcommand{\headrulewidth}{0pt}
\pagenumbering{arabic}

\usepackage{graphicx}
\usepackage{picinpar}

\setlength{\parindent}{0in}
\setlength{\parskip}{2.5ex plus4mm minus3mm}
% \setlength{\parskip}{.3\baselineskip}

\setcounter{section}{0}

% bibtex
\usepackage[super]{natbib}
% \usepackage[globalcitecopy,labelstoglobalaux]{bibunits}
% \makeatletter
% \renewcommand{\@bibunitname}{\jobname.\the\@bibunitauxcnt}
% \makeatother

% \renewcommand{\refname}{缩略语}
% \renewcommand{\refname}{}

\makeatletter                   % Reference list option change
\renewcommand\@biblabel[1]{#1:  } % from [1] to 1: 
\makeatother                    %

\newcommand{\citenumfont}[1]{\textsf{#1}}

% \usepackage{tipa}
% \usepackage{phonetic}

\newcommand{\phonetic}[1]{{\fontspec[Colour=CC00CC]{Microsoft Sans Serif} #1}}


%%% Local Variables:
%%% TeX-master: "medical_trivia"
%%% End:

\usepackage{config}

\makeindex
\begin{document}
\date{}
\thispagestyle{empty}
\renewcommand\contentsname{\textsf{Digest of The New York Times}}
{\footnotesize\textsf{\tableofcontents}}
\clearpage
\setcounter{page}{1}

\section{Wartime China's Elegant Enigma}

\lettrine{``T}{he}\mycalendar{Nov.'09}{06} only thing Oriental about me,'' Soong Mei-ling once
wrote, ``is my face.''

Soong Mei-ling, better known to history as Mme. Chiang Kai-shek, was exaggerating only slightly.
Chinese by birth, American by education and cultural inclination, she was a
seductive\cite{seductive} blend of both societies; for a time, no woman in the world was more
powerful.

Mme. Chiang led a long, vastly complicated life, one that is richly detailed in ``The Last
Empress,'' Hannah Pakula's long, vastly complicated new biography. Ms.~Pakula's book is a yeoman
work of historical research, with fact grinding against fact. It is also a monotonous piece of
storytelling, one that has little pliancy or narrative push. Its 681 pages of text are at times as
grueling as a forced march across the Mongolian steppe.

The story of Mme. Chiang's life has lost none of its strange, piquant appeal, however. Born in
Shanghai in 1898, she was the daughter of a peasant who had gone to America at age 12 and found work
on ships and in printing shops. Her father, Charlie Soong, eventually graduated from Vanderbilt
University and returned to China at 20, where he had six children and became rich publishing Bibles.
He raised Soong Mei-ling and her siblings to appreciate almost everything Western, including
mattresses (soft), food (American) and religion (Methodist).

Cutting against the grain of a staunchly patriarchal society, Mr.~Soong expected big things from his
daughters as well as from his sons. Soong Mei-ling's two older sisters traveled to the United States
to attend Wesleyan Female College in Macon, Ga. Soong-Mei-ling arrived in America at age 10,
studying at a boarding school in New Jersey and a public school in Georgia before graduating from
Wellesley College.

When she arrived at Wellesley in 1913, Ms.~Pakula writes, Soong Mei-ling could lay on a ``Scarlett
O'Hara accent'' she'd picked up in Georgia. (``Ah reckon Ah shan't stay aroun' much longer,'' she
reportedly told the freshman dean.) She was also, Ms.~Pakula writes, ``short, chubby, round-faced
and childish in appearance, with a short haircut and bangs over her eyes that did nothing for her
looks.''

By the time she left Wellesley, however, there was a sense of destiny about Soong Mei-ling. ``She
had not been given a Western education,'' Ms.~Pakula observes, ``in order to spend her afternoons at
the mah-jongg table.''

The Soong sisters married well. One, Soong Qing-ling, married Sun Yat-sen, China's first president
after the last emperor was overthrown in 1911. In a lavish ceremony in 1927, Soong Mei-ling married
one of Sun's former military aides, Chiang Kai-shek, a man who would become the head of the
Nationalist government in China from 1928 to 1949, and later its leader while in exile in Taiwan.

He was a hardened soldier who ``dressed simply in a plain cotton uniform with straw sandals,''
Ms.~Pakula writes, and neither drank nor smoked. Mme. Chiang was by now thin, glamorous and wore
form-fitting clothes. Barely five feet tall, she had, Ms.~Pakula declares, ``a near-hypnotic effect
on men.''

Because Chiang Kai-shek spoke virtually no English, Mme. Chiang became his de facto translator and
the face China turned to the Western world. She wrote articles about China for The New York Times
Magazine and The Atlantic Monthly in the early '40s. She appeared on ``Meet the Press'' in 1958. She
was Chiang's closest adviser and she constantly buffed his -- and the country's -- rough edges.

The pair were seen as a modernizing influence in China; Time magazine named them Man and Woman of
the Year in 1938. The peak of Mme. Chiang's fame arrived in 1943, when she toured America in support
of the Nationalist Chinese cause against Japan.

During that tour she was the first private citizen to address the Senate and the House of
Representatives, and in Los Angeles she gave a speech to a packed Hollywood Bowl. (While in America,
Ms.~Pakula suggests, Mme. Chiang continued a romantic involvement she had begun earlier with Wendell
Willkie, the Republican who had lost the 1940 election to Franklin D.~Roosevelt.)

Chiang Kai-shek's government, increasingly besieged by China's Communist Party as the 1940s went on,
was also rotting from within. He was a ruthless, petty man and a dismal leader. As Theodore H.~White
and Annalee Jacoby observed, ``The manners of the Kuomintang'' -- the Nationalist Party -- ``in
public were perfect; its only faults were that its leadership was corrupt, its secret police
merciless, its promises lies, and its daily diet the blood and tears of the people of China.''

Eleanor Roosevelt got a chilling glimpse of Mme. Chiang's own dark side when Mrs.~Roosevelt asked
her how she would deal with a difficult labor leader like John L.~Lewis of the United Mine Workers.
``She never said a word,'' Ms.~Roosevelt wrote, ``but the beautiful, small hand came up and slid
across her throat.''

Chiang Kai-shek and his wife were forced into exile in Taiwan after the Communist victory in 1949;
he presided for decades over Nationalist politics from there. After his death, in 1975, Mme. Chiang
moved to New York City, where she led a reclusive life, dying in 2003 at 105. She had no children.
Her husband had contracted venereal disease before their marriage, Ms.~Pakula writes, and was
probably sterile.

``The Last Empress'' bogs down in overly long discursions into the intricacies of China's political
history. Indeed, Mme. Chiang's own story often recedes far into the background. But Ms.~Pakula's
book comes alive in its pepperings of telling detail about Mme. Chiang's chaotic life.

Ms.~Pakula notes the way Mme. Chiang loved to deploy esoteric words (``indehiscence,''
``ochlocracy'') in her speeches in English, sending reporters scrambling for their dictionaries. She
observes that President Harry S.~Truman, tired of Mme. Chiang's appeals for money, began to refer to
her husband as ``Cash My-check.''

She details Mme. Chiang's final years at 10 Gracie Square, an elegant apartment building on the
Upper East Side of Manhattan. There she kept three dogs (two bichons and a Yorkshire terrier) and
employed 24 servants. There were reports that neighbors complained about the cooking odors and
cockroaches in her 18-room apartment, and that Mme. Chiang kept a closet filled with gold bars.

Ms.~Pakula is also the author of ``The Last Romantic: A Biography of Queen Marie of Roumania''
(1985) and ``An Uncommon Woman: The Empress Frederick: Daughter of Queen Victoria, Wife of the Crown
Prince of Prussia, Mother of Kaiser Wilhelm'' (1997). She views Mme. Chiang's life with interest and
occasionally, when warranted, with sympathy. She is clearly in agreement with Eleanor Roosevelt, who
summed up Mme. Chiang's striding performance on the world stage by remarking that while she could
``talk beautifully about democracy,'' she did ``not know how to live democracy.''

\section{Chinese Trial Reveals Vast Web of Corruption}

\lettrine{C}{hongqing}\mycalendar{Nov.'09}{06}, China -- Wen Qiang had a fondness for Louis Vuitton
belts, fossilized dinosaur eggs and B-list pop stars. For a public employee in charge of the local
judiciary, he also had a lot of money: nearly \$3 million that investigators found buried beneath a
fish pond.

But Mr.~Wen's lavish tastes were nothing compared with the carnal\cite{carnal} appetites of his
sister-in-law, Xie Caiping, known as ``the godmother of the Chongqing underworld.'' Prosecutors say
she ran 30 illegal casinos, including one across the street from the courthouse. She also employed
16 young men who, according to the state-run press, were exceedingly handsome and obliging.

In recent weeks, Ms.~Xie, Mr.~Wen and a cavalcade of ranking officials and lowbrow thugs\cite{thug}
have been players in a mass public trial that has exposed the unseemly relationships among
gangsters, police officers and the sticky-fingered bureaucrats.

The spectacle involves more than 9,000 suspects, 50 public officials, a petulant billionaire and
criminal organizations that dabbled in drug trafficking, illegal mining, and random acts of
savagery, most notably the killing of a man for his unbearably loud karaoke voice.

But like all big corruption cases in China, this one is as much about politics as graft. The
political machine in Chongqing, a province-size mega-city of 31 million people in the southwest, has
been broken up by a new Communist Party boss, Bo Xilai, who is the son of a revolutionary party
veteran and has his eye on higher office.

Mr.~Bo, a former trade minister sent to Chongqing to burnish his managerial credentials, has
conducted the crackdown in a way that appears devised to maximize national attention. The drawn-out
nature of the trial and the release of lurid details of the criminal syndicate have given him a
reputation as a leading corruption fighter, though the inquiry has yet to implicate any really
high-ranking party officials.

So far six people have been sentenced to death. Ms.~Xie got off relatively lightly, receiving an
18-year prison term on Tuesday.

How Mr.~Bo's performance is regarded by the party elite is a matter of speculation. There are some
suggestions that his swagger\cite{swagger}, including boastful comments to the news media, strikes
some fellow officials as excessive. Anticorruption campaigns by China's one-party state are
generally calibrated to show resolution in tackling venality, but also to reassure the public that
whatever problems are uncovered are localized and effectively contained.

``These guys are all for fighting corruption, but they are a little alarmed by the way Bo Xilai has
been going about it and building up his personality,'' said Sidney Rittenberg, one of the few
American citizens to join the Communist Party here and a confidant of Chinese leaders since 1944.
``People I talk to say he's getting too big for his britches.''

A so-called princeling whose father, Bo Yibo, was an economic planner and a onetime ally of the
paramount leader Deng Xiaoping, Mr.~Bo, 60, is already a member of the Communist Party's powerful
Politburo. He is often talked about as a future top leader in Beijing, although in the party's rigid
hierarchy the No.~1 posts in the party and the government have already been assigned to other
younger officials.

Recent statements by Mr.~Bo suggest he understands the perils\cite{peril} of drawing too much
attention. Two weeks ago, he defended the crackdown, saying he was forced to act by the rampant
violence and brazen criminality that had given this perpetually foggy city a reputation for
lawlessness.

``The public gathered outside government offices and held up pictures of bloodshed,'' he said. ``The
gangsters slashed people with knives just like butchers killing animals.''

In the three weeks since trials began, the crowds have continued to come, and their stories of
bloodshed are indeed horrifying. They press outside the gates of the Fifth Intermediate Court,
hoping to glimpse the orange-vested defendants who are paraded through the hearings.

Others desperately seek out reporters willing to hear tales of crimes unpunished. ``The
bandits\cite{bandit} used to live in the mountains; now they live in the Public Security Bureau,''
said Zheng Yi, a vegetable wholesaler.

Unlike past sweeps that brought down crime bosses and their henchmen, the crackdown in Chongqing has
yielded a number of wealthy businessmen and Communist Party officials, exposing the depth of
corruption that has resulted from the mixing of state control and free-market economics in China.

Ko-lin Chin, who studies the intermingling of organized crime and government in China, said the line
between legitimate business and illegal conduct had become increasingly blurred, although most
official corruption involved bribery, not violence.

``As these gangs have become more powerful, their existence depends entirely on the cooperation and
tolerance of the Communist Party,'' said Mr.~Chin, a professor of criminal justice at Rutgers. ``But
when things get out of hand, as they did in Chongqing, the party can really go after these groups
with a vengeance.''

Among those on trial this week is Li Qiang, a local legislator and billionaire who the authorities
say owned a fleet of 1,000 cabs and 100 bus routes. So great was his power, they say, that he
orchestrated a taxi strike last year that brought the city to a standstill. On trial with him are
three government officials suspected of acting as his ``protection umbrellas'' in exchange for
payments of about \$100,000 each.

While Mr.~Li stood in the dock, more than 200 people gathered outside in the rain, including women
who said they were roughed up in October last year when they refused to vacate their homes for a
redevelopment project. One of them, Wu Pinghui, 67, said 40 people were herded into a
government-owned bus and dumped in the countryside. By the time they made it back, their homes were
gone.

``We called 110,'' she said, referring to the Chinese emergency number, ``but the police said they
couldn't get involved in a government affair.''

Hong Guibi also came to the courthouse. She said the Communist Party chief of her village, enraged
when she and her husband refused to give him part of their orchard, watched as thugs attacked the
couple with cleavers. Ms.~Hong, 47, was critically wounded, and her husband was killed. ``The
neighbors heard our screams, but they were afraid to do anything,'' she said.

Although heartened that so many are being prosecuted, Ms.~Hong is still waiting for someone to come
after the village chief. ``If I could just kneel down in front of Bo Xilai,'' she said, ``I'm sure
he would solve my problem.''

\section{South Koreans Struggle With Race}

\lettrine{O}{n}\mycalendar{Nov.'09}{06} the evening of July 10, Bonogit Hussain, a 29-year-old
Indian man, and Hahn Ji-seon, a female Korean friend, were riding a bus near Seoul when a man in the
back began hurling racial and sexist slurs at them.

The situation would be a familiar one to many Korean women who have dated or even -- as in
Ms.~Hahn's case -- simply traveled in the company of a foreign man.

What was different this time, however, was that, once it was reported in the South Korean media,
prosecutors sprang into action, charging the man they have identified only as a 31-year-old Mr.~Park
with contempt, the first time such charges had been applied to an alleged racist offense. Spurred by
the case, which is pending in court, rival political parties in Parliament have begun drafting
legislation that for the first time would provide a detailed definition of discrimination by race
and ethnicity and impose criminal penalties.

For Mr.~Hussain, subtle discrimination has been part of daily life for the two and half years he has
lived here as a student and then research professor at Sungkonghoe University in Seoul. He says
that, even in crowded subways, people tend not sit next to him. In June, he said, he fell asleep on
a bus and when it reached the terminal, the driver woke him up by poking him in the thigh with his
foot, an extremely offensive gesture in South Korea.

``Things got worse for me this time, because I was with a Korean woman,'' Mr.~Hussain said in an
interview. ``Whenever I've walked with Ms.~Hahn or other Korean women, most of the time I felt
hostilities, especially from middle-aged men.''

South Korea, a country where until recently people were taught to take pride in their nation's
``ethnic homogeneity'' and where the words ``skin color'' and ``peach'' are synonymous, is
struggling to embrace a new reality. In just the past seven years, the number of foreign residents
has doubled, to 1.2 million, even as the country's population of 48.7 million is expected to drop
sharply in coming decades because of its low birth rate.

Many of the foreigners come here to toil at sea or on farms or in factories, providing cheap labor
in jobs shunned by South Koreans. Southeast Asian women marry rural farmers who cannot find South
Korean brides. People from English-speaking countries find jobs teaching English in a society
obsessed with learning the language from native speakers.

For most South Koreans, globalization has largely meant increasing exports or going abroad to study.
But now that it is also bringing an influx of foreigners into a society where 42 percent of
respondents in a 2008 survey said they had never once spoken with a foreigner, South Koreans are
learning to adjust -- often uncomfortably.

In a report issued Oct.~21, Amnesty International criticized discrimination in South Korea against
migrant workers, who mostly are from poor Asian countries, citing sexual abuse, racial slurs,
inadequate safety training and the mandatory disclosure of H.I.V. status, a requirement not imposed
on South Koreans in the same jobs. Citing local news media and rights advocates, it said that
following last year's financial downturn, ``incidents of xenophobia are on the rise.''

Ms.~Hahn said, ``Even a friend of mine confided to me that when he sees a Korean woman walking with
a foreign man, he feels as if his own mother betrayed him.''

In South Korea, a country repeatedly invaded and subjugated by its bigger neighbors, people's racial
outlooks have been colored by ``pure-blood'' nationalism as well as traditional patriarchal mores,
said Seol Dong-hoon, a sociologist at Chonbuk National University.

Centuries ago, when Korean women who had been taken to China as war prizes and forced into sexual
slavery managed to return home, their communities ostracized them as tainted. In the last century,
Korean ``comfort women,'' who worked as sex slaves for the Japanese Imperial Army, faced a similar
stigma. Later, women who sold sex to American G.I.'s in the years following the 1950-53 Korean War
were despised even more. Their children were shunned as ``twigi,'' a term once reserved for animal
hybrids, said Bae Gee-cheol, 53, whose mother was expelled from her family after she gave birth to
him following her rape by an American soldier.

Even today, the North Korean authorities often force abortion on women who return home pregnant
after going to China to find food, according to defectors and human rights groups.

``When I travel with my husband, we avoid buses and subways,'' said Jung Hye-sil, 42, who married a
Pakistani man in 1994. ``They glance at me as if I have done something incredible. There is a
tendency here to control women and who they can date or marry, in the name of the nation.''

For many Koreans, the first encounter with non-Asians came during the Korean War, when American
troops fought on the South Korean side. That experience has complicated South Koreans' racial
perceptions, Mr.~Seol said. Today, the mix of envy and loathing of the West, especially of white
Americans, is apparent in daily life.

The government and media obsess over each new report from the Organization for Economic Cooperation
and Development, to see how the country ranks against other developed economies. A hugely popular
television program is ``Chit Chat of Beautiful Ladies'' -- a show where young, attractive, mostly
Caucasian women who are fluent in Korean discuss South Korea. Yet, when South Koreans refer to
Americans in private conversations, they nearly always attach the same suffix as when they talk
about the Japanese and Chinese, their historical masters: ``nom,'' which means ``bastards.'' Tammy
Chu, 34, a Korean-born film director who was adopted by Americans and grew up in New York State,
said she had been ``scolded and yelled at'' in Seoul subways for speaking in English and thus ``not
being Korean enough.'' Then, she said, her applications for a job as an English teacher were
rejected on the grounds that she was ``not white enough.''

Ms.~Hahn said that after the incident in the bus last July, her family was ``turned upside down.''
Her father and other relatives grilled her as to whether she was dating Mr.~Hussain. But when a
cousin recently married a German, ``all my relatives envied her, as if her marriage was a boon to
our family,'' she said.

The Foreign Ministry supports an anti-discrimination law, said Kim Se-won, a ministry official. In
2007, the U.N. Committee on the Elimination of Racial Discrimination recommended that South Korea
adopt such a law, deploring the widespread use of terms like ``pure blood'' and ``mixed blood.'' It
urged public education to overcome the notion that South Korea was ``ethnically homogenous,'' which,
it said, ``no longer corresponds to the actual situation.''

But a recent forum to discuss proposed legislation against racial discrimination turned into a
shouting match when several critics who had networked through the Internet showed up. They charged
that such a law would only encourage even more migrant workers to come to South Korea, pushing
native workers out of jobs and creating crime-infested slums. They also said it was too difficult to
define what was racially or culturally offensive.

``Our ethnic homogeneity is a blessing,'' said one of the critics, Lee Sung-bok, a bricklayer who
said his job was threatened by migrant workers. ``If they keep flooding in, who can guarantee our
country won't be torn apart by ethnic war as in Sri Lanka?''

\section{Qian Xuesen, Father of China's Space Program, Dies at 98}

\lettrine{Q}{ian}\mycalendar{Nov.'09}{06} Xuesen, a brilliant rocket scientist who single-handedly
led China's space and military rocketry efforts after he was drummed out of the United States during
the redbaiting of the McCarthy era, died on Saturday in Beijing. He was 98.

China's state media reported the death. Mr.~Qian had been frail and bedridden in recent years.

In China, Mr.~Qian was celebrated as the father of Chinese rocketry, the leader of the research that
produced the nation's first ballistic missiles, its first satellite and the Silkworm
anti-shipmissile.

But in the United States in the 1930s and 1940s, he was no less valuable, if not so publicly
celebrated, as a pioneer in American jet and rocket technology.

As a student at the Massachusetts Institute of Technology, and later as a scientist and teacher at
the California Institute of Technology, Mr.~Qian, also known as Tsien Hsue-shen, played a central
role in early United States' efforts to exploit jet and rocket propulsion.

As a graduate assistant at Caltech in the late 1930s, Mr.~Qian helped conduct seminal research into
rocket propulsion, and in the 1940s he helped found the Jet Propulsion Laboratory, now one of NASA's
premier space-exploration centers.

Mr.~Qian served on the United States government's Science Advisory Board during World War II. On the
war front in Germany, he advised the Army on ballistic-missile guidance technology. At the war's
end, holding the temporary rank of lieutenant colonel, he debriefed Nazi scientists, including
Werner von Braun, and was sent to analyze Hitler's V-2 rocket facilities.

In the 1940s his mentor and colleague, the Caltech physicist Theodore von Karman, called Mr.~Qian
``an undisputed genius whose work was providing an enormous impetus to advances in high-speed
aerodynamics and jet propulsion.'' In 1949, Mr.~Qian wrote a proposal for a winged space plane that
the magazine Aviation Week and Space Technology, in 2007, called an inspiration for research that
led to NASA's space shuttle.

But by 1950 his American career was over. Shortly after applying for permission to visit his parents
in the newly Communist China, he was stripped of his security clearance by the Federal Bureau of
Investigation and accused of secretly being a Communist. The charge was based on a 1938 document of
the Communist Party of the United States that showed he had attended a social gathering that the
F.B.I. suspected was a meeting of the Pasadena Communist Party.

Mr.~Qian denied the charges, his Caltech colleagues came to his defense, and the university hired a
lawyer to assist him. Mr.~Qian first sought to return to China but was placed under virtual house
arrest by the government; later, he sought to stay and fight the accusations, but the government
sought to deport him.

In 1955, Mr.~Qian was sent back to China, where he was proclaimed a hero and immediately put to work
developing Chinese rocketry. By many accounts, he later became a committed Communist and served on
the party's ruling body, the Central Committee.

The loyalty allegations have never been fully resolved. Aviation Week, which named Mr.~Qian its man
of the year in 2007, quoted Dan Kimball, a former under secretary of the Navy, as calling Mr.~Qian's
deportation ``the stupidest thing this country ever did.'' A 1999 United States Congress report on
Chinese espionage called Mr.~Qian a spy, but critics say the report provides no basis other than a
claim that he passed to China the secrets of the American Titan missile program, which began years
after he had been deported.

Qian Xuesen was born in 1911, as the Chinese imperial government was collapsing, in Hangzhou, in
eastern China. He earned a mechanical engineering degree in 1934 in Shanghai. At the age of 23 he
went to the United States on a scholarship to study aeronautical engineering at M.I.T. Later, at
Caltech's Guggenheim Aeronautical Laboratory, Mr.~Qian met Mr.~von Karman, who recommended him for
the Science Advisory Board and gave him the lead role in research that developed the first American
solid-fuel rocket to be successfully launched.

After his deportation, Mr.~Qian wrote a position paper for Chinese leaders on aviation and defense,
according to the state-run news service Xinhua.

Under his leadership, China developed its first generation of ``Long March'' missiles and, in 1970,
launched its first satellite. Most of China's recent space achievements, like its manned space
program, began long after Mr.~Qian's retirement.

Mr.~Qian never returned to the United States. In a 2002 published reminiscence, a Caltech colleague
and professor, Frank Marble, stated that he believed that Mr.~Qian had ``lost faith in the American
government'' but that he had ``always had very warm feelings for the American people.''

Caltech gave Mr.~Qian its distinguished alumni award in 2001.

\section{China Dismisses Its Minister of Education}

\lettrine{F}{acing}\mycalendar{Nov.'09}{06} rising criticism over the quality of schools and a crush
of jobless college graduates, China's legislature announced Monday that it had removed the minister
of education after six years on the job and replaced him with a deputy.

The minister, Zhou Ji, had become a prime target for critics of China's education system, which has
stumbled during breakneck expansion that was intended to raise literacy rates and build a
world-class university system.

His dismissal follows a corruption scandal involving a university in Wuhan, where Mr.~Zhou had been
mayor and, before that, president of another university. Mr.~Zhou has not been publicly linked to
the corruption charges, which remain under investigation.

The government-run Chinese-language press reported Mr.~Zhou's removal on Monday largely without
comment in summarizing the work of the standing committee of the National People's Congress, the
legislature whose delegates are largely elected by Communist Party leaders and the military.

Mr.~Zhou's dismissal was described more fully in an official English-language newspaper, China
Daily, which said the education system had been ``plagued with problems, such as underfunding of
primary and secondary schools and poor standards in higher education.''

Many of those problems arose well before Mr.~Zhou became education minister in 2003, but he was
widely criticized for moving too slowly to correct them. When all 3,000 delegates of the National
People's Congress voted in March to retain or replace cabinet-level ministers, Mr.~Zhou drew 384 no
votes -- putting him in last place among the 72 ministers who were considered.

But there had been no hint that the government was considering replacing Mr.~Zhou. Indeed, he met
last month with the Russian prime minister, Vladimir V.~Putin, during Mr.~Putin's visit to Beijing,
an indication that he was in good standing with top leaders.

Late Monday, the prestigious Chinese Academy of Social Sciences stated on its Web site that Mr.~Zhou
would join the organization, further indicating his departure was unrelated to corruption charges.
Still, some Chinese newspaper columnists suggested Monday that Mr.~Zhou's departure offered the
government a chance to address broader corruption in academia, in which excellence and the search
for truth had been subverted by politics and the search for a fast buck.

Both basic and higher education have been hindered by corruption, including the selling of degrees
and stellar test scores by administrators as well as cheating among students.

Mr.~Zhou's successor, Yuan Guiren, may offer a preview of the ministry's new direction. Mr.~Yuan,
59, earlier was president of Beijing Normal University, known by educators as one of the most
progressive institutions in a nation where higher education is tightly bound by ideological and
political restraints.

China has poured billions of dollars into education in the last decade. The results are remarkable:
Higher-education enrollment has more than tripled since 2000, and China today awards more college
degrees than the United States and India combined. Annual awards of doctoral degrees rose sevenfold
between 1996 and 2006.

But critics say the quality of teaching has suffered, and in recent years universities have become
more politicized as Communist Party officials began to view a senior academic position as a ticket
to career success.

Elementary and secondary educational institutions have significantly raised literacy rates and
attendance, but schools are hamstrung by financing problems.

\section{Editor Departs China Magazine After High-Profile Tussle}

\lettrine{T}{he}\mycalendar{Nov.'09}{10} pioneering editor of the top Chinese business magazine has
left her post with plans to start anew, after a tussle\cite{tussle} for control involving much the
same mix of political and financial intrigue that she made her mark uncovering.

Hu Shuli, 56, resigned Monday from Caijing, the magazine she built into a thriving print and Web
outlet that specialized in investigating government corruption and corporate fraud, said a Caijing
spokeswoman, Zhang Lihui. Senior editors and most of Caijing's journalists either had already
resigned or were preparing to as well, magazine employees said.

For months, Ms.~Hu, the editor in chief, and the business managers of the magazine had been locked
in a stalemate with the owners of Caijing over the breadth of the magazine's coverage and the
budgeting of its operations, said former employees and current staff members who asked not to be
identified because they feared losing their jobs.

The owners of the magazine had come under pressure from Communist Party officials to rein in
Caijing's aggressive journalism, people at the magazine have said.

Managers at Caijing told staff members that they had been fighting to maintain the magazine's
editorial integrity.

The managers and editors had been seeking to create a more independent publication by changing the
magazine's shareholding structure, courting outside investors and pressing the owners to allow more
employees to own a stake in the magazine.

In a well-publicized exodus earlier this autumn, nearly 70 business employees resigned. Ms.~Hu held
on until Monday.

She has now accepted a new post as the dean of the journalism school at Sun Yat-sen University in
Guangzhou, a job she had been offered before it became clear that she would leave Caijing.

At the same time, she, along with a large contingent of editors and executives departing Caijing,
was working to secure new licenses and open a new venture, said the employees, who had knowledge of
the plans but were not authorized to speak publicly about them.

Caijing's parent company, the Stock Exchange Executive Council, or S.E.E.C., had already recruited a
new team of editors from another progressive publication, The Economic Observer in Beijing, they
said.

In 11 years at Caijing, editorials by Ms.~Hu pinpointed interest groups and bottlenecks that she
said blocked economic overhauls. And exclusives by Caijing hastened the demise of some of the more
notorious felons in China.

But the magazine's own troubles have involved just the sort of topic that Ms.~Hu and Caijing
relished covering.

The political price of success grew in recent years. Ms.~Hu found herself increasingly at odds with
S.E.E.C. bosses and their Communist Party guardians, according to employees and other colleagues
during interviews in recent months.

After a run-in with a Caijing reporter covering the ethnic riots in the western region of Xinjiang
in July, officials leaned harder on Ms.~Hu's superiors to curb\cite{curb} her coverage, the
employees said.

At one point, the S.E.E.C. was ordered to fire Ms.~Hu, they said. The pressures brought the
infighting over editorial and financial control of Caijing to a boil.

Ms.~Hu did not respond to requests for comment Monday.

Known for enforcing a rigid code of conduct, she has been characteristically guarded during the
crisis.

``I am still working on a good result,'' she wrote in an e-mail message to The New York Times late
last month.

Under her current plan, her new publishing sponsor would be the province-level Zhejiang Daily Press,
said the Caijing employees and a Zhejiang Daily editor.

She has been talking with well-known Chinese investors. Her proposed new publication's title has a
familiar ring: ``Caixin,'' short for ``Caijing Newsweek.''

The split reflects the divergence of interests in a media market still governed by party cadres,
said Zhan Jiang, a journalism professor at Beijing Foreign Languages University.

``Some people still stick to their ideals,'' he said. ``But management has become increasingly
concerned with profits, and increasingly conservative.''

Moreover, as the central authorities lavish official Chinese media giants with support to grow and
compete globally, they also have made moves to tighten their chain of command over muckrakers like
Ms.~Hu.

Not that Ms.~Hu is like any other. She has become an unrivaled celebrity, and counts senior economic
officials friends from her reporting days at state-owned newspapers.

At S.E.E.C., she was uniquely insulated. The chairman of the S.E.E.C., Wang Boming, a former New
York Stock Exchange economist, is the son of a former deputy foreign minister.

When Mr.~Wang and Ms.~Hu started Caijing, in 1998, he met her demands to finance the newsroom and
not interfere.

But their ambitions clashed as the influence of Caijing grew. Caijing now generates about half of
the group's revenue, but the S.E.E.C. has reinvested a considerably smaller percentage.

Mr.~Wang has diversified into less daring titles, most of which have struggled.

Members of Ms.~Hu's team, in turn, went their own way, expanding Caijing online. They also tapped
outside partners, like the Hong Kong tycoon Richard Li, with whom Ms.~Hu has been developing a
financial news service.

Behind the scenes, a conservative official named Quan Zhezhu had taken over Communist Party affairs
at the organization that sponsors Caijing's publishing license, the All-China Federation of Industry
and Commerce, in 2007, replacing the son of a liberal ex-party leader. The Federation ramped up
pressure on Mr.~Wang to curb coverage by Ms.~Hu.

``They say she's ungrateful, that without them the magazine would have been closed a long time
ago,'' a friend of Mr.~Wang's said.

An S.E.E.C. executive did not answer requests for comment in recent days.

Ms.~Hu was able to elude\cite{elude} serious trouble through the spring. After Caijing revealed a
corruption investigation into China Central Television earlier this year, government media officials
demanded that the story be recalled, the employees said.

But within a couple days, Caijing reposted the piece online and handed out hundreds of undistributed
magazines to delegates at the annual legislative sessions.

When the ethnic riots broke out in July, Ms.~Hu promptly dispatched three journalists to Urumqi. But
not all of them were able to obtain a permit to be there.

One day, at the official press center, a veteran reporter named Yang Binbin was caught carrying a
credential borrowed from a former coworker. When an official tried to search his laptop, he
resisted, and he and a guard scuffled. The police carted off the reporter for questioning, then sent
him back to Beijing.

To make matters worse, the altercation unfolded in front of a division chief from the party's
central propaganda department.

``Our pressures and conflicts had accumulated over a long time, but this incident was the fuse,''
said a Caijing colleague of Mr.~Yang, who himself declined to comment.

By mid-July, journalists said, the party's powerful Central Commission on Politics and Law discussed
the need to ``rectify'' Caijing. Propaganda authorities have reprimanded the magazine for at least
eight articles this year, including the China Central Television inquiry, and directed it to
``return to positive reporting on finance and economics.''

Under orders from the All-China Federation, the S.E.E.C. demanded the right to prescreen the
magazine before it went to print.

Ms.~Hu resisted the order. But the magazine was still required to cut at least three investigative
features, including one from Urumqi, and the Web site scrapped two new columns and left the
``Politics and Law'' section without new posts for three days in September, to avoid riling
officials.

At a gathering with Mr.~Wang in August, according to a friend of his in attendance, Mr.~Wang said
that officials had pressured him to fire Ms.~Hu. Mr.~Wang said that he would not go so far as to
dismiss the acclaimed newswoman and that, he told friends, the move would cause an international
scandal.

But his perceived failure to stand up to editorial pressures exacerbated the financial infighting
about ownership shares and budgets, to the point that Ms.~Hu and Mr.~Wang, as another journalist put
it, ``couldn't stand each other.''

In late September, Caijing's general manager and other executives led a walkout of more than 60
business staff members. As of last month, dozens of those who resigned had already started working
at what several said were Caixin's new offices.

For weeks, many journalists have been planning to follow Ms.~Hu to the new venture. But Ms.~Hu could
have to wait months for new publishing licenses, if the authorities approve them, the journalist and
others said.

``She hopes that having this new academic position will make it easier for her to negotiate'' to
start the new outlet, said the journalist, who was among those preparing to rejoin Ms.~Hu.

\section{Chinese Agencies Struggle Over Video Game}

\lettrine{I}{t}\mycalendar{Nov.'09}{10} could almost be a World of Warcraft game session -- two
competing titans, plotting against each other, swapping blows, embarked on a quest for a single
prize that only the stronger of them will claim.

But this is not virtual reality. The titans are two agencies of the Chinese government. And their
quest, during which they have traded a few blows in the past week, is for a potentially rich prize:
the power to regulate the real World of Warcraft, among the most popular online games in China.

The background: On Monday, the Chinese General Administration of Press and Publication ordered the
Shanghai-based operator of World of Warcraft, NetEase, to shut down its servers for World of
Warcraft. The agency said that it had rejected the company's application to become the new host of
the game's four million Chinese players.

But by Wednesday, the Ministry of Culture had struck back.

``In regards to the World of Warcraft incident, the General Administration of Press and Publication
has clearly overstepped its authority,'' a ministry official, Li Xiong, was quoted as saying in the
Economic Information Daily, a newspaper in Beijing. ``They do not have the authority to penalize
online gaming.''

The ministry said it had that authority. And it said NetEase was perfectly free to offer the game on
computers in China. The matter now appears destined for settlement by the State Council, the Chinese
government's cabinet.

Such bureaucratic hair-pulling might seem petty, were so much not at stake. Why the authority to
regulate video games should trigger such a fracas is not altogether clear. But on its face, the
defining aspect of the dispute involves money.

The online gaming industry in China is already huge, and growing fast. About 50 million people crowd
the Internet cafes of China on a regular basis to play. Revenues in 2008 rose about 50 percent to at
least \$2.9 billion, according to Alicia Yap, a Hong Kong analyst for Citi Investment Research and
Analysis. That is 10 times the revenue of just five years ago. IDC, a research company, has
predicted that annual revenue will reach \$6 billion by 2013.

In that context, the question of who decides what games go online -- and how they decide -- looms
large. It is perhaps especially important for game makers outside China, who have had trouble
cracking the vast Chinese market.

Of the 10 most popular games in China ranked by MMLC Group, a Beijing intellectual-property
consulting firm, only World of Warcraft, by Blizzard Entertainment, is American-made; two are South
Korean, and the rest were developed in China.

The press and publication administration has taken a hard line against outside involvement in the
industry, stating flatly last month that foreign investment in Chinese online gaming operations,
whether by joint ventures, cooperatives or other means, is forbidden.

The agency did not directly address the origin of the actual games, although it did bar foreigners
from providing technical support to Chinese companies and declared its authority over foreign
``service packs'' and other improvements to existing online games.

The Ministry of Culture was quick to fire back. Speaking at a conference of computer game
developers, its deputy minister, Tuo Zuhai, condemned the publication administration's ``surly
interference in domestic online game enterprises,'' and said his ministry was the foremost regulator
of online gaming, according to the Shanghai-based blog web2asia.com.

``This is a major issue of right and wrong, and is a matter of principle,'' the blog stated, quoting
a transcript of Mr.~Tuo's remarks. ``We will never compromise on such a matter of principle.''

In practice, some experts say, at least some of those bans could be easily sidestepped. The
proclamation may be part of a larger feud within the government, and perhaps in the business world
as well, over parceling out regulation of the booming industry.

Historically, the publication administration has had the power to censor and ban virtually anything
published, whether a book, a DVD or an online game. The Ministry of Culture has policed film and
other performing arts, including literary and audiovisual works.

The State Council sought to redefine this overlap in 2008, essentially giving the publication agency
the power to approve online games before they are made public, and assigning the Ministry of Culture
to police them once they appear on the Internet.

World of Warcraft fell between the cracks. Long popular among Chinese gamers, the role-playing game
hit a snag in June, when Blizzard dropped the previous operator of the game's Chinese franchise in
favor of NetEase. NetEase shut the game down while it reapplied for permission from the Ministry of
Culture and the publication agency.

The ministry swiftly approved the game, while the publications agency lagged. In September, after
the State Council issued a statement reaffirming the ministry's authority over games already online,
NetEase restarted World of Warcraft -- and drew the publication agency's wrath.

Which agency will win the regulatory battle remains unclear, although the Ministry of Culture, with
allies among other ministerial-level offices, is said to enjoy an edge. Regardless, there appears to
be much for both offices to do. The government this summer proclaimed its desire to clean up the
Internet, ridding it of pornography, gambling, violence and seditious material.

The ministry dived further into that Herculean task in the past week, announcing sanctions against
188 companies that it said were running unlicensed, vulgar or overly violent online games. NetEase
and World of Warcraft were conspicuously absent from the list.

\section{A Squeeze on Customers Ahead of New Rules}

\lettrine{B}{anks}\mycalendar{Nov.'09}{10} are struggling to make money in the credit card business
these days, and consumers are paying the price. Interest rates are going up, credit lines are being
cut and a variety of new fees are being imposed on even the best cardholders.

One recipient of new credit card terms is Anita Holaday, a 91-year-old in Florida, who received a
letter last month from Citibank announcing that her new interest rate was 29.99 percent, an increase
of 10 percentage points.

``I think it's outrageous they pursue such a policy,'' said Susan Holaday Schumacher, Ms.~Holaday's
daughter, who pays her mother's bills. ``That rate is shocking under any circumstances.''

While the average interest rates charged by banks are lower than Ms.~Holaday's, her situation is not
all that unusual. The higher rates and fees reflect the grim new realities of the credit card
industry -- the percentage of uncollectible balances has hit a record even as a new law may further
limit the cards' profitability.

Banks began raising interest rates and pulling back credit lines about a year ago as delinquencies
crept upward and regulators discussed reforms. As banks have become more aggressive in making
changes, lawmakers have accused them of trying to impose rate increases before many of the new rules
take effect in February.

On Monday, the Federal Reserve provided new evidence of the banks' actions. About 50 percent of the
banks responding to the Fed's survey said they were increasing interest rates and reducing credit
lines on borrowers with good credit scores. About 40 percent said they were imposing higher fees.
The banks also said they were demanding higher minimum credit scores and tightening other
requirements.

A study by the Pew Charitable Trusts, released late last month, concluded that the 12 largest banks,
issuing more than 80 percent of the credit cards, were continuing to use practices that the Fed
concluded were ``unfair or deceptive'' and that in many instances had been outlawed by Congress.

In response to voter complaints, the House of Representatives voted last week to make the law
effective immediately. The bill now goes to the Senate, where a vote has not been scheduled. The
Senate Banking Committee chairman, Christopher J.~Dodd, Democrat of Connecticut, meanwhile, is
pushing legislation that would freeze interest rates on existing credit card balances until the law
takes effect.

Whatever the starting date, the law makes it much harder for banks to change interest rates on
existing balances, and requires more time and notice before a new rate can go into effect.

In their defense, banking officials say they have no choice but to raise rates and limit credit.
Because of the new rules and the prolonged economic malaise, they say it is now far riskier to issue
credit cards than it was just a few years ago.

``We sell credit; we don't sell sweaters,'' said Kenneth J.~Clayton, senior vice president for card
policy at the American Bankers Association. ``The only way to manage your return is through the
price of the product or the availability.''

The nation's largest banks are scrambling to figure out a new business model that fits within the
new rules and current economic conditions. Those banks made handsome profits over the last decade by
charging high interest rates and penalty fees to a small group of customers who routinely paid late
or exceeded their balances.

Already, banks are shifting to a model in which a smaller pool of Americans will be eligible for
credit cards, and customers with cards will probably pay more for the privilege through annual fees
and higher interest.

Meanwhile, the banks are in the process of shedding customers considered too risky. That means tens
of thousands of Americans will no longer be able to splurge on Nike gym shoes or flat-screen
televisions unless, of course, they have enough cash to pay for them.

Still, even consumer advocates have said that the banks were too quick in the past to give out
credit. ``You know, it doesn't take a rocket scientist to figure out that if you keep borrowing and
borrowing in order to consume now, eventually you crash and burn,'' said Martin Eakes, chief
executive for the Center for Responsible Lending. ``That's what we're facing.''

In the 12 months that ended in September, the number of Visa, MasterCard, American Express and
Discover card accounts in the United States fell by 72 million, according to David Robertson,
publisher of The Nilson Report, an industry newsletter. There are 555 million accounts still in the
marketplace, he said.

In roughly the same time period, banks lowered credit limits by 26 percent, to \$3.4 trillion, from
\$4.6 trillion, according to an analysis of government data by Foresight Analytics.

Interest on credit card accounts, meanwhile, has increased to an average of 13.71 percent, up from
11.94 percent a year ago, according to federal records.

As to credit card charge-offs -- industry lingo for uncollectible balances -- the number tracks the
unemployment rate and, therefore, is hovering at around 10 percent.

For the banks, this is uncharted territory. In the modern financing era, credit cards were long a
profit center, producing tens of billions in annual profits with a default rate that hovered around
4 percent until the recession.

``We know we are going to lose a lot of money next year in cards, and it could be north of \$1
billion in both the first quarter and the second quarter. And that number will probably only start
coming down as you see unemployment and charge-offs come down,'' Jamie Dimon, chief executive of
JPMorgan Chase, said in an earnings call last month.

Banking officials said that because the new law limits their ability to reprice credit as a
customer's risk profile changes, they will instead have to price for future risk at the start, when
a cardholder applies for a new card.

That means fewer applicants will be approved for new credit cards, and those who are accepted will
increasingly be charged annual fees or variable interest rates, rather than fixed rates. Currently,
about 20 percent of credit cards charge annual fees, a percentage that is rising, said Bill
Hardekopf, chief executive of LowCards.com. Current cardholders, too, will be affected.

Asked to explain its rate increases, Citibank issued a statement saying the ``actions are necessary
given the losses across the industry from customers not paying back their loans and regulatory
changes that eliminate repricing for that risk.''

Ms.~Holaday Schumacher did not accept that explanation. She said she haggled with Citibank to try to
get her mother's bills forwarded to her house in Washington and, during the process, two bills were
inadvertently paid late, resulting in the rate increase.

``How unbelievably unfair for an older person who might not understand what this is all about,'' she
said. Citibank declined to comment on the account.

Still, many of the nation's banks are trying to repair their tarnished reputations with consumers.

American Express and Discover Financial, for instance, have vowed to stop charging fees when
cardholders exceed their credit limits. JPMorgan has started a program that can help consumers
categorize their spending and pay down their balances more quickly.

And Bank of America is promoting a line of consumer products so simple that the terms and conditions
fit on one page. The BankAmericard Basic Visa, for instance, has no rewards and a single interest
rate.

Andrew Rowe, Global Card Services strategy executive at Bank of America, said the new products
represented a sea change in the bank's attitude toward consumer products. Instead of benefiting from
consumers who displayed risky behavior by penalizing them with fees, the bank is now trying to help
them break those bad habits, he said.

``We succeed if our customers succeed,'' he said. ``That's the paradigm shift.''

Treasury Secretary Timothy F.~Geithner, for one, said he would welcome consumer products that were
simpler and less risky. But, he added in an interview with the PBS documentary program
``Frontline'': ``It's a bit of a late conversion. It would have been nice to happen earlier.''

\section{Burberry Looks Online for Ways to Gain Customers}

\lettrine{A}{ngela}\mycalendar{Nov.'09}{10} Ahrendts still remembers when she bought her first
Burberry trench coat. She was 21 years old, had just finished her studies in Indiana and was looking
for a smart but warm coat to wear for her first job at a small men's wear firm in New York.

Those raincoats, a 95-year-old fashion icon, remain Burberry's best-selling item, and Ms.~Ahrendts
-- who now runs the company -- is hoping to move the quintessentially British brand into the age of
the Internet to attract a new generation of shoppers.

On Monday, Burberry introduced a social networking site, artofthetrench.com, to encourage people to
share their own trench coat stories. It is the latest step by Ms.~Ahrendts and her creative
director, Christopher Bailey, to build on the brand's British heritage and trademark plaid with a
more modern twist.

``It's our differentiator,'' Ms.~Ahrendts said. ``It's not so different from what competitors do.
Maybe one was born from shoes and another from luggage; we come from a coat. It's our job to keep
that category hot and cool and relevant for all ages.''

The step reflects a broader move by luxury goods companies, which have generally failed to figure
out how to sell their wares online. Indeed, many have shunned the Web, seeing it as mostly a place
for bargain hunters to search for knock-offs or counterfeits.

As highly affluent, but aging, customers in developed countries cut back on their purchases of
luxury goods, the Web represents the prospect for growth.

Salvatore Ferragamo announced plans for an online store in October, following its Italian rivals
Prada and Bulgari to the Net. Faberg\'e, creator of the legendary Imperial Easter eggs, started
selling its new jewelry collection almost exclusively on the Web in September.

Ms.~Ahrendts, who is tall, charismatic and often dresses head to toe in Burberry, used to work at
Donna Karan and is bringing an American touch to the company based in London.

Burberry has been on the upswing for much of the decade, but since taking over three years ago,
Ms.~Ahrendts, 49, has cut \$82 million in costs, expanded in the United States and Asia, and opened
separate children's fashion stores.

She also added more shoes and handbags, including the \$6,200 ``Beaton'' bag worn by Victoria
Beckham, and took greater control of pricing by opening more Burberry stores.

As a result, Burberry has done better than most of its competitors in steering through the Great
Recession. Sales rose 4.6 percent to \textsterling343 million, or \$568 million, in the three months
to the end of September, beating analyst expectations. The price of its shares has doubled during
the past year.

For years, Burberry was seen as a limited brand, a potential takeover target for the likes of Coach,
Richemont and LVMH Mo\"et Hennessy Louis Vuitton. But the increase in its share price during the
past year has acted as a good defense.

Today, Burberry faces twin challenges: to maintain momentum amid a difficult economic recovery, and
to keep the expansion drive from diluting the brand.

``The biggest thing that keeps me up at night is how can we continue to evolve this organization in
order to stay ahead of the curve,'' Ms.~Ahrendts said during an interview at the company
headquarters last month. ``My job is to always look two to three years ahead and look round the
corner and see what's coming.''

The recession split luxury goods makers into different camps. Those more reliant on sales from
jewelry and watches, like Bulgari, were hit harder than those selling clothes and less expensive or
less ostentatious items. At the same time, companies that depended more on department stores
suffered most of all, largely because they had less control over pricing and inventories.

After a roughly 8 percent decline this year, the \$226 billion global market for luxury goods is
expected to grow again next year as younger consumers and working women replace retiring baby
boomers as the dominant consumer group, according to consulting firm Bain \& Company.

That's one of the main reasons Burberry is now focusing on the Internet. Ms.~Ahrendts said she gets
a lot of inspiration from her three children, who spend time surfing the Web and buy most of their
clothes online.

Ms.~Ahrendts said she is proud that Burberry has more than 699,000 Facebook fans. The company,
founded as a maker of outdoor wear in 1856 by the British draper's apprentice, Thomas Burberry, is
also attracting customers via Twitter and Youtube.

Burberry gets about two-thirds of its revenue from clothing. Ms.~Ahrendts has sharply reduced
Burberry's reliance on department stores; this year, direct retail overtook wholesale as the
company's biggest sales channel. Burberry, which produces the trench coat in Britain, also benefited
from a weaker pound that attracted euro-wielding tourists from mainland Europe.

But Burberry's major dependence on tourists and the benefit it has derived lately from a favorable
exchange rate worries some analysts.

To strengthen the brand, Ms.~Ahrendts divided Burberry's clothing lines more clearly. There's
Prorsum, which includes high-end couture and evening wear worn by celebrities like Gwyneth Paltrow
and Liv Tyler. The classic Burberry London clothes are designed to be office and everyday wear,
while Burberry Brit offers a more casual weekend look.

The separate children's line was almost too successful for its own good: Burberry received so much
demand this year that it ran into production and shipping problems.

Indeed, the trickiest balancing act is to manage further growth.

``They've got the balance sheet to expand when others in the industry can't, which means they can
really benefit once the market recovers,'' said Katharine Wynne, an analyst at Investec in London.
``But you have to make sure you keep that design heritage and retain that exclusivity.''

Its trench coat -- originally designed for the British army in 1914 and later associated with Audrey
Hepburn's character in ``Breakfast at Tiffany's'' and Humphrey Bogart's in ``Casablanca'' -- has
come a long way since then.

Thanks to Mr.~Bailey, the 38-year old British designer who also worked with Ms.~Ahrendts at Donna
Karan, there are now trench evening and mini dresses.

In Burberry's new headquarters along the Thames, its heritage is still visible through the
ultra-modern white-brown interior. Ms.~Ahrendts recently revived a 108-year old company trademark of
an equestrian knight carrying a flag with the words Prorsum -- Latin for forward. The image now
features prominently in the headquarters' glass-walled entrance hall.

Inside, a photo of Emma Watson, the young British actress starring in the Harry Potter movies,
modeling a \$1,300 trench coat covers an entire wall. Next to it, a glass cube showcases a
collection of the best-selling trench coat designs over the years.

Seeking to clearly separate Burberry's different designs, the company is preparing to open its first
stand-alone Burberry London and Burberry Brit stores this month, flanking the entrance to the
company's new Manhattan offices on Madison Avenue. It is also in talks with a licensing partner to
add a make-up and beauty range next year, and is planning to offer more men's accessories, such as
bags and scarves.

Ms.~Ahrendts recently renewed her contract and said she plans to be around for a while. Just like
that first trench coat, which she still wears from time to time.

\section{Justices Hear Patent Case on Protecting the Abstract}

\lettrine{S}{upreme}\mycalendar{Nov.'09}{10} Court justices took up a case on Monday that could
reshape the realm of what can be patented, and expressed skepticism about giving protection to
abstract business innovations.

During arguments, several justices, including the court's newest member, Sonia M.~Sotomayor, seemed
to disagree with arguments advocating a patent for a method of hedging.

The case, Bilski and Warsaw v. Kappos, concerned a business method patent that had been denied to
Bernard Bilski and Rand Warsaw, who in 1997 applied for a process that could help institutions like
utilities, or even factories and schools, have more predictable energy costs.

The justices peppered J.~Michael Jakes, a lawyer for Mr.~Bilski and Mr.~Warsaw, with hypothetical
patents that they clearly found ludicrous.

Justice Antonin Scalia suggested that under Mr.~Jakes's argument, a patent for ``somebody who writes
a book on how to win friends and influence people'' might be allowed, while Justice Sotomayor
suggested a ``method of speed dating.''

Justice Stephen G.~Breyer set off a ripple of laughter when he brought up his ``great, wonderful,
really original method of teaching antitrust law'' -- one in which 80 percent of students actually
stayed awake -- and asked if that could be patented.

Mr.~Jakes argued that some of the examples were potentially patentable, though other considerations
would be brought to bear by examiners, including the question of whether the method was obvious.

The patent process is a balancing act with origins in the United States Constitution, which called
for legal protection for authors and inventors ``to promote the progress of science and useful
arts'' for a limited time.

Initially, patents chiefly concerned tangible things like new machines and novel chemical compounds.
Over time, the rules about what could be patented have come to include increasingly abstract
business methods -- roughly, ways of doing things.

The high-water mark for such patents was a 1998 ruling by the United States Court of Appeals for the
Federal Circuit, State Street Bank v. Signature Financial Group, which has led to the issuance of
thousands of business method patents. In that case, the court ruled that a method of processing
mutual fund data could be patented.

In the Bilski case, the Patent and Trademark Office, pulling back from the realms of abstraction,
denied the patent. The United States Court of Appeals for the Federal Circuit upheld the patent
office's decision, writing that such patents should be ``tied to a particular machine or apparatus''
or transform something ``into a different state or thing.''

The case has drawn intense interest, and nearly 70 amicus briefs by interested parties, including
Microsoft, Google, Bank of America and the American Civil Liberties Union.

``I think this case is the case of the century for patent law,'' said John F.~Duffy, a professor at
the George Washington University Law School. Mr.~Duffy, who was co-author of a brief on behalf of
several technology companies, favors a broad reading of patent law for financial engineering tools
and other emerging technologies.

``There's a tremendous public benefit that could come from encouraging innovation in this space,''
he said. Similar views have been advocated in briefs from the biotechnology and pharmaceutical
industries, as well as technology companies like Yahoo.

Pamela Samuelson, a professor at the University of California, Berkeley, School of Law, said, ``It's
not very often that some obscure issue of patent law can excite so much attention.''

Professor Samuelson, who was an author of a brief on behalf of the Electronic Frontier Foundation,
an online civil liberties group, and others, said it was time for the court to tap the brakes on the
business patents rush. The earlier State Street decision, her brief stated, had the effect of
``knocking patent law loose from its historical moorings and improperly injecting patents into
business areas where they were neither needed nor wanted.''

Briefs from technology companies like Microsoft and Google also recommended greater restrictions on
business method patents.

During the one-hour oral argument on Monday afternoon, the Bilski patent and its ilk also seemed to
sit poorly with Justice Anthony M.~Kennedy and Chief Justice John G.~Roberts Jr., who raised
questions that suggested an interest in narrowing the scope of patents to more squarely focus on
physical inventions and not abstract ideas and processes.

Justice Kennedy described the beginnings of the insurance industry in the late 17th century, thanks
to the development of calculus and the ability to create actuarial tables. ``It's difficult for me
to think Congress would have wanted to give only one person the capacity to issue insurance,'' he
said.

Chief Justice Roberts wrangled with the lawyer for the government, Deputy Solicitor General Malcolm
L.~Stewart, over the final footnote in the government's brief, which conceded\cite{concede} that a
business method might be fit for patenting if it was tied to a computer. Chief Justice Roberts said
that that footnote ``takes away everything you spent 53 pages establishing.''

Mr.~Stewart noted that the government had actually argued against having the court take up the case,
since there were ``difficult problems out there'' in areas like software innovation and medical
diagnostics that have yet to be satisfactorily worked through.

Justice Kennedy broke in and joked that the government ``thought we would mess it up,'' eliciting
laughter from the gallery. Mr.~Stewart said instead that the patent at issue in this case was simply
``an unsuitable vehicle for resolving the hard questions,'' and that those were likely to remain
unresolved at the end of the day.

``But this case could be decided without making any bold steps'' that would complicate future
decisions, Justice Ruth Bader Ginsburg interjected.

\section{Obama Says He Will Raise Yuan Issue With China}

\lettrine{U}{}\mycalendar{Nov.'09}{10}.S. President Barack Obama said on Monday that he plans to
raise the issue of the yuan currency with Chinese officials when he meets with them in Beijing next
week, a potentially disruptive topic for foreign exchange markets.

``Currency, along with a host of other issues, will come up, and I'm confident that both the United
States and China can arrive at a broad set of policies that encourages trade that benefits both
countries, that allows ongoing economic growth,'' Mr.~Obama told Reuters in an interview.

U.S.~manufacturers complain that Beijing artificially holds the value of the yuan down to make its
exports cheaper, and American goods more expensive for Chinese consumers.

Economists say this has led to imbalances in the world economy by contributing to big trade deficits
in the United States and trade surpluses in China. Group of 20 leaders pledged at a summit in
Pittsburgh in September to aim for policies to ease these imbalances.

``That broader conversation will be at the center of our conversations with the Chinese
delegation,'' Mr.~Obama said.

Mr.~Obama has so far resisted domestic pressure to brand China as a currency manipulator, which
could anger a crucial U.S.~creditor, and when asked about this in the interview he
parried\cite{parry} the question.

``I think it is important as president of the United States, as we enter into these discussions,
that we are looking at all the issues involved and not just one,'' Mr.~Obama said.

He did, however, say that the two countries share a common interest in delivering sustainable growth
that will help rebalance the global economy.

``They have a huge amount of U.S.~dollars that they are holding, so our success is important to
them. The flip side of that is that if we don't solve some of these problems, then I think both
economically and politically it will put enormous strains on the relationship,'' he said.

Cheap Chinese exports to the United States have resulted in a massive trade imbalance between the
two countries. But the relationship is further complicated because Beijing has also invested the
dollars it earns from these exports in U.S.~government bonds.

But Mr.~Obama did acknowledge that the complaints about access to Chinese markets had some validity.

``Our manufacturers, I think, would have legitimate concerns about our ability to sell into China,''
he said, and emphasized that boosting U.S.~sales oversees was a crucial part of his economic agenda.

``It is particularly important for us when it comes to Asia as a whole to recognize that in the
absence of a more robust export strategy it is going to be hard for us to rebuild our manufacturing
base and employment base in this country.''

Mr.~Obama spoke to Reuters before his departure on Thursday for Japan, Singapore, China and South
Korea.

\section{Pigs Prove to Be Smart, if Not Vain}

\lettrine{W}{e}\mycalendar{Nov.'09}{10}'ve all heard the story of the third Little Pig, who foiled
the hyperventilating wolf by building his house out of bricks, rather than with straw or sticks as
his brothers had done. Less commonly known is that the pig later improved his home's safety profile
by installing convex security mirrors at key points along the driveway.

Well, why not? In the current issue of Animal Behaviour, researchers present evidence that domestic
pigs can quickly learn how mirrors work and will use their understanding of reflected images to
scope out their surroundings and find their food. The researchers cannot yet say whether the animals
realize that the eyes in the mirror are their own, or whether pigs might rank with apes, dolphins
and other species that have passed the famed ``mirror self-recognition test'' thought to be a marker
of self-awareness and advanced intelligence.

To which I say, big squeal. Why should the pigs waste precious mirror time inspecting their teeth or
straightening the hairs on their chinny-chin-chins, when they could be using the mirror as a tool to
find a far prettier sight, the pig heaven that comes in a bowl?

The finding is just one in a series of recent discoveries from the nascent study of pig cognition.
Other researchers have found that pigs are brilliant at remembering where food stores are cached and
how big each stash is relative to the rest. They've shown that Pig A can almost instantly learn to
follow Pig B when the second pig shows signs of knowing where good food is stored, and that Pig B
will try to deceive the pursuing pig and throw it off the trail so that Pig B can hog its food in
peace.

They've found that pigs are among the quickest of animals to learn a new routine, and pigs can do a
circus's worth of tricks: jump hoops, bow and stand, spin and make wordlike sounds on command, roll
out rugs, herd sheep, close and open cages, play videogames with joysticks, and more. For better or
worse, pigs are also slow to forget. ``They can learn something on the first try, but then it's
difficult for them to unlearn it,'' said Suzanne Held of the University of Bristol. ``They may get
scared once and then have trouble getting over it.''

Researchers have also found that no matter what new detail they unearth about pig acumen, the public
reaction is the same. ``People say, 'Oh yes, pigs really are rather clever, aren't they?' '' said
Richard W.~Byrne, a professor of evolutionary psychology at the University of St.~Andrews. ``I would
recommend that somebody study sheep or goats rather than pigs, so that people would be suitably
impressed to find out your animal is clever.'' His feigned frustration notwithstanding, he added,
``if you want to understand the evolution of intelligence and social behaviors, it's important to
work on animals like pigs that are not at all closely related to us'' but rather are cousins of
whales and hippos.

So far, and yet so near. Last week, an international team of biologists released the first draft
sequence of the pig genome, the complete set of genetic instructions for making the ruddy-furred
Duroc breed of Sus scrofus. Even on a cursory glance, ``the pig genome compares favorably with the
human genome,'' said Lawrence Schook of the University of Illinois at Urbana-Champaign, one of the
team leaders.

``Very large sections are maintained in complete pieces,'' he said, barely changed in the
100-million-plus years since the ancestors of hogs and humans diverged.

Dr.~Schook is particularly eager to see if the many physiological and behavioral parallels between
humans and pigs are reflected in our respective genomes. Pig hearts are like our hearts, he said,
pigs metabolize drugs as we do, their teeth resemble our teeth, and their habits can, too. ``I look
at the pig as a great animal model for human lifestyle diseases,'' he said. ``Pigs like to lie
around, they like to drink if given the chance, they'll smoke and watch TV.''

Pigs have been a barnyard staple for at least 8,000 years, when they were domesticated from the wild
boar in Asia and Europe. Domestication was easy, given that they loved to root around in dump sites.
``The pigs were hard to hunt, but if you put the garbage out, a lot of them would be drawn out from
the woods,'' Dr.~Schook said. ``After a while, people realized, we don't have to hunt them. All we
have to do is put a fence around our garbage.''

Pigs were tireless composting machines. ``They fed on our scraps,'' Dr.~Byrne said. ``Everything we
produced, they turned into good meat.'' Pork is among the world's most popular meats; in many
places, pigs are a valuable form of currency. ``In parts of New Guinea, they're so important to
villages that they're suckled by people,'' he said.

Of course, pigs aren't always handled so lovingly, and these researchers denounced factory farms.
``I'm German and I love sausage, but I would never eat pork that isn't free range,'' Dr.~Held said.

Even in domesticity, pigs have retained much of their foreboar's smarts. Dr.~Byrne attributes pig
intelligence to the same evolutionary pressures that prompted cleverness in primates: social life
and food. Wild pigs live in long-term social groups, keeping track of one another as individuals,
the better to protect against predation. They also root around for difficult food sources, requiring
a dexterity of the snout not unlike the handiness of a monkey.

Because monkeys had been shown to use mirrors to locate food, Donald M.~Broom of the University of
Cambridge and his colleagues decided to check for a similar sort of so-called assessment awareness
in pigs. They began by exposing seven 4-to-8-week-old pigs to five-hour stints with a mirror and
recording their reactions. The pigs were fascinated, pointing their snouts toward the mirror,
hesitating, vocalizing, edging closer, walking up and nuzzling the surface, looking at their image
from different angles, looking behind the mirror. When the mirror was placed in their pen a day
later, the glass-savvy pigs greeted it with a big ho-hum.

Next, the researchers put the mirror in the enclosure, along with a bowl of food that could not be
directly seen but whose image was reflected in the mirror. They then compared the responses of the
mirror-experienced pigs with a group of mirror-na\"ive pigs. On spotting the virtual food in the
mirror, the experienced pigs turned away and within an average of 23 seconds had found the food. But
the na\"ive pigs took the reflection for reality and sought in vain to find the bowl by rooting
around behind the mirror. No doubt the poor frustrated little pigs couldn't wait to get home, crack
open a beer and turn on the TV.

\section{Through 1 Woman, 20 Views of Life's End}

\lettrine{W}{hen}\mycalendar{Nov.'09}{10} Lance Armstrong learned he had testicular cancer, he
assembled a top team of doctors for what he decided would be the race of his life. If he did not
win, he would die.

After Mr.~Armstrong recovered and began racing again, he was a changed man, dominating the grueling
Tour de France -- unstoppable.

``Maybe the ultracompetitiveness came from the disease,'' he says, ``because the stakes were so
high.''

The story of Mr.~Armstrong's extraordinary comeback has been told many times, of course. But this
telling is different. The words are his, but they are spoken by the playwright and actress Anna
Deavere Smith in her one-woman show ``Let Me Down Easy,'' now playing Off Broadway.

The play consists of interviews by Ms.~Smith with 20 people -- from the famous, like Mr.~Armstrong,
to the unknown, like her aunt. Ms.~Smith, a tall, imposing woman, takes on the character of her
subjects by repeating their exact words, complete with regional accents, ums and uhs and stammers.
Her subject is medical care and end-of-life issues in America. It is not always a pleasing picture.

Some of her subjects deliver scathing indictments of disparities in care. A doctor at Charity
Hospital in New Orleans tells of being abandoned by the government in Hurricane Katrina, her
patients left behind while those at other hospitals were evacuated.

The doctor, Kiersta Kurtz-Burke, says that she had a veneer of innocence -- ``that you can take care
of poor people, and you can take care of them as well as you take care of rich people'' -- and that
``all those walls came down.''

But the play goes beyond stories about the nation's medical system or fighting battles with disease
and death. As its dramatic arc builds from Mr.~Armstrong's winning strategy to those who did not
win, the play illustrates much more than exhortations to fight. Instead, it shows why and how that
message can go awry -- and how people cope when they know they are dying.

In an interview with The New York Times, Ms.~Smith said she saw her play, 9 years and more than 300
interviews in the making, as part of a journey to listen to America.

``I try to embody America by embodying its words,'' she said. In this case, they are words about the
deepest human experiences: mortality and nobility of character, perseverance, hope and acceptance.

The play originated a decade ago, when the chairman of the department of medicine at Yale invited
her to spend time interviewing doctors and patients and then perform their words for the medical
school staff.

``I was a little timid about accepting,'' Ms.~Smith said. ``I kept putting it off.''

When she finally came, she was struck by the sort of conversations she had with her subjects. She
had just finished a play, ``House Arrest,'' involving interviews with politicians. Those subjects
were guarded, careful with their words, aware of their message.

``They were not able to speak to me from the heart,'' she said.

Patients were entirely different. ``All I had to do was turn on the tape recorder and say, 'What
happened?' '' she said. The patients were eloquent, their stories ``so powerful.''

Some were a little nervous. Ruth Katz, a breast cancer patient who is now chief public health
counsel for the House Energy and Commerce Committee, worked at Yale when Ms.~Smith was doing
interviews. The playwright ``would be performing me in front of people I work with,'' Ms.~Katz said
in an interview with The Times, adding, ``I would have to see them the next day.''

The play begins with the idea of a dedicated team that will get you through your illness.
Mr.~Armstrong boasts that he put together two such teams: the top doctors who cured him of cancer,
and the talented cyclists whose goal was to support him in his bid to win the Tour de France. Coming
in second was not an option.

The supermodel Lauren Hutton proudly states that she too had the best doctors -- and the best
doctors all refer patients to one another. ``The best always knows the best,'' she says.

Too bad everyone cannot have that. But what about Ms.~Katz, a professional, privileged woman who
should have had the best that medicine has to offer?

Ms.~Smith tells her story, in Ms.~Katz's eastern Pennsylvania accent: One day, she came to the Yale
hospital because she had gotten a fever after chemotherapy. An oncology fellow, a doctor in
training, told her that unfortunately they had lost her records. ``That happens here quite a bit,''
he said to her. So they would have to start all over and try to recreate the records from the
beginning. What kind of cancer do you have? Do you work? Full time? What do you do?

`` `I'm associate dean at the medical school,' '' she replied. ``Now he looks up, and he said, 'At
this medical school?' I said, 'At the Yale School of Medicine.'

``He found my files within half an hour.''

And what about the battles that are half won -- an athlete, for example, who regains health but can
no longer perform?

Michael Bentt, a champion heavyweight boxer, looked up and saw a light and realized he must be lying
on his back -- lights are on the ceiling, he reasoned. He must have been knocked out. Then he
realized the light was from a doctor's penlight. A neurologist was examining him.

It turned out he had been unconscious for four days. ``You can't fight anymore,'' the doctor said.

Mr.~Bentt thought about what being a champion meant to him. The fans, the women who, he recalled,
would fawn over him. ``You know, although I was conflicted about fighting, it defined me,'' he said.

Or what about the athlete who, inevitably, has to face the toll of aging?

``An athlete has two deaths,'' says Sally Jenkins, a sports columnist for The Washington Post.
``When someone like Lance reaches the end of the road, it's a form of death.'' Even when he finishes
third instead of first, people say, ``I don't want to remember him that way.''

Everyone has trouble with death, says Dr.~Philip Pizzo, dean of the medical school at Stanford.
Doctors are reluctant to have end-of-life discussions with patients, and not just because they are
uncomfortable with the subject or because they are afraid they may be taking away hope. Another
reason, Dr.~Pizzo admits, ``which I think is a hard one for me to say -- but I think nonetheless
that, uh, an honest reflection is that it takes a lot of time.''

Death is always going to be difficult, says Dr.~Eduardo Bruera, a palliative care specialist at the
University of Texas M.~D. Anderson Cancer Center in Houston. ``We're not built to assume with a
smile at the end of our lives,'' Dr.~Bruera says.

Some people, knowing they are going to die, sell all they own and sail off into the Caribbean,
Dr.~Bruera says. But that is not what most people do. They go on with the mundane activities of
their lives, as best they can, suffused with ``a profound existential sadness.''

People who were angry when healthy are angry in their final illness, those who were depressed are
depressed, those who were in denial are in denial.

And those who are entertainers entertain.

``I'm always acting,'' Joel Siegel, a movie critic for ABC News before his death from cancer, told
Ms.~Smith. ``It's like a character I play.''

But the anxiety never leaves. Every time he goes for a medical test, he says, ``it's like you're a
little kid going to school, and there's someone in the bathroom waiting to beat you up, and you have
to go to the bathroom.''

He discovered he was a fighter. But being a fighter is not always enough. Even a Lance Armstrong
cannot win every time.

Athletes, Ms.~Jenkins says, are like a cup falling to the ground and smashing: it will never fly
back up to the table.

In Buddhism, death is gentler, says Matthieu Ricard, a Buddhist monk. It is symbolized by a teacup
that is not smashed -- just turned upside down.

\section{A Dream Interpretation: Tuneups for the Brain}

\lettrine{I}{t}\mycalendar{Nov.'09}{10}'s snowing heavily, and everyone in the backyard is in a
swimsuit, at some kind of party: Mom, Dad, the high school principal, there's even an ex-girlfriend.
And is that Elvis, over by the pi\~{n}ata?

Uh-oh.

Dreams are so rich and have such an authentic feeling that scientists have long assumed they must
have a crucial psychological purpose. To Freud, dreaming provided a playground for the unconscious
mind; to Jung, it was a stage where the psyche's archetypes acted out primal themes. Newer theories
hold that dreams help the brain to consolidate emotional memories or to work though current
problems, like divorce and work frustrations.

Yet what if the primary purpose of dreaming isn't psychological at all?

In a paper published last month in the journal Nature Reviews Neuroscience, Dr.~J.~Allan Hobson, a
psychiatrist and longtime sleep researcher at Harvard, argues that the main function of
rapid-eye-movement sleep, or REM, when most dreaming occurs, is physiological. The brain is warming
its circuits, anticipating the sights and sounds and emotions of waking.

``It helps explain a lot of things, like why people forget so many dreams,'' Dr.~Hobson said in an
interview. ``It's like jogging; the body doesn't remember every step, but it knows it has exercised.
It has been tuned up. It's the same idea here: dreams are tuning the mind for conscious awareness.''

Drawing on work of his own and others, Dr.~Hobson argues that dreaming is a parallel state of
consciousness that is continually running but normally suppressed during waking. The idea is a
prominent example of how neuroscience is altering assumptions about everyday (or every-night) brain
functions.

``Most people who have studied dreams start out with some predetermined psychological ideas and try
to make dreaming fit those,'' said Dr.~Mark Mahowald, a neurologist who is director of the sleep
disorders program at Hennepin County Medical Center, in Minneapolis. ``What I like about this new
paper is that he doesn't make any assumptions about what dreaming is doing.''

The paper has already stirred controversy and discussion among Freudians, therapists and other
researchers, including neuroscientists. Dr.~Rodolfo Llin\'as, a neurologist and physiologist at New
York University, called Dr.~Hobson's reasoning impressive but said it was not the only physiological
interpretation of dreams.

``I argue that dreaming is not a parallel state but that it is consciousness itself, in the absence
of input from the senses,'' said Dr.~Llin\'as, who makes the case in the book ``I of the Vortex:
From Neurons to Self'' (M.I.T., 2001). Once people are awake, he argued, their brain essentially
revises its dream images to match what it sees, hears and feels -- the dreams are ``corrected'' by
the senses.

These novel ideas about dreaming are based partly on basic findings about REM sleep. In evolutionary
terms, REM appears to be a recent development; it is detectable in humans and other warm-blooded
mammals and birds. And studies suggest that REM makes its appearance very early in life -- in the
third trimester for humans, well before a developing child has experience or imagery to fill out a
dream.

In studies, scientists have found evidence that REM activity helps the brain build neural
connections, particularly in its visual areas. The developing fetus may be ``seeing'' something, in
terms of brain activity, long before the eyes ever open -- the developing brain drawing on innate,
biological models of space and time, like an internal virtual-reality machine. Full-on dreams, in
the usual sense of the word, come much later. Their content, in this view, is a kind of crude test
run for what the coming day may hold.

None of this is to say that dreams are devoid of meaning. Anyone who can remember a vivid dream
knows that at times the strange nighttime scenes reflect real hopes and anxieties: the young teacher
who finds himself naked at the lectern; the new mother in front of an empty crib, frantic in her
imagined loss.

But people can read almost anything into the dreams that they remember, and they do exactly that. In
a recent study of more than 1,000 people, researchers at Carnegie Mellon University and Harvard
found strong biases in the interpretations of dreams. For instance, the participants tended to
attach more significance to a negative dream if it was about someone they disliked, and more to a
positive dream if it was about a friend.

In fact, research suggests that only about 20 percent of dreams contain people or places that the
dreamer has encountered. Most images appear to be unique to a single dream.

Scientists know this because some people have the ability to watch their own dreams as observers,
without waking up. This state of consciousness, called lucid dreaming, is itself something a mystery
-- and a staple of New Age and ancient mystics. But it is a real phenomenon, one in which Dr.~Hobson
finds strong support for his argument for dreams as a physiological warm-up before waking.

In dozens of studies, researchers have brought people into the laboratory and trained them to dream
lucidly. They do this with a variety of techniques, including auto-suggestion as head meets pillow
(``I will be aware when I dream; I will observe'') and teaching telltale signs of dreaming (the
light switches don't work; levitation is possible; it is often impossible to scream).

Lucid dreaming occurs during a mixed state of consciousness, sleep researchers say -- a heavy dose
of REM with a sprinkling of waking awareness. ``This is just one kind of mixed state, but there are
whole variety of them,'' Dr.~Mahowald said. Sleepwalking and night terrors, he said, represent
mixtures of muscle activation and non-REM sleep. Attacks of narcolepsy reflect an infringement of
REM on normal daytime alertness.

In study published in September in the journal Sleep, Ursula Voss of J.~W. Goethe-University in
Frankfurt led a team that analyzed brain waves during REM sleep, waking and lucid dreaming. It found
that lucid dreaming had elements of REM and of waking -- most notably in the frontal areas of the
brain, which are quiet during normal dreaming. Dr.~Hobson was a co-author on the paper.

``You are seeing this split brain in action,'' he said. ``This tells me that there are these two
systems, and that in fact they can be running at the same time.''

Researchers have a way to go before they can confirm or fill out this working hypothesis. But the
payoffs could extend beyond a deeper understanding of the sleeping brain. People who struggle with
schizophrenia suffer delusions of unknown origin. Dr.~Hobson suggests that these flights of
imagination may be related to an abnormal activation of a dreaming consciousness. ``Let the dreamer
awake, and you will see psychosis,'' Jung said.

For everyone else, the idea of dreams as a kind of sound check for the brain may bring some comfort,
as well. That ominous dream of people gathered on the lawn for some strange party? Probably
meaningless.

No reason to scream, even if it were possible.

\section{Stuck by a Needle, Not by a Decision}

\lettrine{I}{}\mycalendar{Nov.'09}{10} didn't think it would happen to me so soon, just a few months
after beginning my second career as a nurse. I stuck my thumb with a large-bore needle filled with
the blood of a patient with hepatitis C who had come to the emergency room with abdominal pain.

This happened with a 10-milliliter syringe I was using to transfer blood from one tube into another.
I was trained not to do this; it was a bad idea. But I put my patient's comfort above my own safety:
when I learned an extra blood test had been ordered, I hoped to save him the pain of a second needle
stick.

I thought better of my decision to make the transfer and hesitated -- just as I noticed the needle
bending while I struggled to pierce the tough rubber top of the specimen tube. Two drops of blood
came out of the tip of the needle. Afterwards, I saw one bloody smudge on my glove; I feared the
other drop had gone into my thumb.

I froze, breathless.

With some prompting from my colleagues, protocol kicked in: I sped to the sink, bled the wound,
washed it with soap and hot water and checked the patient's record for H.I.V. status. There was
nothing. I asked him if he would mind being tested; taking pity on me, he agreed, though he assured
me he was negative. I ran to the H.I.V. counselor in a nearby unit and told him what had happened.
Within five minutes, the news was in: a rapid H.I.V. test was negative. For the moment, anyway.

Before I knew it, there I was in the employee health office with the physician who recognized me
from my physical on my first day of work 10 months earlier, wiggling his eyebrows and tapping his
pen, waiting for a yes or no: Did I want PEP treatment or not?

The letters stand for the post-exposure prophylaxis: a month of treatment that would, on one hand,
afford me some peace of mind but, on the other, make me sick as a dog. It is a decision many unlucky
health care workers have to make, and it is one that can't wait.

To be effective, PEP treatment for the AIDS virus must begin immediately after injury and must
continue, uninterrupted, for about four weeks. It involves a combination of multiple antiretroviral
drugs, whose side effects routinely include bouts of abdominal cramps, diarrhea, nausea and
vomiting.

The doctor encouraged me to consent to treatment. ``You'll feel sick,'' he said, ``but you can still
work.''

I listened to him talk, but I knew from the first word that I would refuse. All I could think of
were two nurses that, in my work as a writer, I had interviewed for an article about this very
subject.

One opted for PEP after being stuck by a bloody syringe a medical resident had left in the bed
linens of a patient dying from AIDS in a city hospital. By her own account, and those of the
physicians and lawyers involved in her (victorious) negligence lawsuit against the city, she
suffered six months of physical and psychological side effects from the ordeal, despite never
testing positive.

The other nurse was not so lucky. She opted against PEP after being stuck with a needle of a patient
whose H.I.V. status was unknown at the time. She subsequently tested positive for the virus and is
now an outspoken advocate for the use of safer hospital equipment.

Like me, both had been stuck with big, bloody needles. And despite their experiences, statistics
show that the risk of contracting H.I.V. from a needle stick is very low. What matter are the
patient's viral level and the depth of the wound: the higher the level and the deeper the wound, the
higher the risk.

In my case, I calculated, even if the rapid test were a false negative (meaning the patient actually
had the virus but the test didn't yet show it), the patient's viral load couldn't have been very
high. And the needle didn't go in very far. So I refused PEP, confident that the odds were in my
favor.

As the ink dried on the paperwork, the employee health nurse prepared supplies for the first of six
blood tests to which I would submit at prescribed intervals over the next year (three so far have
all been negative). I expect not to feel out of the woods until the six-month test comes out clean,
but I don't worry.

In fact, this essay is the first anyone -- except my co-workers and my husband -- will know about
this. I've kept it to myself, knowing that in a conversation with loved ones about risk, emotions
reign supreme over reason. Statistics are meaningless to a mother worried that one of her worst
nightmares is coming true. For her, an average 1.8 percent risk of contracting hepatitis C and an
0.3 percent risk for H.I.V. adds up to 100 percent sleepless nights.

\section{Obama Presses Senate to Pass Its Health Bill}

\lettrine{T}{he}\mycalendar{Nov.'09}{10} White House, growing concerned that the Congressional
timetable for passing a health care overhaul could slip into next year, is stepping up pressure on
the Senate for quick action, with President Obama appearing Sunday in the Rose Garden to call on
senators to ``take up the baton and bring this effort to the finish line.''

Mr.~Obama's remarks came just 14 hours after the House narrowly approved a landmark plan that would
cost \$1.1 trillion over 10 years and extend insurance coverage to 36 million uninsured Americans;
the president called it ``a courageous vote.'' But the votes had barely been counted when the White
House began turning its attention to an even bigger hurdle: getting legislation passed in the
Senate.

In the Senate, where proposals differ substantially from the House-passed measure on issues like a
government-run plan and how to pay for coverage, the bill is stalled while budget analysts assess
its overall costs. The slim margin in the House -- the bill passed with just two votes to spare, and
39 Democrats opposed it -- suggests even greater challenges in the Senate, where the majority
leader, Harry Reid of Nevada, is struggling to hold on to all 58 Democrats and two independents in
his caucus.

Mr.~Obama has staked his domestic agenda on passing comprehensive health legislation, a goal that
has eluded\cite{elude} presidents for decades. While Democrats were forced to make major concessions
on insurance coverage for abortions to win House passage of the bill, they were nonetheless
ebullient on Sunday, with many saying the vote gave them momentum to push the bill forward.

``For years we've been told that this couldn't be done,'' Mr.~Obama said in the Rose Garden. Of the
American people, he said, ``Moments like this are why they sent us here.''

But for all the exultation, there was a sense inside the White House and on Capitol Hill that the
hardest work is yet to come. The House debate highlighted the pressures that will come to bear on
senators as they weigh contentious issues like federal financing for abortion, coverage for illegal
immigrants and the ``public option,'' a government-backed insurance plan to compete with the private
sector.

In the Senate, Mr.~Reid has merged two bills into one. The fine print is not public, but the broad
outlines are known. Unlike the House bill, which pays to extend coverage by taxing individuals who
earn more than \$500,000 a year and couples who earn more than \$1 million, the Senate bill imposes
a 40 percent excise tax on so-called Cadillac plans that cost more than \$8,000 a year for an
individual or \$21,000 for a family.

And unlike the House bill, which includes a national public plan, the Senate measure would allow
states to opt out. But even that is too much government involvement for moderates like Senator
Joseph I.~Lieberman of Connecticut, a Democrat-turned-independent, who pledged Sunday to wage a
filibuster to block any plan with a public option in it.

``If the public option plan is in there,'' Mr.~Lieberman said on ``Fox News Sunday,'' ``as a matter
of conscience, I will not allow this bill to come to a final vote.''

Apart from substantive hurdles, the Senate bill faces procedural ones; Mr.~Reid cannot bring it to
the floor for debate until he gets an analysis, or ``score,'' from the nonpartisan Congressional
Budget Office, expected later this week. The delay could push Senate consideration of the bill until
after Thanksgiving, which could in turn make it very difficult for Congress to meet Mr.~Obama's goal
of signing a health bill into law by the end of this year.

The timing is crucial. Administration officials say Mr.~Obama wants to wrap up work on health care
so that he can turn his attention to other legislative priorities, including passing an energy bill
and revamping financial regulations. But White House officials also know that the closer the final
vote comes to the November 2010 midterm Congressional elections, the more difficult it will be to
pass legislation.

Sending members of Congress home over an extended Christmas break without a health care bill in hand
could prove disastrous politically. Democrats remember well the setbacks they suffered over the
August recess when the Senate Finance Committee failed to meet Mr.~Obama's deadline for finishing
its measure, and lawmakers were pummeled in town-hall-style meetings around the country.

``The holiday break is viewed the same way as the August break,'' said one Democrat close to the
White House, speaking anonymously to discuss strategy. ``We don't want a repeat. We could probably
survive it, but why take the chance?''

The White House began prodding Mr.~Reid to move quickly even before Saturday's House vote. In a
private meeting with Mr.~Obama this year, Mr.~Reid pledged to work to finish the measure by the end
of December. But last Tuesday, Mr.~Reid said the Senate was ``not going to be bound by any
timelines.''

On Wednesday, Rahm Emanuel, the White House chief of staff, visited Mr.~Reid. The two met on Capitol
Hill to ``continue the discussion on ways to get a bill done by the end of the year,'' said
Mr.~Reid's press secretary, Jim Manley, adding that Mr.~Reid intends to bring the bill to the floor
``as quickly as possible.''

In case the leader did not get the message, Mr.~Obama reinforced it Saturday night. In a statement
after the House vote, he said he looked forward to signing comprehensive health legislation ``by the
end of the year.''

A big question is whether Mr.~Reid has the 60 votes that will almost certainly be necessary to
permit debate to begin. Mr.~Manley said Democrats hoped ``the momentum from the House bill will make
everyone realize that the Senate should at least have a chance to begin debate.'' Still, he conceded
that there was ``no glide path'' toward getting the Senate to actually pass the measure.

As the Senate vote draws closer, the fight on the airwaves, where groups for and against the health
bill are already spending millions of dollars on advertising, will only intensify. Republicans are
also intensifying their opposition as they try to cast Democrats as tax-and-spend liberals, a theme
echoed Sunday by Senator Mitch McConnell of Kentucky, the Republican leader.

``Soon, Senate Democrats will propose their own version,'' Mr.~McConnell said in a statement. ``We
don't know how big it will be or how expensive, but we do know with certainty that it will mean
higher premiums, higher taxes and massive cuts to Medicare to create even more government programs.
That's not reform.''

Representative Pete Sessions of Texas, the chairman of the National Republican Congressional
Committee, issued a fund-raising appeal within hours of the vote. In an e-mail message to
supporters, he said of Democrats, ``Their unprecedented power grab will further bankrupt America
while destroying the finest medical system in the world.''

But Representative Chris Van Hollen of Maryland, the chairman of the Democratic Congressional
Campaign Committee, had his own fund-raising appeal, an e-mail message sent to supporters shortly
after midnight Sunday that said ``thank you for helping to make this historic day.'' Below the text
was a bold blue icon, a link to the committee's fund-raising site, that proclaimed in block letters,
``Contribute Now.''

\section{China's Tough Flu Measures Appear to Be Effective}

\lettrine{C}{hanggang}\mycalendar{Nov.'09}{12}, China -- Few farmers in this southern village gave
much thought to the epidemic that had begun spreading rapidly in the United States early this summer
until the authorities sealed its 100 residents off from the outside world for about a week. It
turned out that a visitor from California had shown symptoms of the swine flu virus, or H1N1, when
he arrived for a funeral.

Quarantines\cite{quarantine} and medical detentions are among the aggressive measures that Chinese
officials have taken to slow the transmission of H1N1, which quickly spread worldwide after being
diagnosed first in North America.

To protests from around the world, China isolated entire planeloads of people entering the country
if anyone on the plane exhibited flulike symptoms. Local authorities canceled school classes at the
slightest hint of the disease and ordered students and teachers to stay home. China was virtually
alone in taking such harsh measures, which continued throughout most of the summer.

Now, Chinese and foreign health officials say that some of those contested measures -- more easily
adopted by an authoritarian state -- may have helped slow the spread of the disease in the world's
most populous country. China has not had to cope with a crush of cases, and it began administering a
vaccine for swine flu in early September, the first country to do so.

Foreign officials also say that China demonstrated an unusual openness to sharing information about
H1N1 with its citizens and other governments, in contrast to its secretive approach to the near
pandemic\cite{pandemic} of severe acute respiratory syndrome, or SARS, a few years ago.

That is not to say that China has been spared. On Tuesday, Health Ministry officials reported an
``explosive'' growth of H1N1 on the mainland because of the onset of winter, with 5,000 new cases in
the previous three days pushing the total to more than 59,000.

At least 30 people have died here after contracting H1N1.

Exact data on the virus are hard to pin down; many more cases are suspected than confirmed, and
countries often use different methods to identify cases. Still, the indications in China are much
more positive than those in India. Like China, India has more than a billion people, many living in
poor, rural conditions, and was exposed to the virus after it had been identified in the West. The
Indian Health Ministry has reported 505 deaths.

The United States, where the virus was spreading before it was identified in the spring, has
reported more than two million cases and about 4,000 deaths in a population of 300 million.

``I think there were a variety of measures put in place by different countries, and it's difficult
to say what worked best and what didn't, but China's has worked very well,'' said Dr.~Michael
O'Leary, the director of the Beijing office of the World Health Organization.

As of August, 56 million people had been screened for flulike symptoms at China's borders, said Feng
Zijian, director of the emergency office of the Chinese Center for Disease Control and Prevention.
Mr.~Feng said he did not know the number of travelers who had been quarantined. The United States
Embassy in Beijing said that 2,046 American citizens had been quarantined by the end of October,
with 215 testing positive for H1N1.

``If these strict measures had not been taken, and if there had been a sudden outbreak of the
disease, there would have been a huge panic among the Chinese population,'' Mr.~Feng said.
``Although there were many criticisms from outside, people should understand China's
considerations.''

But Mr.~Feng and Dr.~O'Leary also say that the social and financial costs of China's tough measures
will have to be evaluated to see whether they were worth the benefits. And it is unclear how
decisive those actions were in slowing the transmission of H1N1 -- the summer heat in much of China
was a critical factor in slowing the spread, and most schools were out of session at the time.

Some foreign officials are still skeptical of the need for the strict quarantine measures, saying
that China should have re-evaluated its policies by June, when it was apparent that the disease was
not as lethal as initially feared. The State Council, China's cabinet, did not decide to relax the
quarantine policy until July.

From the beginning, the W.H.O. has said that tightening borders would not keep the disease out, and
that closing borders or automatically quarantining specific groups of travelers -- as China did for
a brief period with holders of Mexican passports -- would have no benefit.

Quarantines of entire school groups from overseas ignited outrage and led some American officials to
complain to Beijing. The State Department implicitly criticized the Chinese policies by issuing
travel warnings on the quarantine procedures.

In July, a group of 65 students and seven chaperons from St.~Mary's School in Medford, Ore., was
quarantined twice. The first time came after a girl pulled aside at the airport in Beijing tested
positive for H1N1. Then in Henan, a boy running a high fever also tested positive, leading to the
second quarantine. During that time, a dozen students tested positive for H1N1. Most of the students
and chaperons flew back to the United States on July 31, having spent 12 of 17 days of their trip in
quarantine.

``At the time, it seemed extreme, and it seemed restrictive, because I had never experienced an
infectious disease outbreak,'' said Scott Dewing, director of technology at the school and one of
the trip chaperons. ``Now, looking back and seeing some of the measures that are being taken now in
the U.S., the Chinese measures don't seem so extreme.''

Chinese and Western officials say that Chinese leaders put in place a comprehensive plan for a
pandemic outbreak after the disastrous experience of SARS. This includes, at least in the first
stages, some of the stringent quarantine measures taken for SARS but also emphasizes education. A
red banner hanging from the balcony of a rural school building here in Guangdong Province says,
``H1N1 flu is preventable, controllable and curable, and not terrifying.''

The government was so anxious to stay ahead of H1N1 that officials decided in June to start
developing a vaccine even though testing kits for measuring the dosage of the agent in the
experimental vaccines had not arrived from the W.H.O., said Zhao Kai, a virologist who advises the
government. It was an unusual step, but on Sept.~5 China became the first country to declare that it
had developed a vaccine, and by late October it had produced nearly 53 million doses.

\section{Japan Cools to America as It Prepares for Obama Visit}

\lettrine{P}{resident}\mycalendar{Nov.'09}{12} Obama will arrive in Tokyo on Friday, at a time when
America's relations with Japan are at their most contentious since the trade wars of the 1990s --
and back then, the fights were over luxury cars and semiconductors, not over whether the two
countries should re-examine their half-century-old strategic relationship.

When Japan's Democratic Party came to power in September, ending 50 years of largely one-party
government, Obama administration officials put on an outwardly positive face, congratulating the
newcomers. But quietly, some American officials expressed fears that the blunt criticism that the
Japanese had directed at the United States during the political campaign would translate to a more
contentious relationship.

Within weeks, those fears started to play out. The new Japanese government said the country would
withdraw from an eight-year-old mission in the Indian Ocean to refuel warships supporting American
efforts in Afghanistan.

The government also announced that it planned to revisit a 2006 agreement to relocate a Marine
airfield on Okinawa to a less populated part of the island, and to move thousands of Marines from
Okinawa to Guam.

And Japanese government officials have suddenly lost their shyness about publicly sparring with
American officials, as evident in a dispute in September between Japan's ambassador to the United
States, Ichiro Fujisaki, and the Pentagon.

Meanwhile, Japan's new prime minister, Yukio Hatoyama, has called for a more equal relationship with
the United States, and his government wants a review of the status of forces agreement, which
protects American troops from Japanese legal prosecution. Japanese citizens, and Okinawans in
particular, have demanded such a review for years.

When Mr.~Hatoyama met Mr.~Obama in New York during the United Nations General Assembly in September,
the conditions seemed ripe for a kiss-and-make-up session. At their initial meeting, Mr.~Obama
congratulated Mr.~Hatoyama ``for running an extraordinary campaign'' and complimented his party for
``leading dramatic change in Japan.''

Mr.~Hatoyama responded with the usual diplomatic niceties, telling reporters after the meeting that
``I told President Obama that the Japan-U.S.~alliance will continue to be the central pillar, key
pillar of the security of Japan and Japanese foreign policy.''

But there were also a few awkward moments. Mr.~Hatoyama and his wife, Miyuki, were the last to
arrive at a leaders' dinner at the Phipps Conservatory on the margins of the Group of 20 economic
summit meeting in Pittsburgh later that week in September. Mr.~Obama and his wife, Michelle, had
been greeting arriving guests for almost two hours. ``I'm sorry we were late,'' Mrs.~Hatoyama
apologized.

A few days later, after the Obamas and the Hatoyamas flew to Copenhagen to lobby the International
Olympic Committee for the 2016 Olympics, Tokyo beat out Chicago in the first round of voting, then
was bumped as Rio de Janeiro took the prize.

But all of that paled in comparison with the uproar that erupted in Japan after Defense Secretary
Robert M.~Gates visited Tokyo in October. Mr.~Gates, known for speaking bluntly, pressed
Mr.~Hatoyama and Japanese military officials to keep their commitment on the military agreements.

``It is time to move on,'' Mr.~Gates said, calling Japanese proposals to reopen the base issue
``counterproductive.'' Then, adding insult to injury in the eyes of Japanese commentators, Mr.~Gates
turned down invitations to attend a welcoming ceremony at the Defense Ministry and to dine with
officials there.

In the weeks since, in advance of Mr.~Obama's visit, both countries have taken pains to tone down
the rancor. The Japanese government has sent several high-level officials, including members of
Parliament, to Washington to take the political temperature. Besides meeting with Obama
administration officials, the Japanese representatives have spoken with members of research and
policy groups based in Washington, particularly experts on foreign policy issues related to Japan.

``The feelers they've been putting out is, `Please don't push us to make a decision because if you
do, you'll hear what you don't want to hear,' '' said Andrew L.~Oros, a professor at Washington
College and the author of ``Normalizing Japan: Politics, Identity and the Evolution of Security
Practice.''

Japan's new government is ``trying to backtrack from some of their campaign rhetoric, but it's too
soon,'' Mr.~Oros said.

``This was a historic election,'' he added. ``They overturned 50 years of conservative rule. They
can't do everything at once.''

Indeed, the new government is under political pressure at home. More than 20,000 Okinawa residents
held a protest rally against the base last week, and residents have been vociferous in letting the
government know that they expect it to keep its campaign promises.

Administration officials said they had no intention of letting the relationship slide. Mr.~Obama
will be ``looking to build his relationship and his personal ties with the new D.P.J.~government
there,'' Jeffrey A.~Bader, Mr.~Obama's senior director for East Asian affairs, told reporters on
Monday, using the initials for the Democratic Party of Japan. ``This government is looking for a
more equal partnership with the United States. We are prepared to move in that direction.''

But the United States, while tamping down the tone of the discussion, is still pressing Japan,
particularly on the Okinawa base issue. Mr.~Obama, in an interview on Tuesday with NHK television of
Japan, said Japan must honor the agreement.

While ``it's perfectly appropriate for the new government to want to re-examine how to move
forward,'' Mr.~Obama told NHK, he added that he was ``confident that once that review is completed
that they will conclude that the alliance we have, the basing arrangements that have been discussed,
all those things serve the interest of Japan and they will continue.''

In an effort to defuse tensions and perhaps make up for saying it would not refuel the Indian Ocean
warships, Japan said Tuesday that it would sharply increase its nonmilitary aid to Afghanistan,
pledging \$5 billion for a variety of projects that include building schools and highways, training
police officers, clearing land mines, and rehabilitating former Taliban fighters.

But even if the military squabble is eventually resolved, Japan's economic relationship with the
United States is being altered. China has now surpassed the United States as Japan's major trading
partner, a switch that economists expect to continue as China's economy grows.

``Japan sees its future more within Asia,'' said Eswar S.~Prasad, an Asia specialist and professor
at Cornell University. ``They feel that they owe a lot less to the U.S.~right now. U.S.~economic
policy is hurting them in a lot of ways, particularly with the decline in the value of the dollar
versus the yen.''

\section{Video Bingo Has Alabamians Yelling Everything But}

\lettrine{E}{verybody}\mycalendar{Nov.'09}{12} knows what this is: dozens of people, mostly
retirees, hunched over paper grids in a smoke-filled American Legion hall on a Sunday evening
listening eagerly to a woman recite numbers.

But what about this: a dim warehouse of flashing, jingling video terminals with names like Boomtown
Bonanza where, early on a weekday morning, people sit on stools pushing buttons and watching
cherries and 7s reel by.

Whether these people are all essentially doing the same thing -- playing bingo, that is -- is now
the subject of a fierce legal battle across Alabama, a fight with strange bedfellows and billions of
dollars at stake.

An alliance of district attorneys, preachers and the state government say that video bingo is
gambling and should be banned, to the strong protests of charities, small-town officials and
big-money gambling interests.

Gambling on games of chance is illegal under the Alabama Constitution, but paper bingo has been
around for decades, in church recreation rooms and American Legion halls.

Still, its legality was always questionable, and when authorities began to take notice in the early
1980s, bingo-friendly counties managed to have constitutional amendments passed protecting the game.

These allowed for bingo under a variety of conditions, like a limit on the hours of operation and a
requirement that charities be the sole beneficiaries.

Then, in the 1990s, the video machines showed up. They appeared first on Indian tribal land, where
the state's reach is questionable. Thousands were installed at two greyhound racing tracks, which
happen to be run by some of the state's most generous campaign donors.

But in the past few years, machines began popping up elsewhere.

More than three dozen electronic bingo halls sprouted up along a highway in Walker County; nearly 45
showed up in strip malls and old movie theaters in Jefferson County, which includes Birmingham; and
outside the little town of Ashville, American Legion Post 170 ran an electronic bingo hall for 10
lucrative days -- the city made \$13,000 -- before it was shut down.

By using a rotating cast of charities, churches and fire departments, the bingo halls say, they
comply with the mandated limits on hours of operation. Ten hours a week, for instance, a hall might
operate on behalf of the Tutwiler Fire Department, then another 10 hours for the Cordova Youth and
Sports Association, and so on, until there is bingo at the hall 24 hours a day.

Supporters say the bingo halls have produced a welcome source of money for struggling
municipalities. Walker County used \$800,000 in bingo money to build temporary classrooms for a
school that was hit by a tornado, and the debt-laden City of Fairfield was able to pay off a new
fire truck.

``It has been a very welcome development for us,'' said Vincent Smith, a Fairfield city councilman.

But with lax enforcement, no one knows what percentage of the money being made on bingo has gone to
charity and how much to the game suppliers and bingo hall owners and workers. The president of a
company that makes bingo machines testified in court that \$2 billion was wagered annually just in
Walker County.

``The charity was just the loophole that they needed to set up this operation,'' said Brandon Falls,
the district attorney in Birmingham, who ordered all video bingo halls in his jurisdiction shut down
in October. The suppliers of the machines alone are paid at least a fifth of the take, even though
every cent is legally obligated to go to charity, Mr.~Falls said.

Last year, Gov. Bob Riley, a Republican, set up a task force to crack down on electronic bingo,
insisting that the terminals were essentially illegal slot machines. District attorneys began
questioning the legality of video bingo, several halls were raided, and judges began shutting them
down.

The terminals certainly look like slot machines, but supporters say there are differences. First,
there has to be more than one player. The machines are connected to a single server, which acts like
a bingo caller; at the push of a button, the server calls a whole game of bingo at once, and the
machine marks the virtual card for the player. If the server's combination matches the pattern on a
player's card, the player wins.

The cherries and 7s are just eye candy. It all takes seconds, and the players do nothing more than
push a button. But, supporters say, it is still bingo.

Judges, in lengthy decisions that consider the history of bingo since the 16th century, have tended
to disagree.

``Most critically,'' wrote one judge, ``the game is not 'bingo' because it eliminates the
requirement that players manually daub their own cards or call out 'bingo!' when they have won.''

Though the issue seems like a simple divergence of opinion on issues moral and legal, it is crosscut
with competing interests.

Supporters of bingo, including Bill Johnson, a former member of Mr.~Riley's cabinet who is running
for the open governor's seat next year, have raised suspicions about Mississippi gambling interests,
whose history of opposing gambling in Alabama is well documented.

More than once, Mr.~Riley has been accused of benefiting from Mississippi gambling money --
including a comment reportedly made by the incarcerated lobbyist Jack. Abramoff that a casino-owning
Indian tribe spent \$13 million to help win his election -- and several members of Mr.~Riley's staff
have left to join lobbying firms that work with casino groups.

Mr.~Riley has repeatedly denied accusations that he has received gambling money or been influenced
by gambling lobbyists.

``Every question's been answered,'' said Todd Stacy, a spokesman for the governor. ``It is organized
gambling trying to throw a smokescreen.''

Nevertheless, Mississippi casinos are apparently enjoying the current scrap, running advertisements
in Alabama enticing people to ``come visit the real casinos.''

The most pointed complaints in Alabama come from small-time bingo operators, who ask why they have
been closed while the greyhound tracks, with high-rise hotels and rows upon rows of electronic bingo
machines, have somehow remained untouched.

``We have to shut down, and the rich folks get to continue on,'' said Mayor Robert McKay of
Ashville, who also helps run the American Legion hall here. ``It's not right the way they did us.''

The Ashville case is before the State Supreme Court, which, if it rules broadly enough, could shut
down all electronic bingo in the state, ending the discretion of local prosecutors on whether to
shut down the bingo halls. The governor has filed a brief in the case, noting that a definitive
ruling could go far beyond Ashville.

Bamaco Bingo, in Fairfield, sits just down the road from the Payday Loan Store, in a warehouse that
housed a call center before it relocated overseas. On a recent Monday morning, there was a sparse
crowd of older men and women sitting in the glow of the terminals, several of whom expressed
frustration over what they called state interference in their preferred method of entertainment.

``This is the Bible Belt,'' explained Roosevelt Harvey, 73, a retired government worker who comes to
Bamaco a few mornings a week with friends. ``But you cannot save my soul. You got to worry about
your own soul.''

Smoking a cigarette and standing before a machine called Bama Hot, Rodney Miller, 40, a
groundskeeper at a Birmingham cemetery, said he would be upset if the hall closed. Then,
acknowledging that he had dropped \$3,000 in the two months since Bamaco opened, he hedged.

``On the other hand, I may be happy,'' he said. ``I'm losing too much money.''

\section{Job Woes Exacting a Heavy Toll on Family Life}

\lettrine{P}{aul}\mycalendar{Nov.'09}{12} Bachmuth's 9-year-old daughter, Rebecca, began pulling out
strands of her hair over the summer. His older child, Hannah, 12, has become noticeably angrier,
more prone to throwing tantrums.

Initially, Mr.~Bachmuth, 45, did not think his children were terribly affected when he lost his job
nearly a year ago. But now he cannot ignore the mounting evidence.

``I'm starting to think it's all my fault,'' Mr.~Bachmuth said.

As the months have worn on, his job search travails have consumed the family, even though the
Bachmuths were outwardly holding up on unemployment benefits, their savings and the income from the
part-time job held by Mr.~Bachmuth's wife, Amanda. But beneath the surface, they have been a family
on the brink. They have watched their children struggle with behavioral issues and a stress-induced
disorder. He finally got a job offer last week, but not before the couple began seeing a therapist
to save their marriage.

For many families across the country, the greatest damage inflicted by this recession has not
necessarily been financial, but emotional and psychological. Children, especially, have become
hidden casualties, often absorbing more than their parents are fully aware of. Several academic
studies have linked parental job loss -- especially that of fathers -- to adverse impacts in areas
like school performance and self-esteem.

``I've heard a lot of people who are out of work say it's kind of been a blessing, that you have
more time to spend with your family,'' Mr.~Bachmuth said. ``I love my family and my family comes
first, and my family means more than anything to me, but it hasn't been that way for me.''

A recent study at the University of California, Davis, found that children in families where the
head of the household had lost a job were 15 percent more likely to repeat a grade. Ariel Kalil, a
University of Chicago professor of public policy, and Kathleen M.~Ziol-Guest, of the Institute for
Children and Poverty in New York, found in an earlier study that adolescent children of low-income
single mothers who endured unemployment had an increased chance of dropping out of school and showed
declines in emotional well-being.

In the long term, children whose parents were laid off have been found to have lower annual earnings
as adults than those whose parents remained employed, a phenomenon Peter R.~Orszag, director of the
White House Office of Management and Budget, mentioned in a speech last week at New York University.

A variety of studies have tied drops in family income to negative effects on children's development.
But Dr.~Kalil, a developmental psychologist and director of the university's Center for Human
Potential and Public Policy, said the more important factor, especially in middle-class households,
appeared to be changes in family dynamics from job loss.

``The extent that job losers are stressed and emotionally disengaged or withdrawn, this really
matters for kids,'' she said. ``The other thing that matters is parental conflict. That has been
shown repeatedly in psychological studies to be a bad family dynamic.''

Dr.~Kalil said her research indicated that the repercussions\cite{repercussion} were more pronounced
in children when fathers experience unemployment, rather than mothers.

She theorized that the reasons have to do with the importance of working to the male self-image, or
the extra time that unemployed female breadwinners seem to spend with their children, mitigating the
impact on them.

Certainly, some of the more than a dozen families interviewed that were dealing with long-term
unemployment said the period had been helpful in certain ways for their families.

Denise Stoll, 39, and her husband, Larry, 47, both lost their positions at a bank in San Antonio in
October 2008 when it changed hands. Mrs.~Stoll, a vice president who managed a technology group,
earned significantly more than her husband, who worked as a district loan origination manager.

Nevertheless, Mr.~Stoll took unemployment much harder than she did and struggled to keep his spirits
up, before he landed a new job within several months in the Kansas City area, where the family had
moved to be closer to relatives. He had to take a sizable pay cut but was grateful to be working
again.

Mrs.~Stoll is still looking but has also tried to make the most of the additional time with the
couple's 5-year-old triplets, seeking to instill new lessons on the importance of thrift.

``Being a corporate mom, you work a lot of hours, you feed them dinner -- maybe,'' she said. ``This
morning, we baked cookies together. I have time to help them with homework. I'm attending church.
The house is managed by me. Just a lot more homemaker-type stuff, which I think is more nurturing to
them.''

Other families, however, reported unmistakable ill effects.

Robert Syck, 42, of Fishers, Ind., lost his job as a call-center manager in March. He has been
around his 11-year-old stepson, Kody, more than ever before. Lately, however, their relationship has
become increasingly strained, Mr.~Syck said, with even little incidents setting off blowups. His
stepson's grades have slipped and the boy has been talking back to his parents more.

``It's only been particularly in the last few months that it's gotten really bad, to where we're
verbally chewing each other out,'' said Mr.~Syck, who admitted he had been more irritable around the
house. ``A lot of that is due to the pressures of unemployment.''

When Mr.~Bachmuth was first laid off in December from his \$120,000 job at an energy consulting
firm, he could not even bring himself to tell his family. For several days, he got dressed in the
morning and left the house as usual at 6 a.m., but spent the day in coffee shops, the library or
just walking around.

Mr.~Bachmuth had started the job, working on finance and business development for electric
utilities, eight months earlier, moving his family from Austin. They bought something of a dream
home, complete with a backyard pool and spa.

Although she knew the economy was ultimately to blame, Mrs.~Bachmuth could not help feeling angry at
her husband, both said later in interviews.

``She kind of had something in the back of her mind that it was partly my fault I was laid off,''
Mr.~Bachmuth said. ``Maybe you're not a good enough worker.''

Counseling improved matters significantly, but Mrs.~Bachmuth still occasionally dissolved into tears
at home.

Besides quarrels over money, the reversal in the couple's roles also produced friction.
Mrs.~Bachmuth took on a part-time job at a preschool to earn extra money. But she still did most, if
not all, of the cooking, cleaning and laundry.

Dr.~Kalil, of the University of Chicago, said a recent study of how people spend their time showed
unemployed fathers devote significantly less time to household chores than even mothers who are
employed full-time, and do not work as hard in caring for children.

Mr.~Bachmuth's time with his girls, however, did increase. He was the one dropping off Rebecca at
school and usually the one who picked her up. He began helping her more with homework. He and Hannah
played soccer and chatted more.

But the additional time brought more opportunities for squabbling. The rest of the family had to get
used to Mr.~Bachmuth being around, sometimes focused on his search for a job, but other times
lounging around depressed, watching television or surfing soccer sites on the Internet.

``My dad's around a lot more, so it's a little strange because he gets frustrated he's not at work,
and he's not being challenged,'' Hannah said. ``So I think me and my dad are a lot closer now
because we can spend a lot more time together, but we fight a lot more maybe because he's around
24-7.''

When Rebecca began pulling her hair out in late summer in what was diagnosed as a stress-induced
disorder, she insisted it was because she was bored. But her parents and her therapist -- the same
one seeing her parents -- believed it was clearly related to the job situation.

The hair pulling has since stopped, but she continues to fidget with her brown locks.

The other day, she suddenly asked her mother whether she thought she would be able to find a ``good
job'' when she grew up.

Hannah said her father's unemployment had made it harder for her to focus on schoolwork. She also
conceded she had been more easily annoyed with her parents and her sister.

At night, she said, she has taken to stowing her worries away in an imaginary box.

``I take all the stress and bad things that happen over the day, and I lock them in a box,'' she
said.

Then, she tries to sleep.

\section{Bush Emerging for Speech to Kick Off Public Policy Institute}

\lettrine{N}{early}\mycalendar{Nov.'09}{12} 10 months after leaving office, former President George
W.~Bush plans to emerge from self-imposed political hibernation on Thursday as he starts a new
public policy institute to promote some of the domestic and international priorities of his
presidency.

In a speech at Southern Methodist University, home of his future library and museum, the former
president will kick off the new George W.~Bush Institute as a forum for study and advocacy in four
main areas: education, global health, human freedom and economic growth. Advisers said he hoped his
institute would be more focused on producing results than many research organizations are.

Mr.~Bush will announce the appointment of the first five of two dozen scholars to be affiliated with
the institute, which has already scheduled a half-dozen conferences for next year, according to
organizers. The former first lady, Laura Bush, will also speak at Thursday's event to discuss how
women's issues will be injected into all the institute's program areas, including sponsorship of a
conference on the education of women in Afghanistan.

``The president has been working with these ideas for a long time now,'' said James K.~Glassman, a
former top State Department official now serving as the institute's founding executive director.
``He wanted to do something very different from other former presidents, and that is to create a
research institute that's independent, nonpartisan and scholarly and that will have an impact on the
real world.''

Although Mr.~Bush has given several speeches since leaving office, most have been out of the
country, closed to the news media and reminiscent in tone. This is the first event where he has
invited reporters to announce a new venture. Organizers expect as many as 1,500 students, donors and
others to attend.

The former president's approach contrasts with that of his vice president, Dick Cheney, who has
become a vigorous public critic of President Obama's national security policies. The institute will
be a vehicle for Mr.~Bush to re-enter the national conversation and advocate ideas in a less
politicized way, advisers said. The goal, Mr.~Glassman said, is to ``extend principles and work that
were accomplished during the administration.''

The institute will be housed along with the presidential library and museum in a building on the
S.M.U. campus to be completed by 2013. The Bush foundation has secured \$212 million in pledges and
contributions toward its goal of \$300 million by next year's groundbreaking, according to two
people who were privy to the results but not authorized to discuss them on the record.

Mr.~Bush is also completing a book describing some of the most important decisions in his life. He
is about five-sixths through with the manuscript, which he is drafting with Chris Michel, a former
White House speechwriter, for publication next fall, according to his office.

Former aides said Mr.~Bush remained attentive to news developments, even if publicly quiet. ``I get
a lot of e-mails from him now,'' said Michael J.~Gerson, his former senior adviser. ``He responds to
the news. He's very engaged.''

But the Bushes keep a deliberately low profile. After a shooting rampage left 13 dead at Fort Hood,
about 150 miles south of Dallas, the former president and first lady paid a visit in secret to
wounded soldiers and family members.

\section{Hewlett-Packard to Acquire 3Com}

\lettrine{H}{ewlett}\mycalendar{Nov.'09}{12}-Packard said on Wednesday that it had reached an
agreement to acquire 3Com, a provider of computer network equipment, for \$2.7 billion in a deal
that H.P. plans as a springboard for an assault on the market leader in networking, Cisco Systems.

In an interview, Ann M.~Livermore, an executive vice president of H.P., described computer
networking as a \$40 billion-a-year market with high profit margins that is growing briskly and
dominated by Cisco, which has so far had little head-to-head competition.

``H.P. is eager and now positioned to disrupt the networking industry,'' Ms.~Livermore said.

The offer price, \$7.90 a share, is at a premium of about 39 percent to 3Com's closing price on
Wednesday. The deal was announced after the market closed.

Under Mark V.~Hurd, its chief executive, Hewlett-Packard has been beefing up its network equipment
offerings to compete more aggressively against Cisco to sell data center equipment to corporations.

But so far, H.P. has mostly supplied smaller equipment used in office networks, both wired and
wireless.

``What we've been missing is networking for the core of the data center,'' Ms.~Livermore said.
``That's where Cisco is strong, and before, H.P. couldn't attack that.''

This year, Cisco went after one of H.P.'s core businesses by entering the market for server
computers used in data centers. Cisco also recently teamed with EMC, another technology giant, to
sell data center equipment to businesses.

H.P. also provided investors with an early look at its earnings for its fiscal fourth quarter, which
ended in October, and raised its estimates for the next year. The company reported a profit,
excluding some items, of \$1.14 a share in the fourth quarter, which slightly beat expectations of a
\$1.12 a share, as compiled by Thomson Reuters.

Revenue in the quarter was \$30.8 billion, down 8 percent from a year ago. But the company's sales
performance beat analysts' forecasts by \$1 billion.

H.P. joins a series of technology companies that have reported better-than-expected results
recently, including Intel and I.B.M.

``Solid execution drove exceptional performance for H.P. this quarter, fueled by significant growth
in China,'' Mr.~Hurd said in a statement.

H.P. also raised its forecast for its 2010 fiscal year. The company said it expected revenue for
fiscal 2010 to be \$118 billion to \$119 billion, up somewhat from its previous estimate of \$117
billion to \$118 billion. It forecast earnings per share, excluding some items, to be \$4.25 to
\$4.35, up from the earlier forecast of \$4.20 to \$4.30.

In 3Com, H.P. is acquiring a company with a rich heritage in network technology and a solid product
portfolio, analysts say. But it lacks Cisco's size and credibility in the data center market, where
large corporate customers look for strong suppliers who can provide a full range of products and
services.

H.P. may well be able to fill those gaps, with its strong services, storage and server computer
businesses. ``3Com, with H.P.'s backing, is capable of making a real run at Cisco,'' said Rob
Enderle, an independent technology analyst.

But analysts said that while 3Com would give H.P. crucial technology to broaden its reach in
networking, further product development and marketing investment would be needed to compete directly
with Cisco in large data-center accounts. ``It will be a two-year process to roll this out,'' said
Jeffrey Evenson, an analyst at Bernstein Research. ``This is a longer-term play.''

While trailing well behind Cisco in most markets, 3Com is strong in China. Half of its \$1.3 billion
in sales comes from the Chinese market. Three hundred of the largest 500 companies in China, and 70
percent of government agencies, use 3Com equipment, Ms.~Livermore said.

Last year, a deal to sell 3Com for \$2.2 billion to an alliance of Bain Capital, a private equity
firm, and Huawei Technologies, a Chinese maker of networking equipment, fell apart after a federal
review panel, the Committee on Foreign Investment in the United States, expressed national security
concerns about the transaction.

3Com shares were up more than 30 percent, at \$7.67, in after-hours trading on Wednesday. H.P.
shares were down slightly. Trading in 3Com call options -- contracts to buy 3Com shares at a fixed
price -- surged Wednesday to their highest level since September 2007, suggesting that word of the
deal may have leaked before the announcement.

\section{Among Late-Night Writers, Few Women in the Room}

\lettrine{I}{n}\mycalendar{Nov.'09}{12} many ways, television today is about women more than men.

More women watch television than men; female producers and writers have had huge success in prime
time and daytime; in January, women will occupy two of the three seats as anchors of network evening
newscasts.

But there is one glaring exception: very few women make it inside the writing rooms for late-night
television hosts, despite that women make up a larger proportion of their audience than men.

There are no female writers on the new ``The Jay Leno Show,'' none on ``Late Show with David
Letterman,'' none on ``The Tonight Show with Conan O'Brien.''

The lack of women in late-night writing positions was pushed again to the forefront by David
Letterman's confession of sexual relationships with his staff members.

In an article for Vanity Fair's Web site, a former writer for Mr.~Letterman, Nell Scovell, described
her experiences in the late 1980s, citing what she called a ``hostile work environment,'' including
Mr.~Letterman's romantic relationships with women on the staff.

``Writing for late-night talk shows is a great entry-level TV job, and if you deny women that
opportunity it reduces the chance for them to pursue careers in comedy,'' she said in an e-mail
message. Her article was cited in motions filed with the court this week by the lawyer defending the
man accused of blackmailing Mr.~Letterman.

In the 1980s, Mr.~Letterman pioneered the kind of college-age male humor that dominates late night.
But now, his audience is almost 55 percent women; Mr.~Leno's is more than 53 percent, and
Mr.~O'Brien's just over one half. Yet the writing room and sensibilities of the show itself remain
largely male.

Steve Bodow, head writer for ``The Daily Show,'' conceded that the preponderance of male writers had
not changed much in recent years, although his show hired two women writers in September.

``We shook the trees a little,'' Mr.~Bodow said of the special efforts the show made. ``Women have a
different perspective, which we like on the show.''

But perspective is not the same as sensibility. Some women in the business argue that, as long as
the hosts remain almost exclusively male, so will the writers.

``When you're writing for late night, you're writing through one person's prism, and that person at
the shows you're looking at is always a dude,'' said Hallie Haglund, one of the new writers on ``The
Daily Show. `` Allison Silverman, who has served as the only woman writer first on ``The Daily
Show'' then on Conan O'Brien's staff, said she had no trouble assuming the voices of male hosts. But
she added, ``I don't think the issue of sensibility is off base. The hosts and the staffs I worked
on often resembled one another. Have you seen how many tall Irish people are on Conan's staff?''

In what seems like a paradox, ``The Daily Show'' was in fact created by two women, Lizz Winstead and
Madeleine Smithberg, in 1996 for Comedy Central. And one of the most influential writers in the
history of late night is Merrill Markoe , Mr.~Letterman's previous longtime girlfriend who was also
his full partner in the invention of his breakthrough late-night show in the 1980s.

Those women acknowledged that recruiting and hiring female writers was a daunting challenge for them
at the time.

``It's the law of averages,'' Ms.~Winstead said in a telephone interview. ``More guys than women are
in comedy.'' When she and Ms.~Smithberg sought submissions for writers on ``The Daily Show,'' more
than 100 came in and ``only about three or four were from women.''

Ms.~Markoe recalled the pressure she was under to hire writers for Mr.~Letterman. ``I didn't have
any leeway to put bullets in the gun that were going to misfire,'' she said in an e-mail interview.
She, too, had hardly any submissions from women.

``Bad odds do not help the situation,'' she wrote, adding: ``Back then it simply never occurred to
me that there would be a pattern of all guys because I wasn't a guy. Instead I was focused on: Let's
get this thing running so Dave will relax a little.'' That, she said, ``never happened of course.''

Ms.~Markoe said she believed what she called ``an odd shift toward more boys' humor'' in the '90s
might have kept some women from landing late-night jobs. ``The massive popularity of Howard Stern
might have had something to do with that,'' she said.

Whatever the reason, late-night seemed to take on the flavor -- some say aroma -- of a boys' club.

``I would walk into Lizz's office, where the writers were assembled to hear the day's jokes, and
would want to exercise my executive producer privileges by sending half of them home to shower,''
Ms.~Smithberg said in an e-mail message. ``I wonder if the corollary we should be examining is
between body odor and humor rather than gender and humor.''

Writers' rooms have always had a reputation for crude behavior. Ms.~Silverman described both ``The
Daily Show'' and Mr.~O'Brien's show as ``really good experiences'' -- though rife with crudeness.

``There were plenty of penis jokes,'' she said. ``On occasion I've told them myself, though my penis
jokes were rarely good. I don't quite have the feel.''

Ms.~Winstead scoffed at the idea that a strong woman writer would be offended by the writers' room.
``I have no sensitivity to off-color humor,'' she said. ``I only have a sensitivity to bad humor.''

Those who dispute a deliberate tilt in late-night television point out that women occupy positions
of real power. Mr.~Kimmel and Mr.~Leno both have female executive producers -- Jill Leiderman in
Mr.~Kimmel's case, and Debbie Vickers in Mr.~Leno's -- and Mr.~Letterman's show has three, Jude
Brennan, Barbara Gaines and Maria Pope.

Ms.~Silverman was co-head writer and executive producer of ``The Colbert Report'' until stepping
down in August.

But Ms.~Smithberg called the women-as-producers argument an inadequate defense.

``Vis-à-vis women in producing rather than writing roles,'' Ms.~Smithberg said, ``there is an
inherent dynamic not dissimilar to a marriage: male writers and late-night hosts tend to need the
care of female nurturers, whom they routinely undervalue.''

One woman who does have a late-night show, Chelsea Handler on the E Channel, has five women writers
on her staff of 10.

Craig Ferguson of CBS found an ideal female match for his comic sensibility: his sister. Jimmy
Fallon on NBC hired three women on his initial staff of about a dozen. ``The Colbert Report''
currently has one.

There also is only one woman currently writing for Jimmy Kimmel -- Molly McNearney, the co-head
writer. Ms.~Markoe, who is regarded as one of the best writers to ever work in late-night
television, said she never experienced any resentment toward her on the Letterman show.

``In my weird crossover capacity I may have added a level of communication between the host and
writers that hasn't taken place since,'' Ms.~Markoe said. ``I can remember a few times following him
into the bathroom, postrehearsal, preshow, while he showered and got ready. I stood there, reading
him lists of jokes or rewritten ideas for things to do on the show.''

She added, reflecting on the recent headlines: ``I guess over the years there has been an assortment
of other women who continue to have that sort of access to him. But from what I read, it doesn't
appear that their duties include reading him joke rewrites.''

\section{Motorola Said to Explore Dividing Into 3 Companies}

\lettrine{M}{otorola}\mycalendar{Nov.'09}{12}, which has said it wants to split into two separate
companies, is exploring a three-way split instead in order to raise cash and pay down debt.

The pioneering maker of radios and cellphones has hired JPMorgan Chase, Centerview Partners and
Goldman Sachs to seek buyers for its division that manufactures set-top boxes for cable television
companies and radios to go into cellphone transmission towers, according to people briefed on the
matter. These people, who requested anonymity because they were not authorized to discuss the
process, cautioned that it was still in its early stages and may not lead to a deal.

The businesses up for sale could be worth as much as \$5 billion, said Philip Cusick, an analyst for
Macquarie Capital.

Last year, Motorola, based in Schaumberg, Ill., said it wanted to separate its struggling cellphone
handset unit from the rest of the company, which includes a third business unit that sells two-way
radios to businesses and government agencies.

Even if it sells the set-top box and communications equipment unit, the company would still likely
separate the two remaining divisions, Mr.~Cusick said. Motorola is not expected to begin the
separation process, however, until the finances of the handset unit stabilize, perhaps in mid-2010.
A final split-up is not expected until 2011 at the earliest, the analyst said.

Jennifer Weyrauch-Erickson, a Motorola spokeswoman, declined to comment on the potential
transaction, which was first reported by The Wall Street Journal's Web site. She did confirm that
Motorola still wants to split into two independent publicly traded companies.

Maynard Um, an analyst with UBS Securities, said in a note to clients that the cash from selling the
set-top box unit could help Motorola pay down its \$3.9 billion in long-term debt and put its
handset business on a firmer footing.

After a long slide, Motorola has high hopes for turning around its cellphone business. It recently
introduced two smartphones, including the Droid, which went on sale Friday and is being heavily
promoted by Verizon Wireless. Motorola promises dozens more smartphones in the next year.

Motorola got into the cable box business in 2000, when it acquired General Instruments for \$17
billion. Prospects for that business may be dimming, however, in part from more aggressive
competition from its main rival, Cisco Systems, and others. Moreover, the rise of video delivered
over the Internet may undercut the traditional cable TV model.

Motorola's telecom equipment business has also had its share of troubles. It primarily supplies
radios to wireless carriers that use an older technology.

The unit up for sale, which Motorola calls home and network mobility, posted an operating profit of
\$199 million in the quarter ended Oct.~3, down 24 percent from the year earlier. The unit's revenue
fell 15 percent, to \$2 billion.

Mr.~Cusick said it is not clear that there are other companies that would want to own both the
cellular and cable product lines. Other telecommunications equipment providers, like Ericsson, might
be interested in the wireless products. And there are some smaller vendors to the cable industry,
like Arris Group and Pace, that might try to put together a deal to buy the set-top box business, he
said.

Motorola's bankers are also approaching private equity firms, including TPG and Silver Lake
Partners, the people briefed on the deal said.

\section{How to Market Your Business With Facebook}

\lettrine{B}{usiness}\mycalendar{Nov.'09}{12} owner, you might want to friend Facebook.

A growing number of businesses are making Facebook an indispensible part of hanging out their
shingles. Small businesses are using it to find new customers, build online communities of fans and
dig into gold mines of demographic information.

``You need to be where your customers are and your prospective customers are,'' said Clara Shih,
author of ``The Facebook Era'' (Pearson Education, 2009). ``And with 300 million people on Facebook,
and still growing, that's increasingly where your audience is for a lot of products and services.''

Start Small

For most businesses, Facebook Pages (distinct from individual profiles and Facebook groups) are the
best place to start. Pages allow businesses to collect ``fans'' the way celebrities, sports teams,
musicians and politicians do. There are now 1.4 million Facebook Pages and they collect more than 10
million fans every day, according to the site.

Businesses can easily create a Web presence with Facebook, even if they don't have their own Web
site (most companies still should maintain a Web site to reach people who don't use Facebook or
whose employers block access to the site). Businesses can claim a vanity address so that their
Facebook address reflects the business name, like www.facebook.com/Starbucks. Facebook pages can
link to the company's Web site or direct sales to e-commerce sites like Ticketmaster or Amazon.

Facebook offers an array of tools and networks, and it's easy to wander down too many paths.
Ms.~Shih recommends that newcomers start by asking themselves a simple question: What is your basic
objective? Is it getting more customers in the door? Building brand awareness? Creating a venue for
customer support? Once you have set your goal, you can strategize accordingly.

``You can waste a lot of time on Facebook,'' said Ms.~Shih, founder of Hearsay Labs, a Facebook
marketing software company. ``But if you're a business, you don't have any time to waste. Figure out
your objectives first, start small and do things that help you accomplish your objectives.''

Ms.~Shih suggests that businesses ask friends and family to become fans of their pages so that they
display a respectable crowd of supporters when they debut. Pages can grow organically by word of
mouth -- the average Facebook user has 130 friends on the site -- or by advertising or promotion.

You can enliven your page with photos, comments and useful information. As you grow more
comfortable, you can add videos or business applications. Flaunt your personality. The page of an
ice cream parlor should feel different than that of a funeral parlor. ``The pages that are most
successful,'' said Tim Kendall, the director of monetization at Facebook, ``are the ones that really
replicate the personality of the business.''

It's Not All About Selling

Art Meets Commerce, a New York marketing firm, has struck up a never-ending conversation with fans.
The company uses Facebook as a crucial part of its publicity campaigns for theatrical productions.
Its Facebook page for the show ``Rock of Ages,'' for example, has more than 13,000 fans.

Staff members constantly update the page with new photos, videos and quotes from the cast. They've
also learned what not to do: Once they posted a video of Paris Hilton plugging the show and got
negative feedback from fans who professed to be sick of her.

But it's not just about marketing -- or, at least, it's not just about selling. ``You end up moving
away from being an Internet marketer and go into almost customer service,'' said Jim Glaub, creative
director at the agency. ``A lot of times people use Facebook to ask questions: What's the student
rush? How long is the show? Where's parking? You have to answer.''

Some basic rules: Buy-buy-buy messages won't fly. The best practitioners make Facebook less about
selling and more about interacting. Engage with fans and critics. Listen to what people are saying,
good and bad. You may even pick up ideas for how to improve your business. Keep content fresh. Use
status updates and newsfeeds to tell fans about specials, events, contests or anything of interest.

These interactions can take a vast amount of time -- the ``Rock of Ages'' page has 300 to 600
interactions every week -- but they can also provide a big payoff. Facebook is one of the show's top
sources of new ticket sales.

Last year, Art Meets Commerce introduced a Facebook ad campaign to promote an Off Broadway run of
the musical ``Fela!'' The campaign aimed at Facebook users with interests like theatrical shows or
Afro beat. According to the company, it generated 18 million impressions, more than 5,700 clicks and
\$40,000 in ticket sales -- all for \$4,400 spent on advertising.

``We can advertise all day, but if we don't give them what they want they will not be a fan
anymore,'' said Mark Seeley, a marketing associate at Art Meets Commerce. ``Even though we represent
the shows as marketers, we don't want to constantly tell people to buy tickets. You talk to them
like you talk to your friends on Facebook.''

Aim at Potential Customers Only

Some guys use Facebook to find single women. Chris Meyer used it to find women who are already
engaged.

Mr.~Meyer, a wedding photographer in Woodbury, Minn., had had little luck with traditional
advertising. A full-page ad in a bridal magazine generated zero leads and a trade show yielded only
four bookings, barely covering the cost of his booth. But Facebook proved a digital bonanza.

Mr.~Meyer aimed at women ages 22 to 28 who listed their martial status as engaged in the
Minneapolis-St.~Paul metropolitan area. He estimates that he has spent about \$300 on Facebook ads
in the last two years and has generated more than \$60,000 in business. He says about three-quarters
of his clients now come to him through Facebook, either from ads or recommendations from friends.

``I'd be out of business if I didn't have Facebook,'' Mr.~Meyer said. ``Especially with this
economy, I need to stretch each marketing dollar as much as I possibly can.''

Facebook enables small businesses to engage in targeted marketing that they only could have dreamed
about a few years ago. Facebook users fill out profiles with information like hometown, employer,
religious beliefs, interests, education and favorite books, movies and TV shows -- all of which can
help advertisers deliver messages to specific demographic slices.

As you create an ad, you can add demographic criteria and keywords and see how many Facebook users
fall into your target audience and modify it accordingly to get the most bang for your buck.
Advertisers can elect to pay per impression or per click, set maximum budgets and schedule the ad to
run on specific dates.

Thus a coffee shop in San Francisco can display advertisements only to local people whose profiles
or group affiliations suggest they like coffee. According to Mr.~Kendall, Facebook's director of
monetization, ads can also aim at people based on social exchanges, like a person who sends a
message to a friend, ``let's get together for coffee'' or who posts a status update about just
having awakened and needing some java.

``We can help you find customers before they even think about searching for you,'' Mr.~Kendall said.
``We're very, very well-positioned to generate demand, based on the fact that we know a tremendous
amount about a user.''

The Facebook ad system provides instant feedback with metrics like the number of impressions and
clicks-through. This reporting allows Mr.~Meyer to improve his advertising; if one ad doesn't
generate enough hits within 24 hours, he pulls it and tries something new.

Give Away Cupcakes!

Charles Nelson has an M.B.A. and is a former investment banker who owns a growing national chain of
stores. Yet this 40-year-old entrepreneur checks Facebook with the frequency of a college student.
Up to 30 times a day, he logs onto the social networking site via his laptop or Blackberry.

For Mr.~Nelson, this is serious business. He and his wife, Candace, own Sprinkles, a cupcake bakery
that relies on social media in lieu of traditional advertising. Mr.~Nelson considers Facebook
marketing essential. ``People are out there talking about your business everyday, whether you're
looking or not,'' he said. ``This gives people a place to come and speak directly to us.''

Sprinkles uses Facebook to give customers a whiff of what's cooking. Every day it posts a password
on Facebook that can be redeemed for a free cupcake. Since April, its fan base has risen tenfold to
70,000.

Mr.~Nelson and his wife previously worked as investment bankers in the technology sector and were
keenly aware that, even for a traditional business like a bakery, social media is a crucial
ingredient. His advice: make it relevant to the customer, keep it fresh and remember that the return
on investment may come slowly.

``Be patient with it,'' Mr.~Nelson advised. ``People are not going to flock to your social media
site overnight. Technology is about the network effect. It takes time for those connections to
build.''

\section{I'm Innocent. Just Check My Status on Facebook.}

\lettrine{T}{he}\mycalendar{Nov.'09}{12} message on Rodney Bradford's Facebook page, posted at 11:49
a.m. on Oct.~17, asked where his pancakes were. The words were typed from a computer in his father's
apartment in Harlem.

At the time, the sentence, written in street slang, was just another navel-gazing, cryptic Facebook
status update -- meaningless to anyone besides Mr.~Bradford. But when Mr.~Bradford, 19, was arrested
the next day as a suspect in a robbery at the Farragut Houses in Brooklyn, where he lives, the words
took on greater importance. They became his alibi.

His defense lawyer, Robert Reuland, told a Brooklyn assistant district attorney, Lindsay Gerdes,
about the Facebook entry, which was made at the time of the robbery. The district attorney
subpoenaed Facebook to verify that the words had been typed from a computer at an apartment at 71
West 118th Street in Manhattan, the home of Mr.~Bradford's father. When that was confirmed, the
charges were dropped. ``This is the first case that I'm aware of in which a Facebook update has been
used as alibi evidence,'' said John G.~Browning, a lawyer in Dallas who studies social networking
and the law. ``We are going to see more of that because of how prevalent social networking has
become.''

With more people revealing the details of their lives online, sites like Facebook, MySpace and
Twitter are providing evidence in legal battles.

Up to now, social networking activity has mostly been used as prosecutorial evidence, Mr.~Browning
said. He cited a burglary case in September in Martinsburg, Pa., in which the burglar used the
victim's computer to log on to Facebook and forgot to log off. The police followed the digital trail
to Jonathan G.~Parker, 19, who was arrested.

As part of his defense, a suspect in an Indiana murder case, Ian J.~Clark, claimed he was not the
kind of man who could kill his girlfriend's child. But remarks he was found to have posted on
MySpace left him vulnerable to character examination, Mr.~Browning said, contributing to his
conviction and a sentence of life in prison without parole.

In civil cases, too, online communications have helped strengthen evidence, especially in divorce
cases, where they are often used as proof of cheating.

And postings by a probationary sheriff's deputy, Brian Quinn, 26, of Marion County, Fla., on his
MySpace page led to his firing in June 2006 for ``conduct unbecoming an officer.''

Such cases are becoming more prevalent in part because Congress in 2006 mandated changes to the
federal rules of civil procedure, expanding the acceptance of electronically stored information as
evidence.

With the use of a Facebook update as an alibi, such communications may also be used to prove
innocence, Mr.~Browning said.

Mr.~Bradford's arrest was for the mugging at gunpoint of Jeremy Dunklebarger and Rolando
Perez-Lorenzo at 11:50 a.m. on Oct.~17, according to Mr.~Reuland, Mr.~Bradford's lawyer.

Mr.~Bradford, who was facing charges in a previous robbery, contended he was in Harlem at the time
of the Oct.~17 robbery -- a claim supported by Mr.~Bradford's father, Rodney Bradford Sr., and his
stepmother, Ernestine Bradford, Mr.~Reuland said.

Mr.~Reuland acknowledged that, in principle, anyone who knew Mr.~Bradford's user name and password
could have typed the Facebook update, but he regards it as unlikely.

``This implies a level of criminal genius that you would not expect from a young boy like this; he
is not Dr.~Evil,'' Mr.~Reuland said, adding that the Facebook entry was just ``icing on the cake,''
since his client had other witnesses who provided an alibi.

Jonah Bruno, a spokesman for the Brooklyn district attorney, Charles J.~Hynes, said he could not
discuss details of the case because it was sealed. But he acknowledged that Facebook was crucial to
the charges' being dropped.

But Joseph A.~Pollini, who teaches at the John Jay College of Criminal Justice, said prosecutors
should not have been so quick to drop the charges.

``With a user name and password, anyone can input data in a Facebook page,'' Mr.~Pollini said.

``Some of the brightest people on the Internet are teenagers,'' he said. ``They know the Internet
better than a lot of people. Why? Because they use it all the time.''

\section{A Holiday Gift of Some Free Wi-Fi}

\lettrine{G}{oogle}\mycalendar{Nov.'09}{12}, Yahoo, eBay and Microsoft all have the same marketing
idea this holiday season: temporarily providing free Wi-Fi access in airports, airplanes and public
places.

Google announced Tuesday that it would provide free Wi-Fi access in 47 airports across the country
-- including Boston, Houston and Seattle -- through Jan.~15. The airports handle about 35 percent of
American air travelers, the company said.

Travelers who connect to a wireless hot spot at one of the airports will see a browser page that
gives them the chance to donate to three charities and have the donation matched by Google.

The 47 airports include some, such as McCarran International in Las Vegas, that already provide free
Wi-Fi. Sponsorships help the airport keep the service free.

Google is also providing free Wi-Fi on Virgin America flights for the same period.

EBay, the Web auction powerhouse, will provide free Wi-Fi on some Delta Air Lines planes during the
week of Thanksgiving. Logging on will take users first to eBay's holiday page, but they will be free
to roam from there.

Since September, Microsoft has given away Wi-Fi access at some hotels and airports, encouraging
users to make a query on the company's new search engine, Bing.

On Monday, Yahoo said it was giving away Wi-Fi access in Times Square in New York.

Boingo Wireless, which operates for-pay hot spots in airports and is part of Google's project, said
it had been getting good results with free Wi-Fi campaigns sponsored by hotel chains. The campaigns
typically offer users 15 to 20 minutes of access after they watch a 30-second video advertisement.

\section{Train the Mind to Run Right Through Winter}

\lettrine{I}{n}\mycalendar{Nov.'09}{12} late summer, Sharon Henderson, the manager of the Lululemon
athletic clothing store in my town, started organizing Saturday morning group runs. People had two
options: three miles at a slower pace or six miles at a faster one.

There was a good turnout at first -- more than two dozen people, most of them slower runners, showed
up.

Then they stopped coming. Was it the string of gray, rainy Saturdays?

Granted, it is difficult to get up and be at Lululemon by 8:30 a.m. when the temperature is dropping
and a steady rain is falling. But, still. One recent Saturday, it was just me and my friend Claire
Brown running on the slick streets in the rain.

Very few studies have asked whether people exercise less in inclement weather and, if so, which ones
are more likely to slack off or forge ahead. Maybe that's because the results of the studies are not
exactly surprising.

``Why do people work out more in San Diego than in Michigan?'' asked James Pivarnik, an exercise
physiologist at the Michigan State University. ``Gee, I can't imagine.''

HIS study of Michigan residents found that people expended 15 to 20 percent more calories a week
exercising in the spring and summer than they did in the fall and winter.

Something similar seems to happen in Columbus, Ohio, said Janet Buckworth, an exercise physiologist
at Ohio State University.

She found that college students lost cardiovascular fitness in winter but maintained their strength,
indicating that while some of them did not want to go outside and run, at least they may have been
going to the gym.

``Columbus is incredibly dreary in the winter,'' Dr.~Buckworth said. ``It is wet and cold, and we
get snow.''

So maybe the question is not, ``Why do people stay home in dreary weather?'' as much as, ``Why do
some go out and exercise anyway?''

Dr.~Buckworth said that, in her experience, it was the people who were new to exercise who gave up
in bad weather.

``If you are beyond the point that you are learning how to exercise, you can't imagine not running
in bad weather,'' she said. Her advice to people who want to keep exercising all year: find
something you can do indoors, plan to exercise with a friend or do something -- like update your
playlist -- that can make your workout more fun.

Dr.~Pivarnik tells people they need to make up their minds that they will have a regular exercise
routine, no matter what. ``If you are one of those people who are going to back off, you are just
going to have to find something to make you do it,'' Dr.~Pivarnik said. ``It has to be a behavioral
thing in your head. It's not going to happen just because the weather is nice, you have to think
about it.''

My friend Jen Davis, a physical chemist, uses a term from chemistry: Running on dreary days requires
high activation energy, she says. In chemistry, activation energy is what must be added to start a
reaction.

But those of us who exercise in all sorts of weather will attest that there is a certain thrill that
can come from terrible conditions. ``It makes us tough,'' Jen said. She calls our runs in horrendous
conditions ``epic runs.'' And she's right. They are truly memorable, ones we actually recall fondly.

There also are epic bike rides, as Richard Armington will attest.Rich, a software engineer in
Montgomery, N.J., rode 200 miles over two days in a cold rain recently. It was a fund-raising trip
for Battle Against Hunger, and his group had been training all summer.

Last year, the group rode in a hurricane, but that proved too much -- the bikers had to stop at
lunchtime on the second day, three quarters of the way through the trip.

``Why do I do this?'' he said. ``For me, it's two challenges: the athletic challenge and the
challenge of getting others to sponsor and give to the cause.''

Glenn Swan, a cyclist in Ithaca, N.Y., says his area has some of the worst weather in the country,
but he does not let a little rain or snow stop him. Mr.~Swan, a research technician at Cornell and
owner of a bike shop called Swan's Cycles, said, ``Our phrase is, 'We ride even if the sun shines.'
''

His epic ride took place with friends in Virginia. They started at the bottom of a mountain on a
sunny morning. Soon it started to drizzle. ``We said, 'At least it's not raining,' '' Mr.~Swan said.
Then, as they ascended, it started to rain.

``We said, 'At least it's not snowing.' '' Then it started to snow.

``We said, 'At least the snow is not sticking.' '' Then it started to stick.

By the time they got to the top of the mountain, they were in a blizzard. They eventually made it to
a lodge, 20 miles away, where they spent the night. And they have been talking about the trip every
since.

But the problem with epic runs or rides is that each one ups the ante. A day with just ordinary bad
weather simply is not memorable after a while.

Jen and I noticed that recently on a dark, rainy, windy night. We had planned to run after work but
-- just this once -- we thought that maybe we could do one of those mind-numbingly dull treadmill
runs in the gym.

I called my coach, Tom Fleming, and told him our plans. He hates treadmills, thinking that that if
you want to train for road races, you have to run on roads. Treadmills, he says, are ``propelling
you over the running surface.'' When you run, he adds, ``you propel yourself over the surface,''
which can include hills, flat areas, and places where the surface is uneven. ``That's a harder
effort for sure,'' Tom said.

So, Tom told me: Don't go to the gym. Run outside.

So we did, and it was fine. Fun, actually.

But epic? No.~We have had much tougher runs than that.

\section{In China, Obama to Press for Tough Stance on Iran}

\lettrine{P}{resident}\mycalendar{Nov.'09}{16} Obama, fresh from making progress in his efforts to
get Russia on board for possible tough new sanctions against Iran, arrived in China on Sunday, where
he will attempt the even more difficult task of prodding China's leaders to get tough on Iran.

Making his first trip to China, Mr.~Obama landed in Shanghai during a late-night downpour and was
set to begin three days of meetings to discuss climate change, North Korea and the global economic
crisis with President Hu Jintao. But atop that list will be an effort to try to get Mr.~Hu close to
the spot where Russia, the other permanent member of the United Nations Security Council that is
reluctant to impose sanctions, appears to have arrived.

After an hourlong meeting in Singapore on Sunday afternoon, Mr.~Obama managed to get President
Dmitri A.~Medvedev of Russia to express dissatisfaction with Iran's response to a nuclear offer made
by world powers, raising the prospect that sanctions may be the next step in the West's continuing
effort to rein in Tehran's nuclear ambitions.

The leaders, meeting on the sidelines of an Asia-Pacific summit meeting in Singapore, also made
progress in efforts to negotiate a replacement for a key arms control treaty between the United
States and Russia that is set to expire in December, administration officials said.

While White House officials acknowledged on Sunday that there was no way a new agreement to replace
the Strategic Arms Reduction Treaty, or Start, would be ratified by the legislatures of either
country even if it was completed by the end of the year, they said that they expected to at least
reach an interim ``bridge'' agreement that would preserve the status quo until a new agreement was
reached.

Sunday afternoon's session was the fifth bilateral meeting with Mr.~Medvedev since Mr.~Obama took
office vowing to repair America's relationship with Russia, and administration officials expressed
satisfaction with their progress so far.

``I have found, as always, President Medvedev frank, constructive and thoughtful,'' Mr.~Obama said
after the meeting.

``The reset button has worked,'' he added, alluding\cite{allude} to the administration's early
promise to press a ``reset'' button in America's relationship with Russia after several years of
bickering over a host of issues like missile defense and the status of Kosovo.

With Start set to expire soon, the Obama administration is searching for ways to keep inspectors in
Russia -- or it risks losing American eyes on the world's second most formidable nuclear arsenal for
the first time in decades.

In the absence of a treaty or an ad hoc but legally binding ``bridge'' authority, American
inspectors would be forced to leave Russia when the treaty expired, and Russian inspectors would
have to leave the United States.

Under Start, the United States is allowed a maximum of 30 inspectors in Russia to monitor compliance
with the treaty. Russia likewise has interests in finding a bridge mechanism to continue its similar
right to inspections in the United States.

On Iran, Mr.~Obama and Mr.~Medvedev discussed a timetable for imposing sanctions if Tehran and the
West do not reach an accord soon on a proposal in which Iran would send its stockpile of enriched
uranium out of the country for either temporary safekeeping or reprocessing into fuel rods,
administration officials said.

``Unfortunately, so far at least, Iran appears to have been unable to say yes to what everyone
acknowledges is a creative and constructive approach,'' Mr.~Obama said. ``We are running out of time
with respect to that approach.''

More significant, Mr.~Medvedev also alluded to running out of patience. He said that while the
negotiation with Iran was continuing, ``we are not completely happy about its pace. If something
does not work, there are other means to move the process further.''

Robert Gibbs, the White House spokesman, said that the United States had set an internal deadline of
the end of the year.

During Mr.~Obama's trip, his first to Asia as president, he has taken to referring to himself as
``America's first Pacific president,'' a term he first used during a speech in Tokyo on Saturday
morning.

Mr.~Obama drew some fire from conservative American bloggers who accused him of going too far to
reassure Asian leaders: they complained that he should not have bowed to Emperor Akihito of Japan
when he went to the emperor's residence for lunch.

``During his meetings and his speech in Tokyo, the president observed protocol and enhanced the
status of American interests in Japan and across Asia,'' said an administration official traveling
with the president, who spoke on the condition of anonymity according to protocol. ``Those who
suggest otherwise are way off base and only looking to score political points.''

Throughout his trip, Mr.~Obama is seeking to ensure that the fast-growing region remains firmly
cemented to the United States, despite disparities in economic expansion and the rising influence of
China.

On Sunday, he became the first American president to meet with Myanmar's military leaders when he
attended a summit meeting of the Southeast Asian group Asean, held on the sidelines of the
Asia-Pacific Economic Cooperation talks in Singapore. Mr.~Obama, who has made his willingness to
engage with adversaries one of his foreign policy hallmarks, sat four places away from Gen. Thein
Sein, the prime minister of Myanmar, at the meeting table on Sunday afternoon.

After the talks, the group issued a joint statement that called for Myanmar's elections scheduled
for next year to be free and fair. But the statement did not call for the release of the Myanmar
opposition leader Daw Aung San Suu Kyi. White House officials said that Mr.~Obama made a point of
demanding her release when he made a speech to Asean members.

Generally, joint statements out of meetings that encompass many countries -- like Asean and the
Asia-Pacific Economic Cooperation forum -- rarely make news or carry much weight, in part because
these organizations operate by consensus. Since Myanmar is a member of Asean, there was never much
chance that the organization's joint statement would call for the release of Mrs.~Aung San Suu Kyi
or other political prisoners unless Myanmar's military leaders agreed to that first.

On Monday, the president is scheduled to meet with Shanghai's leaders and then hold what is being
billed as a town hall meeting with ``future leaders of China,'' mostly university students.

Mr.~Obama plans to take questions from the youths but also to field questions submitted through the
Internet. The meeting is expected to be broadcast live inside China, according to several Chinese
journalists, and also on the White House's Web site.

\section{China's Role as Lender Alters Obama's Visit}

\lettrine{W}{hen}\mycalendar{Nov.'09}{16} President Obama visits China for the first time on Sunday,
he will, in many ways, be assuming the role of profligate spender coming to pay his respects to his
banker.

That stark fact -- China is the largest foreign lender to the United States -- has changed the core
of the relationship between the United States and the only country with a reasonable chance of
challenging its status as the world's sole superpower.

The result: unlike his immediate predecessors, who publicly pushed and prodded China to follow the
Western model and become more open politically and economically, Mr.~Obama will be spending less
time exhorting Beijing and more time reassuring it.

In a July meeting, Chinese officials asked their American counterparts detailed questions about the
health care legislation making its way through Congress. The president's budget director, Peter
R.~Orszag, answered most of their questions. But the Chinese were not particularly interested in the
public option or universal care for all Americans.

``They wanted to know, in painstaking detail, how the health care plan would affect the deficit,''
one participant in the conversation recalled. Chinese officials expect that they will help finance
whatever Congress and the White House settle on, mostly through buying Treasury debt, and like any
banker, they wanted evidence that the United States had a plan to pay them back.

It is a long way from the days when President George W.~Bush hectored China about currency
manipulation, or when President Bill Clinton exhorted the Chinese to improve human rights.

Mr.~Obama has struck a mollifying note with China. He pointedly singled out the emerging dynamic at
play between the United States and China during a wide-ranging speech in Tokyo on Saturday that was
meant to outline a new American relationship with Asia.

``The United States does not seek to contain China,'' Mr.~Obama said. ``On the contrary, the rise of
a strong, prosperous China can be a source of strength for the community of nations.''

He alluded to human rights but did not get specific. ``We will not agree on every issue,'' he said,
``and the United States will never waver in speaking up for the fundamental values that we hold dear
-- and that includes respect for the religion and cultures of all people.''

White House officials have been working for months to make sure that Mr.~Obama's three-day visit to
Shanghai and Beijing conveys a conciliatory image. For instance, in June, the White House told the
Dalai Lama that while Mr.~Obama would meet him at some point, he would not do so in October, when
the Tibetan spiritual leader visited Washington, because it was too close to Mr.~Obama's visit to
China.

Greeting the Dalai Lama, whom China condemns as a separatist, weeks before Mr.~Obama's first
presidential trip to the country could alienate Beijing, administration officials said. Every
president since George H.~W. Bush in 1991 has met the Dalai Lama when he visited Washington, usually
in private encounters at the White House, although in 2007 George W.~Bush became the first president
to welcome him publicly, bestowing the Congressional Gold Medal on him at the Capitol. Mr.~Obama met
the Dalai Lama as a senator.

Similarly, while he was campaigning for the presidency, Mr.~Obama several times accused China of
manipulating its currency, an allegation that the current Treasury secretary, Timothy F.~Geithner,
repeated during his confirmation hearings. But in April, the Treasury Department retreated from that
criticism, issuing a report that said China was not manipulating its currency to increase its
exports.

While American officials said privately that they remained frustrated that China's currency policies
lowered the cost of Chinese goods and made American products more expensive in foreign markets, they
said that they were relieved that China was fighting the global recession with an enormous fiscal
stimulus program to spur domestic growth, and added that now was not the time to antagonize Beijing.

China is not viewed as a trouble spot for the United States. But this administration, like its
predecessor, has had difficulty grappling with a rising power that seems eager to avoid direct
clashes with the United States but affects its interests in many areas, including currency policy,
nuclear proliferation, climate change and military spending.

In that regard, two members of Mr.~Obama's foreign policy team said that the United States'
interactions with the Chinese had been far too narrow in past years, focusing on counterterrorism
and North Korea. Too little was done, they said, to address China's energy and environmental
policies, or its expansion of influence in Southeast Asia, South Asia and Africa, where China has
invested heavily and used billions of dollars in aid to advance its political influence.

One hint of the Obama administration's new approach came in a speech this fall by James
B.~Steinberg, the deputy secretary of state, who has deep roots in China policy. He argued that
China needed to adopt a policy of ``strategic reassurance'' to the rest of the world, a phrase that
appeared intended to be the successor to the framework of the Bush era, when China was urged to
embrace a role as a ``responsible stakeholder.''

``Strategic reassurance rests on a core, if tacit, bargain,'' Mr.~Steinberg said. ``Just as we and
our allies must make clear that we are prepared to welcome China's 'arrival,' '' he argued, the
Chinese ``must reassure the rest of the world that its development and growing global role will not
come at the expense of security and well-being of others.''

The Chinese reaction has been mixed, at best. The official China Daily newspaper ran a column just
before Mr.~Obama's arrival suggesting that the United States needed to provide some assurance of its
own -- to ``respect China's sovereignty and territorial integrity,'' code words for entirely backing
away from the issues of how China deals with Taiwan and Tibet.

In the United States, the phrase ``strategic reassurance'' has been attacked by conservative
commentators, who argue that any reassurance that the United States provides to China would be an
acknowledgment of a decline in American power.

In an op-ed article in The Washington Post, the analysts Robert Kagan and Dan Blumenthal argued that
the policy had echoes of Europe ``ceding the Western Hemisphere to American hegemony'' a century
ago. ``Lingering behind this concept is an assumption of America's inevitable decline,'' they wrote.
White House officials shot back, insisting that it is China that needs to do the reassurance, not
the United States.

In China, Mr.~Obama will meet with local political leaders and will host an American-style town hall
meeting with students in Shanghai. He will then spend two days in Beijing meeting with President Hu
Jintao.

It seems unlikely that Mr.~Obama will get the same celebrity-type reception in Beijing that he
received in Cairo, Ghana, Paris and London. China seems mostly immune to the Obama fever that swept
other parts of the world, and the Chinese are growing more confident that their country has the
wherewithal to compete with the United States on the world stage, analysts say.

``Obama is still a positive guy, and all over the world most people think he's more energetic, more
sincere, than Bush, more a reformist,'' said Shi Yinhong, a professor and an expert on United
States-China relations at People's University in Beijing. ``But in China, Obama's popularity is less
than in Europe, than Japan or Southeast Asia.'' In China, he said, ``there is no worship of Obama.''

For instance, during the Bush and Clinton years, China might release a few political dissidents on
the eve of a visit by the president as a good-will gesture. This time, American officials say, they
do not expect any similar gestures, although they say that Mr.~Obama will raise human rights issues
privately with Mr.~Hu.

``This time China will agree to have a human rights dialogue with the U.S.~on some cases,'' Mr.~Shi
said, but ``the arguments have changed compared to the past. Now we say, 'We are a different
country, we have our own system, our own culture.' ''

\section{Obama Wades Into Internet Censorship in China Address}

\lettrine{H}{e}\mycalendar{Nov.'09}{17} didn't explicitly call on China's leaders to lift the veil
of state control that restricts Internet access and online social networking here. But President
Obama did tiptoe -- ever so lightly -- into that controversial topic on Monday when he told students
in Shanghai that a free and unfettered Internet is a source of strength, not weakness.

For Mr.~Obama, who has been taking pains to strike a conciliatory note during his first visit to
China, it was a rare challenge to Chinese authorities, but expressed in Mr.~Obama's now familiar
nuance. Responding to a question that came via the Internet during a town hall meeting with Shanghai
students -- ``Should we be able to use Twitter freely?'' -- Mr.~Obama firstly started to answer in
the slightly off-the-point manner which he often uses when he is gathering his thoughts.

``Well, first of all, let me say that I have never used Twitter,'' he said. ``My thumbs are too
clumsy to type in things on the phone.''

But then he appeared to gather confidence. ``I should be honest, as president of the United States,
there are times where I wish information didn't flow so freely because then I wouldn't have to
listen to people criticizing me all the time,'' he said. But, he added, ``because in the United
States, information is free, and I have a lot of critics in the United States who can say all kinds
of things about me, I actually think that that makes our democracy stronger and it makes me a better
leader because it forces me to hear opinions that I don't want to hear.''

On a trip where he has gone out of his way to present a kinder and gentler image of America --
bowing before Emperor Akihito in Japan (which raised the ire of right-wing bloggers back home),
meeting with one of the military rulers of Myanmar, reassuring China that America doesn't seek to
contain the rising economic giant -- the Twitter question, and Mr.~Obama's answer, stood out as a
stark snapshot of a young American president's efforts to reach China's youth while not offending
its authorities.

``I will no forget this morning,'' one Chinese Twitterer said. ``I heard, on my shaky Internet
connection, a question about our own freedom which only a foreign leader can discuss.''

Interestingly, China's government itself demonstrated some restraint, and allowed the Twitter
question and Mr.~Obama's answer to stay up on websites hours after the town hall meeting.

That restraint, however, apparently only went so far. The students --some 500 --in the audience
seemed handpicked by the government and many were members of the Communist Youth League, which is
closely affiliated with President Hu Jintao.

That could explain some of the questions, like this one, offered by a young man who said the
question came in from the Internet from a Taiwan businessman worried that some people in America
were selling arms and weapons to Taiwan. ``I worry that this may make our cross-straits relations
suffer,'' the questioner said. ``So I would like to know if, Mr.~President, are you supportive of
improved cross-straits relations?''

Mr.~Obama grabbed the out that the questioner gave him and ran with it. Making no mention of the
part about arms sales to Taiwan, he instead offered up the standard American talking point on
Taiwan. ``My administration fully supports a one-China policy, as reflected in the three joint
communiqu\'es that date back several decades, in terms of our relations with Taiwan as well as our
relations with the People's Republic of China,'' he said.

Unlike previous town hall gatherings in China with other American presidents, Mr.~Obama's
question-and-answer session was not broadcast live on China's official state network. Instead,
according to the official Chinese news agency Xinhua, live broadcasts inside China were carried on
the agency's Web site and on local Shanghai stations.

The White House streamed the event live on its Web site, , which is not blocked or censored in
China, and a simultaneous Chinese translation was offered. The feed also was available through the
White House page on Facebook.

Unlike American town hall events, where speakers blast campaign songs while the audience chatters
loudly, you could almost hear a pin drop as the students waited for Mr.~Obama in an auditorium at
the Museum of Science and Technology.

Qian Yu, a student from East China Normal University, said she was impressed with Mr.~Obama but not
happy about the limited number of questions he took. ``I wish it had been a longer time,'' she said
after. ``I had lots of questions I'd have liked to ask.''

Mr.~Obama left Shanghai immediately after the town hall meeting, and flew to Beijing where he has a
packed schedule: several meetings with China's leaders, two dinners with Mr.~Hu, including an
elaborate state dinner Tuesday night, and tours of the Great Wall and the Forbidden City.

\section{Is Doomsday Coming? Perhaps, but Not in 2012}

\lettrine{N}{asa}\mycalendar{Nov.'09}{17} said last week that the world was not ending -- at least
anytime soon. Last year, CERN, the European Center for Nuclear Research, said the same thing, which
I guess is good news for those of us who are habitually jittery. How often do you have a pair of
such blue-ribbon scientific establishments assuring us that everything is fine?

On the other hand, it is kind of depressing if you were looking forward to taking a vacation from
mortgage payments to finance one last blowout.

CERN's pronouncements were intended to allay concerns that a black hole would be spit out of its new
Large Hadron Collider and eat the Earth.

The announcements by the National Aeronautics and Space Administration, in the form of several Web
site postings and a video posted on YouTube, were in response to worries that the world will end on
Dec.~21, 2012, when a 5,125-year cycle known as the Long Count in the Mayan calendar supposedly
comes to a close.

The doomsday buzz reached a high point with the release of the new movie ``2012,'' directed by
Roland Emmerich, who previously inflicted misery on the Earth from aliens and glaciers in
``Independence Day'' and ``The Day After Tomorrow.''

In the movie, an alignment between the Sun and the center of the galaxy on Dec.~21, 2012, causes the
Sun to go berserk with mighty storms on its surface that pour out huge numbers of the elusive
subatomic particles known as neutrinos. Somehow the neutrinos transmute into other particles and
heat up the Earth's core. The Earth's crust loses its moorings and begins to weaken and slide
around. Los Angeles falls into the ocean; Yellowstone blows up, showering the continent with black
ash. Tidal waves wash over the Himalayas, where the governments of the planet have secretly built a
fleet of arks in which a select 400,000 people can ride out the storm.

But this is only one version of apocalypse out there. In other variations, a planet named Nibiru
crashes into us or the Earth's magnetic field flips.

There are hundreds of books devoted to 2012, and millions of Web sites, depending on what
combination of ``2012'' and ``doomsday'' you type into Google.

All of it, astronomers say, is bunk.

``Most of what's claimed for 2012 relies on wishful thinking, wild pseudoscientific folly, ignorance
of astronomy and a level of paranoia worthy of 'Night of the Living Dead,' '' Ed Krupp, director of
the Griffith Observatory, in Los Angeles, and an expert on ancient astronomy, wrote in an article in
the November issue of Sky \& Telescope.

Personally, I have been in love with end-of-the-world stories since I started consuming science
fiction as a disaffected child. Scaring the pants off the public has been pretty much the name of
the game ever since Orson Welles broadcast ``War of the Worlds,'' a fake newscast about a Martian
invasion of New Jersey, in 1938.

But the trend has gone too far, suggested David Morrison, an astronomer at the NASA Ames Research
Center in Moffett Field, Calif., who made the YouTube video and is one of the agency's point people
on the issue of Mayan prophecies of doom.

``I get angry at the way people are being manipulated and frightened to make money,'' Dr.~Morrison
said. ``There is no ethical right to frighten children to make a buck.''

Dr.~Morrison said he had been getting about 20 letters and e-mail messages a day from people as far
away as India scared out of their wits. In an e-mail message, he enclosed a sample that included one
from a woman wondering if she should kill herself, her daughter and her unborn baby. Another came
from a person pondering whether to put her dog to sleep to avoid suffering in 2012.

All of this reminded me of the kinds of letters I received last year about the putative black hole
at CERN. That too was more science fiction than science fact, but apparently there is nothing like
death to bring home the abstract realms of physics and astronomy. In such situations, when the Earth
or the universe is trying to shrug you and your loved ones off this mortal plane, the cosmic does
become personal.

Dr.~Morrison said he did not blame the movie for all this, as much as the many other purveyors of
the Mayan prediction, as well as the apparent failure of some people, reflected in so many arenas of
our national life, to tell reality from fiction. But then, he said, ``my doctorate is in astronomy,
not psychology.''

In an e-mail exchange, Dr.~Krupp said: ``We are always uncertain about the future, and we always
consume representations of it. We are always lured by the romance of the ancient past and by the
exotic scale of the cosmos. When they combine, we are mesmerized.''

A NASA spokesman, Dwayne Brown, said the agency did not comment on movies, leaving that to movie
critics. But when it comes to science, Mr.~Brown said, ``we felt it was prudent to provide a
resource.''

If you want to worry, most scientists say, you should think about global climate change, rogue
asteroids or nuclear war. But if speculation about ancient prophecies gets you going, here are some
things Dr.~Morrison and the others think you should know.

To begin with, astronomers agree, there is nothing special about the Sun and galactic center
aligning in the sky. It happens every December with no physical consequences beyond the
overconsumption of eggnog. And anyway, the Sun and the galactic center will not exactly coincide
even in 2012.

If there were another planet out there heading our way, everybody could see it by now. As for those
fierce solar storms, the next sunspot maximum will not happen until 2013, and will be on the mild
side, astronomers now say.

Geological apocalypse is a better bet. There have been big earthquakes in California before and
probably will be again. These quakes could destroy Los Angeles, as shown in the movie, and
Yellowstone could erupt again with cataclysmic force sooner or later. We and our works are indeed
fragile and temporary riders on the Earth. But in this case, ``sooner or later'' means hundreds of
millions of years, and there would be plenty of warning.

The Mayans, who were good-enough astronomers and timekeepers to predict Venus's position 500 years
in the future, deserve better than this.

Mayan time was cyclic, and experts like Dr.~Krupp and Anthony Aveni, an astronomer and
anthropologist at Colgate University, say there is no evidence that the Mayans thought anything
special would happen when the odometer rolled over on this Long Count in 2012. There are references
in Mayan inscriptions to dates both before the beginning and the ending of the present Long Count,
they say, just as your next birthday and April 15 loom beyond New Year's Eve, on next year's
calendar.

So keep up those mortgage payments.

\section{After Microsoft, Bringing a High-Tech Eye to Professional Kitchens}

\lettrine{I}{nside}\mycalendar{Nov.'09}{17} a nondescript warehouse on a nondescript street of this
Seattle suburb is a research laboratory that looks like it came out of a James Bond movie -- had Q
the gadget master been a gastronome.

Here Nathan Myhrvold, a former chief technology officer at Microsoft, and his company, Intellectual
Ventures, pursue an eclectic array of speculative and potentially world-changing ideas -- inventing
a new battery, taming hurricanes, defeating disease. And here, along with the laser designed to
shoot mosquitoes out of the air (a high-speed camera counts the rate of wing-flapping to ensure that
innocent insects are not vaporized), is the best-equipped restaurant kitchen anywhere that never
serves any customers.

Dr.~Myhrvold exuded a Willy Wonka enthusiasm as he talked of the foods that came out of his
industrial food dehydrator. ``Raw lobster tail, freeze dried, is amazing,'' he said.

At another machine, rose petals spun inside a glass globe. ``This is basically a still,'' he said.
``You could crank the temperature up and distill alcohol. What we're trying to do here is get an
essence of rose petals.''

The yield would be a few fragrant tablespoons of liquid.

Around the corner, he pointed to two machines side by side. ``Here's our ice cream machine, and
here's our ultrasonic welder,'' he said.

Had he used the welder as a cooking appliance? ``Not yet,'' he said, earnestly,'' but we're going to
try it out.''

After all, an autoclave designed to sterilize lab equipment has proven culinarily productive --
``It's basically the pressure cooker from hell,'' Dr.~Myhrvold said -- as has a 100-ton hydraulic
press, for beef jerky.

All of this high-tech kitchen tinkering feeds another of Dr.~Myhrvold's projects: a cookbook.

The book, still untitled, intends to be the authoritative reference for chefs wishing to employ
so-called molecular gastronomy -- adapting food industry technologies to restaurant cooking.

Dr.~Myhrvold, who once presided over Microsoft Windows, did not undertake this endeavor as a lonely
intellectual pursuit. He hired 15 people, including 5 professional chefs, a photographer, an art
director and writers and editors, to create it. They included Christopher Young, a
biochemistry-graduate-student-turned-chef who headed the research kitchen at the Fat Duck near
London, one of the most innovative restaurants in the world.

Dr.~Myhrvold has long pursued a Renaissance man portfolio of interests. While still at Microsoft, he
showed that sauropod dinosaurs might have been able to accelerate the tips of their tails to
supersonic speeds, like cracking a whip. More recently, he has been proselytizing among
paleontologists, urging them to hunt for fossilized dinosaur vomit. Owls and some other birds of
prey regurgitate the bones of what they eat, and Dr.~Myhrvold surmises that dinosaurs, as the
ancestors of birds, might have done the same thing.

Every month or so, the cookbook team gathers in a conference room to review their progress.
Dr.~Myhrvold scans each page, points out glitches and sketches how he wants a chart to look.

``It's basically like a software project,'' Dr.~Myhrvold said. ``It's very much like a review we
would do at Microsoft.''

The project has grown in size and scope. Originally planned as a 300-page discussion of sous vide,
an increasingly popular restaurant technique of cooking food in vacuum-sealed bags in warm water
baths, the book has swelled to 1,500 pages that will also cover microbiology, food safety, the
physics of heat transfer on the stove and in the oven, formulas for turning fruit and vegetable
juices into gels, and more.

``And they're big pages,'' Dr.~Myhrvold said.

Because he is self-publishing the book, Dr.~Myhrvold does not have to convince a publisher or anyone
else that such a huge book aimed primarily at a narrow of audience of restaurant chefs makes
economic sense. He said the book would be out in a year, although he admitted that was also what he
said a year ago.

``There's not a chef on Earth who won't learn something from this,'' Dr.~Myhrvold said.

At least some chefs are taking interest.

``I think there are parts of it that are definitely new to me,'' said Wylie Dufresne, the chef and
owner of WD-50 in Manhattan, who visited the kitchen laboratory. ``It's a cookbook that's going to
be in its own category.''

In September, Dr.~Myhrvold, Mr.~Young and two of the other chefs gave a presentation at the
StarChefs.com International Chefs Congress, an annual Manhattan trade show for restaurant
professionals.

They demonstrated how to encrust a pork loin within what was essentially a large crispy pork rind,
how to make stewed prunes look like coal and how to make a ``constructed cream'' -- breaking apart a
fat and a liquid into tiny droplets and mixing them together into something that had the fluidity of
heavy cream.

``For example, why not make a pistachio cream where, instead of milk fat, we use pistachio oil?''
Mr.~Young said. ``If you can get the droplets small enough, if you can coat them in proteins, you
can create a dairy-free pistachio gelato that's 100 percent pistachio oil.''

They also demonstrated cryoseared duck breast, a technique that calls for implements not typically
found in a kitchen: a small satchel of loose metal, dry ice, dog hair brush.

``We do have to perforate it to get the fat out,'' Mr.~Young said. ``The easiest way is a stainless
steel dog hair brush. It will poke a lot of little holes that aren't going to show up.''

``For God's sake,'' Dr.~Myhrvold interjected, ``buy a new one for this.''

The duck breast was placed skin down on the dry ice before being seared, weighed down by the
satchel. The cold froze not only the skin, but also a thin layer of the meat next to the skin, which
acted as a cold barrier to prevent overcooking. ``Until that melts, heat won't go above it,''
Mr.~Young said.

The result was a crispy skin while the meat remained tender and juicy.

The book presents some concepts like wet bulb temperature that will be new to cooks of all skill
levels. The usual temperature set in an oven is scientifically known as the dry bulb temperature.
But for cooking, the wet bulb temperature, which is essentially a measure of the evaporation rate of
water and depends on the humidity, is often more important.

``Nathan would make the point that food is water with a bunch of impurities in it,'' Mr.~Young said.

Dr.~Myhrvold said the wet bulb and dry bulb temperatures could sometimes differ by 30 degrees, which
might be the difference between a souffl\'e that rises and one that collapses. (One advantage of
sous vide is that because the food is immersed in water, the wet and dry bulb temperatures are the
same.)

In another discovery of culinary heat transfer physics, Dr.~Myhrvold said the bulbous shape and
black color of Weber grills were wrong. To achieve an even cooking temperature across the cooking
grate, the inside of the grill should be vertical and shiny to reflect the heat.

That can be fixed by adding an aluminum insert to the grill. ``So we have directions for that,''
Dr.~Myhrvold said.

The conclusions have often been backed up by careful scientific exploration. For example, confit,
the French technique of cooking slowly in fat, is supposed to impart a unique taste and texture as
the fat penetrates the meat.

But Dr.~Myhrvold said: ``There's no way it could penetrate. The molecules are too big.''

He said double-blind taste tests proved that the same tasty results could be achieved by steaming
and then rubbing some of the fat on the outside.

\section{In China, Obama Pushes Need for `Strong Dialogue'}

\lettrine{P}{resident}\mycalendar{Nov.'09}{17} Obama and President Hu Jintao of China met in private
off Tiananmen Square here on a frigid Tuesday morning to discuss issues like trade, climate change
and the nuclear programs of Iran and North Korea, in a session that signaled the central role of
China on the world stage.

The leaders told reporters afterward that the United States and China were in agreement on a range
of issues, but they spoke only in general terms.

At a news conference where both presidents appeared, neither took questions from reporters, staying
in line with the minutely stage-managed atmosphere of Mr.~Obama's first visit to China. They said in
separate speeches that the two nations would work together to stabilize the teetering world economy,
contain the dangers of climate change and prevent nuclear proliferation.

The public pronouncements were full of familiar rhetoric. At the start of their first meeting,
Mr.~Obama told Mr.~Hu: ``We believe strong dialogue is important not only for the U.S.~and China,
but for the rest of the world.''

The leaders greeted each other at the door of the Great Hall of the People after Mr.~Obama's
motorcade slithered its way past thousands of onlookers crowding around Tiananmen Square, in front
of the giant portrait of Mao, to catch a glimpse of the American president.

The leaders shook hands and walked up the red carpet, Chinese military leaders facing them. At the
conference table where the first bilateral meeting was held, Mr.~Obama sat flanked by senior cabinet
members.

The meeting came the day after Mr.~Obama tried to hold a frank and public discussion with Chinese
students in Shanghai. The event was called a town hall, but Mr.~Obama's meeting with about 500
students had little in common with the sometimes raucous exchanges that have become a fixture of
American politics.

It was, instead, an example of Chinese stagecraft. Most of those who attended the event at the
Museum of Science and Technology turned out to be members of the Communist Youth League, an official
organization that grooms obedient\cite{obedient} students for future leadership posts.

Some Chinese bloggers whom the White House had tried to invite were barred from attending. Even
then, the Chinese government took no chances, declining to broadcast the event live to a national
audience -- or even mention it on the main evening newscast of state-run China Central Television.

The scripted interaction underscored the obstacles Mr.~Obama faces as he tries to manage the
American relationship with an authoritarian China, whose wealth and clout have surged as its economy
has weathered the global downturn far better than the United States' or Europe's.

It remained unclear whether the United States would make progress on several issues on this trip,
including on the management of its tightly controlled currency, the renminbi, or on how to keep Iran
from developing nuclear weapons. China has rejected American pressure to allow the renminbi to float
freely and has opposed tougher sanctions on Iran.

The degree of control exercised over the most public event of Mr.~Obama's three-day stay in China
suggests that Chinese leaders are less willing to make concessions to American demands for the
arrangements of a presidential visit than they once were.

The White House spent weeks wrangling with the Chinese authorities over who would be allowed to
attend the Shanghai town hall meeting, including how much access the media would have and whether it
would be broadcast live throughout the country. In the end, Mr.~Obama had little chance to promote a
message to the broader Chinese public.

One student who participated in the meeting said she was trained for four days by the Chinese
government, and told not to ask questions about Tibet or human rights and to be respectful of
President Obama and consider the implications any question would have on U.S.-China relations. She
asked not to be identified for fear of being punished by her university.

The event in some respects signaled a retreat from the reception given at least two earlier American
presidents, Bill Clinton and George W.~Bush, who both asked for, and were granted, the opportunity
to address the Chinese people and answer their questions in a live national broadcast.

One local television station broadcast Mr.~Obama's session live. But the official Xinhua news agency
offered only a transcript of the exchange on its Web site instead of the live Webcast it had
promised. The White House streamed the event live on its Web site, which did not appear to be
blocked inside China. But that site is not a common destination for most Chinese looking for
breaking news.

Although it was carefully choreographed, the event gave Mr.~Obama a little room to prod the Chinese
authorities toward more openness. In his initial remarks at the forum, Mr.~Obama said the United
States was not seeking to impose its political system on other countries, but he called freedom of
expression and worship among the ``universal rights'' common to all people.

He did, however, steer clear of the most delicate human rights topics, like the recent unrest in the
Chinese regions of Tibet and Xinjiang, and he focused most of his comments on the need for China and
the United States to become partners instead of rivals.

His tone reflected the fact that China had become the largest foreign lender to the United States at
a time when America's total public debt is surging and its economy is still trying to claw its way
out of a deep slump. Mr.~Obama said the two countries carried a ``burden of leadership'' on issues
like climate change and nuclear nonproliferation, and said they needed to work more closely on
matters of mutual concern.

``I will tell you, other countries around the world will be waiting for us,'' Mr.~Obama said at the
town hall meeting. He later flew to Beijing for a dinner and full state visit hosted by Mr.~Hu.

At the Shanghai forum, Mr.~Obama was asked only one question -- ``Should we be able to use Twitter
freely?'' -- that delved into an area the Chinese government considers controversial.

His cautious answer stood out as a sign that he hopes to reach China's youth without offending its
increasingly influential leaders. He delivered an oblique critique of China's rigid controls and
restrictions on the Internet and free speech without mentioning that China practices online
censorship as a matter of policy.

``I have a lot of critics in the United States who can say all kinds of things about me,'' he said.
But, he added, ``I actually think that that makes our democracy stronger, and it makes me a better
leader because it forces me to hear opinions that I don't want to hear.''

That snippet, at least initially, captured the attention of Chinese netizens. It was a topic of
discussion on Web sites for a couple of hours after Mr.~Obama spoke, before being deleted or removed
from prominent positions. According to several Web snapshots in the hours after the meeting,
``What's Twitter?'' and ``Obama Shanghai'' shot up to the list of top 10 Chinese Google searches.

``I will not forget this morning,'' one Chinese Twitter user posted on the Internet, apparently
using software to get around the government firewall. ``I heard, on my shaky Internet connection, a
question about our own freedom which only a foreign leader can discuss.''

But most of the questions appeared to reflect the careful vetting of the crowd by the Chinese.
Beijing vetoed the White House's attempt to invite a group of popular bloggers, an audience
component that administration officials hoped would make the session more authentic, according to
several people who were asked to participate in the forum.

``I was invited, but then a few days ago I was told we can't go,'' said Michael Anti, a popular
blogger who formerly worked as a research assistant at The New York Times and was a Nieman fellow at
Harvard last year. ``I don't know why.''

\section{During Visit, Obama Skirts Chinese Political Sensitivities}

\lettrine{W}{hether}\mycalendar{Nov.'09}{18} by White House design or Chinese insistence, President
Obama has steered clear of public meetings with Chinese liberals, free press advocates and even
average Chinese during his first visit to China, showing a deference to the Chinese leadership's
aversions to such interactions that is unusual for a visiting American president.

Mr.~Obama held a ``town hall'' meeting with students on Monday. But the students were carefully
vetted and prepped for the event by the government, participants said. And the Chinese authorities,
wielding a practiced mix of censorship and diplomatic pressure, succeeded in limiting Mr.~Obama's
exposure to a point where a third of some 40 Beijing university students interviewed Tuesday were
unaware that he had just met in Shanghai with their peers.

Some students who were aware cast him in terms rarely applied to American leaders, like ``rather
humble'' and ``bland.''

``Is America being capricious because their economic difficulties force them to be nicer to China
and other countries, or is this a genuine change?'' asked Liu Ziqi, 18, a freshman at the University
of International Business and Economics. ``I don't know.''

This is no longer the United States-China relationship of old but an encounter between a weakened
giant and a comer with a bit of its own swagger. Washington's comparative advantage in past meetings
is now diminished, a fact clearly not lost on the Chinese.

Human rights is the prime example. In 1998, President Bill Clinton staged a nationally broadcast
discussion with the president at the time, Jiang Zemin, about human rights, the Dalai Lama and
perhaps China's most taboo topic, the 1989 Tiananmen Square protests. In 2002, President George
W.~Bush stressed liberty, the rule of law and faith in a speech to university students broadcast
across China.

When Mr.~Obama visited Moscow in July, he met with opposition political activists and journalists,
and he publicly questioned the prosecution of an anti-Kremlin businessman.

In China, by contrast, Mr.~Obama, in nuanced references to human rights, has shied away from citing
China's spotty record, even when offered the chance. Asked Monday in Shanghai to discuss China's
censorship of the Internet, the president replied by talking about America's robust political
debates.

American scholars and activists, who requested anonymity for fear of damaging relations with the
White House, said the administration rejected proposals for brief meetings in Beijing with Chinese
political activists, and then with lawyers.

American officials did consider organizing meetings between Mr.~Obama and Chinese lawyers,
university students in Beijing and Hu Shuli, a well-known Chinese journalist who recently ceded
control of Caijing, one of the nation's most respected and independent magazines. But officials say
time constraints, not political considerations, sidelined those options, although the sightseeing
agenda remained intact.

One prominent defense lawyer, Mo Shaoping, said Tuesday that an American official called this month
to ask if he would meet with Mr.~Obama but never called back. ``The U.S.~should be the safeguard of
universal values,'' Mr.~Mo said, but Mr.~Obama ``actually didn't make it a very high priority.''

For its part, the Chinese government made sure Mr.~Obama did not bump into protesters by placing
well-known activists under tighter security. Chinese Human Rights Defenders, a local organization,
said 20 people were detained, placed under house arrest or prohibited from traveling before
Mr.~Obama's visit.

Zhang Zuhua, once a Communist Party official and now among China's most influential civil rights
activists, said that additional police officers were watching his apartment and that he had been
warned to avoid political activity.

Mr.~Zhang expressed concern over what he called America's growing reluctance to criticize China on
human rights, saying ``the Communist Party can pay even less regard to it and tighten up.''

But an alternative explanation for Mr.~Obama's comparatively low profile here, curiously, is the
very insecurity of China's autocratic government.

In contrast to Mr.~Jiang, who sparred openly with President Clinton over human rights, President Hu
Jintao is a cautious politician whose tenure has been marked by an obsession with stability. In
Mr.~Obama's case, for example, Chinese officials hamstrung negotiations over items like the national
broadcast of Shanghai's town hall meeting until they achieved most of their objectives to limit its
exposure.

In China, Mr.~Obama does not enjoy the matinee-idol status that has followed him elsewhere. But the
Chinese are curious about the young president, and in some cases, they clearly find him a refreshing
contrast to their own retirement-age, shoe-black-haired leadership.

A topic of awe on Chinese chat sites this week was the image of Mr.~Obama descending from Air Force
One into rainy Shanghai, holding his own umbrella, without an aide's assistance.

In a Nov. 11 Internet poll, people were asked to say what was most memorable about Mr.~Obama. A
majority noted his Nobel Peace Prize. No.~2, improbable to foreigners, was a Chinese report that the
president had insisted on paying for his own hamburger at a Washington restaurant.

In this basketball-crazy nation, Mr.~Obama might single-handedly have remade America's image by
showing up on one of the city's many outdoor courts for a few rounds of hoops. Instead, he tiptoed
around fractious issues like human rights, as Chinese authorities took extra steps to ensure that
the state media not project any hint of disharmony.

One state newspaper editor said his newsroom now was more tense even than in June, when China passed
the 20th anniversary of the 1989 Tiananmen crackdown.

Late Monday, he said, after a Foreign Ministry official combed over the paper, editors scrapped two
articles scheduled for publication on Tuesday, including one straightforward news article on the
value of China's currency.

Mr.~Obama's trip, journalists at the paper joked, ``had driven the homeless from Beijing and brought
more censorship to China.''

``It's as if they think he'd read the paper and it would offend him and trigger an international
uproar,'' the editor said. ``As it is now, it would only trigger a snore.''


\section{China Holds Firm on Major Issues in Obama's Visit}

\lettrine{I}{n}\mycalendar{Nov.'09}{18} six hours of meetings, at two dinners and during a stilted
30-minute news conference in which President Hu Jintao did not allow questions, President Obama was
confronted, on his first visit, with a fast-rising China more willing to say no to the United
States.

On topics like Iran (Mr.~Hu did not publicly discuss the possibility of sanctions), China's currency
(he made no nod toward changing its value) and human rights (a joint statement bluntly acknowledged
that the two countries ``have differences''), China held firm against most American demands.

With China's micro-management of Mr.~Obama's appearances in the country, the trip did more to
showcase China's ability to push back against outside pressure than it did to advance the main
issues on Mr.~Obama's agenda, analysts said.

``China effectively stage-managed President Obama's public appearances, got him to make statements
endorsing Chinese positions of political importance to them and effectively squelched discussions of
contentious issues such as human rights and China's currency policy,'' said Eswar S.~Prasad, a China
specialist at Cornell University. ``In a masterstroke, they shifted the public discussion from the
global risks posed by Chinese currency policy to the dangers of loose monetary policy and
protectionist tendencies in the U.S.''

White House officials maintained they got what they came for -- the beginning of a needed
give-and-take with a surging economic giant. With a civilization as ancient as China's, they argued,
it would be counterproductive -- and reminiscent of President George W.~Bush's style -- for
Mr.~Obama to confront Beijing with loud chest-beating that might alienate the Chinese. Mr.~Obama,
the officials insisted, had made his points during private meetings and one-on-one sessions.

``I do not expect, and I can speak authoritatively for the president on this, that we thought the
waters would part and everything would change over the course of our almost two-and-a-half-day trip
to China,'' said Robert Gibbs, the White House spokesman. ``We understand there's a lot of work to
do and that we'll continue to work hard at making more progress.''

Several China experts noted that Mr.~Obama was not leaving Beijing empty-handed. The two countries
put out a five-point joint statement pledging to work together on a variety of issues. The statement
calls for regular exchanges between Mr.~Obama and Mr.~Hu, and asks that each side pay more attention
to the strategic concerns of the other. The statement also pledges that they will work as partners
on economic issues, Iran and climate change.

But despite a conciliatory tone that began weeks ago when Mr.~Obama declined to meet the Tibetan
spiritual leader, the Dalai Lama, before visiting China to avoid offending China's leaders, it
remains unclear whether Mr.~Obama made progress on the most pressing policy matters on the American
agenda in China or elsewhere in Asia.

The president has had to fend off criticism from American conservatives that he appeared to soften
the American stance on the positioning of troops on the Japanese island of Okinawa, and for bowing
to Japan's emperor.

At a regional conference in Singapore, Mr.~Obama announced a setback on another top foreign policy
priority, climate change, acknowledging that comprehensive agreement to fight global warming was no
longer within reach this year.

Past American presidents have usually insisted in advance on some concrete achievements from their
trips overseas. President Bush received vigorous endorsements of his top foreign policy priority,
the global war on terrorism, during his visits to Beijing, and President Bill Clinton guided China
toward joining the World Trade Organization after prolonged negotiations. When either of those
presidents visited the country, China often made a modest concession on human rights as well.

This time, Mr.~Hu declined to follow the lead of President Dmitri A.~Medvedev of Russia, who, after
months of massaging by the Obama administration, now says that he is open to tougher sanctions
against Iran if negotiations fail to curb Iran's nuclear program. The administration needs China's
support if tougher sanctions are to be approved by the United Nations Security Council. But during
the joint appearance in Beijing on Tuesday, Mr.~Hu made no mention of sanctions.

Rather, he said, it was ``very important'' to ``appropriately resolve the Iranian nuclear regime
through dialogue and negotiations.'' And then, as if to drive home that point, Mr.~Hu added,
``During the talks, I underlined to President Obama that given our differences in national
conditions, it is only normal that our two sides may disagree on some issues.''

White House officials acknowledged that they did not get what they wanted from Mr.~Hu on Iran but
said that Mr.~Obama's method would yield more in the long term. ``We're not looking for them to lead
or change course, we're looking for them to not be obstructionist,'' one administration official
said.

In a meeting in Beijing with a senior Chinese official on Wednesday morning, Secretary of State
Hillary Rodham Clinton again pressed China on Iran. She told the official, Dai Bingguo, that even if
China had not decided what sanctions on Iran it would accept, ``you need to send a signal,'' said a
senior American official, who spoke on condition of anonymity so he could describe the exchange.

Mr.~Obama did not appear to move the Chinese on currency issues, either. China has come under heavy
pressure, not only from the United States but also from Europe and several Asian countries, to
revise its policy of keeping its currency, the renminbi, pegged at an artificially low value against
the dollar to help promote its exports. Some economists say China must take that step to prevent the
return of large trade and financial imbalances that may have contributed to the recent financial
crisis.

Mr.~Obama on Tuesday could only cite China's ``past statements'' in support of shifting toward
market-oriented exchange rates, implying that he had not extracted a fresh commitment from Beijing
to move in that direction soon.

There are many reasons the White House may have heeded China's clear desire for a visit free of the
polemics that often accompany meetings between leaders of the two countries. Mr.~Obama's foreign
policy is rooted in recasting the United States as a thoughtful listener to friends and rivals
alike. ``No we haven't made China a democracy in three days -- maybe if we pounded our chest a lot
that would work,'' Mr.~Gibbs said in an e-mail message on Tuesday night. ``But it hasn't in the last
16 years.''

Kenneth Lieberthal, a Brookings Institution scholar who oversaw China issues in President Clinton's
White House, agreed. ``The United States actually has enormous influence on popular thinking in
China, but it is primarily by example,'' he said. ``If you go to the next step and say, 'You guys
ought to be like us,' you lose the impact of who you are.''

The National Security Council's spokesman, Michael A.~Hammer, added, ``What we did come to do is
speak bluntly about the issues which are important to us, not in an unnecessarily offensive manner,
but rather in the Obama style of showing respect.''

Mr.~Obama, even as he projected a softer image, did nudge the Chinese on some delicate issues.

On Tuesday, standing next to Mr.~Hu, Mr.~Obama brought up Tibet, where Beijing-backed authorities
have clamped down on religious freedom. ``While we recognize that Tibet is part of the People's
Republic of China, the United States supports the early resumption of dialogue between the Chinese
government and representatives of the Dalai Lama to resolve any concerns and differences that the
two sides may have,'' he said.

\section{Those Soccer Plays, in Context}

\lettrine{N}{early}\mycalendar{Nov.'09}{18} two weeks later, the University of New Mexico soccer
player Elizabeth Lambert said she still could not fully explain what led her to yank an opponent
from Brigham Young down by her ponytail in what has become a highly publicized incident of violent
behavior.

Her action was indefensible, Lambert said Tuesday in her first interview since the incident occurred
Nov. 5 in a Mountain West Conference semifinal game and led to her indefinite suspension from the
New Mexico team.

``I still deeply regret it and will always regret it and will carry it through the rest of my life
not to retaliate\cite{retaliate},'' said Lambert, a 20-year-old junior on scholarship.

She has watched the video a handful of times and does not recognize herself pulling down Brigham
Young's Kassidy Shumway, Lambert said.

``I look at it and I'm like, 'That is not me,' '' said Lambert, a defender and an all-conference
academic player. ``I have so much regret. I can't believe I did that.''

At the same time, she said other moments of aggressive play -- in which Lambert elbowed a Brigham
Young player in the back, received a yellow card for tripping, seemed to throw a punch at an
opponent's head and made a hard tackle from behind -- came during the forceful, insistent play that
routinely occurred in women's soccer but might be misunderstood by casual fans.

Some of her actions -- like the apparent punch, which she said was inadvertent -- were
misinterpreted or taken out of context on a condensed video, Lambert said. And she said she believed
that the incident was blown out of proportion because it occurred in a women's game. She said it was
wrongly reported to be her when it was actually a teammate who tried to clear a ball and
accidentally kicked it into the face of a B.Y.U. player.

``I definitely feel because I am a female it did bring about a lot more attention than if a male
were to do it,'' Lambert said. ``It's more expected for men to go out there and be rough. The
female, we're still looked at as, Oh, we kick the ball around and score a goal. But it's not. We
train very hard to reach the highest level we can get to. The physical aspect has maybe increased
over the years. I'm not saying it's for the bad or it's been too overly aggressive. It's a game.
Sports are physical.''

She added: ``I think the way the video came out, it did make me look like a monster. That's not the
type of player I am. I'm not just out there trying to hurt players. That's taking away from the
beauty of the game. And I would never want to do that.''

Lambert said she was shaken and appalled by some of the responses she received in e-mail messages,
telephone messages and on blogs, which included the publishing of her parents' home phone number in
Southern California and one suggestion that ``I should be taken to a state prison, raped and left
for dead in a ditch.''

She said she felt conflicting emotions and sometimes still woke up in a sweat.

``I'll be angry with myself that I did this, to my team, my university, that I did this to women's
soccer, a sport that many females have worked very hard to get respect for,'' Lambert said. ``And
I'll be sad that people want to see me suffer.''

She said she was taken aback at how the incident had been perceived by some as sexy catfighting
between two women. She said she was aghast that some men had sent her messages saying, ``Hey, we
should meet up some time.''

``That appalled me,'' Lambert said. ``A lot of people think I have a lot of sexual aggression. I was
like, 'Whoa, no, I don't feel that way at all.' That's bizarre and shocking to me.''

The game against Brigham Young began with familiar passion and intensity, Lambert said. Emotions
escalated after Brigham Young took a 1-0 halftime lead, given that a defeat could mean the end of
the season and a failure to qualify for the N.C.A.A. tournament, she said.

Opposing fans were mockingly chanting her name, she said, and players on both teams were playing
aggressively. She said she was called names and taken down to the ground with cheap shots. On video,
a B.Y.U. player can be seen elbowing Lambert in the stomach before she shoves the opponent in the
back in retaliation. Shumway can be seen tugging on Lambert's shorts before she is yanked down by
her ponytail.

If the referee Joe Pimentel had issued more yellow cards or a red card, Lambert said, ``It would
have been a very different game.''

Still, Lambert said that she did not want to throw Pimentel ``under the bus'' and that she did not
consider the game to be out of control.

Her coach, Kit Vela, never instructed her to ``take anybody out,'' Lambert said, adding that the
B.Y.U. players also did not appear to have malign intent.

Lambert said she eventually grew frustrated, as much with herself as with the opponent, saying she
had often struggled with self-confidence and with feeling ``that I'm accepted playing at this
level.''

Lambert said of the match: ``I've never been in a situation like that, where I was out of my
element. There were times in the game where I was literally like, 'All right, Elizabeth, you've got
to get control' of myself.''

In each of her two previous matches, Lambert had received a yellow-card warning, but those were the
only cautions in more than 2,500 minutes of play at New Mexico, a university official said.

Lambert said she did not consider herself a dirty player. Yet in the second half, she yanked Shumway
down by her ponytail and assumed widespread villainy.

``In that one moment, I let it all get into my head,'' Lambert said of the emotion of the game.

Later in the match, Lambert received a yellow card for tripping.

She is seeing a clinical psychologist on campus to better understand what caused the hair-pulling
incident. It is one of several steps she is taking, along with speaking to youth players about
acceptable behavior, so she can seek reinstatement to the team in the spring.

``I'm working on my mental game to never let that happen again,'' Lambert said. ``That's
unacceptable in any sport to get to that point where you feel it's necessary that you have to
retaliate in a dirty manner.''

\section{Blair's Chances for Europe Presidency Said to Fade}

\lettrine{B}{ritain}\mycalendar{Nov.'09}{18} is under pressure to abandon its campaign to install
former Prime Minister Tony Blair in a new European Union presidential post and indicate whether it
would instead put up a candidate for the job of foreign policy chief for the bloc.

With Herman Van Rompuy, the prime minister of Belgium, emerging as the front-runner for the post of
president, the search continued for a senior figure from the center-left -- to give political
balance -- for the foreign affairs post.

The center-left foreign minister of Spain, Miguel \'Angel Moratinos, has emerged as a potential
candidate for the foreign policy job, a post for which the former prime minister of Italy, Massimo
D'Alema, had been mentioned.

Despite speculation that the foreign secretary of Britain, David Miliband, its business minister,
Peter Mandelson, or the E.U. trade commissioner, Catherine Ashton, might be contenders, the British
government is pushing Mr.~Blair for the presidential post.

E.U. heads of government meet Thursday night to fill the top jobs created by the Lisbon Treaty,
which is designed to streamline decision-making and increase the bloc's global clout. But the task
has proved complex -- and a failure to agree smoothly about these top spots could prove
embarrassing.

``There is a near consensus -- though one in which London does not want to participate -- that it
cannot be Blair for the Council President job,'' said Thomas Klau of the European Council on Foreign
Relations.

Mr.~Klau argued that there was solid support for the post to be filled from the center-right to
reflect the fact that most E.U. governments were politically conservative.

Also on Tuesday, Sweden, which holds the rotating E.U. presidency and which is negotiating a jobs
package, nominated its European minister, Cecilia Malmstrom, for the European Commission. That is
significant because it means that the foreign minister of Sweden, Carl Bildt -- one of the strongest
contenders for the foreign policy post -- will not be a candidate for the foreign policy post, which
also takes up a country's seat in the European Commission.

Meanwhile, the prime minister of Spain, Jos\'e Luis Zapatero, said that Mr.~Moratinos was in the
running for the foreign policy role, Reuters reported.

That was seen by some diplomats as a warning to Britain to make clear whether it wanted the job, if
it fails to secure Mr.~Blair the presidential post. ``The post is Britain's if Britain wants it,''
said one official speaking on condition of anonymity, ``but if it is unable to field somebody it
will go elsewhere.''

\section{Younger Buyers Challenge Luxury Retailers in Asia}

\lettrine{W}{hen}\mycalendar{Nov.'09}{18} Ermenegildo Zegna opened a 7,300 square-foot shop in Hong
Kong last month, it did so with the elegant fanfare and glamour that befits the 99-year-old Italian
luxury men's wear company.

But the canap\'es and wine, the woollen suits and male mannequins were framed in a setting different
from the more traditional, sedate, salon-like Zegna establishments in Milan, Paris or New York: The
new 678-square-meter store boasts spacious rooms; a huge, shiny facade that is brilliantly lit at
night; a sleek, modern design; and a highly visible location in Tsim Sha Tsui, one of Hong Kong's
main shopping districts.

It is, after all, pitched toward a luxury shopping crowd whose tastes and expectations, and even
average age, are different from those in Europe or America.

``Asia has an incredible thirst for fashion and quality; the region is very important for us,''
Gildo Zegna, the company's chief executive and the grandson of its founder, said during an interview
at what is one of a handful of the ``global stores'' that house all of Zegna's brands -- formal,
casual and sportswear. ``Greater China has been our fastest-growing market for the past three
years.''

But the trick to succeeding here, Mr.~Zegna said, is to know your local customers -- and adapt to
them.

A study published by the consulting firm Bain \& Co.~last month highlighted the importance of
emerging Asian nations, with their fast growth and increasingly affluent populations, for luxury
retailers who face falling revenue in other parts of the world.

The global luxury market is expected to shrink 8 percent this year, to \texteuro153 billion, or
\$229 billion, as sales in the Americas plummet 16 percent, in Europe 8 percent and in Japan 10
percent, according to the Bain report.

In China, by contrast, luxury sales are forecast to grow 12 percent, to \texteuro6.6 billion, this
year, lifting the pace of growth in Asia over all to 10 percent.

But starting up or expanding in Asia is not merely a matter of replicating the tried-and-tested
models used back home -- not least because the sort of shopper who can afford top-of-the-market
Zegna suits -- or other luxury items -- tends to be younger than elsewhere.

This is especially noticeable in mainland China.

The ranks of the wealthy there are mostly composed of first-generation entrepreneurs, said Rupert
Hoogewerf of Hurun Report, a publishing house in Shanghai that compiles information on China's
millionaires and billionaires.

Typically, a Chinese individual worth \$150 million or more is about 50 years old -- about 15 years
younger than someone in that category in Britain or the United States, Mr.~Hoogewerf said.

The average age of someone with 100 million yuan, or about \$15 million, is 43. The approximately
825,000 Chinese with personal wealth of 10 million yuan are on average as young as 39, according to
Hurun's data -- again, about 15 years younger than their counterparts in America or Europe.

``Much of this wealth has only been created since the 1980s -- in other words, a solid generation
later than in Hong Kong or Taiwan. You'd have to look back to the late nineteenth century in the
United States or to the industrial revolution in Britain, to find anything comparable to the wave of
entrepreneurs who are now starting up in China,'' Mr.~Hoogewerf said.

``Many luxury brands are having quite a bit of trouble with that,'' he added. ``A lot of Western
brands are trying to apply the same model in China as they have back home, and are thus potentially
targeting too old a group, when really they need to be more youthful and dynamic to attract these
sorts of people.''

In addition to the age issue, luxury retailers have to deal with regional cultural complexities,
Claudia D'Arpizio, the author of the Bain study, said by telephone from Milan.

Even within China, for instance, there are differences between the cities on the coast and towns in
the sprawling interior.

Those who shop for luxury apparel in coastal cities like Shanghai, for instance, prefer
Western-style formal wear for business occasions. Those in the country's interior tend to dress down
and wear more casual clothes, according to Ms.~D'Arpizio.

``Companies need to have a completely different merchandising plan for different parts of the
country,'' she said.

For Ermenegildo Zegna, China, or indeed the rest of Asia, is not uncharted waters. The company was
among the first luxury retailers to open a shop in mainland China, in 1991 in Beijing.

Of the \texteuro870.6 million in revenue Zegna generated last year, 88 percent came from abroad and
36 percent of total revenue came from the Asia-Pacific region.

The company now has 73 outlets in the China region, including 10 in Hong Kong, and one in Macao, the
former Portuguese colony that is now China's gambling hub. (A second Macao outlet is to open in
December.) Zegna has more than 550 outlets in 86 countries.

Last month also saw Zegna's debut in Mongolia and two store openings in Singapore. Zegna has three
shops in India and aims to open more in the next few years.

``You have to constantly fine-tune,'' Mr.~Zegna said. ``We've become much more scientific and
analytical. We seek constant feedback from the customer, and monitor how particular items are doing
-- we take that into account in our store planning, we adjust our marketing efforts accordingly, the
look of each store, the product mix.''

In China, for instance, Mr.~Zegna and his team learned that shoppers like stores to be ``more
glamorous, with bigger facades,'' he said, with more space and very high service levels.

Taking account of the lower age of its customers, much of the space in the new Hong Kong store is
devoted to the sports apparel and ``upper casual'' ranges that appeal to younger shoppers.

Introduced only three years ago, the upper casual range accounts for about 30 percent of revenue in
Greater China -- two to three times as much as in Zegna's traditional markets.

Mr.~Zegna's appraisal: ``Thank God we had that!''

Zegna's marketing in Asia is less geared toward magazine advertising and more toward in-store
educational events for customers.

Zegna realized that Asian shoppers like in-depth information about the craftsmanship that goes into
its suits and top-quality textiles.

And so it arranges for its craftsmen to tour the region, giving talks about the processes of making
ties or shoes. Store staff members also have to be more versed in such information and are trained
accordingly.

The global economic turmoil has caused double-digit percentage drops in Ermenegildo Zegna's sales in
countries like Japan, Russia and the United States, leading it to reduce capital and media spending
this year in all areas except Asia, which now has a larger portion of such expenditures than before.
Sales in Greater China are booming -- up about 30 percent this year.

``China can't quite compensate for the double-digit drop-offs we've seen elsewhere,'' Mr.~Zegna
said. ``But at least it means that, over all, our revenue decline will be in the single digits, and
not in double digits.''

\section{Library in a Pocket}

\lettrine{W}{ith}\mycalendar{Nov.'09}{18} Amazon's Kindle, readers can squeeze hundreds of books
into a device that is smaller than most hardcovers. For some, that's not small enough.

Many people who want to read electronic books are discovering that they can do so on the smartphones
that are already in their pockets -- bringing a whole new meaning to ``phone book.'' And they like
that they can save the \$250 to \$350 that they would otherwise spend on yet another gadget.

``These e-readers that cost a lot of money only do one thing,'' said Keishon Tutt, a 37-year-old
pharmacist in Texas who buys 10 to 12 books a month to read on her iPhone, from Apple. ``I like to
have a multifunctional device. I watch movies and listen to my songs.''

Over the last eight months, Amazon, Barnes \& Noble and a range of smaller companies have released
book-reading software for the iPhone and other mobile devices. One out of every five new
applications introduced for the iPhone last month was a book, according to Flurry, a research firm
that studies mobile trends.

All of that activity raises a question: Does the future of book reading lie in dedicated devices
like the Kindle, or in more versatile gadgets like mobile phones? So far, e-book software for phones
does not appear to have cut into demand for single-function e-readers. According to the Codex Group,
a consultant to the publishing industry, about 1.7 million people now own one, and that number could
rise to four million by the end of the holiday season.

But there are already 84 million smartphones that can run applications in the United States alone,
according to IDC, a research firm. Apple has sold more than 50 million iPhones and iPod Touches,
which both run e-book software.

Apple itself doesn't see the iPhone as the ultimate reading device. Next year, it is likely to
further stir up the e-book market if, as expected, it introduces a tablet computer -- a device
bigger than a phone that will most likely run e-reader software along with other programs intended
for the iPhone.

People once scoffed\cite{scoff} at the idea of reading a book on a 3.5-inch mobile screen. For many
readers, though, sheer convenience trumps everything else.

``The iPod Touch is always at hand,'' Shannon Stacey, who has written several romance e-novels,
said. ``It's my calendar, it's my everything, so my books are always with me.'' Ms.~Stacey, who also
owns an early Sony Reader model, said she had now bought twice as many e-books for her iPod Touch as
for her Sony.

While the Kindle, the Reader and the Nook, the Barnes \& Noble device that will reach the market
later this month, feature screens that use very little power and are close to the size of a page in
a trade paperback, they have comparatively limited features, like gray-and-white reading displays
and either partial Internet access or none at all.

Ian Freed, vice president for the Kindle division at Amazon, said customers still bought more books
to read on the Kindle than they did for its iPhone application, though he declined to disclose
figures. Amazon is working on e-reading software for the BlackBerry and for Macintosh computers; it
introduced software for Windows PCs last week.

``It's a surprisingly pleasant experience to read on a small screen,'' said Josh Koppel, a founder
of ScrollMotion, a New York company that has made some 25,000 e-books available through Apple's App
Store and has sold more than 200,000 copies.

Companies like ScrollMotion and BeamItDown sell books in the form of individual applications, so
novels like ``Twilight'' by Stephenie Meyer show up right in the App Store. Amazon and Barnes \&
Noble give away e-reading software instead; users buy the actual books through the browser on a
phone or PC.

Publishers are now rushing to develop new forms of books to cater to readers who will see them on
smartphones -- books that will not work on today's stand-alone e-readers. When Nick Cave, the rock
musician, wrote his second novel, ``The Death of Bunny Munro,'' he and his British publisher,
Canongate, worked with a multimedia company to develop an app for the iPhone that incorporated not
just the text but also videos, music composed by Mr.~Cave and audio of the author reading the book.

``What you can do with graphics and moving images creates a lot of possibilities for a publisher
that have never existed before,'' said Jamie Byng, Canongate's publisher.

Of course, e-readers like the Kindle and the Nook will evolve as well, most likely adding color to
their reading screens.

But in the meantime, Amazon executives say that the limitations of the Kindle actually make it more
attractive for reading.

``The Kindle is for people who love to read,'' Mr.~Freed of Amazon said. ``People use phones for
lots of things. Most often they use them to make phone calls. Second most often, they use them to
send text messages or e-mail. Way down on the list, there's reading.''

Indeed, Sarah Wendell, an administrative assistant in Manhattan who blogs about romance novels, said
that although she used the iPhone to read while on a coffee or lunch break, she still used her
Kindle during her one-hour commute from New Jersey.

For long reading sessions, she said, the iPhone is ``a small screen, and my eyes would start to
hurt, even though I crank the text up to grandma or great-grandma size.''

Travis Bryant, director of digital products for Keen Communications, a small publisher in
Birmingham, Ala., said he had gotten a surprising amount of reading done while waiting in lines.
Mr.~Bryant said he had recently read ``The Shack,'' the best-selling Christian allegorical novel, as
well as ``The Templar Legacy,'' a historical thriller by Steve Berry, on his iPhone.

But Mr.~Bryant acknowledged that the iPhone, while convenient, did not serve every reading purpose.
``I've got a 3-year-old at home, and he really digs books,'' Mr.~Bryant said. ``I remembering
pilfering my parents' shelves, and if everything is on the iPhone, he's just not going to have that
visual temptation. So we keep the shelves loaded.''

\section{A Case in Antiquities for 'Finders Keepers'}

\lettrine{Z}{ahi}\mycalendar{Nov.'09}{18} Hawass regards the Rosetta Stone, like so much else, as
stolen property languishing in exile. ``We own that stone,'' he told Al Jazeera, speaking as the
secretary general of Egypt's Supreme Council of Antiquities.

The British Museum does not agree -- at least not yet. But never underestimate Dr.~Hawass when it
comes to this sort of custody dispute. He has prevailed so often in getting pieces returned to what
he calls their ``motherland'' that museum curators are scrambling to appease him.

Last month, after Dr.~Hawass suspended the Louvre's excavation in Egypt, the museum promptly
returned the ancient fresco fragments he sought. Then the Metropolitan Museum of Art made a
pre-emptive display of its ``appreciation'' and ``deep respect'' by buying a piece of a shrine from
a private collector so that it could be donated to Egypt.

Now an official from the Neues Museum in Berlin is headed to Egypt to discuss Dr.~Hawass's demand
for its star attraction, a bust of Nefertiti.

These gestures may make immediate pragmatic sense for museum curators worried about getting
excavation permits and avoiding legal problems. But is this trend ultimately good for archaeology?

Scientists and curators have generally supported the laws passed in recent decades giving countries
ownership of ancient ``cultural property'' discovered within their borders. But these laws rest on a
couple of highly debatable assumptions: that artifacts should remain in whatever country they were
found, and that the best way to protect archaeological sites is to restrict the international trade
in antiquities.

In some cases, it makes aesthetic or archaeological sense to keep artifacts grouped together where
they were found, but it can also be risky to leave everything in one place, particularly if the
country is in turmoil or can't afford to excavate or guard all its treasures. After the Metropolitan
Museum was pressured to hand over a collection called the Lydian Hoard, one of the most valuable
pieces was stolen several years ago from its new home in Turkey.

Restricting the export of artifacts hasn't ended their theft and looting any more than the war on
drugs has ended narcotics smuggling. Instead, the restrictions promote the black market and
discourage the kind of open research that would benefit everyone except criminals.

Legitimate dealers, museums and private collectors have a financial incentive to pay for expert
excavation and analysis of artifacts, because that kind of documentation makes the objects more
valuable. A nation could maintain a public registry of discoveries and require collectors to give
scholars access to the artifacts, but that can be accomplished without making everything the
property of the national government.

The timing of Dr.~Hawass's current offensive, as my colleague Michael Kimmelman reported, makes it
look like retribution against the Westerners who helped prevent an Egyptian from becoming the leader
of Unesco, the United Nation's cultural agency. But whatever the particular motivation, there is no
doubt that the cultural-property laws have turned archeological discoveries into political weapons.

In his book ``Who Owns Antiquity?'', James Cuno argues that scholars have betrayed their principles
by acquiescing to politicians who have exploited antiquities to legitimize themselves and their
governments. Saddam Hussein was the most blatant, turning Iraqi archeology museums into propaganda
for himself as the modern Nebuchadnezzar, but other leaders have been just as cynical in using
antiquities to bolster their claims of sovereignty.

Dr.~Cuno advocates the revival of partage, the traditional system in which archeologists digging in
foreign countries would give some of their discoveries to the host country and take others home.
That way both sides benefit, and both sides have incentives to recover antiquities before looters
beat them to it. (To debate this idea, go to nytimes.com/tierneylab.)

As the director of the Art Institute of Chicago, Dr.~Cuno has his own obvious motives for acquiring
foreign antiquities, and he makes no apology for wanting to display Middle Eastern statues to
Midwesterners.

``It is in the nature of our species to connect and exchange,'' Dr.~Cuno writes. ``And the result is
a common culture in which we all have a stake. It is not, and can never be, the property of one
modern nation or another.''

Some of the most culturally protectionist nations today, like Egypt, Italy and Turkey, are trying to
hoard treasures that couldn't have been created without the inspiration provided by imported works
of art. (Imagine the Renaissance without the influence of ``looted'' Greek antiquities.) And the
current political rulers of those countries often have little in common culturally with the creators
of the artifacts they claim to own.

Dr.~Hawass may consider the Rosetta Stone to be the property of his government agency, but the
modern state of Egypt didn't even exist when it was discovered in 1799 (much less when it was
inscribed in 196 B.C., during the Hellenistic era). The land was under the rule of the Ottoman
Empire, and the local historians were most interested in studying their Islamic heritage.

The inscribed stone fragment, which had been used as construction material at a fort, didn't acquire
any significance until it was noticed by Napoleon's soldiers and examined by the scholars on the
expedition.

When the French lost the war, they made a copy of the inscriptions before surrendering the stone to
the English victors, who returned it to the British Museum. Eventually, two scholars, working
separately in Britain and in France, deciphered the hieroglyphics.

This all happened, of course, long before today's nationalistic retention laws and the United
Nations' Convention on the Means of Prohibiting and Preventing the Illicit Import, Export and
Transfer of Ownership of Cultural Property. But what if the Rosetta Stone were unearthed in modern
times?

Were the Rosetta Stone to appear on the art market without the proper export permits and documented
provenance, Dr.~Cuno says, a museum curator who acquired it would risk international censure and
possible criminal charges. Scholars would shun it because policies at the leading archeological
journals would forbid the publication of its text.

``Not being acquired or published, the Rosetta Stone would be a mere curiosity,'' Dr.~Cuno writes.
``Egyptology as we know it would not exist, and modern Egyptians would not know to claim it as
theirs.''

The Supreme Council of Antiquities wouldn't even know what it was missing.

\section{3 Democrats Could Block Health Bill in Senate}

\lettrine{S}{enator}\mycalendar{Nov.'09}{18} Ben Nelson, Democrat of Nebraska, says he is not sure
he is ready to help a Democratic health care proposal clear even the most preliminary hurdle:
gaining the 60 votes his party's leaders need to open debate on the measure later this week.

Two of his fellow Democrats, Senators Mary L.~Landrieu of Louisiana and Blanche Lincoln of Arkansas,
are proving tough sells as well, raising the prospect that one or perhaps all three of them could
scuttle the bill before the fight over it even begins on the Senate floor.

``I think what is most important for me is to take a look at what is presented on behalf of
Arkansans and figure out whether it is something that really makes sense,'' Mrs.~Lincoln said
Tuesday. ``I am responsible to the people of Arkansas, and that is where I will take my direction.''

Typically routine, the procedural approval needed to begin consideration of a bill looms as anything
but routine in this instance. Instead, the vote is fast becoming a test of the leadership abilities
of Senator Harry Reid of Nevada, the majority leader. It will also decide at least the near-term
prospects of President Obama's top domestic priority. And it is providing a case study of the
Democrats' difficulties in managing high expectations fueled by large Congressional majorities.

Inability by the Democrats to advance their emerging plan could require them to regroup and redraw
the measure or even switch to a more contentious procedural shortcut around the need for a 60-vote
majority.

Given what is riding on the vote, party leaders have been busy talking to holdouts, negotiating
deals in an effort to get them on board as wavering lawmakers exert their leverage at a critical
moment. In response to demands from Mr.~Nelson, for instance, the leaders appear willing to drop
plans to use the bill to strip health insurance companies of their antitrust exemption.

Though he lacks 60 firm commitments, Mr.~Reid said Tuesday that he remained cautiously optimistic
that he could get the Senate's 58 Democrats and 2 independents to vote to thwart a filibuster on
what is known as a motion to proceed, the initial step in any debate.

But the fact that the outcome remains in doubt is frustrating to some lawmakers, given the sentiment
in favor of moving forward among an overwhelming number of Senate Democrats. Others say they have
their own reservations about the bill but are eager to have the floor debate as a way to seek
changes in the measure.

``We've all got problems with the bill,'' said Senator Bob Casey, Democrat of Pennsylvania. ``I
would hope we could all support the effort to get the bill to the floor.''

Some senators who were previously considered potential obstacles to starting the debate have made
clear that they will vote with the party at this point, drawing a distinction between this initial
advance and later votes of more consequence.

``I think we ought to begin a dialogue on this and see where it leads,'' said Senator Evan Bayh,
Democrat of Indiana, adding that this position was no guarantee he would support the measure in the
end.

On Tuesday, Republicans demonstrated the risks some Democrats face in agreeing to even open debate.
The Republican National Committee attacked Senator Mark Pryor of Arkansas, a state that went heavily
Republican in last year's presidential election, for agreeing to side with fellow Democrats on the
initial vote, equating it with a vote for the measure itself.

There is no Republican on whom Mr.~Reid can count for support in the procedural vote. ``We're going
to do everything we can to defeat this monstrosity and actually get to a meaningful debate about how
we genuinely have health care reform in this country,'' Senator John Thune of South Dakota, a member
of the Republican leadership, said Tuesday.

Senators Landrieu, Lincoln and Nelson have all said they cannot commit to backing the preliminary
step until they see the final legislation, which is being written by Mr.~Reid and his lieutenants
and could be unveiled as early as Wednesday.

``I'm anxious to see the bill,'' Mr.~Nelson said.

The three lawmakers have all been skeptical of a public health insurance option, and Mr.~Reid's
proposal includes a public plan, though it would give states the opportunity to opt out of it.

All three also represent states won handily last year by Senator John McCain, the Republican
presidential nominee, and one of them, Mrs.~Lincoln, is up for re-election next year. She must run
while under fire from both left and right in Arkansas, where Mr.~Obama had one of his worst
showings.

Worried about getting the votes needed to open debate, some Democratic senators are already looking
for alternatives to the government-run plan proposed by Mr.~Reid.

Senator Thomas R.~Carper, Democrat of Delaware, said he was trying to devise such an alternative to
meet ``centrist concerns about the public option.'' Over and over, Mr.~Carper said, the centrists
have made clear that they do not want to create an insurance plan that is ``government-run or
government-funded.''

While acknowledging the reservations of the three Democrats, senior party officials said they
expected that none of the three would want to be responsible for stalling such a major party
initiative in a procedural vote. The leaders predicted that Senators Landrieu, Lincoln and Nelson
would grudgingly join ranks.

Mr.~Reid said that once his colleagues saw his final measure, the decision would be made easier.

``I am going to have a bill that's fiscally responsible,'' he told reporters. ``Of all the bills
we've seen, it'll be the best: saves more money, is more protective of Medicare, is a bill that's
good for the American people.''

\section{Medical Schools Quizzed on Ghostwriting}

\lettrine{S}{enator}\mycalendar{Nov.'09}{18} Charles E.~Grassley wrote to 10 top medical schools
Tuesday to ask what they are doing about professors who put their names on ghostwritten articles in
medical journals -- and why that practice was any different from plagiarism by students.

Mr.~Grassley, of Iowa, the ranking Republican on the Senate Finance Committee, sent the letters as
part of his continuing investigation of so-called medical ghostwriting. The term refers to
publication of medical journal articles in which an outside writer -- sometimes paid by a drug or
medical devices company whose product is being studied -- has done extensive work on the article
without being named on the publication. Instead, one or more academic researchers may receive author
credit.

Mr.~Grassley said ghostwriting had hurt patients and raised costs for taxpayers because it used
prestigious academic names to promote medical products and treatments that might be expensive or
less effective than viable alternatives.

``Any attempt to manipulate the scientific literature, which can in turn mislead doctors to
prescribe treatments that may be ineffective and/or cause harm to their patients, is very
troubling,'' the senator wrote.

Some journals, medical associations, writers' and editors' groups and pharmaceutical companies
themselves have called for crackdowns on ghostwriting. But some universities that employ the
professors who put their names on the articles have been slow to respond. Merck, Wyeth (now part of
Pfizer), GlaxoSmithKline and AstraZeneca are among the companies accused by lawyers and
investigators of providing ghostwriters for research papers.

Mr.~Grassley asked the universities to describe their policies on both ghostwriting and plagiarism
and to enumerate complaints and describe investigations into both practices since 2004.

Dr.~Ross McKinney Jr., director of the Trent Center for Bioethics at Duke University, said faculty
who took credit for a ghostwritten paper should suffer the same penalties as students who
plagiarized.

``But it is a very, very difficult thing to prove, just as it turns out that plagiarism is hard to
prove,'' he said in an interview.

Mr.~Grassley's letters went to the top medical schools for research as ranked by U.S.~News and World
Report this year, in order: Harvard, Johns Hopkins, the University of Pennsylvania, Washington
University in St.~Louis, University of California, San Francisco, Duke, Stanford, the University of
Washington, Yale and Columbia.

Most of them already have policies against ghostwriting or honorary authorship of research papers, a
review of their Web sites shows.

Harvard Medical calls the practices ``deplorable.'' Duke says, ``Severe and/or repeated offenses
will result in formal disciplinary action.''

Arthur L.~Caplan, director of the Penn's Center for Bioethics, said there was a difference in
degree, if not in kind, between ghostwriting and plagiarism. Faculty members who sign their names to
ghostwritten papers for research credit usually have some agreement with the paper, he said, even
if, improperly, they did not write it. Students who plagiarize a paper may know nothing about the
subject.

``Ghostwriting and plagiarism, they're on a continuum,'' Mr.~Caplan said. ``They're related. I
wouldn't say they're twins, but they're cousins.''

Mr.~Grassley's letter highlighted the disparate treatment of students and professors who claimed
authorship of a paper that was not their own.

``Students are disciplined for not acknowledging that a paper they turned in was written by somebody
else,'' Mr.~Grassley wrote. ``But what happens when researchers at the same university publish
medical studies without acknowledging that they were written by somebody else?''

The medical schools were asked to answer the questions by Dec.~8.

\section{Obama's Pacific Trip Encounters Rough Waters}

\lettrine{S}{eoul}\mycalendar{Nov.'09}{19}, South Korea -- For all of President Obama's laying claim
to the title of ``America's first Pacific president,'' Asia was always going to be a tough nut for
him to crack.

Without the first lady at his side, he would not have the kind of round-the-clock coverage the first
couple got during their inaugural tour of Europe.

Without a popular gesture like elevating the plight of the Palestinian people to equal status of the
Israelis, he would not be showered with the kind of praise he got for his speech to the Muslim world
in Cairo.

And without a stop in Indonesia, his boyhood home, he would not bask in the kind of adulation he
received in Accra, Ghana.

Instead, with the novelty of a visit as America's first black president having given way to the
reality of having to plow through intractable issues like monetary policy (China), trade (Singapore,
China, South Korea), security (Japan) and the 800-pound gorilla on the continent (China),
Mr.~Obama's Asia trip has been, in many ways, a long, uphill slog.

So it is no wonder that on the last day of the toughest part of his trip -- the China part --
Mr.~Obama took a hike: a brisk, bracing 30-minute climb up the Great Wall. Around 3:30 Wednesday
afternoon, Mr.~Obama's mile-long motorcade arrived at the Badaling section of the Great Wall, which
snakes over jagged, rocky mountains.

Visitors to that touristy section of the wall generally encounter a cacophonous melee of vendors,
but on this day, the place was like a ghost town, courtesy of the Chinese authorities who had shut
it down. (The same thing happened Tuesday when Mr.~Obama sped through an empty-but-for-his-entourage
Forbidden City.)

Even the two sightseeing trips did not offer a total respite, however, as they were prominent,
well-publicized examples of what Mr.~Obama did not do in China. He steered clear of public meetings
with Chinese liberals, free press advocates and even average Chinese, with his aides citing
scheduling conflicts. Mr.~Obama did, though, give an interview on Wednesday morning to Southern
Weekly, one of China's most popular newspapers, sometimes known for poking the authorities by
breaking news on delicate subjects.

Still, for an American president who has tried to make openness a hallmark of his public persona, it
was a departure, made more stark since Chinese authorities largely hijacked Mr.~Obama's one other
attempt at a give and take with Chinese students, a town hall meeting in Shanghai, by stuffing the
auditorium with young Communist Party aspirants.

A week ago, when Mr.~Obama kicked off his trip in Japan, things were not so grim. Tokyo welcomed him
as much as a celebrity as a world leader, with cries of ``Obama-san!'' from the people who gathered
in the rain to watch his motorcade pass. Local newspapers gushed about how he told his Japanese
hosts that he wanted to eat tuna and Kobe beef. Even the ballyhoo from right-wing bloggers back at
home over Mr.~Obama's deep bow to Emperor Akihito did not seem to dent Mr.~Obama's image in Japan;
his aides said he was unfazed by the criticism.

But Mr.~Obama quickly discovered that popularity on the Asian streets did not necessarily translate
into policy successes behind closed doors in the Kantei, the Japanese White House, let alone in the
Great Hall of the People in Beijing.

Political analysts in Japan gave Mr.~Obama high marks for what was one of his principal goals:
improving communication with Japan's outspoken new leaders.

But the trip managed only to paper over some of the recent differences between the sides, like the
contentious issue of the relocation of an unpopular Marine air base in Futenma, on the southern
island of Okinawa. Mr.~Obama and Prime Minister Yukio Hatoyama could not solve that issue, instead
merely deferring a tough decision by agreeing to form a working group to look at the relocation
problem.

One former Japanese diplomat praised the president for showing patience and avoiding mishaps that
would have further tarnished the relationship. The former diplomat, Kunihiko Miyake, who now teaches
international affairs at Ritsumeikan University in Kyoto, said the United States and Japan still did
not see eye to eye on their single biggest bilateral issue: how to make their cold-war-era alliance
relevant in a region where the balance of power had been upset by China's rise.

``The two countries are in the same bed, but dreaming different dreams,'' Mr.~Miyake said. ``The
Americans want the alliance to be stronger, but the Japanese seem to want to do less.''

Mr.~Obama's next stop was Singapore for the Asia-Pacific Economic Cooperation conference, best known
for its quaint custom of making all the leaders wear the same style of colorful shirt, helpfully
supplied by the host country. Mr.~Obama, in blue, wore a brave grin in the group photo, flanked by
the red-shirted Singaporean prime minister and an identical blue-shirted Indonesian president.

This year, APEC made headlines, though not the sort Mr.~Obama might have liked. With a deadline
looming for a big climate change conference in Copenhagen, the leaders convened a hastily called
breakfast meeting to acknowledge that they would not be able to resolve entrenched differences in
time.

And then, Mr.~Obama departed for China, where the authorities stage-managed and restricted access to
his town hall meeting in Shanghai. He did offer a nuanced, oblique critique of China's rigid
controls and restrictions of the Internet and free speech without mentioning, let alone condemning,
China's government.

Mr.~Obama and President Hu Jintao presented their two days of talks as substantive, even though they
did not appear to make much progress on issues like Iran, China's currency or human rights. Robert
Gibbs, the White House spokesman, took the unusual step of sending a statement to reporters --
something he did not do for either stop in Japan or Singapore -- saying the China trip went well.

In Seoul, where Mr.~Obama ends his trip, he will have perhaps his easiest leg. South Korea is a
longtime ally that has been cooperating with the United States on vital issues like North Korea and
does not appear to have any big ax to grind with the United States.

\section{Questioning a Korean Wedding Tradition}

\lettrine{W}{hen}\mycalendar{Nov.'09}{19} a daughter of Kim Jong-chang, South Korea's top financial
regulator, got married last June, Mr.~Kim did something unusual: He eliminated the cashier and the
cash-filled envelopes.

These are fixtures of a South Korean wedding, as much so as the wedding officiant. Before entering
the wedding hall, guests line up in front of the cashier's table to hand over an envelope stuffed
with cash. The cashier opens the envelope and registers the guest's name, and the amount given, in a
velvet-covered ledger -- often while the guest is still standing there.

``The problem with this tradition is that it can be abused for bribery,'' said Mr.~Kim, governor of
the Financial Supervisory Service, which regulates the South Korean banking and securities
industries. ``In my case, many banking officials would have shown up with cash gifts. They would
have wondered whether I was annoyed that they didn't put enough in the envelope.''

Chipping in to help friends defray wedding or funeral expenses is an old custom here. But in recent
months, it has been criticized as wasteful, and sometimes even as a conduit for vote-buying and
bribery.

In May, after some critical news stories about extravagant weddings being held at five-star hotels
during the economic downturn, President Lee Myung-bak exhorted South Korea's rich and powerful to
set an example in fighting the ``vain and extravagant'' wedding culture.

Mr.~Kim is one of a small but growing number of people, from ordinary families to dignitaries, who
are joining this campaign, refusing to accept cash gifts and keeping their guest lists relatively
short. Ban Ki-moon, the South Korea-born secretary general of the United Nations, invited only a few
close friends and relatives to the wedding of his son in May, as did Foreign Minister Yu Myung-hwan
when his daughter married in April. In October, Chung Jung-kil, Mr.~Lee's chief of staff, followed
suit.

Still, these low-key weddings were considered such oddities that they made the news.

In South Korea, where ``face'' is famously cherished, the measure of a family's social standing is
seen in the number of guests at weddings, as well as the amount of money given and the sumptuousness
of the banquet. At funerals, the number of wreaths presented by friends, business associates and
local politicians is a comparable social metric.

``Here, a wedding is less a celebration than an occasion for a family to show off,'' said Lee
Yoon-ji, who runs a wedding management agency and photo studio in Seoul's upscale Kangnam district.
``For instance, if the bride's family finds its guests are much fewer than the groom's, it's
humiliating.''

Some families send out thousands of wedding invitations. A bank account number is sometimes included
so people who can't attend can still send money.

Often, the decision of whether to attend is based on whether the couple, or their relatives,
attended weddings or funerals in one's own family -- or might be expected to. Families keep records
of how much they receive and from whom so that they can reciprocate. Failure to do so can ruin a
friendship.

``Sometimes you even get invitations from people you don't know very well,'' Mr.~Kim said. ``They
arrive like tax bills or I.O.U.'s.''

Every year, the roughly 330,000 South Korean couples who get married spend an average of 15 million
to 20 million won, or \$13,000 to \$17,000, in wedding expenses, said Lee Woong-jin, head of Sunoo,
a matchmaking company that conducts an annual survey on wedding expenses. The cost can exceed 50
million won for hotel weddings.

Much of that is covered by the cash gifts. Last year, South Koreans gave out 8 trillion won, or
524,500 won for each household, in cash gifts for weddings and funerals, according to the National
Statistical Office.

``This is a 'you-help-me, I-help-you' tradition. I don't see anything wrong with it. You chip in and
you get help in return,'' said Han Seung-ho, 33, a photographer whose wedding in October attracted
370 guests. ``Without their cash gifts, my wedding would have been a serious financial burden for
me.''

But these envelopes also reflect a culture in which giving cash is considered so natural that people
sometimes call it a ``greeting'' -- and, in some cases, use it as a cover for bribery. When South
Korea's election laws were revised in 2004, they banned politicians from giving cash envelopes,
except at the weddings and funerals of close relatives.

Three candidates running for election at provincial farmers' and fisheries' cooperatives were
indicted in September and October on charges of giving cash gifts at voters' weddings. A provincial
education chief was widely criticized in the media in April after he reportedly invited 2,000 people
-- including the principals of all 460 schools under his jurisdiction -- to his son's wedding.

Chung Woo-jin, 50, president of Q\&Q Medi, a medical supplies company, said many wedding guests show
up ``reluctantly,'' fearing they might lose out on business contracts or promotions if they don't.
``So they show up to prove that they were there, give the envelope and hurry off to have the meal,
without even taking a look at the bride or groom,'' he said.

Mr.~Chung refused to accept cash envelopes at his mother's funeral in June. But he said he still
felt compelled to attend 40 to 50 weddings or funerals a year for friends, employees and business
acquaintances, each time donating an average of 100,000 won.

Meanwhile, some younger couples are rebelling against what they call a ``commercial'' wedding
culture controlled by parents. It is generally the parents who send out invitations, collect the
cash and pay for the wedding, and by and large, more guests are there for the parents than for the
couple getting married.

``Some of my friends feel frustrated, wondering if their wedding is for them or for their parents,''
said Lee Eun-jeong, 35, who works at a publishing company in Seoul. She limited her wedding in June
to 135 guests and did not accept envelopes. ``We also hate it when a friend who hasn't contacted us
for years suddenly gets in touch with us before her wedding, obviously with our envelopes in mind,''
she said.

South Korea has seen campaigns for wedding frugality before. In 1973, the late military strongman
Park Chung-hee tried to ban written invitations, flowers and gifts from weddings and funerals, in
the belief that such customs were wasteful and detracted from his campaign to build and modernize
the economy.

But enforcement was sporadic at best, and experts say weddings grew more extravagant after 1999,
when the restrictions were lifted and five-star hotels and wedding agencies entered the market.

Mr.~Kim, the financial watchdog chief, predicted that it would be some time before the cash envelope
tradition faded.

``Frankly, I found myself thinking, 'I've given out all these envelopes over the years. Why
shouldn't I get them once for my daughter's wedding?''' he said. ``It's not always easy in our
weddings to tell the difference between bribes and genuine gifts.''

\section{No 'System,' but Music Housewide}

\lettrine{``M}{ultiroom}\mycalendar{Nov.'09}{19} sound system? There's an app for that.''

That's not Sonos's slogan for its new ZonePlayer S5 wireless speaker consoles, but it could be. Now,
in addition to tracking your weight, calculating recipes and identifying bird calls, your iPhone or
iPod Touch can serve as a touch-screen remote control for the music playing in different rooms of
your home.

It's a pretty giddy feeling. You stroll magisterially through your dwelling, browsing your music
collection on your phone. You can start playing some Broadway in the bedroom, Latin in the living
room and rock in the rec room. All with a few taps on the screen of your iPod/iPhone.

You can adjust the volume all at once or room by room. If you like, you can sync everything up so
that the same music is playing in every room.

In short, you can pretend to be a well-wired rich person -- without having had to pay \$10,000 for
somebody to install a ``system.''

This, of course, has always been the Sonos specialty: multiroom stereo system components that
communicate wirelessly, operate effortlessly and require very little technical ability.

In the past, though, Sonos provided only electronic receiver boxes (\$350 apiece); you had to supply
your own powered speakers to plug into them. You still undercut the price of a professional sound
system, but for a wireless system, you wound up with a lot of wires.

The S5 (\$400 per unit) is different; it has both the wireless receiver and the speakers built into
a single box. Five speakers, actually: two tweeters, two midrange speakers and a woofer for bass,
all built into a 14 x 8 x 4-inch white cabinet that looks a lot like one of those white Bose speaker
docks for the iPod.

This time, though, there's no socket for the iPod. Instead, like the speakerless Sonos boxes before
it, the S5 speaker is designed to play the music collection you keep on your Mac or PC.

But already, this system, famous for its simplicity, is starting to sound complicated. And it's true
that there are a bunch of different pieces to it:

1. Some music to play. For most people, this means a collection of music files. It might be your
iTunes collection on a Mac or PC, or even a bunch of MP3 files sitting in folders.

(If you have some of the older, copy-protected iTunes songs in your stash, they won't play. When
Apple eliminated copy protection from the iTunes Store, I'll bet nobody whooped louder than the
Sonos gang.)

So what happens if the computer is turned off or, worse, gone (because it's a laptop and you took it
with you -- hey, it could happen)? In that case, nobody else at home can play any music until you
return. It's nice to be missed, of course, but not after you've just spent \$400 on a music system.

In those situations, you can copy your music collection to a network-attached hard drive (or Apple
Time Capsule) either attached to the S5 or somewhere on the network.

There's also an audio input on the back of each S5 speaker (and a headphone jack). That's so you can
connect an iPod, CD player or some other directly connected audio source.

2. A home wireless network, complete with a router. Presumably you already have this part.

3. An S5 speaker unit, or more than one.

You're supposed to plug the first S5 directly into your router. Additional S5s don't have to be
connected to anything but a power outlet; that's the beauty of Sonos's ``mesh network.'' Each S5
repeats the wireless signal from the first one, rebroadcasting the music up to 100 feet at a time.

Ah, but what if the router is someplace awkward -- like a utility closet, where you presumably don't
listen to music very often?

In that case, you'll have to spend another \$100 for something called a ZoneBridge. It's a little
white box that connects to your router, wherever it is, and sends its signal to the S5s, wherever
they are.

4. Playback software. The only buttons on the S5 speaker console itself are volume and mute. So you
need a way to start and stop the music, choose which music you want (by song, band, genre, playlist
or whatever) and so on.

For this, you have three options. First, you can download Sonos's free iPhone or iPod Touch app,
which is beautifully designed and wicked fun to use. (My only criticism: Too many confirmations. If
I tap the Pause button, I don't need the software to ask, ``Do you really want to Pause?'')

Second, you can use the free Desktop Controller software on your Mac or Windows machine.

Third, you can buy Sonos's hand-held, touch-screen remote for \$350. (Hint: If you buy an iPod Touch
instead, you'll save \$130 and have a lot more fun after hours.)

All three of these playback tools can operate both S5s as well as Sonos's earlier wireless music
boxes.

Sonos sent me two S5s and a bridge for testing. My router isn't where I wanted the music, so I
connected the bridge thing to my router, and parked the S5 speakers in two different rooms.

To get started, the instruction card says to open the Controller software on the Mac or PC, and
then, while it's ``listening,'' press the volume and mute keys together on the S5. That's how you
introduce the first S5 to the computer.

It didn't work. The computer didn't find the S5.

I called the help line. A friendly technician told me to introduce the computer to the bridge first,
and then the S5's. That worked, but why on earth wasn't it in the instructions to begin with?

But never mind; I was flying. I was a multiroom maestro, starting and stopping music in two distant
rooms, all from the privacy of my breakfast nook.

Better yet, two taps let you tap into the Sonos's menu of 25,000 podcasts and Internet radio
stations (including Pandora and Last.fm). Christmas music, anyone? Kiddy tunes? Neo-grunge? It's all
there.

The Sonos can also play Internet music from paid services like Rhapsody, Napster and Sirius.

The music sounds fantastic. Obviously, there's not much sense of stereo-channel separation unless
you have a very skinny head. But holy cow, the bass, the distinct instruments, the clarity -- it's
all there. And with serious power. The higher volume settings are literally ear-splitting indoors.
One S5 could fill a very large backyard with sound, and probably a school gym, without distortion or
skipping.

This all sounds great, and it is great. But you hecklers in back are no doubt thinking: ``Well, duh!
Why not just buy a \$95 AirPort Express pocket Wi-Fi base station, connect speakers to it and then
control playback using Apple's free Remote app on your iPhone/iPod Touch?''

This is true. That's a wireless music system for a lot less money. There is, however, a caveat or
five: the price doesn't include speakers. That system doesn't work when the computer is off or
iTunes isn't running. It doesn't let you control the volume of each room. It doesn't let you pipe
different music to each room. It's not nearly as easy to grab by the back-panel handle and carry out
to the patio for a party. And the music sometimes drops out because it's using Wi-Fi instead of
Sonos's much more reliable, stutter-free music signal.

So no, the Sonos S5 may not be just what the doctor ordered for everyone, especially in recessionary
times. But \$400 is a reasonable price for the S5's size, sound, portability, control, refinement
and pure pleasure.

And yes, now there's an app for that.

\section{Retailers Extend Deals Beyond Black Friday}

\lettrine{A}{ttention}\mycalendar{Nov.'09}{27} shoppers: It might pay to just sleep in this Black
Friday.

The conventional wisdom is that the most stupendous bargains of the year are to be had on the Friday
after Thanksgiving. But the marketplace has become so packed on that crowded shopping day that some
retailers are shifting their strategy.

Deals on certain products are likely to be just as good, perhaps even better, in the days and weeks
after Friday. In this economy, retailers need to stand out -- and some of them are betting they can
do so by offering bargains later in the season. Also, while chains are not discounting as deeply as
last year, they know the primary way to get penny-pinching consumers to spend is to keep the deals
coming all season long.

Exactly which strategy retailers are pursuing this year differs not only among shopping chains but
among categories of merchandise. That means the best time to shop for the season could hinge on
which items are on your list.

``Black Friday is about cheap stuff at cheap prices,'' said Daniel de Grandpre, the editor in chief
of DealNews.com, which tracks such sales each year. ``That means that high-end stuff is not on sale
on Black Friday. It just isn't.''

That is not to say consumers who brave the nation's stores on Friday will not find deals on
flat-screen televisions and fluffy ear muffs. But the products on sale that day, particularly
electronics, generally are lower-end products without many extras, or they are older models on the
verge of being discontinued. That is, of course, a reason stores are able to offer them at low
prices.

``It looks like a real mixed bag of deals and duds,'' said Andrew Eisner, director of content for
Retrevo, a Web site that reviews consumer electronics and recommends where and when to buy them.

Many of the gadgets on sale this Friday will be outdated models, he said, like navigation devices
without speech capability, Blu-ray players without Internet connections and digital cameras without
face-recognition technology.

Manish Rathi, a co-founder of Retrevo, cited some ``over-the-hill'' products, like a Nikon CoolPix
digital camera being sold at Target for \$88, reduced from \$140.

Mr.~de Grandpre said luxury retailers tended to stay out of the Black Friday fray because they would
rather not associate with bargain-basement shopping. To participate in the nation's über-shopping
day in a way that is befitting their status, luxury chains do offer deals, but only in certain
popular holiday and seasonal categories, like coats and home d\'ecor. And they do not bother to open
at 5 a.m., a common opening time on the day after Thanksgiving for the lower-end retailers known as
big-box stores.

The luxury chain Saks, for instance, is offering 40 percent off already reduced merchandise, but not
its newest collections. The chain does not bring in merchandise specifically for the day after
Thanksgiving. And the doors open at the relatively late hour of 8 a.m.

``The key difference is we don't run a strategy of these key items that we buy thousands of and that
we mark down to these low, low amounts,'' said Kimberly Grabel, senior vice president for marketing
at Saks. ``That is the big-box mentality.''

Stores have greatly reduced their inventories since last year, when the economic downturn forced
them into panic selling. So while there will be sales, as there are every year in any economy, it is
possible stores might run out of certain products or sizes. ``For the best selection you are going
to need to shop early this year,'' Ms.~Grabel said.

Her best advice? ``Stay home and shop online.''

Indeed, on Wednesday, many stores had begun holiday promotions on their Web sites. On Wednesday,
clothing was up to 40 percent off at Bloomingdale's, Saks, Neiman Marcus and Barneys. At Nordstrom,
merchandise was up to 50 percent off. At Lord \& Taylor, some items, like women's coats, were half
off.

Saks planned to offer its early deals online at midnight Wednesday. ``We figure some people will
need a break from their family on Thursday,'' Ms.~Grabel said.

To distinguish themselves from discounters, the likes of Saks and Neiman Marcus also offer gift
cards for spending a certain amount of money within a designated time period on Black Friday, which
got its name because it was thought to be the day when retailers often shifted into the black, or
became profitable. In the coming months, the chains will probably continue borrowing the idea of
private, timed online sales.

Made popular by members-only shopping Web sites like Gilt, Rue La La and HauteLook, the sales give
consumers a now-or-never incentive to buy, and are a less conspicuous way for stores to sell
designer merchandise at a discount. Consumers must sign up for e-mail alerts on retailers' Web sites
to be in the know.

In general, chain stores will not divulge their promotion schedules for the rest of the shopping
season, but some Web sites also track sales throughout the holiday season, including Deal- News.com,
BlackFriday.info, and Bfads.net.

Expect some notable deals online on Monday, now called Cyber Monday, when consumers return to their
offices and go holiday shopping on the Web.

This year, more e-commerce sites plan to offer bargains to win over reluctant consumers. For
instance, Blue Nile, the online jeweler, has long avoided participating in promotional days. But no
more.

``Even the wealthy have become more value-conscious,'' said John Baird, a spokesman for Blue Nile,
noting a shift in the behavior of its more affluent customers.

Blue Nile plans to offer two promotions on Monday -- free overnight shipping on orders placed by 6
p.m. Eastern time, and discounts on a different item each day through Dec.~23. (On one of the days,
a 15-carat diamond eternity bracelet will cost \$27,000, down from \$36,500. Blue Nile would not say
which day the bracelet would be on sale.)

``We want to be very careful when we offer special promotions like this, to make sure they are
brand-appropriate,'' Mr.~Baird said. ``We're not going to be like the mall jeweler, so you're not
going to see 50 percent off the entire store.''

Across-the-board sales are the specialty of the nation's biggest big-box stores and Web sites, and
they, too, began discounting before consumers roasted their turkeys.

Kmart had a ``Better Than Black Friday Sale.'' Best Buy offered ``Early Black Friday Prices'' on
some televisions, and on Nov. 11 the chain offered its lowest-ever advertised price on a laptop,
\$250. Toys ``R'' Us had early blockbuster deals known as doorbusters. And for weeks, Wal-Mart, the
nation's largest retailer, has been offering deep discounts.

Consumers who missed those bargains need not fret.

``Once you've started discounting, it's hard to stop midseason,'' said Maggie Taylor, a vice
president and senior credit officer for Moody's Investors Service.

Some chains are even getting a jump on Monday, re-imagining it as Cyber Sunday. J.C. Penney, for
instance, plans to offer its Monday sales beginning Sunday, for two days. Staples also plans to
offer its Web deals on both Sunday and Monday.

Consumers can track hourly specials on CyberMonday.com, which will continue posting news about deals
at the nation's major online retailers long after Monday.

\section{China Seen Facing Protectionist Backlash Next Year}

\lettrine{C}{hina}\mycalendar{Nov.'09}{27} faces a protectionist backlash next year because its
manufacturers are saddled with overcapacity and are disposing of its excess output on world markets,
the European Union Chamber of Commerce in China said Thursday.

Jörg Wuttke, the head of the business group, said it would take about 12 months to prepare a case
alleging dumping, the practice of selling goods for less than it costs to produce them.

``This lead time would indicate to me that in the second half of 2010, there will be far more
dumping cases against China, unfortunately,'' Mr.~Wuttke said.

He was speaking at the introduction of a study into industrial overcapacity in China, a longstanding
situation that the chamber believes has grown more serious as a result of the global financial
meltdown and Beijing's aggressive response to it.

``The crisis has throttled demand for exports from China at a time when even more investment, in the
form of the Chinese government's massive stimulus package, is being pumped into building new plants
and adding unnecessary capacity,'' the report said.

``As a result, the problem is actually getting worse in many industries,'' it said.

In a survey of the chamber's members, 56 percent of respondents identified local government policies
aimed at luring investment as the main macroeconomic reason for overcapacity; loose lending was the
second most frequently cited cause.

Mr.~Wuttke welcomed efforts by the central government to curb new capacity but said it was often
powerless in the face of local governments that craved new factories for the tax revenue and jobs
they could generate and that have done everything to keep existing plants from going under.

``Local protectionism kicks in,'' Mr.~Wuttke said. ``So even if Beijing sees a problem and wants to
tackle it, they are very often derailed by local politics.''

Apart from generating trade friction, rampant overcapacity would weigh on foreign direct investment
into China.

``Why would you invest if that market is already oversupplied?'' he asked.

By wasting resources and eroding profits, overcapacity deters research and development and
encourages companies to cut corners on health and safety standards as well as environmental
protection.

Further, the creation of unneeded capacity raises the risk of nonperforming loans for banks that
finance the investment. It also generates trade tension as producers sell their surplus production
overseas at cut-rate prices, the study said.

The State Council, China's cabinet, recently singled out the iron and steel, cement, electrolytic
aluminum, glass, coal, chemical, polysilicon and wind power equipment sectors as the worst offenders
when it came to overcapacity and announced steps to rein in their expansion.

The chamber applauded the cabinet's actions but said the fundamental answer was to shift from
investment-led and export-led growth and to focus more on domestic consumption and services.

The report made a series of recommendations that amounted to a root-and-branch overhaul of China's
economic model.

The recommendations included the redistribution of national income from companies to households. It
said state-owned enterprises should disgorge dividends for the government to spend on social
security, health and education instead of plowing profits back into fresh investment.

It also recommended a sharp reduction in corporate capital expenditure in coming years.

In addition, the report said that China should remove subsidies for energy and other inputs, which
are provided indirectly by households. Claiming that the manufacturing sector has become addicted to
these subsidies, it recommended increasing resource and environmental charges.

\section{The Biology Behind the Milk of Human Kindness}

\lettrine{A}{s}\mycalendar{Nov.'09}{27} the festival of mandatory gratitude looms into view, allow
me to offer a few suggestions on what, exactly, you should be thankful for.

Be thankful that, on at least one occasion, your mother did not fend off your father with a pair of
nunchucks, but instead allowed enough contact to facilitate your happy conception. Be thankful that
when you go to buy a pale, poultrylike entity, the grocery clerk will accept your credit card in
good faith and even return it with a heroic garble of your last name. Be grateful for the empathetic
employee working the United Airlines ticket counter the day after Thanksgiving, who understands why
you must leave town today, this very minute, lest someone pull out the family nunchucks.

Above all, be thankful for your brain's supply of oxytocin, the small, celebrated peptide hormone
that, by the looks of it, helps lubricate our every prosocial exchange, the thousands of acts of
kindness, kind-of kindness and not-as-nakedly-venal-as-I-could-have-been kindness that make human
society possible. Scientists have long known that the hormone plays essential physiological roles
during birth and lactation, and animal studies have shown that oxytocin can influence behavior too,
prompting voles to cuddle up with their mates, for example, or to clean and comfort their pups. Now
a raft of new research in humans suggests that oxytocin underlies the twin emotional pillars of
civilized life, our capacity to feel empathy and trust.

Reporting this month in The Proceedings of the National Academy of Sciences, researchers found that
genetic differences in people's responsiveness to the effects of oxytocin were linked to their
ability to read faces, infer the emotions of others, feel distress at others' hardship and even to
identify with characters in a novel or ``Doonesbury.'' ``I came into this research as a big
skeptic,'' said Sarina M.~Rodrigues of Oregon State University, an author of the new report, ``but
the results had me floored.''

Oxytocin may also be a capitalist tool. In a series of papers that appeared in Nature, Neuron and
elsewhere, Ernst Fehr, director of the Institute for Empirical Research in Economics at the
University of Zurich, and his colleagues showed that the hormone had a remarkable effect on the
willingness of people to trust strangers with their money. In the Nature study, 58 healthy male
students were given a single nasal squirt of either oxytocin or a placebo solution and, 50 minutes
later, were instructed to start playing rounds of the Trust Game with each other, using monetary
units they could either invest or withhold.

The researchers found that the oxytocin-enhanced subjects were significantly more likely than the
placebo players to trust their financial partners: whereas 45 percent of the oxytocin group agreed
to invest the maximum amount of money possible, just 21 percent of the control group proved so
amenable. Moreover, the researchers showed that the oxytocin boost didn't simply make subjects more
willing to take risks and throw their money around. When participants knew they were playing against
a computer rather than a human being, there was no difference in investment strategy between the
groups. Trust, it seems, is a strictly wetware affair.

Yet the hormone doesn't turn you into a sucker. In the Nov. 1 issue of Biological Psychiatry, Simone
Shamay-Tsoory of the University of Haifa and her colleagues reported that when participants in a
game of chance were pitted against a player they considered arrogant, a nasal spritz of oxytocin
augmented their feelings both of envy whenever the haughty one won and of schadenfreudian gloating
when their opponent lost.

As a rule, though, oxytocin is a joiner not a splitter. Analogues of the molecule are found in fish,
perhaps to help facilitate the delicate business of fertilization, by inhibiting a fish's natural
tendency to flee from other fish. The more elaborate grew the social demands, the more roles
oxytocin assumed, reaching its apotheosis in mammals. If you're going to give birth to a litter of
needy young, why not let the same signal that helped push those mewlers into the world give you tips
on their care and feeding? And if you're a human, bent on turning everything into an extended family
affair, here is oxytocin again to cheerlead and teleprompt. C.~Sue Carter of the University of
Illinois at Chicago, a pioneer in the study of oxytocin, suspects that the association between the
hormone and childbirth long kept scientists from taking it seriously. ``But now that it's been
brought into the world of economics and finance,'' Dr.~Carter said, ``suddenly it's very hot.''

Oxytocin acts as a hormone, traveling through the bloodstream to affect organs far from its origin
in the brain, and as a kind of neurotransmitter, allowing brain cells to communicate. Unlike most
neurotransmitters, oxytocin seems to deliver its signal through just one receptor, one protein
designed to recognize its shape and shudder accordingly when clasped; dopamine and serotonin, by
contrast, each have five or more receptors assigned to their care. Yet the precise contours of
oxytocin's hardworking receptor differ among individuals, to apparently noticeable effect.

In their new study, Dr.~Rodrigues and Laura R.~Saslow and Dacher Keltner of the University of
California, Berkeley, looked at how two variants in the genetic code for the receptor might
influence a person's capacity for empathy, as measured by a standard empathy questionnaire (``I
really get involved with the feelings of the characters in a novel'') and a behavioral task called
``Reading the Mind in the Eyes.'' In it, participants looked at 36 black-and-white photographs of
people's eyes and were asked to choose the word that best described each subject's mood. Uneasy,
defiant, contemplative, playful? In a related measure of oxytocin's presumed calming effects,
subjects were also tested for how strongly they reacted to the stress of hearing a series of loud
noises.

In their sample of 192 male and female college students, the researchers found that those carrying
the so-called A version of the oxytocin receptor, which previous reports had associated with autism
and poor parenting skills, scored significantly lower on the eye-reading task and higher on the
stress-prone test than did subjects with the G variant of the receptor.

``We're all different, and that's a good thing,'' Dr.~Rodrigues said. ``If everyone were gooey and
lovey-dovey, it would be an obnoxious world.'' As she drolly admitted, she herself is Type A.

\section{Obama Backs Senate on Health Bills' Disparities}

\lettrine{T}{wo}\mycalendar{Nov.'09}{27} of President Obama's senior health care advisers said
Wednesday that a proposed tax on high-cost insurance plans and a new commission to control Medicare
spending were among ``four pillars'' essential to major health care legislation.

Their remarks firmly aligned the White House with the Senate on two major disagreements facing
Democratic Congressional leaders trying to pass a bill.

The health care legislation that the House adopted on Nov. 7 did not include the excise tax on the
so-called Cadillac insurance plans. But the Senate majority leader, Harry Reid of Nevada, included
the levy in the bill that the Senate will begin debating in earnest next week, and many health
economists say it is vital to slowing long-term costs.

The House legislation also did not include the proposed Medicare commission, a central component of
the Senate measure.

The White House and Democratic lawmakers have been under pressure in recent weeks over the issue of
whether the legislation moving through Congress does enough to contain health care costs. Republican
critics, and even some Democrats, have said the legislation does not go far enough to slow medical
spending.

In a conference call with reporters Wednesday to highlight cost-containment provisions, the two
senior presidential advisers, Peter R.~Orszag, White House budget director, and Nancy-Ann DeParle,
director of health reform, cited a letter to Mr.~Obama by a group of respected economists who
strongly endorsed the Medicare commission and the tax on high-priced plans.

The letter also backed the two other pillars of the four embraced by the White House: that the
legislation not add to the federal budget deficit and that it promote changes, known as ``delivery
system reforms,'' aimed at rewarding high-quality care rather than high-quantity care.

``We are in favor of a fiscally responsible reform bill that includes these four pillars,''
Mr.~Orszag said in the conference call.

Mr.~Orszag declined to say whether the White House had told the House speaker, Representative Nancy
Pelosi of California, that to win Mr.~Obama's signature, any final bill must include the provisions
for a Cadillac tax and a Medicare commission.

Aides to Ms.~Pelosi pointed to comments she made last week indicating that she was prepared to work
out any differences between the two houses. ``I like what we see in their bill,'' she said of the
Senate measure, adding that the bills ``have a great deal in common'' and that when the Senate
adopts its version, ``we'll be prepared to go to the table as soon as possible to pass this
important legislation.''

Mr.~Obama has long signaled a preference for both an excise tax on high-end plans, which was first
proposed in a version of the legislation developed by the Senate Finance Committee, and the
independent commission, which some House leaders view as encroachment on Congressional authority
over Medicare.

The president made reference to both proposals in his speech to Congress just after Labor Day as he
sought to refocus the health care debate.

``This reform will charge insurance companies a fee for their most expensive policies, which will
encourage them to provide greater value for the money -- an idea which has the support of Democratic
and Republican experts,'' Mr.~Obama said in that address. ``And according to these same experts,
this modest change could help hold down the cost of health care for all of us in the long run.''

Of the Medicare panel, he said, ``We will also create an independent commission of doctors and
medical experts charged with identifying more waste in the years ahead.''

``The commission,'' he added, ``can help encourage the adoption of these common-sense best practices
by doctors and medical professionals throughout the system -- everything from reducing hospital
infection rates to encouraging better coordination between teams of doctors.''

The fate of the excise tax in particular is uncertain, partly because it is opposed by many labor
unions, an important component of the Democratic Party's base of support. In a nod to that
opposition, Mr.~Reid took steps to reduce the impact of the tax when he put together the Senate's
version of the legislation.

The Senate bill would impose an excise tax of 40 percent on the cost of employer-sponsored insurance
policies above \$8,500 for individuals and \$23,000 for families. It provides for increasing those
thresholds by \$1,350 for individuals and \$3,000 for families in cases where workers are in ``a
high-risk profession or employed to install electrical or telecommunications lines.'' And there
would be an additional increase in the thresholds, by the same amounts, in the 17 states where
health insurance is most expensive.

The excise tax is the single biggest revenue-raiser in the Senate bill. The nonpartisan
Congressional Budget Office has projected that over 10 years, it would generate \$149 billion toward
the \$848 billion 10-year cost of the legislation.

The main revenue-raiser in the House bill is an income surtax, on individuals earning more than
\$500,000 and couples earning more than \$1 million. It would generate \$460 billion over 10 years
toward the \$1.05 trillion cost of that measure.

The budget office has projected that both bills would reduce federal deficits to levels below what
they would be under current law, because the legislation's cost would be more than offset by new
taxes and fees and by reductions in government spending, mostly through slowing the growth of
Medicare.

The differing approaches in paying for the legislation are among the biggest divisions that Senate
and House leaders will face, assuming the Senate's passage of its bill. Those differences would have
to be worked out, with the resulting new version of the legislation approved in each chamber before
the measure could be sent to Mr.~Obama for his signature.

\section{Budget Hawks Have a Buffet of Options With Health Bill}

\lettrine{I}{f}\mycalendar{Nov.'09}{27} the Senate were going to write a new rule for Medicare
payments meant to slow the growth of medical costs, you might think that the rule would apply to
hospitals and doctors. A fair amount of medical care is, after all, provided by hospitals and
doctors.

But the Senate's health reform bill exempts hospitals from just such a rule until 2019. Doctors,
meanwhile, are likely to be effectively exempt, the Congressional Budget Office says. And come 2019,
a separate part of the rule will change, making the entire thing mostly moot.

This odd situation is just one of the opportunities that awaits those senators who don't yet seem to
have a firm position on health reform. All along, these centrist senators -- Blanche Lincoln of
Arkansas, Ben Nelson of Nebraska, Olympia Snowe of Maine and others -- have claimed that containing
costs was among their highest priorities. Yet it's never been clear how serious how they were. Most
have not been aggressively pushing cost containment proposals.

The bill brought to the floor last week by Harry Reid, the Senate leader, has offered them a chance
to make good on their rhetoric.

On the positive side, the bill includes nearly every big idea that health economists and medical
researchers have for slowing cost growth -- as well as for improving the patchwork quality of
American health care.

But many of the ideas, like the rule on Medicare reimbursement, have been at least partly neutered.
A provision to punish hospitals for infecting their patients, for example, would cut payments for
the related treatments by a mere 1 percent. A provision meant to help people who don't like the
insurance options offered by their employer would apply to only a tiny fraction of them. A provision
to encourage more cooperation among doctors would not apply to the areas where it is needed the
most: chronic diseases like diabetes and congestive heart failure.

``There is a lot to like in the bill,'' Dr.~Alan Garber of the Stanford School of Medicine says,
``but it needs to go further.''

Thus the opportunity for those centrist senators: to achieve their stated goal, they don't suddenly
need to turn themselves into health care wonks and rewrite the bill. They just need to improve
what's already there.

•

Complaining that Congress and the White House aren't doing enough to reduce the deficit is always a
popular pundit game. So it's no surprise that the last few weeks have been filled with knowing
claims that health reform will fail to control spiraling health costs.

Sometimes, however, Washington really does succeed in reducing the deficit. Presidents Harry Truman
and Dwight Eisenhower both did it. President Bill Clinton and Congress eliminated the deficit. Their
1993 budget bill was derided by some of the same people now criticizing health reform as an economy
wrecker. Instead, that budget bill created the first significant surpluses since the late 1940s (and
helped make possible the 1990s economic boom).

Health reform, done right, also has enormous potential. And the Senate bill isn't all that far from
doing it right. Ronald Brownstein has written a detailed essay for The Atlantic's Web site, in which
he breaks down the bill and quotes mostly pleased health experts, like Robert Reischauer, a former
head of the Congressional Budget Office whose skepticism of Mr.~Clinton's reform helped kill it.
Jonathan Gruber of M.I.T. calls himself ``a known skeptic on this stuff'' before adding, ``I can't
think of a thing to try that they didn't try.''

It is the execution of the ideas that can be problematic.

In a column two weeks ago, I laid out six issues that would determine whether health reform would
change health care (beyond expanding insurance, which clearly remains crucial). That is, would the
bill encourage the best treatments? Would it discourage errors and needless infections? Would it
bring costs under control so that most workers could again receive a decent raise?

By these measures, Mr.~Reid's bill is vastly better than the bill the House passed this month. The
Reid bill, however, is not as strong as the bill that the Senate Finance Committee wrote last month,
which itself was good but not ideal.

On only one of the six areas has Mr.~Reid seemed to have come up with an improvement: insurance
choice. His bill would give a small number of people with high premiums the ability to opt out of
their employer-provided plan and choose a plan on the insurance exchange for the uninsured and small
businesses. This provision is one of the few in any bill that would immediately help people who
don't like their employer coverage.

On three of the six areas -- preventing medical errors, financing more research into what kinds of
care actually work and creating pilot programs to reward doctors who provide good, less expensive
care -- the bills from Mr.~Reid and the finance committee are similar. On a fifth area -- a tax on
the costliest insurance plans -- the Reid version is only slightly worse. The tax is still big
enough to bring down insurance costs and raise incomes.

The real problem with the Reid bill comes in the sixth area, the one I alluded to at the start of
this column. The finance committee bill called for the creation of an independent commission with
the power to suggest changes to Medicare payment rates, based on available evidence. The goal would
be to steer the medical industry away from treatments that brought big profits and no net health
benefit.

Economists put the idea near the top of their wish list, as has President Obama. It has the
potential to bend the curve of Medicare spending, as the experts say, and eventually spread to the
rest of medicine. When Medicare policy has changed in the past, private insurers have often followed
suit.

Even as is, the finance panel's version of the commission had its flaws. It specifically exempted
hospitals and hospices until 2019, presumably for political reasons. The Congressional Budget Office
also thinks the commission will be too weak to prevent Congress from overriding any ideas that apply
to doctors and their powerful lobbyists. The commission's work would then be restricted to home
health care agencies and a few other parts of the medical industry.

Yet, as worrisome as these provisions were, they still left a promising commission in place after
2019. The Reid bill does not.

The details are technical but important. Beginning in 2020, the commission would be told to reduce
the growth rate of Medicare's budget merely to the growth rate of non-Medicare health spending. So
if per capita health spending kept rising 6 percent a year, as it has recently, Medicare's budget
could keep rising at that rate -- eventually making the program insolvent or requiring huge tax
increases.

Instead of turning Medicare into a role model, the Reid bill says it simply must be no more wasteful
than the rest of the health care system. Lewis Carroll would be proud.

What will those centrist senators do?

A cynic may suggest they were never as fiscally conservative as they let on. Some have preferred
repeating bromides about fiscal conservatism to engaging in the details of cost-cutting. Others,
like Mary Landrieu of Louisiana and many Blue Dog Democrats in the House, have seemed mainly
interested in securing Medicare dollars for rural areas. Still others -- including Joe Lieberman,
who represents a state, Connecticut, that is home to many insurance executives -- have seemed
obsessed with the so-called public insurance option. None of these positions will fix the health
care system.

Fortunately, though, the senators still have plenty of opportunity to prove the cynics wrong.

\section{At Obama's First State Dinner, the First Crashers}

\lettrine{T}{his}\mycalendar{Nov.'09}{27} much is known: About 7:15 Tuesday night, a glittering
blonde, decked out in a red and gold sari, holding the hand of her black-tuxedoed escort, swept past
the camera crews and reporters camped out to catch the red-carpet arrivals for the first state
dinner given by President Obama.

``Hey, that's a Desperate Housewife!'' one reporter yelled out, alluding to the TV series named for
a collection of them.

In fact, the couple -- Michaele Salahi and her husband, Tareq -- are Virginians who have been
auditioning for a possible role in a different housewives TV franchise: ``The Real Housewives of
Washington.''

They swept past the camera crews and followed the trail of other bigwigs attending the dinner.

But neither Mr.~nor Mrs.~Salahi, best known in the Washington area for promoting wine and polo in
Virginia, were on the guest list for the event, a fact first reported Wednesday morning on the
Washington Post Web site.

A White House official confirmed Wednesday that the Salahis had not been invited nor seated for
dinner.

It was not clear Wednesday night how close the Salahis got to Mr.~Obama and his wife, Michelle, or
to the guests of honor, Prime Minister Manmohan Singh of India, and his wife, Gursharan Kaur.

Edwin M.~Donovan, a spokesman for the Secret Service, said Wednesday night that ``initial findings
identified a Secret Service checkpoint which did not follow proper procedures,'' allowing the couple
to gain access to the festivities, even though their names were not on the invitation list.

``It is important to note that these individuals went through magnetometers and other levels of
security, as did all guests attending the dinner,'' Mr.~Donovan said in an e-mail message.

He declined to offer details, including what kind of procedures were not followed or whether the
couple had been interviewed. He said the investigation was continuing.``The bottom line is that they
should not have gotten in,'' he said.

Mr.~Donovan also said he was unaware of any other examples of this kind of party-crashing at the
White House. ``I'm not aware of any other incidents like this,'' Mr.~Donovan said.

The Salahis posted photographs of themselves at the dinner on their Facebook page: ``Honored to be
at the White House for the state dinner in honor of India with President Obama and our First Lady!''

As word of the incident spread on Wednesday night, Michaele Salahi's Facebook page was flooded with
comments from friends and acquaintances.

``Glad it was you, and not someone with malicious intent, who tested and beat White House
security,'' wrote Joe Galietta. ``Hopefully the officers who failed to stop you will learn from
this.''

Another commenter, Katie Paige, wondered how the couple pulled it off, asking ``What did you do
during the dinner when you weren't seated?''

Ms.~Salahi's profile also revealed a possible link to one of the dinner guests. In one of her
photographs, Ms.~Salahi and her husband are shown drinking wine and celebrating at the Indian
embassy with Arun K.~Singh -- the Deputy Chief of Mission at the embassy -- for the kickoff of a polo
match between India and the United States earlier this year. Ambassador Singh was on the guest list
for the state dinner Tuesday night.

\section{China's Impolitic Artist, Still Waiting to Be Silenced}

\lettrine{A}{i}\mycalendar{Nov.'09}{28} WEIWEI is perhaps China's most famous living artist and its
most vociferous domestic critic, titles of a sort this committed iconoclast disdains. Which is not a
bad thing, considering that recently, he very nearly lost them both.

Mr.~Ai was in Chengdu, the capital of Sichuan Province, preparing to testify at the trial of a
fellow political activist. ``By 3 a.m., we heard a very strong noise in the hallway, very brutal,
much like a Hollywood movie -- knocking on every door: 'Open it up -- we are the police!' '' he
said.

``They kicked open the door. I said, 'How do I know you are the police?' They said, 'I'll show you,'
and punched me here.'' Mr.~Ai pointed to the right side of his forehead. ``It was a very solid
punch.''

A month later, at an art exhibition in Munich, Mr.~Ai went to a doctor with a pounding headache and
was rushed into surgery to drain a pool of blood from his brain.

Mr.~Ai nearly died. Three months later, he says, his memory still fails him. On the other hand, ``I
don't have so many good memories anyway.''

That seems an exaggeration. At 52, Mr.~Ai, a beefy, bearded man with an air of almost monastic
composure, is an international figure in the art world, successful beyond what anyone might have
predicted even a decade ago. He is a celebrated architect, a co-designer of Beijing's landmark
Bird's Nest Olympic stadium, an installation artist and a documentary filmmaker with a 100-member
staff.

Artistically, he can do almost anything he wishes, like personally shipping 16 40-foot containers,
including 9,000 custom-made children's backpacks, from Beijing for his recent exhibition in Munich.

Yet clearly, all is not rosy in Mr.~Ai's world. In one of his early acclaimed works, a series of
three photographs called ``Dropping a Han Dynasty Urn,'' he dispassionately shatters a priceless
ancient Chinese vase, striking a theme -- destruction and recreation -- that runs through much of
his art.

Other works employ Ming and Qin period urns, furniture and architecture, assembled into haunting new
creations, or painted over, Warhol-style, with the Coca-Cola logo, or speared by wooden beams. A
series of photographs depicts global icons -- the Forbidden City, the White House, the Eiffel Tower
-- interrupted by Mr.~Ai's hand, middle finger raised.

Then there are his politics, an in-your-face criticism of China's leaders that, given Beijing's
limited tolerance for dissent, seems almost suicidal. Long before the Olympics, Mr.~Ai disavowed his
role in designing the Bird's Nest, saying the government had transformed the Olympics into a
patriotic celebration instead of using them to create a more open society.

In a 90-minute interview in his minimalist studio in north Beijing, Mr.~Ai called the government
unimaginative, prevaricating, suspicious of its own people and utterly focused on self-preservation.

``They don't believe in liberty. They don't believe in China before the Communists,'' he said.
``There is only one simple, clear task: to protect their control, to maintain their governing. Which
is such a pity.''

All of this he has said many times before. China's nationalists often accuse him of shilling for the
West, and in fact, Mr.~Ai ended his chat with a plea to President Obama to call for greater freedom
in China, saying ``we still need the moral support of the Western leaders'' to press for more
uncontrolled space in a still-closed society.

With or without help, Mr.~Ai is pressing hard. His most provocative art, as well as his latest
cause, concerns the question of why the May 2008 Sichuan earthquake killed thousands of children in
their classrooms -- and why the government has refused to give the public an official explanation.

After the quake, Mr.~Ai used the Internet to assemble scores of volunteers who combed the disaster
area, compiling a list of more than 5,000 dead children, organized by age and school, that now
covers one wall of his studio. ``The picture became clear. All of them belonged to about 20 schools,
and those schools, the buildings collapsed to dust,'' he said. ``Why did those buildings collapse,
and the ones next to it are standing?''

THE citizens' inquiry has produced a detailed list of questions, sent to government agencies, which
were supposed to be answered by this past Tuesday under law, but have yet to be addressed. In
December it will yield a documentary film on the disaster.

In Munich, the inquiry produced Mr.~Ai's most arresting work of art to date: those 9,000 children's
backpacks, covering one exterior wall of the Haus der Kunst. Against a blue background, colored bags
form the Chinese characters for the message, ``She lived happily on this earth for seven years,'' a
quotation from a mother of one earthquake victim.

The Munich exhibit is titled ``So Sorry,'' a caustic comment on the government's near-silence on the
schools disaster.

Mr.~Ai's beating in August occurred as he was preparing to testify at the trial of Tan Zuoren, a
Sichuan writer and activist who was trying to investigate the same issue. Mr.~Tan was accused of
inciting the subversion of state power. Mr.~Ai was blocked from testifying at his trial, which has
yet to produce a verdict. This week, though, another activist, Huang Qi, received a three-year
sentence for encouraging parents to press their grievances.

A disquieting sense of foreboding accompanies these jousts with the all-powerful state.

Ai Weiwei's father, Ai Qing, was both an artist and one of China's most revered contemporary poets,
who as a young man studied Baudelaire and Mayakovski in Paris. When he returned to Shanghai in 1932,
the ruling Kuomintang party jailed and tortured him, calling him a leftist. It was right: in 1941,
Ai Qing joined the Communist Party.

BUT 17 years later, in the infancy of Mao's new People's Republic, he ran afoul of the Communist
Party for subtly criticizing its suppression of free speech. The party exiled him, first to
Manchuria, then to remotest northwest China; Siberia, essentially.

Mr.~Ai and his family lived in a hut dug into the ground. His job for the next 16 years was to clean
out the village's public toilets.

``He was 60 years old. He had never done physical work in his life and he had to start doing it,''
his son said. ``Every night, he comes home very, very dirty. But he says, 'For 60 years, I don't
know who cleans my toilets. So now I do something for them.'

``That's something I learned from him. He became very powerful in terms of his thinking. He made the
toilet so clean, he would see it as a work of art -- like a museum, like MoMA.''

The family returned to Beijing in 1976, with the end of the Cultural Revolution. In 1985, the elder
Mr.~Ai, now rehabilitated, would receive a literary award from President François Mitterrand of
France. His son, on the other hand, could hardly wait to flee China.

Young Ai Weiwei studied at the Beijing Film Academy but in 1981 left for the United States. In New
York, Mr.~Ai said, he was in the city's art scene, not of it. He held temporary jobs and moved 10
times, throwing out his canvases each time for lack of storage room.

When his father fell ill in 1993, he agonized over returning to his homeland, which harbored such
painful memories. But after 1989, and the silencing of protesters at Tiananmen Square, he had
decided that ``the world became different.'' And so he returned to China in 1993, reckoning that one
day he might face something like his father's fate.

Lately, there are indeed signs that the government is reaching its limit. His blogs on Chinese Web
sites, about issues political and otherwise, have been shut down. Someone has installed two video
cameras outside his studio. The police are said to be scrutinizing his finances, an ominous
development in a state where other political critics have been prosecuted for what appear to be
concocted fiscal misdeeds.

``He has never done anything illegal,'' said his lawyer and friend, Liu Xiaoyuan. ``But if he
continues on his current path, getting involved in some very high-profile cases, I will get worried.
Some government departments are already very annoyed about him.''

Mr.~Ai says he is ready for whatever comes. ``I came to art because I wanted to escape the other
regulations of the society. The whole society is so political,'' he said. ``But the irony is that my
art becomes more and more political.''

\section{Uninvited Pair Met Obama; Secret Service Offers Apology}

\lettrine{P}{resident}\mycalendar{Nov.'09}{28} Obama and his wife, Michelle, had a face-to-face
encounter with the couple who sneaked into a state dinner at the White House this week, White House
officials acknowledged on Friday. The revelation underscored the seriousness of the security breach
and prompted an abject apology from the Secret Service.

A White House spokesman said that the couple, Michaele and Tareq Salahi of Virginia, met and shook
hands with the president and the first lady in the receiving line in the Blue Room, as the Obamas
greeted each of their 400 invited guests Tuesday night before moving to a tent on the South Lawn for
dinner.

That disclosure coincided with a statement from the director of the Secret Service, Mark Sullivan,
saying that his agency was ``deeply concerned and embarrassed'' by the events. Secret Service
officials said the agency wanted to interview everyone connected with the episode, including the
Salahis, and had not ruled out criminal charges.

``The preliminary findings of our internal investigation have determined established protocols were
not followed at an initial checkpoint, verifying that two individuals were on the guest list,''
Mr.~Sullivan said.

``Although these individuals went through magnetometers and other levels of screening, they should
have been prohibited from entering the event entirely,'' Mr.~Sullivan said. ``That failing is
ours.''

On Friday night, the White House released a photograph of the couple in the receiving line, being
greeted by Mr.~Obama. In the photo, a smiling Mrs.~Salahi, wearing a red and gold sari, is clasping
Mr.~Obama's hand with both of hers, as her husband looks on. The prime minister of India, Manmohan
Singh, is standing next to Mr.~Obama.

Mrs.~Salahi's Facebook page has photographs of the couple with officials including Vice President
Joseph R.~Biden Jr.~and Rahm Emanuel, the White House chief of staff. It is not known who took the
pictures.

The couple's lawyer, Paul W.~Gardner of Baltimore, asserts that the Salahis had been ``cleared, by
the White House,'' to be at the event, and so were not really ``crashing.'' Mr.~Gardner declined a
request to elaborate on his assertion.

But a Secret Service spokesman, James Mackin, said he had no idea how Mr.~Gardner could make such a
claim.

For their part, White House officials took pains to publicly refrain from criticizing the Secret
Service.

``The men and women of the Secret Service put their lives on the line every day to protect us; they
are heroes, and they have the full confidence of the president of United States,'' said Nick
Shapiro, a White House spokesman. ``The White House asked the United States Secret Service to do a
full review, and they are doing that. The United States Secret Service said they made a mistake, and
they are taking action to identify exactly what happened, and they will take the appropriate
measures pending the results of their investigation.''

Domestic security experts said that the episode raised serious questions about protection for the
president. Ronald Kessler, author of ``In the President's Secret Service: Behind the Scenes With
Agents in the Line of Fire and the Presidents They Protect'' said threats against the president had
increased 400 percent since Mr.~Obama took office.

On Friday, agents went to a Virginia winery founded by Mr.~Salahi's father in search of the couple,
one administration official said. A phone call to the establishment, Oasis Winery in Hume, Va., was
not returned. The winery's ownership has been the subject of extensive litigation involving disputes
within the family. The Salahis and the winery were saddled with hundreds of thousands of dollars of
unpaid bills, according to bankruptcy and legal findings in Virginia.

Mr.~and Mrs.~Salahi, who are known in the area to have a taste for polo and fine wine, are aspiring
reality-show celebrities. For months, the couple have been trailed by camera crews with the cable
channel Bravo, as it prepared for a new show, ``The Real Housewives of D.C.''

Seemingly distancing itself from the Salahis' actions, Bravo said Friday that it would not comment
about ``ongoing investigations.'' Earlier, the channel said that while its cameras were filming the
Salahis before the dinner, producers were told by the couple that they had been invited to it.

The Salahis have not officially been selected for the ``Houswives'' show.

``The decision as to who will be included in the series will not be made for several months,'' Bravo
said Friday.

By the sometimes twisted logic of reality TV, it is hard to tell whether the publicity surrounding
the state dinner would help or hinder the couple's goals.

For the breakout stars of the multicity ``Housewives'' franchise, reality television is a full-time
-- and lucrative -- job. Between per-episode fees, endorsements and public appearances,
``Housewives'' cast members can net six-figure salaries, said a producer with knowledge of the
franchise.

In a brutal economy, fortune is an even bigger motivator than fame for many aspiring reality
television contestants, producers and researchers say.

Financial hardships were clearly a factor for the Heene family of Fort Collins, Colo., who gained
nationwide attention last month when they claimed that their child had floated away on a homemade
flying saucer. According to a police affidavit, the Heenes were working with a production company to
pitch a television series to networks at the time. The parents pleaded guilty in the case this
month.

Both reality television and the Internet have trained people ``to brand themselves, to distribute
themselves, to get themselves out there,'' said James Hay, a communications professor at the
University of Illinois and a co-author of the book ``Better Living Through Reality TV.''

\section{Some Indians Find It Tough to Go Home Again}

\lettrine{W}{hen}\mycalendar{Nov.'09}{28} 7-year-old Shiva Ayyadurai left Mumbai with his family
nearly 40 years ago, he promised himself he would return to India someday to help his country.

In June, Mr.~Ayyadurai, now 45, moved from Boston to New Delhi hoping to make good on that promise.
An entrepreneur and lecturer at the Massachusetts Institute of Technology, with a fistful of
American degrees, he was the first recruit of an ambitious government program to lure talented
scientists of the so-called desi diaspora back to their homeland.

``It seemed perfect,'' he said recently of the job opportunity.

It wasn't.

As Mr.~Ayyadurai sees it now, his Western business education met India's notoriously inefficient,
opaque government, and things went downhill from there. Within weeks, he and his boss were at
loggerheads. Last month, his job offer was withdrawn. Mr.~Ayyadurai has moved back to Boston.

In recent years, Mother India has welcomed back tens of thousands of former emigrants and their
offspring. When he visited the United States this week, Prime Minister Manmohan Singh personally
extended an invitation ``to all Indian-Americans and nonresident Indians who wish to return home.''
But, like Mr.~Ayyadurai, many Indians who spent most of their lives in North America and Europe are
finding they can't go home again.

About 100,000 ``returnees'' will move from the United States to India in the next five years,
estimates Vivek Wadhwa, a research associate at Harvard University who has studied the topic. These
repats, as they are known, are drawn by India's booming economic growth, the chance to wrestle with
complex problems and the opportunity to learn more about their heritage. They are joining
multinational companies, starting new businesses and even becoming part of India's sleepy government
bureaucracy.

But a study by Mr.~Wadhwa and other academics found that 34 percent of repats found it difficult to
return to India -- compared to just 13 percent of Indian immigrants who found it difficult to settle
in the United States. The repats complained about traffic, lack of infrastructure, bureaucracy and
pollution.

For many returnees the cultural ties and chance to do good that drew them back are overshadowed by
workplace cultures that feel unexpectedly foreign, and can be frustrating. Sometimes returnees
discover that they share more in their attitudes and perspectives with other Americans or with the
British than with other Indians. Some stay just a few months, some return to the West after a few
years.

Returnees run into trouble when they ``look Indian but think American,'' said Anjali Bansal,
managing partner in India for Spencer Stuart, the global executive search firm. People expect them
to know the country because of how they look, but they may not be familiar with the way things run,
she said. Similarly, when things don't operate the way they do in the United States or Britain, the
repats sometimes complain.

``India can seem to have a fairly ambiguous and chaotic way of working, but it works,'' Ms.~Bansal
said. ``I've heard people say things like 'It is so inefficient or it is so unprofessional.' '' She
said it was more constructive to just accept customs as being different.

Sometimes, the better fit for a job in India is an expatriate\cite{expatriate} who has experience
working in emerging markets, rather than someone born in India who has only worked in the United
States, she said.

While several Indian-origin authors have penned soul-searching tomes about their return to India,
and dozens of business books exist for Western expatriates trying to do business here, the
guidelines for the returning Indian manager or entrepreneur are still being drawn.

``Some very simple practices that you often take for granted, such as being ethical\cite{ethical} in
day to day situations, or believing in the rule of law in everyday behavior, are surprisingly absent
in many situations,'' said Raju Narisetti, who was born in Hyderabad and returned to India in 2006
to found a business newspaper called Mint, which is now the country's second-biggest business paper
by readership.

He said he left earlier than he expected because of a ``troubling nexus'' of business, politics and
publishing that he called ``draining on body and soul.'' He returned to the United States this year
to join The Washington Post.

There are no shortcuts to spending lots of time working in the country, returnees say. ``There are
so many things that are tricky about doing business in India that it takes years to figure it out,''
said Sanjay Kamlani, the co-chief executive of Pangea3, a legal outsourcing firm with offices in New
York and Mumbai. Mr.~Kamlani was born in Miami, where his parents emigrated from Mumbai, but he has
started two businesses with Indian operations.

When Mr.~Kamlani started hiring in India, he met with a completely unexpected phenomena: some new
recruits would not show up for work on their first day. Then, their mothers would call and say they
were sick for days in a row. They never intended to come at all, he realized, but ``there's a
cultural desire to avoid confrontation,'' he said.

The case of Mr.~Ayyadurai, the M.I.T. lecturer, illustrates just how frustrating the experience can
be for someone schooled in more direct, American-style management. After a long meeting with a top
bureaucrat, who gave him a handwritten job offer, Mr.~Ayyadurai signed on to the Council of
Scientific and Industrial Research, or C.S.I.R., a government-financed agency that reports to the
ministry of science.

The agency is responsible for creating a new company, called C.S.I.R.-Tech, to spin off profitable
businesses from India's dozens of public laboratories. Currently, the agency, which oversees 4,500
scientists, generates just \$80 million in cash flow a year, even though its annual budget is the
equivalent of half a billion dollars.

Mr.~Ayyadurai said he spent weeks trying to get answers and responses to e-mail messages,
particularly from the person who hired him, the C.S.I.R. director general, Samir K.~Brahmachari.
After several months of trying to set up a business plan for the new company with no input from his
boss, he said, he distributed a draft plan to C.S.I.R.'s scientists asking for feedback, and
criticizing the agency's management.

Four days later, Mr.~Ayyadurai was forbidden from communicating with other scientists. Later, he
received an official letter saying his job offer was withdrawn.

The complaints in Mr.~Ayyadurai's paper could be an outline for what many inside and outside India
say could be improved in some workplaces here: disorganization, intimidation, a culture where top
directors' decisions are rarely challenged and a lack of respect for promptness that means meetings
start hours late and sometimes go on for hours with no clear agenda.

But going public with such accusations is highly unusual. Mr.~Ayyadurai circulated his paper not
just to the agency's scientists but to journalists, and wrote about his situation to Prime Minister
Manmohan Singh. India is ``sitting on a huge opportunity'' to create new businesses and tap into
thousands of science and technology experts, Mr.~Ayyadurai said, but a ``feudal culture'' is holding
the country back.

Mr.~Brahmachari said in an interview that Mr.~Ayyadurai had misunderstood nearly everything -- from
his handwritten job offer, which he said was only meant to suggest what Mr.~Ayyadurai could receive
were he to be hired, to the way Mr.~Ayyadurai asked scientists for their feedback on what the
C.S.I.R. spinoff should look like.

To prove his point, Mr.~Brahmachari, who was two hours late for an interview scheduled by his
office, read from a government guide about decision-making in the organization. Mr.~Ayyadurai didn't
follow protocol, he said. ``As long as your language is positive for the organization I have no
problem,'' he added.

As the interview was closing, Mr.~Brahmachari questioned why anyone would be interested in the
situation, and then said he would complain to a reporter's bosses in New York if she continued to
pursue the story.

\section{Seeking the Best Medical Care Prices}

\lettrine{C}{an}\mycalendar{Nov.'09}{28} you really shop for doctors and hospitals the way you would
for airfares and flat-screen TVs?

Health care consumers are encouraged to comparison-shop on things like doctor's fees and heart
surgery rates. But unfortunately, most of us have little clear or useful information to go shopping
with.

``When you go to the doctor, how much you fork over when all is said and done is often just a
mystery,'' said Dr.~Anthony P.~Geraci, a Manhattan neurologist who is trying to buck that trend by
posting his prices on his Web site.

With the growing number of uninsured people, the increase in high-deductible insurance plans and big
jumps in co-payments, just about everybody is paying more out of pocket for health care nowadays. An
estimated 15 percent of adults younger than 65 now pay with their own money medical costs greater
than 5 percent of their annual household income, according to the Center for Studying Health System
Change, a nonpartisan research group in Washington.

So the typical person is probably far more motivated to ask how much an M.R.I. or a hip replacement
costs. And just as often, people are asking -- or should be -- ``How can I get a better price?''

Take Katie Kyser, 30, the mother of a year-old daughter, who lives north of Seattle. She and her
husband, Jason, who works in construction, recently moved from California. They have no health
insurance, so they pay all costs out of pocket.

When Ms.~Kyser needed a routine gynecological exam, she called a handful of local doctors, all of
whom were charging \$200 or more. ``There's no way we could pay that,'' Ms.~Kyser said. ``I had to
find another way.''

Having seen an ad for PriceDoc.com, a new Web site that lists doctors throughout the country who are
willing to post their prices and negotiate with patients, she decided to try it. Ms.~Kyser found a
nearby clinic where doctors charged only \$75 for the exam.

``I was a little nervous at first because the price was so cheap, but when I got there, it was
wonderful,'' Ms.~Kyser said. ``Everyone was so professional and helpful.''

The crucial part of shopping wisely for health care (or anything else for that matter) is comparing
prices the way Ms.~Kyser did.

But that is also where problems arise. Medical pricing is a quagmire, oozing with jargon and current
procedural terminology codes. Just look, if you dare, at your latest ``explanation of benefits''
from your insurer.

What's more, rarely is there one standard price for a medical treatment. Prices vary based on
geography and type of provider -- whether hospital, stand-alone clinic or any alternative.

Then, doctors, hospitals and other providers may negotiate different rates with different insurers.
It is not unusual for a provider to have 10 or more different prices for the same procedure,
depending on who is paying. Providers often charge a completely different rate for people paying on
their own, which is almost always much more expensive than the discounted rate that insurers pay.

``It's a challenge for consumers to sift through these different price structures,'' said Ha T.~Tu,
senior health researcher at the Center for Studying Health System Change. And there is no one place
to go for good information, she added.

Despite the challenges, here are several steps consumers can take to make health care shopping a bit
more manageable:

CHECK WITH YOUR INSURER Many insurance companies have begun posting provider prices on their Web
sites so enrollees can access cost information. These tools allow you to compare prices among
network doctors (not all network doctors are paid the same) and check on the price of diagnostic
tests and other treatments.

``This is especially helpful if you're in a high-deductible plan,'' said Ms.~Tu, ``because you can
see how much you'll pay out of pocket.''

USE THE INTERNET A few Web companies have tried to fill the price information gap online, all with
varying approaches.

On PriceDoc.com, the site Ms.~Kyser used, you plug in your ZIP code to find a list of providers in
your area who have posted their prices. You can also plug in the price you're willing to pay.
Providers will then respond if they are willing to accept that price.

HealthcareBlueBook.com compiles prices paid for specific treatments and procedures in ZIP codes
throughout the country, then lists what the site determines is a range of fair prices. Consumers can
then use these ranges as a jumping-off point for negotiating with their providers, said Dr.~Jeffrey
Rice, the chief executive of the concern.

Dr.~Rice tells the story of a woman in northern Ohio who had been quoted a price at a local hospital
of \$2,500 for an M.R.I. of her knee. When she looked up the test on the site, she found the fair
price in that area was more like \$500.

She went back to the hospital where she had been quoted the high price and started asking questions.
The clerk told her it would be much less expensive if she went to the clinic down the street instead
of the hospital. The woman followed that advice and paid \$300 for her M.R.I.

Another Web site, OutOfPocket.com, combines price information that users send in to determine a
going rate for specific health care costs throughout the country.

None of these sites are comprehensive, although all of them are easy to use and are expanding their
listings. It's worth taking a look to see if you can glean any useful information from them.

BROWSE STATE DATA If you're checking out hospitals, you will want to see what information your state
government offers. At least 33 states mandate that hospitals make their prices public, Ms.~Tu says.

But there are caveats. Often, only the most expensive, nondiscounted prices are listed.

Moreover, on most sites, costs are not bundled, so you may find the price of a general surgery, for
example, but it would not include the surgeon's or anesthesiologist's fees.

Some states offer more information than others, points out Ms.~Tu. Minnesota, for example, uses
average prices for some procedures, and New Hampshire and Maine have some bundled prices. To see
what information, if any, is available in your state, you can use the links on
healthcarebluebook.com.

PICK UP THE PHONE ``The most important thing to do if you're looking for price information is call
your doctor,'' said Jonathan Weiner, professor of health policy and management at Johns Hopkins
University.

``This is still an awkward discussion for most doctors,'' Professor Weiner said. ``But if you sit
down and talk about money, it almost always leads to discounts, particularly for self-paid people.''

If your doctor balks\cite{balk} at having this conversation, ask to speak to the office worker in
charge of billing, who will know the prices your doctor charges and can at least estimate what you
will be paying. Then, when you do collect price information, you can return to that person to
negotiate a better price.

\section{At Odds Over Land, Money and Gas}

\lettrine{C}{hris}\mycalendar{Nov.'09}{28} and Robert Lacey own 80 acres of idyllic\cite{idyllic}
upstate New York countryside, a place where they can fish for bass in their own pond, hike through
white pines and chase deer away.

But the Laceys hope that, if all goes well, a natural gas wellhead will soon occupy this bucolic
landscape.

Like many landowners in Broome County, which includes the town of Chenango, the Laceys could
potentially earn millions of dollars from the natural gas under their feet. They live above the
Marcellus Shale\cite{shale}, a subterranean layer of rock stretching from New York to Tennessee that
is believed to be one of the biggest natural gas fields in the world.

As New York environmental officials draft regulations to allow drilling in the shale as early as
next year, thousands of residents like the Laceys in upstate counties have banded together in
coalitions to sign leases with gas companies for drilling on their land -- for \$5,000 to \$6,000 an
acre for a term of five years, and royalties of up to 20 percent on whatever gas is found.

``When I heard about drilling, what came to mind was `Thank you,' '' said Mrs.~Lacey, 58, who has
lived on her property here for 27 years with her husband, Robert, 68, a commercial insurance agent.
``Finally our community can recover, and our children don't have to leave the state to find jobs.''

In New York City, natural gas exploration is largely seen as a threat to the drinking water the city
gets from watersheds to the north in the Catskills. But in the rural communities above the shale,
the reaction has been far more mixed -- and far more contentious.

Some residents welcome the drilling as a modern-day gold rush and salvation from the economic
doldrums that they say have chased jobs and young people away from their area. Others express
concerns about the environment and quality-of-life issues like noise and heavy-truck traffic.

In some cases, the issue has pitted neighbor against neighbor or spouse against spouse.

Exploring the shale involves a drilling method called hydraulic fracturing that requires pumping
huge volumes of water laced with benzene and other chemicals into the rock to break it and extract
gas. The process raises issues about the use and disposal of wastewater, and the danger of leaks,
spills and other contamination. It has been linked to contamination of water wells in Pennsylvania
and Wyoming and to the death of livestock in Louisiana.

Mark Dunau, an organic vegetable farmer with 50 acres in the town of Hancock in Delaware County,
says he is passing up any potential rewards from drilling because of worries about the pollution of
water and air and the cumulative impact of thousands of wells. ``That water is my resource,'' he
said.

Mr.~Dunau, 57, and his wife, Lisa Wujnovich, 55, said that they were holdouts not only among their
neighbors but also among their friends, and that the character of their community was already
changing. Mr.~Dunau said he knew people who said they would take the money and move away, families
who were fighting over whether to sign gas leases, and neighbors who regretted signing too early for
too little money.

``It's a nightmare,'' he said.

One of Mr.~Dunau's neighbors, Grace K.~Kinzer, signed a lease with Chesapeake Energy two years ago
when a representative of the company knocked on her door with an offer of \$25 an acre and royalties
of 12.5 percent, she said. Ms.~Kinzer, 83, said she needed money to pay her taxes so she signed,
getting about \$2,750 for 110 acres that would now fetch more than half a million dollars.

Ms.~Kinzer said she felt as if ``I got taken for a ride.''

But Brad Gill, executive director of the Independent Oil and Gas Association of New York, said that
two years ago, gas companies could not justify high leasing costs. Prices have steadily climbed with
the success of drilling the Marcellus Shale in states like Pennsylvania, he said.

``I understand the frustration, but I'd also be fair-minded,'' he said.

Sometimes rifts over drilling are found under the same roof.

Myron and Shirley Ernst, who own two acres in the Broome County town of Vestal, were recently
invited by a neighbor to join a landowners' coalition that is negotiating with gas companies as a
group. The Ernsts, teachers who ran a Montessori preschool on the property until their retirement
five years ago, are at a standoff.

Mr.~Ernst, 72, favors drilling as a matter of survival. ``We're prostrate, and dependent on oil from
enemy countries,'' he said.

His wife, 74, said she was too worried about possible accidents and chemical spills.

``My main concern is the aquifer,'' she said. ``We have our own well. I don't want to have to buy
bottled drinking water.''

She refuses to sign a lease, and, after 46 years of marriage, Mr.~Ernst said, he knows better than
to think he can persuade her. ``That's like asking how do I plan to fly to Pluto,'' he said.

But the holdouts cannot stop the transformation of their surroundings, both good and bad, once
drilling is allowed. The area of the shale in New York that is expected to be the most productive
spans about 3.4 million acres in 10 counties, but lies mostly in Broome, Delaware, Sullivan,
Chenango and Tioga, said officials from the Independent Oil and Gas Association of New York.

Further fostering bad feelings within communities, under a concept known as
``compulsory\cite{compulsory} integration'' gas companies can drill under land without the owner's
consent if they have leases in most of the surrounding area. Those owners without leases would get
royalties of 12.5 percent on gas from their property, the minimum allowed under state law, but many
worry more about exposing their water to pollution.

``They could be drilling directly under your well and threatening your groundwater,'' said Wes
Gillingham, the program director for the environmental group Catskill Mountainkeeper. He owns 100
acres in Sullivan County and said he was trying to keep it off-limits to drilling.

The regulations proposed by New York's Department of Environmental Conservation are under fire as
inadequate from environmentalists around the state. Another concern shared by both the pro- and
anti-drilling sides is whether the environmental agency, which has suffered financial and staff
cuts, will have the resources to monitor the drilling.

But the financial benefits beckon\cite{beckon} strongly in the communities that stand to benefit
from a gas boom at a time when they are suffering economically. Broome County, which has lost jobs
and population over the last decade, would stand to gain \$3.72 billion a year in wages and tax and
retail revenue from up to 4,000 gas wells, a study commissioned by local officials estimated.
(However, a report by Columbia University researchers this year noted that any prediction of
economic results was ``entirely speculative'' because of unknowns like the potential cost of
cleaning up contamination.)

Broome County officials themselves are seeking to lease the rights beneath the county landfill and
airport.

``We cannot afford to chase this industry away,'' Patrick J.~Brennan, deputy county executive, told
state environmental officials this month at a hearing on the proposed rules governing drilling that
drew about 800 people to a school auditorium in the county.

Both to poor residents with small parcels of land and wealthy landowners with vast estates, drilling
represents an economic opportunity tantamount to winning the lottery. Jim Ward, who heads a
landowners' coalition with 138 members whose properties range in size from 650 acres to half an
acre, said the sentiment in his group is: ``I should have a right to prosper from my land.''

Mrs.~Lacey, who also belongs to a landowners' coalition with her husband, said she was satisfied
that the proposed state drilling regulations offered ``reasonable protection.'' She said her group's
lawyers were drawing up contracts that added even more restrictions.

Common provisions in the leases require gas companies to restore the land to its original state
after drilling ceases and to assume all liability for any water contamination.

``If people are concerned, they should hide in a cave,'' Mrs.~Lacey said. ``These are the people who
want to put gas in their car and watch TV, but you can't do all these things if there's no energy.
We have to take a risk here.''

The Laceys would make more than \$400,000 upfront once they signed their lease. Although they are
not struggling financially, Mrs.~Lacey said, she and her husband worry about rising property taxes
biting into their retirement comfort. They also want their two married daughters to move back to the
area with their families, and they long for an economically vibrant community.

``We love it here, and we'd like to stay here after retirement,'' Mrs.~Lacey said. ``It'd be easier
for us to do that if some of this gas would help us defray our property taxes.''

She added, ``It just basically means that you sleep better at night.''

\section{Buying, Selling and Twittering All the Way}

\lettrine{O}{nce}\mycalendar{Nov.'09}{28} upon a time, people mailed their holiday wishes to the
North Pole and hoped for a reply on Christmas Day. Nowadays they are sending their wishes into
cyberspace and are apt to get a reply in minutes.

America's first Twitter Christmas got under way in earnest on Friday. Across the land, retailers and
their customers used the social networking site to talk to one another about bargains, problems,
purchases and shopping strategies.

After buying a new navigation system at 6 a.m. on the most frenzied shopping day of the year, Laura
S.~Kern of Los Angeles could not figure out why it was not giving her traffic updates. She sent a
message to Best Buy's Twitter account and within five minutes not one, but two Best Buy employees
responded with fix-it advice.

In Bloomington, Minn., Mall of America used its Twitter page to tell consumers two of its parking
areas were at capacity and that their best bet was to park near Ikea.

Twitter permits public communication via short, to-the-point messages. Many people use it to send
mundane updates to their friends, but increasingly, the nation's retailers see it as a business
tool.

It gives customers a practical way to cajole a retailer, complain about something or ask questions.

A Twitter post can in theory be seen by millions, and thus packs more punch than an e-mail message
or a phone call to a store. The big retailers are all scrambling this Christmas to come up with
Twitter plans. They are designating tech-savvy employees to respond to the posts, sometimes by
providing up-to-minute inventory information from a sales floor, for example, or by offering help
with some balky gadget.

``It's one of the greatest emerging communication channels out there,'' said Greg Ahearn, senior
vice president of marketing and e-commerce for Toys ``R'' Us. ``This is a way people can stay
connected with the brand in a way they've never been able to before.''

So far this shopping weekend, special deals have been posted on Twitter from stores as varied as
Best Buy, J.C. Penney, Toys ``R`` Us, Staples, Gap, Bloomingdale's, and Barneys. (Links to the
retailing Twitter accounts mentioned in this article can be found in the Web version of the story on
NYTimes.com.)

For the uninitiated, Twitter.com is a Web site where each member has a password-protected page. It
has a blank box for typing in a message of 140 characters or fewer, an act known as tweeting.

To see a retailer's messages, Twitter users ``follow'' the retailer, which means that the chain's
posts show up on their Twitter home page when they log in. And the system allows users to send
messages in the other direction, so that a retailer's employees will see them.

``I think in this economy you need to leverage every asset that you have,'' said James Fielding,
president of Disney Stores Worldwide, who sends messages under the Twitter name, or handle,
DisneyStorePrez.

On Friday morning, as consumers flooded Disney Stores around the country, Mr.~Fielding messaged:
``We have amazing ONE DAY ONLY deals previewing on our Facebook page -- become a fan today and find
out more!''

Retailers hope that if they send Twitter messages, consumers will come. About 47 percent of
retailers said they would increase their use of social media this holiday season, according to a
study by Shop.org, part of the National Retail Federation, an industry group. And more than half of
retailers said they added or improved their Facebook and Twitter pages. There are advantages for
consumers too, like discounts. For instance, those who decided to follow Gap Outlet received an
offer for 15 percent off purchases of \$75 or more.

As shoppers jammed the aisles on Friday at a Best Buy store in Arlington Heights, Ill., an employee,
Jerry DeFrancisco, went up to a computer kiosk and used his Twitter account to tell customers about
Best Buy's home theater deals. Then he resumed his in-store duties, helping a customer decipher a
sales circular.

A few months ago, Best Buy began piloting a Twelpforce -- a Twitter-inspired play on ``help force''
-- of some 2,500 employees that answer consumers' questions in real time.

``It's 24-hour access to our employees,'' said Brad Smith, director of interactive marketing and
emerging media for Best Buy. The Twelpforce had fielded about 25,000 questions even before gearing
up for Thanksgiving weekend.

Ms.~Kern in Los Angeles used the service on Friday. After she could not get her new navigation
system to work, she tried Best Buy's telephone support line, only to receive a warning that her wait
would be an hour. So she posted on Twitter instead, and within minutes, Best Buy employees were
sending her useful links and details about her gadget. ``It's amazing,'' she said later in the day.
(Her interaction with the employees ultimately helped her realize she would need to go back to the
store for help.)

Many retailers will be posting to their Twitter pages throughout the weekend and the entire holiday
season. Some chains have an official Twitter account. Others have many, like one for each store, or
one for each employee who wants to post messages. There are Twitter pages for designers, like Nicole
Miller and Diane von Furstenberg.

Retailers also use Facebook to interact with their customers. But Facebook, with its photo albums
and various applications, does not have the same no-frills immediacy as Twitter -- which is why
Twitter is ideal for instantaneously announcing sales.

In addition to bargains, stores are also using Facebook and Twitter to promote contests and games
that they hope will keep consumers engaged and coming back. Best Buy has an interactive Secret Santa
application on its Facebook page. Gap is using Twitter to inform New York City residents and
visitors where its ``Gap Cheer'' bus (filled with dancers and drummers) will be parked and giving
away sweaters and jeans.

Of course, sometimes retailers simply use their Twitter posts to capture the spirit of the season.
At 3:30 Thursday morning, an employee posted seven words on the Macy's Twitter page, about a
marching band that was practicing hours before the chain's Thanksgiving day parade.

It said: ``Is he really running with a tuba?''

\section{Some Jostling, but Less of a Frenzy Among Shoppers}

\lettrine{A}{}\mycalendar{Nov.'09}{28} year after an unfolding economic crisis sent consumers into
shock, they ventured out Friday and opened their wallets a bit -- still hunting bargains, but no
longer quite so afraid to spend.

On the marathon shopping day known as Black Friday, anecdotal reports across the country suggested
an encouraging opening of the holiday shopping season, with larger, more enthusiastic crowds than
last year. Across many categories of merchandise, retailers noticed that consumers seemed more
willing to spend money on themselves than they had been a year ago.

And thousands of people explored a new way of interacting with merchants, posting messages on
Twitter and getting replies from store employees, often within minutes.

Despite the improved mood, frugality has not gone out of style. Many of Friday's shoppers seemed to
stick to lists and take advantage of discounts, rather than give in to impulse purchases.

``My feeling is that consumers are just more realistic and more accepting of the economy and where
their personal situation is this year,'' said Jim Fielding, president of the Disney Store chain.
``They're coming in with a budget and they're coming in with a list.''

Indeed, veteran retailing executives and analysts said that, while store turnout was improved from a
year ago, they had never seen the American consumer shop so strategically.

``If it's at full price, nobody's even looking at the stuff,'' said Marshal Cohen, chief industry
analyst for the NPD Group, who spent the day surveying malls.

Terry J.~Lundgren, president and chief executive of Macy's, said some 5,000 people lined up outside
Macy's in Midtown Manhattan on Friday morning, slightly more than last year. But he said consumers
were still shopping with caution.

``They've got a budget that they've planned to spend and I think most consumers will stick to
that,'' he said. ``We expect that this will be a market-share holiday season. We don't need
consumers to spend more money. We just need to make them spend more money with us.''

The National Retail Federation, an industry group, reported that high-definition TVs, laptops,
winter coats and Zhu Zhu Pets toy hamsters were among the most sought-after items.

How much consumers actually spent on Friday will not become clear for several days. Gerald
L.~Storch, chairman and chief executive of Toys ``R'' Us, said he thought the day would be a
success. ``For the retailers as a whole this might be the biggest Black Friday in history,'' he
said.

It was also a more orderly day, despite some jostling and angry words here and there. Wal-Mart, the
nation's largest retailer, strengthened security in some stores a year after a temporary store
worker in Valley Stream, N.Y., died after being trampled by an unruly crowd. The biggest disturbance
appeared to be at a Wal-Mart store in Upland, Calif., which closed for a few hours when shoppers
began fighting over merchandise

Black Friday got its name because it was perceived as the day retailers shifted into the black, or
became profitable for the year. While it is always a high-volume shopping day, that turnout does not
necessarily translate into big profits or determine how consumers will behave in the weeks to come.

Chains are hoping this season will be better than last, when retailers posted the worst sales
figures in decades. Industry sales last year fell 2.2 percent for the shopping season (November and
December combined), the biggest decline since at least 1970, according to the International Council
of Shopping Centers, an industry group. This year, sales are expected to be about the same as last
Christmas, give or take a couple of percentage points.

The National Retail Federation, a trade group, predicted sales would decline 1 percent, to \$437.6
billion. Some industry groups and professionals say that figure is low. The International Council of
Shopping Centers, for instance, estimates holiday sales will rise 1 to 2 percent.

One bright spot on Friday was a trend of people giving themselves gifts. Many shoppers, after
denying themselves for the better part of a year, snapped up discounted merchandise like boots and
sweaters. Matthew F.~Katz, a managing director in the retailing practice of AlixPartners, a
reorganization firm, said the shoe department at Nordstrom was packed with such shoppers. Consumers
will wait for better discounts around Christmastime to buy gifts for others, he said.

Madeline Hart, 18, a high school senior, arrived at Woodfield Shopping Center in Schaumburg, Ill.,
at 4:30 a.m. and spent \$200 on clothing and accessories at Express and Victoria's Secret.

``I work part time at a car wash, so I bought this with my own money,'' said Ms.~Hart, peeking into
a pair of shopping bags holding a few shirts, a bra and a shiny red vinyl tote bag, a free extra
from Victoria's Secret.

While the economy is showing signs of recovery, unemployment is the highest in decades and consumers
are still deeply worried, which is likely to affect holiday shopping. The Conference Board said this
week that American households were expected to spend an average of \$390 on Christmas gifts, down
from \$418 last year.

``Job losses and uncertainty about the future are making for a very frugal shopper,'' Lynn Franco,
the director of the board's consumer research center, said in a statement.

At Wal-Mart stores around the country, consumers said they were feeling the pain. At a store in
Columbus, Ohio, Michelle Perry, 33, a waitress at Waffle House, said she moved this year from a nice
rented house in the middle-class suburb of Pataskala to one in a rough neighborhood on the north
side of Columbus.

``I can't even let my kids outside without me going with them,'' she said. ``But we had to do
something to save money, and I'm saving \$300 a month by switching from renting to buying.''

Still, she feels better about her finances this year than she did last year. After Wal-Mart, she was
off to Meijer's to look for an Xbox 360 for her son, and for gifts inspired by the film ``Twilight''
for her 12-year-old daughter.

``She's got to have the Twilight lip gloss,'' Ms.~Perry said. ``Every girl at her school has it, so
she's got to have it, too.''

\section{China Charges 58 With Covering Up Deadly Mine Blast}

\lettrine{T}{en}\mycalendar{Dec.'09}{01} journalists and 48 officials have been charged with taking
bribes to cover up a mining disaster last year, according to a report published on Monday in China
Daily, an official English-language newspaper.

Mine bosses relocated bodies, destroyed evidence and paid the journalists the equivalent of
\$381,000 to cover up the explosion, in which 34 miners and a rescue worker were killed, China Daily
reported. Earlier reports by other news organizations indicated that the bosses also cremated
miners' bodies against the wishes of family members, paid grieving relatives to silence them and
sealed the mine shaft with truckloads of dirt.

The disaster took place on July 14, 2008, almost a month before the Beijing Olympics, at the Lijiawa
mine in Hebei Province, about 100 miles west of Beijing. The cover-up kept the disaster out of the
public eye for 85 days.

In September 2008, someone reported the cover-up on an Internet chat site, and the ensuing clamor
forced the central government in Beijing to step in, firing 25 local officials and putting 22 of
them under criminal investigation. The charges reported by China Daily on Monday were the result of
an investigation by the State Council, China's cabinet.

The report said that 48 officials were being charged, including the mine owners, the county chief,
work safety officials and police officers. The China Daily said none of the 10 journalists charged
in the case had been identified.

Because of low salaries, journalists in China are often tempted to accept bribes, called hongbao, or
red envelopes.

The China Daily said one of the journalists being charged is almost certainly Guan Jian, a reporter
from Beijing working for a newspaper called China Internet Weekly. Mr.~Guan was detained in Shanxi
Province last December and charged in April with taking bribes from officials in Yuxian County,
where the mine is located, in the aftermath of the disaster.

In that case, prosecutors accused Mr.~Guan of receiving \$36,600 from officials under the pretense
of running two pages of advertising in his newspaper, and also receiving a so-called newspaper
subscription fee of \$4,400.

Last year in Shanxi Province, two journalists and 26 people posing as journalists were accused of
taking money to cover up a coal mine accident in which a worker was killed.

Mine fatalities in China, even if underreported, annually rank among the highest in the world. On
average, nine coal miners died each day in China last year -- a rate 40 times that of the United
States, according to statistics from the State Administration of Work Safety. Small mines, legal and
illegal, accounted for three-fourths of the deaths but only a third of China's production.

\section{Obama's Speech on Afghanistan to Envision\cite{envision} Exit}

\lettrine{P}{resident}\mycalendar{Dec.'09}{01} Obama plans to lay out a time frame for winding down
the American involvement in the war in Afghanistan when he announces his decision this week to send
more forces, senior administration officials said Sunday.

Although the speech was still in draft form, the officials said the president wanted to use the
address at the United States Military Academy at West Point on Tuesday night not only to announce
the immediate order to deploy roughly 30,000 more troops, but also to convey how he intends to turn
the fight over to the Kabul government.

``It's accurate to say that he will be more explicit about both goals and time frame than has been
the case before and than has been part of the public discussion,'' said a senior official, who
requested anonymity to discuss the speech before it is delivered. ``He wants to give a clear sense
of both the time frame for action and how the war will eventually wind down.''

The officials would not disclose the time frame. But they said it would not be tied to particular
conditions on the ground nor would it be as firm as the current schedule for withdrawing troops in
Iraq, where Mr.~Obama has committed to withdrawing most combat units by August and all forces by the
end of 2011.

Officials of one allied nation who have been extensively briefed on the president's plan said,
however, that Mr.~Obama would describe how the American presence would be ratcheted back after the
buildup, while making clear that a significant American presence in Afghanistan would remain for a
long while. That is designed in part to signal to Pakistan that the United States will not abandon
the region and to allay\cite{allay} Pakistani fears that India will fill the vacuum created as
America pulls back.

Some leading members of Congress talked publicly Sunday about their hope that the president would
explain an endgame for American involvement in the eight-year war that includes how Afghans will
assume more of their security needs.

But more hawkish Republicans cautioned that setting a deadline for withdrawal could signal a lack of
resolve to allies, including Afghanistan and Pakistan.

``Talk of an exit strategy is exactly the wrong way to go,'' said Senator Jon Kyl, an Arizona
Republican. ``I certainly hope the president doesn't do that, because all that does is signal to the
enemies and also to our allies, to the folks in Pakistan as well as the Afghanis, that we're not
there to stay until the mission is accomplished.'' He spoke on ``Fox News Sunday.''

Senior lawmakers also warned the White House on Sunday that its expected troop buildup in
Afghanistan would fail unless the governments of Afghanistan and Pakistan did more to combat
militants attacking American forces, a concern that administration officials concede is a major
vulnerability in President Obama's new war strategy.

``The key here is an Afghan surge, not an American surge,'' Senator Carl Levin, a Michigan Democrat
who heads the Armed Services Committee, said on CBS's ``Face the Nation.'' ``And if the president
lays out the case for why our combat forces that are going particularly to the south will increase
the speed-up of the Afghan Army, it seems to me that that would be very, very important.''

With the cost of the war rising, some Democrats have even talked of a surtax. And a Republican
senator, Richard G.~Lugar of Indiana, asked: ``If we were talking about several years of time, how
many more years beyond that? What is the capacity of our country to finance this particular type of
situation as opposed to other ways of fighting Al Qaeda and the war against terror?''

At West Point, Mr.~Obama was expected to describe commitments from Afghanistan's president, Hamid
Karzai, and specific benchmarks his government must meet: to crack down on corruption, deploy
well-trained Afghan troops and police officers, and focus on development in one of the world's
poorest nations. Mr.~Obama was expected to be far less specific about Pakistan, where Taliban
leaders are commanding operations across the border against American forces, and where Al Qaeda's
central leadership still lives.

``We agree that no matter how many troops you send, if the safe haven in Pakistan isn't cracked, the
whole mission is compromised,'' said one official who has participated in the debate over the
strategy. ``But if you make too many demands on the Pakistanis in public, it can backfire.''

The problems in Afghanistan have only been compounded by the fragility of Mr.~Obama's partner in
Pakistan, President Asif Ali Zardari, who is so weak that his government seems near collapse. On
Friday, Mr.~Zardari relinquished his position in Pakistan's nuclear command structure, turning it
over to the prime minister, in what appeared to be an effort to avoid impeachment or prosecution,
and retain at least a figurehead post.

On Sunday, one of the Obama administration's staunchest allies, Prime Minister Gordon Brown of
Britain, joined in the campaign to press Pakistan to step up attacks on Al Qaeda's leadership in
Pakistan's unruly tribal areas and other militant groups there. ``People are going to ask why, eight
years after 2001, Osama bin Laden has never been near to being caught,'' Mr.~Brown told Sky News,
``and what can the Pakistan authorities do that is far more effective.''

White House officials have said relatively little about the Pakistan side of the administration's
evolving war strategy, in part because they have so few options and so little leverage. They cannot
send troops into Pakistan, and they cannot talk publicly about one of their most effective measures,
the Central Intelligence Agency's Predator drone strikes in the country.

``Everyone understands this is a complex, nuanced, critical relationship,'' said a senior American
official, who spoke on condition of anonymity because Mr.~Obama's review had not been announced.
``Everyone has their eyes open, and there are genuine concerns. But one focus now is on trying to
expand cooperation. The Pakistanis are doing some positive things in the tribal areas. That presents
opportunities on which to build.''

Mr.~Obama's advisers previously signaled that the president wanted to outline, as he had before,
expectations for the Afghan government. This time, they said, the goals would be more explicit and
demanding, aimed at improving governance and curbing corruption.

But the advisers have been debating whether to put deadlines on those benchmarks, like the pace of
training Afghan security forces to defend their country.

Gen.~Stanley A.~McChrystal, the top NATO and American commander in Afghanistan, is expected to
testify about Mr.~Obama's new strategy on Dec.~8 to the Senate and House Armed Services Committees
in Washington, the official said. His appearance is expected to follow Congressional testimony later
this week by Defense Secretary Robert M.~Gates, Secretary of State Hillary Rodham Clinton and Adm.
Mike Mullen, the chairman of the Joint Chiefs of Staff.

The administration has sought to build consensus among crucial allies to reach this point. In the
last two weeks, Mr.~Obama dispatched two top aides to Pakistan to deliver the same message: Keep the
pressure on.

In separate visits to Islamabad, the capital, the director of the Central Intelligence Agency, Leon
E. Panetta, and the president's national security adviser, Gen.~James L.~Jones, told Pakistani
officials that no matter how many more troops the president sent to Afghanistan, the effort would
fail unless Pakistan increased strikes against Al Qaeda's leadership and Mullah Muhammad Omar and
the leadership of the Afghan Taliban in the southern Pakistani city of Quetta, and the Haqqani
network, militants operating out of North Waziristan who have attacked Afghan and NATO targets in
eastern Afghanistan and Kabul, the Afghan capital.

\section{China's Impolitic Artist, Still Waiting to Be Silenced}

\lettrine{A}{i}\mycalendar{Dec.'09}{01} WEIWEI is perhaps China's most famous living artist and its
most vociferous domestic critic, titles of a sort this committed iconoclast disdains. Which is not a
bad thing, considering that recently, he very nearly lost them both.

Mr.~Ai was in Chengdu, the capital of Sichuan Province, preparing to testify at the trial of a
fellow political activist. ``By 3 a.m., we heard a very strong noise in the hallway, very brutal,
much like a Hollywood movie -- knocking on every door: 'Open it up -- we are the police!' '' he
said.

``They kicked open the door. I said, 'How do I know you are the police?' They said, 'I'll show you,'
and punched me here.'' Mr.~Ai pointed to the right side of his forehead. ``It was a very solid
punch.''

A month later, at an art exhibition in Munich, Mr.~Ai went to a doctor with a pounding headache and
was rushed into surgery to drain a pool of blood from his brain.

Mr.~Ai nearly died. Three months later, he says, his memory still fails him. On the other hand, ``I
don't have so many good memories anyway.''

That seems an exaggeration. At 52, Mr.~Ai, a beefy, bearded man with an air of almost monastic
composure, is an international figure in the art world, successful beyond what anyone might have
predicted even a decade ago. He is a celebrated architect, a co-designer of Beijing's landmark
Bird's Nest Olympic stadium, an installation artist and a documentary filmmaker with a 100-member
staff.

Artistically, he can do almost anything he wishes, like personally shipping 16 40-foot containers,
including 9,000 custom-made children's backpacks, from Beijing for his recent exhibition in Munich.

Yet clearly, all is not rosy in Mr.~Ai's world. In one of his early acclaimed works, a series of
three photographs called ``Dropping a Han Dynasty Urn,'' he dispassionately shatters a priceless
ancient Chinese vase, striking a theme -- destruction and recreation -- that runs through much of
his art.

Other works employ Ming and Qin period urns, furniture and architecture, assembled into haunting new
creations, or painted over, Warhol-style, with the Coca-Cola logo, or speared by wooden beams. A
series of photographs depicts global icons -- the Forbidden City, the White House, the Eiffel Tower
-- interrupted by Mr.~Ai's hand, middle finger raised.

Then there are his politics, an in-your-face criticism of China's leaders that, given Beijing's
limited tolerance for dissent, seems almost suicidal. Long before the Olympics, Mr.~Ai disavowed his
role in designing the Bird's Nest, saying the government had transformed the Olympics into a
patriotic celebration instead of using them to create a more open society.

In a 90-minute interview in his minimalist studio in north Beijing, Mr.~Ai called the government
unimaginative, prevaricating, suspicious of its own people and utterly focused on self-preservation.

``They don't believe in liberty. They don't believe in China before the Communists,'' he said.
``There is only one simple, clear task: to protect their control, to maintain their governing. Which
is such a pity.''

All of this he has said many times before. China's nationalists often accuse him of shilling for the
West, and in fact, Mr.~Ai ended his chat with a plea to President Obama to call for greater freedom
in China, saying ``we still need the moral support of the Western leaders'' to press for more
uncontrolled space in a still-closed society.

With or without help, Mr.~Ai is pressing hard. His most provocative art, as well as his latest
cause, concerns the question of why the May 2008 Sichuan earthquake killed thousands of children in
their classrooms -- and why the government has refused to give the public an official explanation.

After the quake, Mr.~Ai used the Internet to assemble scores of volunteers who combed the disaster
area, compiling a list of more than 5,000 dead children, organized by age and school, that now
covers one wall of his studio. ``The picture became clear. All of them belonged to about 20 schools,
and those schools, the buildings collapsed to dust,'' he said. ``Why did those buildings collapse,
and the ones next to it are standing?''

THE citizens' inquiry has produced a detailed list of questions, sent to government agencies, which
were supposed to be answered by this past Tuesday under law, but have yet to be addressed. In
December it will yield a documentary film on the disaster.

In Munich, the inquiry produced Mr.~Ai's most arresting work of art to date: those 9,000 children's
backpacks, covering one exterior wall of the Haus der Kunst. Against a blue background, colored bags
form the Chinese characters for the message, ``She lived happily on this earth for seven years,'' a
quotation from a mother of one earthquake victim.

The Munich exhibit is titled ``So Sorry,'' a caustic comment on the government's near-silence on the
schools disaster.

Mr.~Ai's beating in August occurred as he was preparing to testify at the trial of Tan Zuoren, a
Sichuan writer and activist who was trying to investigate the same issue. Mr.~Tan was accused of
inciting the subversion of state power. Mr.~Ai was blocked from testifying at his trial, which has
yet to produce a verdict. This week, though, another activist, Huang Qi, received a three-year
sentence for encouraging parents to press their grievances.

A disquieting sense of foreboding accompanies these jousts with the all-powerful state.

Ai Weiwei's father, Ai Qing, was both an artist and one of China's most revered contemporary poets,
who as a young man studied Baudelaire and Mayakovski in Paris. When he returned to Shanghai in 1932,
the ruling Kuomintang party jailed and tortured him, calling him a leftist. It was right: in 1941,
Ai Qing joined the Communist Party.

BUT 17 years later, in the infancy of Mao's new People's Republic, he ran afoul of the Communist
Party for subtly criticizing its suppression of free speech. The party exiled him, first to
Manchuria, then to remotest northwest China; Siberia, essentially.

Mr.~Ai and his family lived in a hut dug into the ground. His job for the next 16 years was to clean
out the village's public toilets.

``He was 60 years old. He had never done physical work in his life and he had to start doing it,''
his son said. ``Every night, he comes home very, very dirty. But he says, 'For 60 years, I don't
know who cleans my toilets. So now I do something for them.'

``That's something I learned from him. He became very powerful in terms of his thinking. He made the
toilet so clean, he would see it as a work of art -- like a museum, like MoMA.''

The family returned to Beijing in 1976, with the end of the Cultural Revolution. In 1985, the elder
Mr.~Ai, now rehabilitated, would receive a literary award from President François Mitterrand of
France. His son, on the other hand, could hardly wait to flee China.

Young Ai Weiwei studied at the Beijing Film Academy but in 1981 left for the United States. In New
York, Mr.~Ai said, he was in the city's art scene, not of it. He held temporary jobs and moved 10
times, throwing out his canvases each time for lack of storage room.

When his father fell ill in 1993, he agonized over returning to his homeland, which harbored such
painful memories. But after 1989, and the silencing of protesters at Tiananmen Square, he had
decided that ``the world became different.'' And so he returned to China in 1993, reckoning that one
day he might face something like his father's fate.

Lately, there are indeed signs that the government is reaching its limit. His blogs on Chinese Web
sites, about issues political and otherwise, have been shut down. Someone has installed two video
cameras outside his studio. The police are said to be scrutinizing his finances, an ominous
development in a state where other political critics have been prosecuted for what appear to be
concocted fiscal misdeeds.

``He has never done anything illegal,'' said his lawyer and friend, Liu Xiaoyuan. ``But if he
continues on his current path, getting involved in some very high-profile cases, I will get worried.
Some government departments are already very annoyed about him.''

Mr.~Ai says he is ready for whatever comes. ``I came to art because I wanted to escape the other
regulations of the society. The whole society is so political,'' he said. ``But the irony is that my
art becomes more and more political.''

\section{HSBC Returns to Its Roots in Hong Kong}

\lettrine{H}{SBC}\mycalendar{Dec.'09}{01} is so tied to Hong Kong that when the bank's shares
dropped sharply earlier this year a local television commentator burst into tears while reporting
the news.

Never mind that for the better part of two decades, HSBC's focus has been less on its 144-year-old
roots in Asia and more on expanding into Europe and the United States -- a shift that produced a
sprawling behemoth\cite{behemoth} with branches in 86 countries.

Now, with huge subprime losses in the United States still a drag on earnings and demand for credit
declining in Europe, the London-based bank is returning its focus to Asia, looking to benefit from
rapid growth in China and other emerging markets. As a first step, its chief executive, Michael
F.~Geoghegan, will move to Hong Kong from London in February with the bank's head of strategy and
about a dozen staff members.

``The center of gravity has shifted significantly towards Asia,'' Christopher White, a fund manager
at Threadneedle in London, said of the global business growth. ``It's been a quantum shift and
Geoghegan moving over to Hong Kong is a statement of that.''

In fact, some critics say that HSBC waited too long to turn to Asia and that its return was too
timid\cite{timid}.

Knight Vinke, an asset management firm based in Monaco that owns a small stake in HSBC, pushed
management as early as June 2007 to abandon its ailing operation in the United States and focus on
Asia instead. Now, the investor says the bank should commit exclusively to the region, as demand for
credit in Western markets is unlikely to recover in the medium term.

``They have something in the market which no one else has: a huge heritage, a footprint, deep
connections in a way that Barclays, JPMorgan and Citi have not, and as an investor I'm looking for a
way to tap into that,'' Glen P.~Suarez, a Knight Vinke executive, said.

HSBC is not the only company eager to tap into the faster-growing Asian markets. Coca-Cola, Wal-Mart
and NYSE Euronext are among firms that recently said they planned to grow in the region. Morgan
Stanley plans to set up its own brokerage there, BBVA of Spain might increase its stake in China's
Citic bank, and Deutsche Bank and Credit Suisse recently moved global division heads to Asia.

A growing middle class, a rising number of high-net-worth individuals and relatively underdeveloped
banking products make China an attractive hunting ground for banks. HSBC's own economists in October
raised the forecast for China's economic growth this year to 8.5 percent from 8.1 percent. HSBC
might be the biggest foreign bank in China but it has only a 1.3 percent share of the highly
fragmented banking market, according to figures by KPMG and Goldman Sachs.

But few banks have the strong ties to local businesses and a rich heritage in the region HSBC can
point to. Founded in 1865 as the Hongkong and Shanghai Banking Corporation by Thomas Sutherland, a
Scot who was then working for the Peninsular and Oriental Steam Navigation, it established itself as
the main financier for British shipping and trade in opium, silk and tea. It handled China's first
loan in 1874 and grew so important to Hong Kong that many consumers hold their entire pension in
HSBC shares.

It was a good business for the better part of a century. But in the 1980s, when financial markets in
the West were booming, HSBC decided to expand into the United States and Europe. It bought Marine
Midland Bank in the United States and, after doubling its assets with the purchase of Midland Bank
of Britain in 1992, HSBC moved its headquarters to London from Hong Kong.

Still hungry for market share in the American market, HSBC bought the subprime mortgage lender,
Household International, for \$14.2 billion in 2003, a purchase it now says it regrets. It was that
business that made HSBC the first major bank in 2007 to sound the alarm about mortgage market losses
after its own piled up.

Almost three years later, provisions for bad loans at the unit are finally starting to decline as
HSBC winds down most of the business. Over the last two years, HSBC has relied on its earnings from
emerging markets to compensate for losses in the United States.

Along the way, it was supported by Hong Kong investors like Li Ka-shing, Hong Kong's richest man,
who helped HSBC raise \textsterling12.5 billion, or \$18.5 billion, in a rights issue earlier this
year while many rivals had to accept government succor.

HSBC already generates 52 percent of its earnings from Hong Kong and the rest of Asia-Pacific but
the company plans to increase that share to 60 percent in the short term.

The bank's leadership denied that the Hong Kong move meant it was turning its back on the United
States, Europe or its London home, which will remain its global headquarters, at least for now.

Douglas Flint, the chief financial officer, said the bank was unlikely to again try to grab domestic
market share in the United States. Instead, it will seek to generate about 10 percent of earnings
from its American business, which is currently unprofitable, by using its large Asian operation to
tap into huge capital flows between the United States and China.

``The incremental demand is increasingly going to be out of Asia,'' Mr.~Flint said in an interview
on Wednesday. ``Our differentiator to other banks is not just that we are big in some emerging
markets, it's that we connect those markets to each other and the rest of the world.''

Mr.~Geoghegan, 56, will be the bank's first chief to be based in Hong Kong in 16 years, but he knows
the Asian market well. He joined HSBC 36 years ago and spent eight of those years in Asia.
Mr.~Geoghegan's wife will move with him to Hong Kong while his two sons, both university students in
Britain, will stay behind. Mr.~Geoghegan plans to spend about one week a month in Britain.

Analysts disagree whether his move will turn out to be more than just symbolic but said it would
probably help the bank with its plan to list shares on the Shanghai stock exchange, which would give
HSBC access to local currency needed for its expansion. HSBC, whose shares are already listed in
London and Hong Kong, plans to hire about 1,000 employees in China next year, open branches and
invest in technology.

Over the last 12 months, its shares have performed better than those of many rivals but have lagged
behind those of BNP Paribas of France, Banco Santander of Spain and Standard Chartered, another
London-based bank that generates most of its profit from emerging markets.

Even as Mr.~Suarez of Knight Vinke expresses concern that HSBC is leaving itself too yoked to the
moribund American market, he also worries about the rapid growth of bank loans in China. He
acknowledges that the billions of dollars the Chinese government pumped into its economy to cushion
the effects of the global economic downturn could pose a problem in the short term by creating a
credit bubble.

In fact, HSBC already started to notice how increased lending by Chinese banks has started to hurt
its loan margins. Falling interest rates have narrowed the difference between what HSBC pays on
deposits and what it can charge for loans.

Encouraged by the government to increase lending to offset the drag on the Chinese economy from
falling exports, Chinese banks lent more money in the first seven months of this year than in the
two previous years combined.

``Whenever people borrowed that much over that short period of time it ended in tears,'' Colin
Morton, a fund manager at Rensburg in Leeds, England, said. But in the long term, he said, exposure
to emerging markets will be beneficial for HSBC.

``The only problem is that it's already reflected in their share price,'' he said. ``The stock is
not cheap anymore.''

\section{In Pursuit of Loyalty}

\lettrine{W}{hen}\mycalendar{Dec.'09}{01} Greg McHale checks into his hotel room after a day of
business travel, he expects what he calls the ``wonderful and bizarre,'' namely complimentary
Snickers bars, Diet Pepsi and, sometimes, a compact disc of his favorite electronic dance music.

For Kimpton Hotels, it is a small price to pay for the loyalty of someone like Mr.~McHale, a Web
entrepreneur who spends 50 or 60 nights a year on the road. And for Mr.~McHale, the personal touches
-- part of the hotel chain's loyalty program -- make it worth his while to seek out Kimpton's
hotels.

``The level of personal attention really blows me away,'' said Mr.~McHale, founder and chief
executive of Good2gether, which connects nonprofit organizations with donors and volunteers. ``So if
there's a Kimpton in town, that's where I'll stay.''

Not all hotels go to such lengths to please their guests, but this year most are stretching their
creativity to attract and, perhaps more important, retain guests. Hotels have been particularly hard
hit by the drop in business travel, and brand loyalty has often given way to practical cost concerns
as companies have cut expenses.

Only 36 percent of business travelers said they were brand loyal this year, compared with 42 percent
two years ago, according to Henry H.~Harteveldt, a travel analyst for Forrester Research. ``And 2010
is likely to be more difficult for hotels because companies are telling their employees that every
penny saved means fewer people laid off or fewer cuts in pay.''

Hotels are responding by offering free nights, upgrades and loyalty points. Many hotels, especially
the high-end chains, are introducing twists to cultivate customers. Amenities including free
breakfasts, no-fee Internet connections, late checkouts and paid parking are being
bundled\cite{bundle} in business traveler packages. Marriott Hotels, for instance, calls its package
``Business Boost,'' while Hyatt Hotels has ``Business Plan'' and Sheraton Hotels ``Road Warrior.''

As part of its package, Hilton's Conrad Chicago Hotel is giving guests their choice among
best-selling books, and a personal shopper is available to help select gifts for those left at home.

Beyond packages, some hotels are trying to make stays more enticing by reducing fees for the
minibar, subsidizing some meals, offering free in-room spa services or free dry cleaning. The hotel
industry is trying to keep room rates stable, which is not easy. The average occupancy rate in
October was down 6.2 percentage points to 58.1 percent, and per-room revenue dropped 13.8 percent to
\$57.57 from the year before -- the worst numbers in more than two decades, according to Smith
Travel Research, in Hendersonville, Tenn.

To try to hold the line on rates, hotels are offering guests more for their money.

``It's about offering added value rather than lowering rates,'' said Sam Shank, chief executive of
DealBase.com, an online search engine for hotels. ``When travel fell after 9/11, hotels dropped
their rates and it took a while for them to bring those prices up again. They don't want to go
through that again.''

Many corporations have room rate arrangements with hotel chains, but their employees typically claim
the loyalty points for personal use, especially for long weekends or upgrading to concierge floors,
where they can have a nicer breakfast, access to snacks and, in the evening, a glass of wine with
hors d'oeuvres.

Loyalty points are a major selling feature for many travelers, said Don Berg, vice president for
loyalty for Intercontinental Hotels Group, the world's largest hotel operator whose brands include
Crowne Plaza and Holiday Inn. The group has four million club members.

About 90 percent of the points are redeemed for personal use, he said. The hotel group, taking a
page from the American Express and Visa rewards programs, also offers members special access to
concerts and sporting events. Starwood Hotels has a similar program.

``People have guilt over being away from home and family, and this is guilt-free currency to make up
for that,'' Mr.~Berg said. ``No expiration on our loyalty points is, by far, our most popular
feature.''

William R.~Snider, a Houston software consultant, was able to use his loyalty points from Holiday
Inns to indulge\cite{indulge} his love of baseball. He used his points to bid on, and win, World
Series and All-Star game packages that provided accommodation, meals and transportation and also
allowed him to mingle with players.

``I had a blast,'' he said.These awards make me want to stay at Holiday Inns, if at all possible.''

Among the most inventive in catering to customers are the high-end properties. Four Seasons Los
Angeles at Beverly Hills, for example, will store a guest's suitcase between visits, the hotel's
general manager, Mehdi Eftekari, said. ``So if you are traveling between Los Angeles and New York or
London, you always come back to freshly washed and ironed clothes packed away in your suitcase.''

An array of exercise gear, including socks and shoes, is available to guests so they do not have to
worry about smelly clothing, he said.

Kimpton offers specially prepared dinners for its most frequent guests, including one recently in
Manhattan for top-tier female travelers. The chain has also introduced weekend trips like the one in
October for its most frequent travelers and their spouses, in Oregon's wine country. The
wine-tasting getaway came with meals made by Kimpton chefs, and a balloon ride over the vineyards.

Paul Seus, a management consultant from Chicago who attended the Oregon weekend with his wife, Amy,
said Kimpton's special treatment cemented his loyalty.

``Kimpton called me and asked me if I would like to do something special,'' Mr.~Seus said.

``I've traveled my whole career, and I used to stay, well, wherever,'' he said. ``Now I'll only stay
somewhere else if I can't find one of their hotels.''

\section{In Japan, an Odd Perch\cite{perch} for Google: Looking Up at the Leader}

\lettrine{I}{n}\mycalendar{Dec.'09}{01} 2001, a fledgling\cite{fledgling} Internet company named
Google opened its first overseas office in Japan, eager to tap a huge technology market.

But after eight years, Japan is one of a few major countries Google has yet to conquer. The Web
giant still trails far behind Yahoo Japan, the front-runner here, operated by the Japanese
telecommunications giant Softbank.

In a reversal of the rivalry in the United States, Yahoo Japan dominates Japan's Web search market
with 56.5 percent of all queries, according to the Internet research company, GA-Pro. Google, at
33.7 percent, is a distant second.

Unaccustomed to being second, Google is bending some of its most time-honored traditions in a
renewed push into the Japanese market. Earlier this year, Google's splash page for Japan abandoned
the company's classic spare design and added links to YouTube, Gmail and other services -- an
attempt to lure Japanese users who favor sites decorated with a cacophony of text and graphics.

And in a first for Google, which is based in Mountain View, Calif., it initiated branding ads for
Japan and staged attention-grabbing publicity stunts, including one in which it invited passers-by
to float into the air with the help of 2,500 balloons.

Google's dogged interest in Japan has partly to do with sheer size. Japan is one of the world's most
wired countries, with more than 90 million regular Internet users -- of which three-quarters use
fast broadband connections and two-thirds also log in from cellphones.

And despite a sluggish economy, Japan's 6.6 trillion yen (\$77 billion) advertising market remains
the world's second-largest, one that an increasingly global advertising force like Google cannot
afford to ignore.

``Japan is absolutely a key market for Google,'' said Koichiro Tsujino, president of Google Japan.
Every day, for example, Japanese view 10 million clips on YouTube, Google's video-sharing site --
and that is just from their cellphones, making them the world's most avid adopters of video
on-the-go. ``Japan leads the world in many ways,'' he said.

That Japanese propensity\cite{propensity} to try new things is the other reason Google is intent on
staying put in Japan. Over the years, Japan has become a testing lab for many of the Web giant's
cutting-edge new ideas, especially in mobile technology. Google's Tokyo-based programmers, immersed
in Japan's mobile and Web culture, have become a valuable source of ideas for the entire company.

Overseas markets now account for half of Google's revenue, and the company is becoming more keenly
aware of the need to tailor its services to local markets, as well as the advantages of absorbing
ideas from outside the United States, company executives say. ``Japan made us realize that
non-U.S.~ideas can go global,'' David Eun, a vice president for Google, said on a recent trip to
Japan, where he closed deals with two Japanese broadcasters to allow YouTube to run some of their
content.

Google Japan's offices occupy several floors in a skyscraper in Shibuya, a Tokyo neighborhood
popular with start-ups that is also a hangout for the city's hippest teenagers. Minutes away from
where Google developers work, young Japanese perch on sidewalks, playing with their Web-enabled
cellphones, thumbs flying and eyes glued to the tiny screens.

But most of those trendsetters do not regard Google as being very Japanese -- a big headache for the
company. Google has never been able to overcome Yahoo's advantage as the first Web-based search
engine. And although 35 percent of Yahoo Japan is owned by Yahoo in Sunnyvale, Calif., it is viewed
as a local company.

``Yahoo Japan is a Japanese company, and most of their employees are Japanese people who fluently
understand how the Japanese mind-set and business work,'' said Nobuyuki Hayashi, a technology
analyst. ``But Google's still a foreigner who's learned how to speak some Japanese.''

Popularizing Google in Japan has been fraught with 21st-century versions of the cultural mishaps
that have long plagued American companies here. In May, Google was forced to reshoot its entire
``Street View'' image stock in Japan -- with a camera positioned to capture views 15 inches lower --
after intense criticism that the service peeked over fences and into people's homes, invading
privacy. The narrower width of Japan's roads made the service especially intrusive, bloggers fumed.

Google Earth also came under fire after posting historical maps that detailed locations of former
communities of an ``untouchable'' caste, still a sensitive topic in Japan. Human rights advocates
were furious that the maps could be used to identify families that had lived in the low-caste
neighborhoods.

But Google keeps trying. After studying feedback from Japanese users, developers designed Google's
maps service here so that a query led users to the town's train station or bus terminal, not the
center of town as it would in the United States, reflecting the way the Japanese, heavily reliant on
public transportation, think of their personal geography.

Programmers based in Tokyo have proposed and developed a line of services and functions, including
``emoticons'' for Gmail -- a particular Japanese obsession -- and a function allowing users to add
photos to Google Maps. It created ``Spellmeleon,'' to correct misspelled queries. It took developers
based in Tokyo to realize that non-native English speakers, who might not be very good spellers of
English words, could use a little help with queries.

``Part of our job is to think specifically about the Japanese market,'' said Kentaro Tokusei, group
product manager at Google Japan. ``We find whatever we build works globally, too.''

Some services in Japan offer a glimpse into the future. The Japanese version of Google's
photo-sharing service, Picasa, offers quick response, or Q.R., bar codes that contain Web address
information. Scanning a Q.R. bar code with a Japanese cellphone takes the user to a Web site to view
an online photo album.

Japan has been an especially important market for YouTube, with viewers here making up the site's
biggest audience outside the United States. The site's big presence in Japan has put developers here
at the forefront of crucial projects -- for example, a recently announced feature that will bring
text captions to many videos on the site, linked with automatic translation into 51 languages.

The captions will go a long way toward helping videos go viral across language divides, said Hiroto
Tokusei, YouTube product manager in Japan and Kentaro Tokusei's younger brother. (The Tokusei
brothers, both Stanford graduates with experience in Silicon Valley, were brought to Google Japan
with an eye to localizing Google's products while keeping Google at the cutting edge of innovation.)

Next month YouTube will also start a mobile version of its ``Click-to-Buy'' feature, which
identifies songs used in video clips, then lets users download them to their cellphones for use as
ring tones.

``To have an audience so obsessed with video and TV, and with access to broadband, means Japan is
the perfect place to experiment,'' Mr.~Tokusei said.

\section{China Pushes Back Against Calls for Yuan Rise}

\lettrine{C}{hina}\mycalendar{Dec.'09}{01} pushed back Monday against calls to let its currency
rise, with Prime Minister Wen Jiabao warning that an appreciation in the yuan could hobble Chinese
growth.

Speaking to reporters after a summit meeting with the European Union, Mr.~Wen said the demands being
made of Beijing to let the yuan strengthen were not fair.

He reaffirmed China's determination to take its own, gradual steps with regards to the currency but
said that for now the yuan, also known as the renminbi, would be kept steady.

``In this international financial crisis of a kind rarely seen in history, maintaining the basic
stability of the renminbi exchange rate has benefited China's economic development and benefited
world economic recovery,'' Mr.~Wen said.

``Now some countries, on the one hand, want the renminbi to appreciate but, on the other hand,
engage in brazen trade protectionism against China,'' he said. ``This is unfair. In fact, it amounts
to restricting China's development.''

The E.U. had pushed China to rethink its currency position at the summit meeting, saying the weak
yuan, combined with the weak dollar, was hurting European exports. But Mr.~Wen swatted away the
bloc's complaints, defending the policy and attacking critics as promoters of a dangerous
protectionism that threatened the global economic recovery.

The United States in particular has taken a number of steps recently to slow a surge in low-priced
imports from China in sectors including steel, tires and paper. ``Faced with the present complex
economic conditions, we must appropriately handle trade friction and not engage in trade
protectionism,'' Mr.~Wen said.

The contentious question of the yuan's exchange rate has dominated two days of talks with an array
of high-ranking E.U. policy makers in this eastern Chinese city. The E.U. is China's biggest market,
absorbing 20 percent of its exports, and runs a large bilateral trade deficit with China.

The three top economic officials representing the 16 E.U. members that use the euro pleaded Sunday
for a renewed gradual rise in the yuan, which China has virtually pegged to around 6.83 per dollar
since July 2008 to help its exporters weather the global credit crunch.

The European Commission president, Jos\'e Manuel Barroso, also pressed those demands Sunday, telling
Mr.~Wen that the low value of the yuan was hurting parts of the E.U. economy.

But in his remarks to reporters on Monday, Mr.~Barroso did not mention the yuan, and he abruptly
cancelled another news conference about the E.U.-China summit.

A joint statement issued by the two sides was also devoid of any comments on the currency dispute.

While Mr.~Wen refused to yield ground on the yuan, he struck a conciliatory tone about another issue
important to E.U. leaders: global warming.

Mr.~Wen said that China would deliver on a promise to curb carbon dioxide emissions, the main
greenhouse gas emitted by burning fossil fuels.

China's vow last week to cut ``carbon intensity'' -- the carbon dioxide released in generating each
yuan of economic output -- by 40 percent to 45 percent by 2020, compared with 2005 levels, was an
earnest pledge, Mr.~Wen said.

But he also restated China's position that developed countries must lead the way in climate change
talks starting next week in Copenhagen by offering big cuts in carbon emissions. Advanced economies
must also make financing and technology available to developing countries so that they can tackle
global warming, Mr.~Wen added.

\section{Open Source as a Model for Business Is Elusive\cite{elusive}}

\lettrine{I}{n}\mycalendar{Dec.'09}{01} many ways, MySQL embodies the ideals of the populist
software movement known as open source, in which a program's creator releases it to the world free
of charge, and legions of volunteers contribute improvements that are also freely shared.

The start-up company came out of nowhere, building a database application beloved by vibrant, young
Internet companies. Logging in from homes scattered around the globe, its workers seemed more a part
of a virtual commune than a corporate monolith, and they relished taking on proprietary software
giants like Microsoft.

But like most open-source companies, MySQL's sales, tied to support deals, never matched the
astronomical number of downloads for its product, about 60,000 a day. In January 2008, the founders
decided to sell the company for \$1 billion to Sun Microsystems. And this year, Sun agreed to sell
itself to Oracle, which makes database software aimed at larger companies and tougher jobs, for
\$7.4 billion.

Now, disagreement over the value of MySQL -- both as a stand-alone entity and as part of a big
company -- lies at the heart of a bitter public battle between Oracle and the European Union over
the Sun acquisition. The fight illuminates a larger truth about open-source companies: their
societal and strategic importance far exceeds their financial value as operating businesses.

European regulators view MySQL as sort of a database of the people, a low-cost alternative to
Oracle's costly proprietary products. The regulators worry that Oracle may stop improving MySQL in
favor of protecting its core traditional products, and customers will lose an important option in
the database market.

``In the current economic context, all companies are looking for cost-effective I.T. solutions, and
systems based on open-source software are increasingly emerging as viable alternatives to
proprietary solutions,'' said the European Commission's competition chief, Neelie Kroes, in a recent
statement. ``The commission has to ensure that such alternatives would continue to be available.''

Oracle, meanwhile, insists that it will continue to develop MySQL and other Sun technologies.
Oracle's chief executive, Lawrence J.~Ellison, contends that MySQL serves a different part of the
database market than Oracle's main products do -- an assessment supported by many analysts. One main
incentive for Oracle to keep improving MySQL is that the program serves as a bulwark against
Microsoft's SQL Server database, which challenges Oracle's products on the low end.

``The commission's statement of objections reveals a profound misunderstanding of both database
competition and open source dynamics,'' Oracle said in a statement.

To Ms.~Kroes's point, there is an open-source alternative, and usually a pretty good one, to just
about every major commercial software product. In the last decade, these open-source wares have put
tremendous pricing pressure on their proprietary rivals. Governments and corporations have welcomed
this competition.

Whether open-source firms are practical as long-term businesses, however, is a much murkier
question.

The best-known open-source company is Red Hat, which produces a variant of the Linux operating
system for server computers. Like most of its peers, Red Hat offers a free version of its base
product and relies on selling support services and extra tools for revenue. In its last fiscal year,
which ended in March, the company's revenue rose 25 percent to \$653 million, and it reported net
income of \$79 million.

But Red Hat is a rare case. ``There's only one company making real money out of open source, and
that's Red Hat,'' said Simon Crosby, the chief technology officer at Citrix Systems, which acquired
the open-source software maker XenSource for \$500 million in 2007. ``Everyone else is in trouble.''

The enduring appeal of open-source software revolves more around its disruptive nature than
blockbuster sales.

As long as there has been software, there have been some people eager to share and improve it for
the common good. The rise of the Internet made such sharing easier than ever, enabling people the
world over to work together on projects outside the confines of a formal corporate structure.

Open-source software has thrived and played a prominent role in the building of the Internet's
infrastructure. Many companies rely on Linux-based computers and Apache Web server software to
display their Web pages. Similarly, the Mozilla Firefox Web browser has emerged as the most
formidable competitor to Microsoft's Internet Explorer.

The grass-roots nature of open source has led advocates to view the projects as a populist foil to
proprietary software, where a company keeps the inner workings of its applications secret.

But in the last decade, open-source software has become more of a corporate affair than a people's
revolution.

In some cases, dominant technology companies have used open-source projects as pawns. Google, for
example, has needled Microsoft by providing financial support to the nonprofit Mozilla Foundation,
which oversees of the development of Firefox. I.B.M. has been a major backer of Linux, helping to
raise it as a competitor to Microsoft's Windows and other proprietary operating systems.

Many of the top open-source developers are anything but volunteers tinkering in their spare time.
Companies like I.B.M., Google, Oracle and Intel pay these developers top salaries to work on
open-source projects and further the companies' strategic objectives.

In the last three years, there have been five big acquisitions in which a major technology company
bought an up-and-coming open-source company for many times its annual revenue. Sun, for example,
bought MySQL for about 10 times its revenue, while Citrix bought XenSource for more than 150 times
its revenue, according to people familiar with the companies' sales.

Most recently, VMware, the leading maker of virtualization software, brought SpringSource for \$420
million, or about 20 times its annual sales.

``A lot of these guys were getting close to an I.P.O., but they elected to go the acquisition route
instead,'' said Michael Olson, the chief executive of Cloudera, an open-source start-up. ``A lot of
open-source firms are one-product companies, and it's hard to build a long-term, successful business
that way.''

The larger technology companies have tended to buy these one-trick ponies for strategic purposes.
With its core server business declining, Sun hoped it could piggy-back on MySQL's momentum with
Internet companies. In SpringSource, VMware acquired a company that had cultivated deep interest
with software developers and helped VMware diversify beyond its virtualization roots.

``VMware took into consideration that which money can't buy, which is a critical mass of adoption,''
said Peter Fenton, a venture capitalist at Benchmark Capital, who has been involved in some fashion
with many of the large open-source deals. ``SpringSource's main product was the equivalent of a
best-selling novel.''

Citrix took perhaps the biggest risk of all, paying a huge premium for XenSource in the hopes of
disrupting VMware's position in the virtualization market.

``I don't think Citrix would ever say it paid too much,'' Mr.~Crosby said. ``Citrix leaped to the
forefront of a whole software category. The ability to talk credibly about virtualization is worth a
huge amount in its own right.''

Meanwhile, the ideal of an independent open-source giant has faded.

Mr.~Fenton said that many open-source advocates had once hoped Red Hat would scoop up the top
open-source start-ups, keeping these crown jewels out of the hands of proprietary software makers.
But the company failed to go after other open-source companies initially and later could not afford
to pay the high prices offered by larger companies.

``You could make the case there was a window of opportunity to do that three to five years ago,''
Mr.~Fenton said. ``That opportunity has gone away. And it's hard to put Humpty Dumpty back together
again now.''

\section{Upstart Chinese Telecom Company Rattles Industry as It Rises to No.~2}

\lettrine{W}{hen}\mycalendar{Dec.'09}{01} the telecommunications company Telenor decided to upgrade
its mobile phone network in Norway, the job drew bids from the companies that had built its original
grid: Ericsson and Nokia Siemens Networks.

Instead, Telenor this month chose Huawei Technologies, a Chinese equipment maker, to build its
ultrafast network based on a technology called L.T.E., or Long Term Evolution.

Grabbing the contract in the heart of Scandinavia was the latest coup for Huawei, the rising star of
the mobile equipment industry. The company's low-cost, multipurpose networks have catapulted it to
No.~2 in the world, behind Ericsson.

``Huawei has established itself as a serious competitor,'' said Morten Karlsen Sorby, the head of
global business development at Telenor. ``They have been able in a very short period of time to
build the necessary competence and innovation.''

Huawei leapfrogged Alcatel-Lucent and Nokia Siemens in quarterly sales this year, according to
Dell'Oro, a research firm in Redwood City, Calif. But Huawei, a private company founded in 1988 by
Ren Zhengfei, a former officer of the People's Liberation Army, has fought the perception that it
has ties to China's government and military.

A 2007 report by the RAND Corporation, a policy research institute, for the United States Air Force
said Huawei ``maintains deep ties with the Chinese military, which serves a multifaceted role as an
important customer, as well as Huawei's political patron and research development partner.''

Security concerns motivated the Committee on Foreign Investment, a United States government panel,
to reject Huawei's joint \$2.2 billion bid with Bain Capital in 2008 for 3Com, a communications
equipment maker that also produces security software for the American military.

Huawei's private ownership -- which it has never disclosed in detail -- has brought it scrutiny this
year in India and Australia, as well, as it has bid on contracts.

Edward Zhou, the Huawei marketing director for Europe, dismissed questions surrounding the company's
ownership structure as ``market speculation,'' saying Huawei was owned by its 80,000 employees and
had no links to Chinese officialdom.

``No government or government-linked organizations have any ownership stake,'' Mr.~Zhou said.
``Huawei has no connection to the Chinese military, and none of our directors hold, or has held, any
positions with, or serves or has served as a consultant or adviser to, any Chinese government or
agency.''

A senior executive for one of Huawei's main rivals, who did not want to be identified because he was
not authorized to speak for the company, said he did not know of any government ties. But that does
not mean they do not exist, he added.

``They are very aggressive as a vendor,'' he said. ``But there is a definite lack of transparency,
which makes it very tricky for us.''

Questions surrounding its ownership have not prevented Huawei from supplying 36 of the world's top
50 mobile operators, including Telus in Canada and Cox Communications, Leap and Clearwire, a WiMax
operator majority owned by Sprint Nextel, in the United States.

All were attracted to innovative Huawei products like its SingleRAN, a wireless network that
transmits in second- and third-generation, as well as L.T.E., signals, saving operators the expense
of separate networks. In-Stat, a research firm in Scottsdale, Ariz., says Huawei was the first
company to produce an L.T.E. base station, the fastest in the industry, on a large scale.

``In whatever sector they are competing in, Huawei is a major player,'' said Jeff Heynen, an analyst
at Infonetics in Raleigh, N.C. ``They are now among the top three in almost every market. The future
for them continues to look bright.''

Supplying equipment to the three big Chinese operators-- China Mobile, China Telecom and China
Unicom -- helped Huawei almost double its share of the \$38 billion global mobile equipment market
to 20.1 percent in the third quarter, from 11 percent a year earlier, according to Dell'Oro. Huawei
moved past Nokia Siemens, at 19.5 percent, and trails Ericsson, which has 32 percent.

ZTE, China's second-largest maker of networking equipment, is also growing rapidly as Chinese
operators this year are building the first 3G networks on mainland China. Sales at ZTE, which is
publicly traded, rose 43 percent, to \$2.2 billion, in the third quarter and its profit rose 59
percent, to \$60 million.

Huawei is based in Shenzhen, like ZTE, in mainland China near Hong Kong. The company has grown
beyond China's borders, with foreign orders accounting for 75 percent of its \$18.3 billion in 2008
sales, up 43 percent from a year earlier. Huawei's profit rose to \$1.2 billion, from \$957 million
in 2007.

In Europe, where it has been selling equipment since 2000, Huawei supplies all of the Continent's
major operators, including Vodafone, Deutsche Telekom, France T\'el\'ecom and Telef\'onica. European
sales are about \$3 billion, Mr.~Zhou said.

Scott Siegler, an analyst at Dell'Oro, said inexpensive Chinese labor and Huawei's private
ownership, which spares the company from having to pay regular dividends, let it underbid rivals by
40 percent to 50 percent on average.

That has rattled the industry, Mr.~Siegler said. As Huawei's sales soared in 2008, Ericsson's rose
11 percent, while Nokia Siemens' dipped 0.7 percent. The two companies said they would cut thousands
of jobs this year.

``Huawei is crushing the market,'' Mr.~Siegler said. ``They continue to win market share and are
forcing others to compete on their terms.''

Rather than the underbidding, Mr.~Zhou said it was Huawei's unique products, like the SingleRAN,
short for single radio access network, that were winning contracts. That system lets mobile
operators cut operating costs, which account for 80 percent of expenses over the lifetime of a
network.

``From a cost perspective, we are not the lowest bidder in many projects,'' Mr.~Zhou said. ``But our
focus has been on lowering the total cost of ownership for the network as a whole.''

In Moldova, a former Soviet republic between Romania and Ukraine, France T\'el\'ecom chose Huawei
last year to build a new 3G network, the first in Europe that can transmit voice calls in
high-definition audio. Liudmila Climoc, the chief executive of Orange Moldova, said Huawei made the
best offer to Orange, France T\'el\'ecom's mobile unit.

``Their network was easier to roll out, and it was already H.D. voice-ready,'' Ms.~Climoc said.
``With Huawei, it was much easier to implement this feature. Everything was already in place.''

Mr.~Sorby, the Telenor business development chief, said his company had examined Huawei's private
ownership but decided it was not a factor in its deal. The Norwegian contract will place Huawei on
track to bid on a multibillion-dollar overhaul of Telenor's networks in Eastern Europe, Russia and
Central Asia.

``Their technology is innovative, and their bid provided the lowest cost of ownership for us,'' Mr.
Sorby said. ``And in our business, the lowest cost of ownership is key.''

\section{We May Be Born With an Urge to Help}

\lettrine{W}{hat}\mycalendar{Dec.'09}{01} is the essence of human nature? Flawed, say many
theologians. Vicious and addicted to warfare, wrote Hobbes. Selfish and in need of considerable
improvement, think many parents.

But biologists are beginning to form a generally sunnier view of humankind. Their conclusions are
derived in part from testing very young children, and partly from comparing human children with
those of chimpanzees, hoping that the differences will point to what is distinctively human.

The somewhat surprising answer at which some biologists have arrived is that babies are innately
sociable and helpful to others. Of course every animal must to some extent be selfish to survive.
But the biologists also see in humans a natural willingness to help.

When infants 18 months old see an unrelated adult whose hands are full and who needs assistance
opening a door or picking up a dropped clothespin, they will immediately help, Michael Tomasello
writes in ``Why We Cooperate,'' a book published in October. Dr.~Tomasello, a developmental
psychologist, is co-director of the Max Planck Institute for Evolutionary Anthropology in Leipzig,
Germany.

The helping behavior seems to be innate because it appears so early and before many parents start
teaching children the rules of polite behavior.

``It's probably safe to assume that they haven't been explicitly and directly taught to do this,''
said Elizabeth Spelke, a developmental psychologist at Harvard. ``On the other hand, they've had
lots of opportunities to experience acts of helping by others. I think the jury is out on the
innateness question.''

But Dr.~Tomasello finds the helping is not enhanced by rewards, suggesting that it is not influenced
by training. It seems to occur across cultures that have different timetables for teaching social
rules. And helping behavior can even be seen in infant chimpanzees under the right experimental
conditions. For all these reasons, Dr.~Tomasello concludes that helping is a natural inclination,
not something imposed by parents or culture.

Infants will help with information, as well as in practical ways. From the age of 12 months they
will point at objects that an adult pretends to have lost. Chimpanzees, by contrast, never point at
things for each other, and when they point for people, it seems to be as a command to go fetch
something rather than to share information.

For parents who may think their children somehow skipped the cooperative phase, Dr.~Tomasello offers
the reassuring advice that children are often more cooperative outside the home, which is why
parents may be surprised to hear from a teacher or coach how nice their child is. ``In families, the
competitive element is in ascendancy,'' he said.

As children grow older, they become more selective in their helpfulness. Starting around age 3, they
will share more generously with a child who was previously nice to them. Another behavior that
emerges at the same age is a sense of social norms. ``Most social norms are about being nice to
other people,'' Dr.~Tomasello said in an interview, ``so children learn social norms because they
want to be part of the group.''

Children not only feel they should obey these rules themselves, but also that they should make
others in the group do the same. Even 3-year-olds are willing to enforce social norms. If they are
shown how to play a game, and a puppet then joins in with its own idea of the rules, the children
will object, some of them vociferously.

Where do they get this idea of group rules, the sense of ``we who do it this way''? Dr.~Tomasello
believes children develop what he calls ``shared intentionality,'' a notion of what others expect to
happen and hence a sense of a group ``we.'' It is from this shared intentionality that children
derive their sense of norms and of expecting others to obey them.

Shared intentionality, in Dr.~Tomasello's view, is close to the essence of what distinguishes people
from chimpanzees. A group of human children will use all kinds of words and gestures to form goals
and coordinate activities, but young chimps seem to have little interest in what may be their
companions' minds.

If children are naturally helpful and sociable, what system of child-rearing best takes advantage of
this surprising propensity? Dr.~Tomasello says that the approach known as inductive parenting works
best because it reinforces the child's natural propensity to cooperate with others. Inductive
parenting is simply communicating with children about the effect of their actions on others and
emphasizing the logic of social cooperation.

``Children are altruistic by nature,'' he writes, and though they are also naturally selfish, all
parents need do is try to tip the balance toward social behavior.

The shared intentionality lies at the basis of human society, Dr.~Tomasello argues. From it flow
ideas of norms, of punishing those who violate the norms and of shame and guilt for punishing
oneself. Shared intentionality evolved very early in the human lineage, he believes, and its
probable purpose was for cooperation in gathering food. Anthropologists report that when men
cooperate in hunting, they can take down large game, which single hunters generally cannot do.
Chimpanzees gather to hunt colobus monkeys, but Dr.~Tomasello argues this is far less of a
cooperative endeavor because the participants act on an ad hoc basis and do not really share their
catch.

An interesting bodily reflection of humans' shared intentionality is the sclera, or whites, of the
eyes. All 200 or so species of primates have dark eyes and a barely visible sclera. All, that is,
except humans, whose sclera is three times as large, a feature that makes it much easier to follow
the direction of someone else's gaze. Chimps will follow a person's gaze, but by looking at his
head, even if his eyes are closed. Babies follow a person's eyes, even if the experimenter keeps his
head still.

Advertising what one is looking at could be a risk. Dr.~Tomasello argues that the behavior evolved
``in cooperative social groups in which monitoring one another's focus was to everyone's benefit in
completing joint tasks.''

This could have happened at some point early in human evolution, when in order to survive, people
were forced to cooperate in hunting game or gathering fruit. The path to obligatory cooperation --
one that other primates did not take -- led to social rules and their enforcement, to human altruism
and to language.

``Humans putting their heads together in shared cooperative activities are thus the originators of
human culture,'' Dr.~Tomasello writes.

A similar conclusion has been reached independently by Hillard S.~Kaplan, an anthropologist at the
University of New Mexico. Modern humans have lived for most of their existence as hunter gatherers,
so much of human nature has presumably been shaped for survival in such conditions. From study of
existing hunter gatherer peoples, Dr.~Kaplan has found evidence of cooperation woven into many
levels of human activity.

The division of labor between men and women -- men gather 68 percent of the calories in foraging
societies -- requires cooperation between the sexes. Young people in these societies consume more
than they produce until age 20, which in turn requires cooperation between the generations. This
long period of dependency was needed to develop the special skills required for the hunter gatherer
way of life.

The structure of early human societies, including their ``high levels of cooperation between kin and
nonkin,'' was thus an adaptation to the ``specialized foraging niche'' of food resources that were
too difficult for other primates to capture, Dr.~Kaplan and colleagues wrote recently in The
Philosophical Transactions of the Royal Society. We evolved to be nice to each other, in other
words, because there was no alternative.

Much the same conclusion is reached by Frans de Waal in another book published in October, ``The Age
of Empathy.'' Dr.~de Waal, a primatologist, has long studied the cooperative side of primate
behavior and believes that aggression, which he has also studied, is often overrated as a human
motivation.

``We're preprogrammed to reach out,'' Dr.~de Waal writes. ``Empathy is an automated response over
which we have limited control.'' The only people emotionally immune to another's situation, he
notes, are psychopaths.

Indeed, it is in our biological nature, not our political institutions, that we should put our
trust, in his view. Our empathy is innate and cannot be changed or long suppressed. ``In fact,''
Dr.~de Waal writes, ``I'd argue that biology constitutes our greatest hope. One can only shudder at
the thought that the humaneness of our societies would depend on the whims of politics, culture or
religion.''

The basic sociability of human nature does not mean, of course, that people are nice to each other
all the time. Social structure requires that things be done to maintain it, some of which involve
negative attitudes toward others. The instinct for enforcing norms is powerful, as is the instinct
for fairness. Experiments have shown that people will reject unfair distributions of money even it
means they receive nothing.

``Humans clearly evolved the ability to detect inequities, control immediate desires, foresee the
virtues of norm following and gain the personal, emotional rewards that come from seeing another
punished,'' write three Harvard biologists, Marc Hauser, Katherine McAuliffe and Peter R.~Blake, in
reviewing their experiments with tamarin monkeys and young children.

If people do bad things to others in their group, they can behave even worse to those outside it.
Indeed the human capacity for cooperation ``seems to have evolved mainly for interactions within the
local group,'' Dr.~Tomasello writes.

Sociality, the binding together of members of a group, is the first requirement of defense, since
without it people will not put the group's interests ahead of their own or be willing to sacrifice
their lives in battle. Lawrence H.~Keeley, an anthropologist who has traced aggression among early
peoples, writes in his book ``War Before Civilization'' that, ``Warfare is ultimately not a denial
of the human capacity for cooperation, but merely the most destructive expression of it.''

The roots of human cooperation may lie in human aggression. We are selfish by nature, yet also
follow rules requiring us to be nice to others.

``That's why we have moral dilemmas,'' Dr.~Tomasello said, ``because we are both selfish and
altruistic at the same time.''

\section{Developmental Psychologist Says Teenagers Are Different}

\lettrine{L}{aurence}\mycalendar{Dec.'09}{01} Steinberg, a developmental psychologist at Temple
University in Philadelphia, is one of the leading experts in the United States on adolescent
behavior and adolescent brain biology. Dr.~Steinberg, 57, has won the \$1 million Klaus J.~Jacobs
Research Prize, which will be awarded to him at a ceremony in early December in Switzerland. Here is
an edited version of two conversations with Dr.~Steinberg last month:

Q. YOU HEAR PARENTS SOMETIMES SAY, ``I'M LIVING WITH AN INSANE PERSON. MY CHILD IS A TEENAGER.'' ARE
THEY BEING HYPERBOLIC?

A. I'm not one of those people who labels adolescence as some sort of mental illness. Teenagers are
not crazy. They're different.

When it comes to crime, they are less responsible for their behavior than adults. And typically, in
the law, we don't punish people as much who are less responsible. We know from our lab that
adolescents are more impulsive, thrill-seeking, drawn to the rewards of a risky decision than
adults. They tend to not focus very much on costs. They are more easily coerced to do things they
know are wrong. These factors, under the law, make people less responsible for criminal acts. The
issue is: as a class, should we treat adolescents differently?

Q. IS THE CRIMINAL JUSTICE SYSTEM BEGINNING TO TAKE THESE DIFFERENCES INTO ACCOUNT DURING
SENTENCING?

A. It's been coming up in cases. I went to Washington in November to watch the oral arguments in two
related cases before the Supreme Court that ask: should someone who committed a crime as a teen be
subjected to life imprisonment without a chance for parole, ever?

With these cases, and another in 2005 where the high court threw out the death penalty for
adolescents, I was scientific consultant to the American Psychological Association on its amicus
brief. What we said in the death penalty case -- and now -- was that we have considerable evidence
showing that adolescents are different from adults in ways that mitigate their criminal
responsibility. But since 2005, there's been a lot of new scientific evidence supporting this
position.

Q. WHAT IS THE NEW EVIDENCE?

A. In the last five years, as neuroscience has moved forward with functional magnetic resonance
imaging and with research on animals, there have been dozens of new studies of adolescent brain
development. These show that the brain systems providing for impulse control are still maturing
during adolescence. Neuroscientists have shown that the part of the brain that improves most during
adolescence is the prefrontal cortex, which is involved in complicated decision-making, thinking
ahead, planning, comparing risks and rewards. And the neuroscientific research is showing that over
the course of adolescence and into the 20s, there is this continued maturation of this part of the
brain. So now, we have brain evidence that supports behavioral studies.

Moreover, we're seeing that behavior can change once the brain more fully matures. Take
thrill-seeking, for instance. What happens is that when people move out of adolescence, they become
less interested in it. For example, I can't stand riding on a roller-coaster now. I liked it as a
teenager. I can't stand driving fast now. I liked driving fast when I was a teenager. What has
changed? I'm not as driven today by this thrill-seeking sensation. And in our studies, we've shown
that there is a kind of normative decline in sensation-seeking after middle adolescence. A lot of
adolescent crime is driven by thrill-seeking.

Q. HOW DOES THIS NEW INFORMATION

lead to concluding that the courts shouldn't sentence some adolescents to life in prison without
parole?

A. Given the fact that we know that there will be a developmental change in most people, the science
says that we should give them a chance to mature out of it. No one is saying that kids who commit
horrific crimes shouldn't be punished. But most in the scientific community think that we know that
since this person is likely to change, why not revisit this when he's an adult and see what he's
like?

Q. DO YOU HAVE TEENAGERS AT HOME?

A. We have a son, Ben, who is now 25 and who works at Random House. He did something as a teenager
that led me to a whole program of research. He and some friends went to the window of a girl they
knew and inadvertently set off a burglar alarm. When a police squad car came, they panicked and
fled. When I found out, I said: ``Do you realize that you were running from armed police officers
who thought they were interrupting a break-in. What were you thinking?'' He said: ``Well, that's the
problem. I wasn't.'' I wondered: ``What goes on when kids are in a peer group that pushes them to
make bad decisions?''

Since then, we've had people of different ages come to the lab and bring two friends with them. We
give them computerized risk-taking tests while we image their brains. We compare brain activity when
individuals are watched by their friends and when they are alone. For the adults, the presence of
friends has no effect. But for adolescents, just having friends nearby doubles the number of risks
they take. We've found that a certain part of the brain is activated by the presence of peers in
adolescents, but not in adults.

Q. YOU ADVISED THE DEFENSE TEAM OF OMAR KHADR, THE YOUNGEST DETAINEE AT GUANT\'aNAMO BAY. WHY GET
INVOLVED IN THAT CASE?

A. Because he was 15 when he was captured in a safe house in Afghanistan, where he'd been sent by
his father, who was active in Al Qaeda. There was a battle in 2002 to take this house where American
troops died.

He was interrogated for many hours and admitted to having thrown a grenade that killed an American
soldier. He later recanted. I was asked by his Defense Department counsel to advise on whether what
he said during interrogation was reliable and his degree of culpability, if he did do it.

In my deposition, I said I don't know whether he did it or not, but there are studies that say that
adolescents are more likely than adults to give false confessions. There's the Central Park jogger
case, where it turned out a group of teenagers gave false confessions. Five were convicted. Several
years later, an adult murderer and rapist confessed to the crime.

Q. IT HAS JUST BEEN ANNOUNCED THAT YOU'VE WON THIS \$1 MILLION KLAUS JACOBS PRIZE. WHAT DO YOU
INTEND TO DO WITH THE MONEY?

A. I want to extend our work on adolescent development to teenagers in other cultures so that we can
determine whether the patterns are universal. There's a longstanding debate over how much of
adolescent behavior is biological or cultural. Perhaps this award will lead to more answers.

\section{Story? Unforgettable. The Audience? Often Not.}

\lettrine{I}{f}\mycalendar{Dec.'09}{01} a friend is someone who laughs at our stories, then a good
friend is one who enjoys them even the second time around. But anyone who gasps with delight on
hearing a story for the third time is faking it. Or, it's a relative: some poor nephew Will or aunt
Emily, sitting captive at the holiday table, being polite, perhaps covering a shudder of dread that
life is caught in some endless loop where the punch lines never change.

It is not an entirely irrational fear, either, according to new research published in the journal
Psychological Science.

``You hear people of all ages, not just elderly people, say, 'Stop me if I've told you this before,'
'' said Nigel Gopie, a postdoctoral fellow at the Rotman Research Institute, in Toronto, who has a
paper in the current issue of the journal on these memory lapses.

``We often have a hard time remembering who we told things to, and clearly it starts early.''

In their long study of memory, psychologists have made important distinctions between the short-term
and long-term varieties. They have documented crucial differences between explicit memories, like
for faces and vocabulary, and the implicit kind, like for driving skills. They have published
hundreds of studies on autobiographical memory, false memories and so-called source memory -- the
ability to recall where a fact was learned, whether from the radio or a book, from a work colleague
or the neighborhood gossip.

Yet they have paid little, if any, attention to what Dr.~Gopie and his co-author, Colin M.~MacLeod
of the University of Waterloo, in Ontario, call destination memory: about whose ears information has
landed on. While the source of remembered information can be crucially important (Did I read that in
The Onion or the daily newspaper?), so is its destination. Our stories, our jokes, our gossip form
an important part of our social identity, psychologists say. Repeating oneself is not only
embarrassing; it can be damaging, for diplomats, liars or anyone else trying to guard secrets,
personal or professional.

``I think people simply get a lot more practice monitoring the sources of information, asking
themselves and others, 'Where is that from?' '' said Morris Moscovitch, a psychologist at the
University of Toronto. ``Whereas, it's rare we get any feedback about'' whom we told.

The main finding by Dr.~Gopie and Dr.~MacLeod -- that destination memory is relatively weak -- helps
explain several embarrassing, and annoying, kinds of social interaction. In one experiment, they had
60 University of Waterloo students associate 50 random facts (a shrimp's heart is in its head; 8
percent of men are color blind) with the faces of 50 famous people, like Madonna, Wayne Gretzky and
Oprah Winfrey. Half of the students ``told'' each fact to one of the faces, reading it aloud when
the celebrity's picture appeared on a computer screen. The other half read each fact silently and
saw a different celebrity moments afterward.

The students then took a memory test. They chose from face-fact pairs: those which they remembered
from learning a fact, and those they remembered from reading facts out loud in the first phase of
the study. The students who simulated telling the facts did 16 percent worse on the test than the
students who were fed the facts while seeing celebrity faces. The study authors concluded that
outgoing information ``was less integrated with its environmental context -- i.e., the person --
than was incoming information.''

This makes sense, psychologists say, given what is known about attention: namely, that it is finite.
A person who is conveying information, even trivial facts, will devote some mental resources to
monitoring what is being said. Self-absorption is also a factor. In another study, Dr.~Gopie and
Dr.~MacLeod repeated the famous-face exercise, with one big difference. This time the facts that the
students simulated telling to celebrities were personal (``My zodiac sign is Pisces''). The result
was their destination memory worsened significantly.

``Now, the situation may be reversed entirely for highly emotional personal information,'' like
devastating personal anxieties, Dr.~Gopie said. ``That is, that people are in those cases very aware
of whom they told. We just don't know that yet.''

The results suggest nonetheless that some of people's most intricate, richly detailed stories -- the
most self-distracting to tell -- are at high risk for being met with rolled eyes that say, ``Been
there, heard all that.''

The tendency to blank on who-I-told-what may in fact reflect the workings of a healthy memory.
Psychologists have found evidence that when people reset a password or a new phone number for an old
friend, their brain actively suppresses the out-of-date digits. The old numbers are a competing
memory, and potentially confounding.

Reprised stories aren't always embarrassing or socially redundant, either. If they are repeated
often enough, they become ritual, or, over time, oral history, Dr.~Gobie suggests. Still, it is
telling that people who have the most invested in who hears what -- salesmen, lobbyists -- often
remind themselves whom they are addressing: ``Have I told you, Gail, about the special price we have
on laser printers?'' That may be sucking up, but it may also be a way of keeping tabs on where the
information is going.

That is precisely what the two researchers found in the final experiment reported in their paper.
Saying the recipient's name (``Oprah Winfrey, the United States Postal Service handles 40 percent of
the world's mail volume!'') increased the accuracy of their destination memory.

Researchers say that if destination memory proves significantly weaker in further studies, the next
step will be to find out when the risks of such lapses are highest and in whom. An improved
understanding of destination memory could help doctors detect age-related memory problems earlier,
for instance. It may also be relevant to some models of how memory works.

None of which will bail out the holiday raconteur, caught short in the middle of telling a rerun
story. Unless he or she can reshape that tale on the fly, and pass it off as oral history.

\section{Looking for Balloons and Insights to Online Behavior}

\lettrine{T}{he}\mycalendar{Dec.'09}{01} prize is \$40,000, and it goes to the first person or group
to determine the locations of 10 red balloons that can be anywhere in the continental United States.

The apparent frivolity of the challenge is only on the surface. This is not a game invented by some
eccentric Web Midas. The contest, which takes place on Dec.~5, is being sponsored by Darpa, the
Pentagon's research agency.

The goal is to learn more about social behavior in computer networks and how large
computer-connected teams use their resources and connections to compete.

There is also an invention being celebrated. Peter Lee, a computer scientist and one of the Darpa
directors organizing the contest, said Dec.~5 would be the 40th anniversary of the day when the
first four nodes of the Arpanet -- the experimental military-sponsored computer network that was the
forerunner of today's Internet -- were connected.

Darpa has previously sponsored three ``grand challenges'' in an effort to advance the technology for
autonomous vehicles. In the second one, in 2005, a Stanford University team won \$2 million when its
roboticized Volkswagen Touareg was the quickest to navigate a 131-mile course through California
desert.

The mission of the agency, created in 1958 after the Sputnik satellite's launching, is to guard the
country against technological surprise. But Darpa prompted concerns about privacy after the
Sept.~11, 2001, attacks when it created a program to use data-mining technologies to identify
potential terrorists.

Dr.~Lee said he was not certain what to expect in the tactics that teams might use to track down the
balloons, which will be visible from public roadways for a single day. Some groups are developing
software applications. Dr.~Lee said he also expected large teams of spotters and even the
possibility that some groups might use subterfuge like disseminating false information.

Other groups may try to pay for information, he said, noting that even during a brief experiment the
agency ran with a balloon near its headquarters, information on the location was offered for sale on
Craigslist.

Dr.~Lee said the agency would continue to pursue a number of large and small challenge-style
contests to foster what he described as new ways to tap into pools of talented individuals and
creative groups. Contestants from anywhere in the world may participate in this contest, he said,
and registration will stay open until the contest begins.

\section{The Doctors Were Real, the Patients Undercover}

\lettrine{I}{t}\mycalendar{Dec.'09}{01} had all the markings of a television detective show. Posing
as patients, three undercover observers got themselves admitted as patients to a locked psychiatric
ward to investigate conditions on the inside.

Each undercover patient had rehearsed an extensive back story, and the supposed family members who
visited them were professional actors. A remote team monitored the project via hidden cameras and
microphones from a command center in a nearby hotel.

The project, which took place this spring in De Gelderse Roos, a psychiatric complex about 40 miles
from Amsterdam, was not a sting operation. The staff was told there would be mystery shoppers, of a
sort, in the facility over a couple of months.

``We didn't go in there like cowboys,'' said Menko Soeters, a partner at Clearfields, a consulting
firm that developed the project with De Gelderse Roos. ``But we did use an unorthodox instrument for
psychiatric care.''

Surrounded by manicured greenery, the closed-off ward of the complex, known as De Riethorst, recalls
a suburban dental clinic, and its sunny gymnasium and carpeted hallways do little to suggest that it
houses up to a dozen acute psychiatric patients, many of whom are there involuntarily.

And that is why the undercover participants were all experienced psychiatric nurses. ``You couldn't
have done it otherwise,'' said Edo De Vries, the director of De Gelderse Roos, which released the
results of the project last summer.

It and Clearfields are working on a project for 2010, most likely to involve five to eight
psychiatric hospitals in the Netherlands.

Mr.~De Vries said the impetus for the project came in part from a pair of patient deaths last year
in psychiatric facilities in Amsterdam -- one involving a suicide, the other a man who choked on
food while locked in an isolation cell. Several managers and staff members were fired as a result,
and others were suspended.

``Of course, incidents can happen anywhere,'' Mr.~De Vries said. ``But what if there is something
structurally wrong that we don't know about? We have to be more transparent, and I think this method
is a good tool for that.''

Because this initiative was the first of its kind, the goal in Ede was to see whether such a project
was even practical.

The untested nature of the project meant the organizers prepared as meticulously as possible.

The undercover patients developed their fictive biographies in months of meetings with an acting
coach and a psychotherapist. ``Ronald,'' for example, was a middle-aged man with a history of
aggression problems, and after a supposed suicide attempt he was taken to De Riethorst by an actor
playing his brother. To make their stories believable, the patients memorized details about where
their children went to school or which supermarket they shopped at, and the psychotherapist advised
them on how to present their given mental illness convincingly.

To ensure their safety, the fake patients checked in via text message every three hours, and they
carried letters identifying them as plants. A code word (``fireplace'') was in place if they had to
communicate genuine distress to visitors. The visitors wore hidden cameras and microphones; the
undercover patients did not.

The warning to the staff about the undercover patients made the project more amenable to Martien
Opdam, a psychiatric nurse at De Riethorst who worked while the plants were there. Though he and his
colleagues did wonder who the plants might be, he said: ``You can't keep doing that for two months.
You go back to your routine.''

He said the experience was helpful. ``It taught me not to go too much on autopilot,'' he said.

Among the findings of the project were that patients frequently found it difficult to get
information on their treatment and medications and that the sound of a staff member's key chain
jangling could be jarring to an already anxious patient.

Malingering one's way into a psychiatric ward to report on conditions within has been a journalistic
staple since as far back as 1887, when Nellie Bly got herself admitted to the insane asylum on
Blackwells Island (now Roosevelt Island) in New York City. Similar endeavors have been the subject
of many books and television documentaries over the years.

The best-known scientific example remains the experiments by David Rosenhan, a psychologist. He and
seven other ``pseudopatients'' got themselves admitted to a dozen hospitals by pretending to hear
voices, and the study, published in Science magazine in 1973 as ``On Being Sane in Insane Places,''
is still widely seen as a critique of psychiatric diagnosis.

In the July issue of the journal Psychiatric Services, Arthur Lazarus, a psychiatrist with decades
of experience in the health insurance and pharmaceutical industries, wrote a commentary supporting
the use of mystery shoppers in mental health care, though not in in-patient settings and only if the
staff is informed.

Dr.~Lazarus suggested in an interview that perhaps doctors themselves could play the role of
undercover patients. ``I think there is an opportunity to get beyond the mere fact that the carpet
is dirty or the staff is rude,'' he said.

Likewise, Dr.~James Sabin, a Harvard Medical School professor as well as a medical ethicist for
Harvard Pilgrim Health Care, sees possibilities. ``It's technically much more challenging in a
psychiatric hospital,'' Dr.~Sabin said, ``but it's certainly doable.'' He said it could help doctors
``try to fathom what is it like for the people we're dealing with.''

Opponents of the idea think enough ways already exist. Melissa Miller, a licensed clinical social
worker based in Sarasota, Fla., said mystery shopping in mental-health settings was ``an intrusive
redundancy.''

``The field is really already well covered with ways to have checks and balances,'' Ms.~Miller said,
including surveys and exit interviews with staff members and patients.

She said that having snooping eyes in her practice would violate the implicit trust between doctor
and patient. ``I could imagine having to have my client sign a statement saying, 'I am not a mystery
shopper,' '' she said. ``Otherwise how can I be free to give my professional self to the therapy?''

Last year, the ethics council of the American Medical Association considered issuing guidelines for
the use of ``secret shopper patients,'' but the measure has not advanced.

Many medical professionals object to the use of real doctors and resources by fake patients. For the
project in Ede, the organizers had made provisions to relocate the undercover patients if there were
a sudden shortage of beds.

The potential in mystery shopping for reality television has also been explored.

Richard Bentall, a professor of clinical psychology at the University of Bangor in Wales, took part
in a 2008 BBC program called ``How Mad Are You?'' The program challenged medical experts to
distinguish between ``normal'' volunteers and those with a history of mental illness. Despite the
gimmicky nature of the show, Dr.~Bentall said he was ``cautiously in favor'' of projects like the
recent Dutch example.

``Having covert observation is going to provide you with information you probably wouldn't get in
any other way,'' he said.

But Dr.~Bentall also sees some irony in using proxy mental patients to illuminate the experiences of
real ones. ``Their stories are neglected,'' he said, ``and their understanding of how they got to be
in the hospital is not considered important.''

There are some signs of change in that respect. The National Alliance on Mental Illness, for
example, recently began using mental health care consumers to rate the Web sites and telephone
switchboards of various facilities.

As Dr.~Bentall put it: ``It's common sense, really. Involve the customers in the decision-making,
and you're going to get a better product.''

\section{Obama Adds Troops, but Maps Exit Plan}

\lettrine{P}{resident}\mycalendar{Dec.'09}{03} Obama announced Tuesday that he would speed 30,000
additional troops to Afghanistan in coming months, but he vowed to start bringing American forces
home in the middle of 2011, saying the United States could not afford and should not have to
shoulder an open-ended commitment.

Promising that he could ``bring this war to a successful conclusion,'' Mr.~Obama set out a strategy
that would seek to reverse Taliban gains in large parts of Afghanistan, better protect the Afghan
people, increase the pressure on Afghanistan to build its own military capacity and a more effective
government and step up attacks on Al Qaeda in Pakistan.

``America, we are passing through a time of great trial,'' Mr.~Obama said. ``And the message that we
send in the midst of these storms must be clear: that our cause is just, our resolve unwavering.''

The military escalation Mr.~Obama described and defended in his speech to a national television
audience and 4,000 cadets at the United States Military Academy here, the culmination of a review
that lasted three months, could well prove to be the most consequential decision of Mr.~Obama's
presidency.

In his 33-minute address, he sought to convince an increasingly skeptical nation that the resurgence
of the Taliban in Afghanistan and the continued existence of Al Qaeda across the border in Pakistan
-- what he called a ``cancer'' on the region -- were direct threats to the United States, and that
he could achieve the seemingly contradictory goals of expanding American involvement in the war even
as he sought to bring it to a close.

The scene in the hall was striking and somber: row after row of cadets, in their blue-gray uniforms,
listening intently to a strategy that could put many of them in harm's way. ``If I did not think
that the security of the United States and the safety of the American people were at stake in
Afghanistan, I would gladly order every single one of our troops home tomorrow,'' Mr.~Obama said.
``So no, I do not make this decision lightly.'' He called on foreign allies to step up their
commitment, declaring, ``This is not just America's war.''

He delivered a pointed message to Hamid Karzai, the president of Afghanistan, saying, ``The days of
providing a blank check are over.''

Addressing critics who have likened Afghanistan to Vietnam, Mr.~Obama called the comparison ``a
false reading of history.'' And he spoke directly to the American people about the tough road ahead.

``Let me be clear: none of this will be easy,'' Mr.~Obama said. ``The struggle against violent
extremism will not be finished quickly, and it extends well beyond Afghanistan and Pakistan. It will
be an enduring test of our free society, and our leadership in the world.''

With the economy weak and the issue of jobs foremost on Americans' minds, the president conceded
that the new strategy would carry an expensive price tag, which he put at an additional \$30 billion
in the first year.

Yet with some Democrats talking of a war surtax, Mr.~Obama offered no details of how he intended to
pay for his new policy, saying only that he was ``committed to addressing these costs openly and
honestly.''

White House advisers said they expected the administration would do so in the coming weeks, as
officials including Secretary of Defense Robert M.~Gates and Secretary of State Hillary Rodham
Clinton testify on Capitol Hill starting Wednesday.

The approach laid out by Mr.~Obama -- not so much a new strategy as a doubling down on the one he
embraced earlier this year -- incorporated the basic goals and came close to the force levels
proposed in the counterinsurgency plan that Gen.~Stanley A.~McChrystal, the top military commander
in Afghanistan, put forward in September.

In that report, General McChrystal said, in stark language, that unless significantly more troops
were sent, the war in Afghanistan was likely to be lost.

But by including an explicit timetable to begin a withdrawal, Mr.~Obama highlighted the seemingly
conflicting pressures defining the debate over how to proceed: to do what is necessary to ensure
that the region is not a launching pad for attacks on the United States and its allies, and to
disengage militarily as quickly as possible.

Senior administration officials suggested, however, that any initial withdrawal starting in mid-2011
could be very limited, depending on the military situation at that point.

``The pace, the nature and the duration of that transition are to be determined down the road by the
president based on the conditions on the ground,'' said Mich\`ele A.~Flournoy, under secretary of
defense for policy.

The initial political reactions showed the crosscurrents facing the White House. Republicans
applauded the buildup of troops but questioned the commitment to a timetable for bringing them home.

``Setting a draw-down date before this surge has even begun is a mistake, and it sends a mixed
message to both our friends and our enemies regarding our long-term commitment to success,'' said
Senator John Cornyn, Republican of Texas.

But among many Democrats, the response ranged from noncommittal to outright opposition.

``I see no good reason for us to send another 30,000 or more troops to Afghanistan when we have so
many pressing issues -- like our economy -- to deal with in this country,'' said Representative
Louise M.~Slaughter, Democrat of New York.

Mr.~Obama is calculating, administration officials said, that the explicit promise of a drawdown
will impress upon the Afghan government that his commitment is not open-ended.

Mr.~Obama was less clear publicly on how he planned to address the issue of Pakistan, which many
administration officials say will prove to be a far more intractable problem in the long term than
Afghanistan.

Administration officials said that Mr.~Obama had signed off on a plan by the Central Intelligence
Agency to expand C.I.A. activities in Pakistan. The plan calls for more strikes against militants by
drone aircraft, sending additional spies to Pakistan and securing a White House commitment to bulk
up the C.I.A.'s budget for operations inside the country.

The expanded operations could include drone strikes in the southern province of Baluchistan, where
senior Afghan Taliban leaders are believed to be hiding, officials said.

The new Afghanistan strategy draws heavily on lessons learned from Mr.~Bush's ``surge'' and strategy
shift in Iraq in 2007, which Mr.~Obama opposed.

In addition to the influx of troops and the training of the Afghan Army, administration officials
said they were taking other lessons from the Iraq buildup, like empowering local security forces to
stand up to Taliban militants in their communities and enhancing the training of national forces by
embedding American troops with Afghan counterparts and later pairing American and Afghan units to
fight side by side.

The 30,000 troops that Mr.~Obama is sending are part of what one administration official
characterized as a short-term, high-intensity effort to regain the initiative against the Taliban.

Administration officials said that they were hoping to get a commitment for an additional 5,000 to
8,000 troops from NATO allies -- perhaps as early as Friday at a foreign ministers' meeting in
Brussels -- which would bring the number of additional troops in Afghanistan to close to the 40,000
that General McChrystal was seeking.

Mr.~Obama is sending three of the four brigades requested by General McChrystal. The first Marines
will begin arriving as early as Christmas, and all forces will be in place by May, a senior
administration official said.

The 30,000 new American troops will focus on securing and protecting the country's top population
centers, including Kabul, Khost and Kandahar, the Taliban's spiritual capital. Military officials
said that two brigades would go south, with the third going to eastern Afghanistan.

Military officers said that they could maintain pressure on insurgents in remote regions by using
surveillance drones and reports from people in the field to find pockets of Taliban fighters and to
guide attacks, in particular by Special Operations forces.

The strategy also includes expanded economic development and reconciliation with less radical
members of the Taliban.

In addition, Mr.~Obama is making tougher demands on the Afghan government; he spent an hour on the
phone Monday with Mr.~Karzai, White House officials said, and pressed him on the need to combat the
corruption and drug trafficking, which many Western officials say has fueled the resurgence of the
Taliban.

During the conversation, Mr.~Obama, described by one White House official as ``very explicit,''
pressed Mr.~Karzai on the need to take steps that would show progress. Mr.~Obama congratulated Mr.
Karzai on setting up a corruption task force, but also pressed him on the need to make sure that
officials appointed by the government are untainted by corruption.

\section{Obama Team Defends Policy on Afghanistan}

\lettrine{D}{efense}\mycalendar{Dec.'09}{03} Secretary Robert M.~Gates, Secretary of State Hillary
Rodham Clinton and the nation's top military officer on Wednesday laid out a muscular defense of
President Obama's decision to send 30,000 additional troops to Afghanistan, but they made clear that
his plan to begin withdrawing those forces by July 2011 was flexible.

At two hearings on Capitol Hill, where they faced deep skepticism about different parts of
Mr.~Obama's war plan from both parties, they also said that the arrival of the additional forces,
while speedy, would not be as fast as Mr.~Obama suggested in a speech to the nation on Tuesday
night.

Although most of the additional troops would arrive in Afghanistan in the first six months of next
year, as the president announced, Adm. Mike Mullen, the chairman of the Joint Chiefs of Staff, said
they would not all be in place until the fall of 2010.

At the same time, American diplomats said that NATO allies had expressed surprise at Mr.~Obama's
commitment to begin withdrawing by July 2011. Richard C.~Holbrooke, the administration's special
representative for Afghanistan and Pakistan, met Wednesday with officials from several countries to
explain the president's thinking, similar to what Mr.~Gates, Mrs.~Clinton and Admiral Mullen were
doing before Congress.

The Democrat who leads the House defense appropriations subcommittee, Representative John P.~Murtha
of Pennsylvania, said Wednesday that he expected the White House to seek \$40 billion by the spring
to pay for the additional troops, or \$10 billion more a year than the president estimated in his
speech. Administration officials said they did not support efforts of some Democrats in Congress to
pay for the additional troops with a surtax on incomes, meaning that war costs would continue to add
to the deficit.

In a full day of testimony before the Senate Armed Services Committee and the House Foreign Affairs
Committee, Mr.~Gates, Mrs.~Clinton and Admiral Mullen faced repeated criticism that the target date
to begin withdrawals -- less than a year after all the 30,000 troops arrive -- would be an
invitation to Al Qaeda and the Taliban to prepare and plan.

Senator John McCain, Republican of Arizona, started the day by grilling Mr.~Gates and Admiral
Mullen, and said it made no sense to set an exit date if the withdrawal was also going to be based
on conditions on the ground, as the president said in his speech.

Eight hours later on the House side, Mr.~Gates was still answering the same question when he said:
``I have adamantly opposed deadlines. I opposed them in Iraq, and I oppose deadlines in
Afghanistan.'' In Afghanistan, he said, ``This will be a gradual process.''

The American commander in Afghanistan, Gen.~Stanley A.~McChrystal, made the same point while
speaking to reporters in Kabul, the Afghan capital. He said he was ``absolutely supportive of the
timeline'' laid out by Mr.~Obama, but he also cautioned that the timeline was flexible and ``is not
an absolute.''

``It's not, 'At 18 months, everybody leaves,' '' the general said.

The July 2011 date, Mr.~Gates, Mrs.~Clinton and Admiral Mullen said, was in large part meant to be a
wake-up call to the Afghans that the United States would not be in the country forever. ``How do you
demonstrate resolve and at the same time convey a sense of urgency to the Afghans that they must
step up to the plate and begin to take responsibility for their own security and to protect their
own country against the extremists?'' Mr.~Gates told the House committee.

Although members of both parties expressed doubts about the war plan, there were few political
fireworks. Representative Dana Rohrabacher, a California Republican who opposes the war, dismissed
Mr.~Obama's strategy as ``maybe a different facade, but it's the same old policy.''

Over all, the hearings had little of the drama that characterized Capitol Hill hearings on the war
in Iraq. In the morning, Senator Evan Bayh, Democrat of Indiana, offered his congratulations to
Mrs.~Clinton on the engagement of her daughter, Chelsea.

The hearings did offer a glimpse of a few details and some of the thinking behind the president's
strategy. The testimony from the three officials showed that the White House would keep to its
longstanding goal to build up the Afghan security forces to 240,000 by 2011, and not expand to
400,000 as General McChrystal had proposed.

Senator Carl Levin, the Michigan Democrat who has been a proponent of training more Afghan security
forces, showed his own skepticism on that issue, and questioned whether sending so many additional
troops might keep the Afghans from building up their security forces on their own. ``Where I have
questions is whether the rapid deployment of a large number of U.S.~combat forces, without an
adequate number of Afghan security forces for our troops to partner with, serves that mission,'' he
said.

In his opening statement, Mr.~Gates, who pushed for the 30,000 additional American troops and was
singled out by the White House as influential in Mr.~Obama's decision, sharply differed with some of
Mr.~Obama's advisers who have argued that the United States should focus on rooting out Al Qaeda
from Pakistan, and that the Taliban in Afghanistan do not present a serious long-term threat to the
national security of the United States.

On the contrary, Mr.~Gates said, Al Qaeda and the Taliban are inextricably linked.

``While Al Qaeda is under great pressure now and dependent on the Taliban and other extremist groups
for sustainment, the success of the Taliban would vastly strengthen Al Qaeda's message, to the
Muslim world, that violent extremists are on the winning side of history,'' Mr.~Gates said. He
added, ``The Taliban and Al Qaeda have become symbiotic, each benefiting from the success and
mythology of the other.''

When pressed by Senator Susan Collins, Republican of Maine, on why the United States had to invest
so much military power and money in Afghanistan when Al Qaeda still had the ability to establish
havens in other countries, Mr.~Gates replied that Afghanistan was unique.

Not only was it the place where the 2001 attacks against the United States were planned, he said, it
``is still the wellspring of inspiration for extremist jihadism everywhere.''

He said that the ``guidance and strategic leadership'' for Al Qaeda came from the group's leaders
who were in the border area with Pakistan, and that there was an ``unholy alliance'' that had
developed in the past year between Al Qaeda, the Taliban in Pakistan and the Taliban in Afghanistan.

He added, ``If anything, the situation, I think, is more serious today than it was a year ago
because of the attacks of the Taliban in Pakistan on Pakistan, and the effort of Al Qaeda in
collusion with the Taliban in Pakistan to try and destabilize Pakistan itself.''

\section{Obama's War Speech Wins Over Some Skeptics}

\lettrine{A}{}\mycalendar{Dec.'09}{03} month ago, Donnie Jones, a 40-year-old Republican who lives
outside Dallas, told pollsters that he was not sure President Obama had a plan for the war in
Afghanistan. But after hearing the president speak Tuesday night, Mr.~Jones feels reassured that
Mr.~Obama not only has a plan, but also one he can generally support.

Margaret Gilbert, 62, a Democrat from Portsmouth, Va., told the same pollsters that she did not want
the United States to send more troops to Afghanistan. But after listening to Mr.~Obama, Ms.~Gilbert
now believes that he has no choice.

And Dave Cegledi, a 66-year-old independent from Olmsted Falls, Ohio, says he does not like
Mr.~Obama any more today than he did in November. But Mr.~Cegledi thinks the president gave a good
speech -- good enough, indeed, that he might vote for him for re-election if the strategy for
Afghanistan works.

Mr.~Obama intended his speech on Tuesday at West Point to rally Americans behind his plan to send
30,000 more troops to Afghanistan and to set an 18-month timetable for starting a withdrawal. And
interviews on Wednesday suggested that, while opinions on the war remained wildly diverse, Mr.~Obama
managed to persuade a significant number of people on both sides of the political aisle, though it
was impossible to know how many.

Many Democrats who opposed the war said they now understood the need for escalation, in some cases
to the point of supporting it. And Republicans who had thought Mr.~Obama unwilling or unable to send
more troops praised his decision, though many also criticized him for articulating a timetable for
bringing troops home.

``I like the fact that he's sending more troops,'' said George Bronner, 45, a Republican from
Knoxville, Tenn. ``This speech does change my opinion about Obama, and it changes it for the better.
I didn't think he was going to step up to the plate and get more people over there.''

The people interviewed on Tuesday night and Wednesday were randomly selected from a list of
respondents to a poll conducted last month by CBS News, which is a frequent partner with The New
York Times on national polls. In that poll, 53 percent approved of Mr.~Obama's overall job
performance, compared with 36 percent who did not, but only 38 percent approved of his handling of
the war in Afghanistan. Opinions on troop levels were almost evenly divided, with 39 percent
supporting a decrease, 32 percent calling for an increase and 20 percent saying to keep it the same.

The telephone poll was conducted Nov. 13-16 with 1,167 adults nationwide and has a margin of
sampling error of plus or minus three percentage points.

Clearly, Mr.~Obama's plan poses a risk of alienating supporters who had wanted him to swiftly scale
back the American military presence in Afghanistan. Betty Holstine, a Democrat from Belle, W.Va.,
said that she had turned against the war after deciding that the death toll was too great and that
she was now on the verge of turning against Mr.~Obama as well.

``I was so eager to trust somebody to make things better,'' Ms.~Holstine, 70, said, ``and perhaps I
put too much faith and trust in the man. I just felt really disappointed and betrayed that it simply
isn't happening.''

Similarly, Dorothy Lingenfelser, a Republican from Lawrence, Kan., who a month ago said she approved
of Mr.~Obama's performance, said she was furious with him for sending more troops.

``I think we ought to take care of people in the United States before we go take care of them in
Afghanistan,'' Ms.~Lingenfelser, a 77-year-old retired factory worker, said. ``I liked Obama before,
but now I'm through with him.''

More common were Democrats who said they continued to support Mr.~Obama even as they held their nose
about his plan.

``Sending more troops is a very bad idea,'' said Robert Labar, 60, a Democrat from Modesto, Calif.
``But I know that the right guy is in the White House, and we just have to go along with what he
wants to do.''

And there were several Democrats who said the president had convinced them of the need to dispatch
more troops.

``I believe him; I trust him,'' said Dianna Sampson, 60, a Democrat from Dayton, Ky., who a month
ago said she wanted to reduce troop levels. ``He met in a room with a lot of people, and this is
what they came up with. It wasn't just this idea out of the blue.''

Though most of the Republicans interviewed said they were pleased that Mr.~Obama was sending more
troops, many expressed skepticism about his motives.

``I think it's good Obama is putting extra troops into Afghanistan, but I know he's not doing it
because he wants to,'' said Suzanne Miller, 60, a Republican from La Grande, Ore. ``He is trying to
earn points with Republicans and conservatives who don't agree with his health care program.''

And there was little evidence that most of the people who last month expressed disapproval of
Mr.~Obama were now prepared to embrace him warmly, as most said they continued to dislike his
policies on health care and the economy.

``I'm trying to give the guy the benefit of the doubt, but he's not doing anything that would push
me toward liking his policies,'' said Scott Taylor, 40, an insurance agent from Franklin, Ind., who
supports sending more troops to Afghanistan.

Republicans, independents and even a few Democrats also expressed concern about Mr.~Obama's
announcement that he would start bringing troops home in mid-2011. Some said he had given away too
much information to the enemy; others said 18 months would not be enough time to stabilize the
country.

``I don't think that after 18 months the Afghan army will be large enough or trained enough to
continue the work on its own,'' said Louis Bruso Jr., 58, a computer programmer and independent
voter from Jamaica, Vt. ``My hope is that we leave enough troops there to help the Afghan army.''

Mr.~Obama seemed to do well among independents, several of whom said they felt more comfortable with
him after the speech.

``He did a better-than-fair job,'' said Fran Turner, 75, an independent voter from San Antonio, who
told the pollsters last month that she opposed Mr.~Obama's performance as president. ``Not a perfect
job, but a good job.''

In particular, Ms.~Turner said she was pleased Mr.~Obama had taken his time in making a difficult
decision, something many Republicans found fault with.

``I go back to the old saying, 'Measure twice, cut once,' '' she said. ``Once you've jumped in,
you're in. Bush jumped in, and he didn't know how to get out. Obama is thinking twice and cutting
once. I really like that.''

\section{Immigration Detention System Lapses Detailed}

\lettrine{G}{rowing}\mycalendar{Dec.'09}{03} numbers of noncitizens, including legal immigrants, are
held unnecessarily and transferred heedlessly in an expensive immigration detention system that
denies many of them basic fairness, a bipartisan study group and a human rights organization
concluded in reports released jointly on Wednesday.

Confirmation of some of their critical conclusions came separately from the Homeland Security
Department's inspector general, in an investigation that found detainee transfers by Immigration and
Customs Enforcement were so haphazard that some detainees arrived at a new detention center without
having been served a notice of why they were being held, or despite a high probability of being
granted bond, or with pending criminal prosecutions or arrest warrants in the previous jurisdiction.

The bipartisan group, the Constitution Project, whose members include Asa Hutchinson, a former under
secretary of homeland security, called for sweeping changes in agency policies and amendments to
immigration law, including new access to government-appointed counsel for many of those facing
deportation.

In its report, the human rights organization, Human Rights Watch, revealed government data showing
1.4 million detainee transfers from 1999 to 2008, most of them since 2006. The transfers are
accelerating, the report found, with tens of thousands of longtime residents of cities like
Philadelphia and Los Angeles being sent to remote immigration jails in Texas and Louisiana, far from
legal counsel and the evidence that might help them win release.

``ICE is increasingly subjecting detainees to a chaotic game of musical chairs, and it's a game with
dire consequences,'' said Alison Parker, deputy director in the United States for the human rights
group, and author of its report. The data underlying the report was obtained under the Freedom of
Information Act, and analyzed by the Transactional Records Access Clearinghouse (TRAC) of Syracuse
University, which issued its own report.

The inspector general's investigation found that the consequences of haphazard transfers include a
loss of access to legal counsel and relevant evidence; additional time in detention; and ``errors,
delays and confusion for detainees, their families, legal representatives'' and the immigration
courts.

Some detainees were transferred with files lacking a photo and a security classification, field
inspectors found in work conducted from October 2008 to February.

Officials at Immigration and Customs Enforcement, which is part of Homeland Security, said the
agency would issue advisories reminding field offices of 10-year-old national detention standards
that require a review of a detainee's ``alien file'' before any transfer, and reinforcing the need
to coordinate with immigration courts.

In August, the Obama administration announced ambitious plans to overhaul immigration detention, a
disjointed network that relies heavily on private prisons and county jails. But taken together, the
three reports underscore the gap between the plans and the problems on the ground in a system that,
according to the inspector general, is estimated to be detaining more than 442,000 people a year --
more than double the number in 2003, ICE's first year of operation.

John T.~Morton, director of the immigration agency, envisions a ``truly civil detention system''
shaped by more centralized agency control. In contrast, the Constitution Project recommends
shrinking the use of detention, in part by adding more constitutional safeguards required in the
criminal justice system.

``None of the recommendations being made should in any way compromise national security,''
Mr.~Hutchinson said Wednesday in an interview before he presented the report at the National Press
Club in Washington. ``It simply allows for a more humane and more efficient system.''

Immigration law is complex, and the deprivation of liberty is quite similar to the situation in
other settings that require court-appointed counsel for the indigent. But 60 percent of noncitizens
face deportation without a lawyer, and transfers compound the problem, the reports said.

The immigration agency has said it uses transfers to deal with an imbalance in the number of
detention beds at various locations. But the TRAC analysis shows that the number of transfers has
grown much more rapidly than the detention population. It found that in the first six months of the
2008 fiscal year, 53 percent of detainees were transferred at least once, and that one in four were
transferred multiple times, a fivefold increase since 1999.

Though transfers occur in almost every state, the data show that the jurisdiction receiving the most
transferred detainees is the Federal Court of Appeals for the Fifth Circuit, covering Louisiana,
Mississippi and Texas -- which is widely known for decisions hostile to the rights of noncitizens
and has the worst ratio of immigration lawyers to detainees, the human rights report said.

A strong case against deportation sometimes simply evaporates in such a transfer, the report said.
It cited a Jamaican New Yorker transferred to Texas after three months in detention in New York and
New Jersey.

Immigration authorities contended that he should be deported based on two prior convictions for drug
possession.

In New York, his drug misdemeanors were not considered an ``aggravated felony,'' and based on the
man's 22 years of legal residency and strong family relationships in the United States, he would
have been eligible for ``cancellation of removal,'' a form of relief from deportation. In Texas, he
was barred from relief based on Fifth Circuit rulings, and deported.

The bipartisan group said the agency makes it too hard for people to avoid detention while
challenging deportation. It recommended a significant easing in the burden of proof, and a hardship
waiver from mandatory detention for lawful permanent residents.

In what it called ``an aspirational goal,'' it recommended that where free counsel is not available,
all indigent noncitizens in standard deportation proceedings have access to a government-paid
lawyer. It also urged Congress to give immigration judges discretion to appoint counsel, and to
require a lawyer in certain cases, including those involving unaccompanied children and the mentally
ill.

Mr.~Hutchinson said that the immigration agency could make many other changes immediately, including
some that would ``correct some potential unfairness in the system'' unintentionally left by his own
efforts when he was in office.

According to the Human Rights Watch report, a memorandum Mr.~Hutchinson issued in 2004 is now used
as a loophole to hold detainees for weeks without giving them notice of why the government is
seeking to deport them. ``This can certainly be tightened up and narrowed,'' Mr.~Hutchinson said.

\section{Senate Breaks Health Stalemate; First Votes Today}

\lettrine{A}{t}\mycalendar{Dec.'09}{03} the end of a third day of Senate debate over sweeping health
care legislation, Democrats and Republicans said Wednesday night that they had broken an impasse
over the seemingly simple question of how and when to vote on the first amendments.

But even as lawmakers announced an agreement to begin voting Thursday, Democrats accused Republicans
of stalling debate and obstructing the legislation.

In a closed-door meeting of his caucus, the Senate majority leader, Harry Reid of Nevada, told
Democrats that Republicans were not interested in passing a bill. In effect, he prepared them for
trench warfare, saying the Democrats must stick together and should be ready to work weekends to
finish the bill before Christmas.

For their part, Republicans said it was unrealistic to expect quick action on such a big bill, and
they denied they were stalling.

``That's an odd charge about a bill that would cost \$2.5 trillion, when fully implemented, and
restructure one-sixth of the economy, affecting 300 million people,'' said the Senate's No.~3
Republican, Lamar Alexander of Tennessee.

Noting that the Senate had spent one month on a farm bill, seven weeks on an education bill and
eight weeks on an energy bill, Mr.~Alexander said, ``Surely we can spend at least that much time on
a comprehensive health care bill.''

Democrats need at least 60 votes to limit debate on the measure, whose overarching purpose is to
guarantee insurance for most Americans while curbing the growth of health care costs and federal
health spending.

Senate leaders of both parties said they would vote Thursday on four proposals dealing with two
issues: how to guarantee additional health benefits for women and how to squeeze nearly a
half-trillion dollars from Medicare over 10 years without adversely affecting older Americans.

Democrats offered one proposal to cover a wide range of screenings and preventive health services
for women. Republicans drafted an alternative, intended to outdo the Democrats.

Republicans offered a proposal to strike provisions of the bill that would save money in Medicare.
Democrats came back with an alternative that would protect ``guaranteed Medicare benefits'' -- not
necessarily the extra benefits offered by private Medicare Advantage plans, which cost the
government more, on average, than the traditional Medicare program.

The agreement to hold votes on Thursday followed some testy exchanges on the floor.

``We are going to go into the fourth day of the debate on one of the most important bills in the
history of the United States Senate,'' said Senator Richard J.~Durbin of Illinois, the No.~2
Democrat. ``The fact is we've not had a single vote this week on health care reform amendments,
because of objections from the other side.''

The Senate Republican leader, Mitch McConnell of Kentucky, protested that ``Republicans were
prepared and fully ready and willing to vote.'' Mr.~McConnell said Democrats shared responsibility
for the delay because they had rejected a Republican offer to vote on two of the proposals Wednesday
night rather than all four at once. Democrats feared that such votes would lead the Republicans to
stretch out any voting on the two remaining proposals.

As the skirmishing continued on the floor, the American Medical Association endorsed the thrust of
the Senate bill. In a letter to Mr.~Reid, the doctors' group praised the measure's strict federal
regulation of health insurance and new tax credits to help low- and moderate-income people buy
coverage.

But the medical association objected to some provisions, including a new tax on ``elective cosmetic
medical procedures'' and creation of an Independent Medicare Advisory Board that could ``mandate
payment cuts for physicians.''

The association also objected to antifraud provisions that it said could penalize doctors who make
``an honest mistake,'' with no intent to defraud Medicare or Medicaid. And it complained about a
section that would impose new restrictions on doctor-owned hospitals.

Elsewhere on Capitol Hill, hundreds of supporters of abortion rights held a rally Wednesday to
insist that the legislation allow insurance coverage of abortion. They denounced restrictions on
abortion that were accepted by House Democratic leaders as the price of passing their version of the
legislation last month.

Representative Judy Chu, Democrat of California, said ``women will lose benefits'' under the House
bill, which she described as ``the biggest rollback of reproductive rights in decades.''

Marcia D.~Greenberger, co-president of the National Women's Law Center, an advocacy group, said the
House bill went to ``ridiculous extremes'' in trying to prevent the use of federal money to pay for
abortions.

Senators Ben Nelson, Democrat of Nebraska, and Orrin G.~Hatch, Republican of Utah, said they would
try to add similar restrictions to the Senate bill. Senate supporters of abortion rights said that
they were confident they could defeat such proposals but that they needed a large majority vote to
give the Senate a strong hand in expected negotiations with the House.

The United States Conference of Catholic Bishops said the Senate bill, in its current form, ``does
not live up to President Obama's commitment of barring the use of federal dollars for abortion.''

\section{In Secret, Comcast Wooed G.E. and Won NBC}

\lettrine{T}{he}\mycalendar{Dec.'09}{03} secret meeting was set for an early July afternoon in a
condominium along the ninth hole of a golf course in Sun Valley, Idaho. Jeffrey R.~Immelt, General
Electric's chief executive, arrived first, taking care to avoid being spotted by his own employee,
Jeff Zucker, the chief executive of NBC Universal, who was mingling with other executives nearby.

Ralph J.~Roberts, the 89-year-old co-founder of the cable giant Comcast, and its chief operating
officer, Steve Burke, arrived 15 minutes later.

The gathering, which had been brokered by James B.~Lee Jr., a vice chairman of JPMorgan Chase, was
set up for one purpose: Mr.~Immelt, who had resisted the urge to sell NBC for years, was finally
ready to sell. For months, he had been in discussions with Mr.~Roberts's son, Brian, Comcast's chief
executive. But now, at the investment bank Allen \& Company's annual media conference -- known for
big deal-making -- he wanted to hear it from the mouth of the company's patriarch.

``Do you want to do this?'' Mr.~Immelt, dressed informally in a polo shirt, asked Mr.~Roberts, who
was wearing his trademark bow tie, and Mr.~Burke, who was Mr.~Immelt's classmate at Harvard Business
School.

``Yes,'' Mr.~Burke said.

Mr.~Roberts, who founded Comcast in Tupelo, Miss., in 1963, said: ``I've done a lot of deals in my
life. Every deal has its time. This is the right time.''

On Thursday, G.E. is planning to finally announce what had leaked more than a month ago: it is
selling a controlling stake in NBC Universal to Comcast, a deal that will once again reshape the
media landscape.

The transaction, the largest during Mr.~Immelt's tenure as chief executive, will also reshape G.E.,
refocusing it into an industrial and financial conglomerate without the flash -- and financial
instability -- of a television and film business. And in the process, he has been undoing much of
the legacy of his predecessor, John F.~Welch Jr.

The deal was a long time in the making and was filled with meetings at the Four Seasons hotel in
Philadelphia, in New York City apartments and on helicopter rides. It also featured code names: G.E.
was Green, NBC was Navy, Vivendi was Violet and Comcast was Crimson (because of the Harvard link).

More than a half-dozen executives involved in those discussions, speaking on the condition of
anonymity because the deal had yet to be formally announced and because the negotiations were
considered confidential, helped reconstruct a nearly yearlong dance between G.E. and Comcast.

Mr.~Roberts had long wanted to control not just the pipes into people's homes, but the television
shows and movies that flow over them. But since 2004, when he sought and failed to buy the Walt
Disney Company, the media industry's economics had cratered. Broadcast television was suffering
through ratings declines, and a falloff in DVD sales had dented profits in Hollywood. But cable
channels, of which NBC Universal has many, were flourishing.

The prospect of a deal with G.E. began in earnest in the late afternoon on March 3 on the 48th floor
of JPMorgan, when Mr.~Roberts and Mr.~Burke came to meet with that firm's chief executive, Jamie
Dimon, at the behest of Mr.~Lee.

The meeting began with a general discussion of Comcast's finances, but Mr.~Roberts said the company
did not need a bank to raise money. Instead, he changed direction by saying he had been pursuing
Mr.~Immelt about NBC but felt like he was getting nowhere. He felt that G.E. was in a vulnerable
position and highlighted the fact that when NBC acquired the Weather Channel earlier in the year, it
partnered with private equity instead of buying the network on its own. It was a sign, Mr.~Roberts
believed, that Mr.~Immelt might not be fully committed to the television business. Mr.~Lee said he
was having breakfast the next morning with Mr.~Immelt and agreed to mention Comcast's interest.

A day later, Mr.~Roberts was standing in the lobby of a Marriott hotel in Baltimore, where his
daughter was playing in a squash tournament, when Mr.~Immelt called his cellphone.

``I want you to know that I'm going to study this,'' Mr.~Immelt told Mr.~Roberts. The two agreed
that measures should be taken to ensure secrecy and that only a handful of executives should be
informed. Mr.~Roberts, who had the failed hostile takeover bid for Disney behind him, had one
requirement: he said he would not participate in an auction.

``We've got to be monogamous,'' he said.

Mr.~Immelt's evolution in thinking about NBC had come over the last year as his company's fortunes
were battered during the financial crisis. In the weeks after Lehman Brothers' bankruptcy,
Mr.~Immelt had spent many hours on the phone with the Treasury secretary, Henry M.~Paulson Jr.,
worrying about the conglomerate's fate.

In the beginning of 2009, as the stock market continued to plunge and G.E. hovered as low as \$5.87
a share, Mr.~Immelt listened to presentations about its assets at a management retreat, where his
thoughts began to crystallize. NBC Universal, whose cable channels continued to do well but whose
flagship broadcast network was deteriorating, no longer appeared to be core to the business and he
thought his capital could be redeployed better elsewhere.

Comcast had also undertaken an internal review to consider where the company could grow by
acquisition. It considered buying another cable company, a mobile phone company or even Facebook. At
one point, it considered acquiring Viacom, which owns several cable networks but is unencumbered by
a broadcast network, but Sumner M.~Redstone, the controlling shareholder in Viacom, was not
interested in selling.

As the spring wore on, G.E. and Comcast met repeatedly, trying to come up with a structure for the
deal. By August, the broad points, in which Comcast would acquire 51 percent of the company, with
G.E. holding 49 percent, were agreed upon. G.E. can begin selling its remaining stake back to NBC
three and a half years after the deal closes at a 20 percent premium to the market value. However,
it would have to share 50 percent of any increase in the value of NBC with Comcast.

The deal nearly fell apart several times. Once, when it seemed that it had been derailed over price
and structure, Michael J.~Angelakis, Comcast's chief financial officer, flew to the summer home of
Keith S.~Sherin, G.E.'s chief financial officer, in Cape Cod and took him and his wife out to dinner
to put the deal back on track. By the end of dinner, they had shaken hands.

The largest complication was that Vivendi, the French conglomerate that owned 20 percent of the
company, could force G.E. to pursue an initial public offering if they could not come to terms on a
deal.

Even within the last two weeks amid a constant stream of leaks, it appeared the deal could collapse.
Vivendi wanted to value the business at \$6.1 billion; G.E. wanted to value it at \$5.5 billion.
They ended up at \$5.8 billion, but there was still a worry about what would happen if G.E.'s deal
with Comcast were blocked by regulators.

Mr.~Immelt, after attending the state dinner last month at the White House, flew to Paris to
persuade Vivendi to complete the deal. An agreement was reached over the weekend after he offered to
pay Vivendi \$2 billion even if the Comcast deal collapsed.

For nearly six months, only a small cadre of G.E. and Comcast executives knew about the deal --
nobody at NBC was ever told -- and it had not leaked. On Sept.~30, several hours after the talks
were disclosed to a tiny group of executives at NBC, the blockbuster talks appeared on TheWrap.com,
a Hollywood news site.

``I'm telling you to be prepared for this to leak,'' Mr.~Sherin had told Mr.~Angelakis earlier that
day.

\section{Senator Moves to Hold Up Bernanke Confirmation}

\lettrine{S}{enator}\mycalendar{Dec.'09}{03} Bernard Sanders of Vermont said on Wednesday that he
would try to block the Senate from confirming Ben S.~Bernanke to a second term as chairman of the
Federal Reserve.

The move is unlikely to derail Mr.~Bernanke's reappointment, but it could slow the confirmation
process and give the Fed's critics additional opportunity to press their case. As a practical
matter, it means Senate Democratic leaders will have to line up 60 votes in favor of Mr.~Bernanke
rather than a simple majority at a time when the Federal Reserve is under increasing populist
attacks from lawmakers on both the right and the left.

Mr.~Bernanke will testify on Thursday at his confirmation hearing before the Senate banking
committee. He is expected to face criticism for not doing more to prevent the financial crisis, and
calls by some lawmakers for a sharply reduced regulatory role in the future.

Mr.~Sanders, an independent, is not a member of the Senate banking committee, but he has frequently
accused the Federal Reserve of bailing out Wall Street firms and the banking industry at the expense
of ordinary citizens.

``In this country, there is profound disgust at what happened on Wall Street,'' Mr.~Sanders said in
an interview. ``People want a new direction and people are asking, where was the Fed? How did the
Fed allow this to happen, when one of their mandates is to oversee the safety and soundness of the
banking system?''

Mr.~Sanders said he would place a hold on Mr.~Bernanke's nomination when it reached the Senate
floor. Under Senate rules, lawmakers would need 60 votes to override Mr.~Sanders and proceed with a
vote.

Though the Senate has been paralyzed by similar blocking tactics on countless other issues,
Mr.~Bernanke probably has enough support in both parties to clear the 60-vote hurdle.

The Fed chairman was appointed by President George W.~Bush and took over the central bank in
February 2006. Despite his Republican ties, Mr.~Bernanke forged a close working relationship with
President Obama and his top economic advisers during the financial crisis.

Mr.~Bernanke's four-year term as Fed chairman ends on Jan.~31, 2010. If he were not reconfirmed by
then, he could continue to serve in an acting role until he was confirmed or someone else was
confirmed to succeed him. His supporters include much of the Senate's Democratic leadership, though
some senior Democrats have been grudging.

Senator Christopher J.~Dodd, Democrat of Connecticut and chairman of the banking committee, has said
Mr.~Bernanke was ``probably'' the best person to lead the Fed because he responded valiantly to the
financial crisis when it began two years ago.

But Mr.~Dodd has also proposed stripping the Federal Reserve of virtually all its powers as a
banking regulator, and consolidating all the federal government's bank regulatory efforts in a new
agency. In an Op-Ed article last Sunday in The Washington Post, Mr.~Bernanke sharply criticized
Mr.~Dodd's proposal.

Senator Richard C.~Shelby of Alabama, the top Republican on the Senate banking committee, has also
been sharply critical of the Federal Reserve but has not yet said how he will vote on Mr.~Bernanke's
nomination.

Mr.~Sanders, one of the Senate's most left-leaning lawmakers, had already said he would vote against
Mr.~Bernanke's appointment. He has also sponsored a bill that would allow Congress to ``audit'' all
of the Federal Reserve's activities, including its core mission of steering the economy by setting
interest rates.

Mr.~Sanders's bill is identical to one championed by Representative Ron Paul, Republican of Texas,
which Mr.~Bernanke and other top Fed officials have adamantly opposed on the grounds that it would
undermine the central bank's independence.

But in a surprising setback to the Fed, the House Financial Services Committee voted last month to
approve Mr.~Paul's bill as an amendment to a broader bill on financial regulation.

\section{Organizing the Chaos of Online Travel Tips}

\lettrine{T}{his}\mycalendar{Dec.'09}{03} fall, I was lured to Panama. Friends who settled there a
year ago had regularly regaled me, an avid traveler, with visions of a lively Latin city with
nearby, unspoiled beaches bordered by jungle.

For half of my 10-day trip, I intended to traipse around Panama City. The rest would be a solo
adventure -- but to where? Greedily, I began surfing the Web, exploring my friends' suggestions and
amassing options.

It didn't take long for disorganization to set in. There were so many possibilities and interesting
tidbits in travel articles. How would I keep it together to make a wise destination decision -- or
take the information I'd found with me? I could cut and paste text into a Word document or save Web
pages in a bookmarks folder. But the former seemed tedious, and the latter inadequate.

Then I remembered Gliider, a browser tool that saves and organizes travel research. I downloaded the
free add-on for Firefox, which deposited a small arrow icon on my browser navigation bar that, when
clicked, opened a sort of file box. There I created a ``trip'' to Panama City, and began selecting,
dragging and dropping text and photos from the Web into handy folders with labels like flights,
hotels, see and do.

I could share the stash with friends and e-mail myself PDF dossiers of smartly organized
information, ready to print out as a bespoke travel guide. (In January, iPhone users will be able to
carry their Gliider content via an app.)

Gliider, I learned, is just one of a growing number of new online tools to help travelers plan where
to go and what to do when they get there. While still immature and somewhat buggy, they are making
finding, sorting, saving and organizing nuggets culled from the mountains of online travel content
more efficient and fun.

GET ORGANIZED With the help of Gliider, my Panama vision soon crystallized. I decided to spend two
days in Bocas del Toro on the Caribbean, then I would head over the continental divide and spend two
days in Boquete, a mecca for coffee fiends.

I made Gliider trips for each destination. Local flight confirmation details went into the flights
folder, lodging choices into hotels, surf-lesson information into see and do, and so on.

With a little more luck (or more common destinations), I might have gotten a price break on a hotel.
Gliider queries its partner sites, Hotels.com and TravelZoo, and alerts you to deals at properties
you put in hotel folders and nearby properties of the same quality. If you book a room, Gliider
collects a commission. In the future, it plans to provide deals for flights, rental cars and
adventure excursions specific to trips.

Two other services, TravelMuse and NileGuide, also help users organize trip information, create
their own travel guides and make bookings via partners. These Web sites also have unique travel
content and user-created content, which makes them especially useful for travelers who haven't
decided on destinations yet. In fact, TravelMuse has a feature called Inspire Me that provides
suggestions based on the kind of trip you want -- say, a romantic getaway no more than a two-hour
flight from your home airport. NileGuide offers ideas by region and categories like adventure,
beach, budget, family and honeymoon.

Because the sites are still building their content, they work best for people planning trips to
well-worn destinations. Neither were fully prepared to help me plan my trip to Panama.

TravelMuse provided more overall flexibility because, in addition to having its own content, it lets
you search the Web. The only catch is that you have to search the Web from within TravelMuse, which
can feel a bit cumbersome. I also found the search results were sometimes off. For example, while
looking for a restaurant in Panama City, I got some results in a city of the same name in Florida.
Nevertheless, the approach enables content to be saved in a Web-based Tripfolio that you can share
without having to download software.

NileGuide is a slicker site but limits users to content it creates itself and gets from partners. To
test it out, I created a guide for Cairo, a city I once lived in. The hotel, restaurant and activity
suggestions were good and reasonably comprehensive, with many items helpfully informed by a ``local
expert.'' Once I'd made my own list of items, I could see them on a Google map and create a
logistically sensible day-by-day itinerary for a big, gridlocked city. The site also offers 16 Cairo
guides, including itineraries for lovers, budget travelers and visitors with children.

FIND FUN A spate of new sites have also cropped up that focus on helping travelers find great
activities.

Kijubi (sounds like ``could you be'') is building an online marketplace for tours, adventure trips
and other leisure activities in California, Colorado and Florida, for starters. Travelers can browse
activities like backpacking, dinner cruises and even indoor skydiving and purchase them directly on
the site. Companies that organize the activities pay commissions when Kijubi sends them customers.
The site also includes user ratings and reviews.

Trekaroo focuses on families. Modeled on the consumer-review site Yelp, Trekaroo helps parents find
child-friendly hotels and age-appropriate activities recommended by other parents. It has more than
20,000 reviews covering 65 destinations in the United States, and is expanding steadily.

Parents looking for activities can select a destination and travel radius, specify price and age
ranges, and choose an activity category. I searched for free activities within 10 miles of Manhattan
suitable for children age 10 and under and got 27 results, topped by Central Park and including
ideas like the Socrates Sculpture Park in Long Island City, Queens.

SEARCH MORE EFFICIENTLY Activity-focused search engines, including Goby and UpTake, confine
themselves to searching sites oriented toward travelers. Since they are more focused, they hope to
do a better job delivering tips on great experiences than far-ranging engines like Google do.

On Goby's site, which was started in September, users can choose among 300 categories or type
keywords that express what they'd like to do, type in a destination in the United States (and one
day, the world), and select a timeframe or specific dates.

A search for jazz music within 10 miles of Lower Manhattan this weekend yielded 23 results,
including a Branford Marsalis gig in Newark on Friday night. I clicked ``what's nearby'' and found
27 restaurants only a quarter-mile away. While the first result was a misfire (a fast-food company),
the second made a desirable pairing -- a well-reviewed soul-food restaurant.

This time, I wasn't even planning a vacation -- some of these sites are just as useful for locals.

\section{New Stem Cell Lines Open to Research}

\lettrine{T}{he}\mycalendar{Dec.'09}{03} National Institutes of Health said Wednesday that it had
approved 13 new human embryonic stem cell lines for use by federally financed researchers, with 96
more under review.

The action followed President Obama's decision in March to expand the number of such cell lines
beyond those available under a policy set by President George W.~Bush, which permitted research to
begin only with lines already available on Aug.~9, 2001.

Since that date, biomedical researchers supported by the N.I.H. have had to raise private money to
derive the cells, which are obtained from the fertilized embryos left over from in vitro fertility
clinics.

With federal money banned from being used in any part of the work on the derived lines, researchers
had to divide their laboratories and go to extreme lengths to separate research materials based on
the financing source.

``You can imagine what it meant not to be able to carry a pipette from one room to another,'' said
Ali H.~Brivanlou, a researcher at Rockefeller University. ``They even had to repaint the walls to
ensure no contamination by federal funds.''

Two of the newly approved 13 lines were derived by Dr.~Brivanlou with private financing. The rest
were prepared by Dr.~George Daley of Children's Hospital Boston.

Dr.~Daley said that private financing had been drying up and that he was eager to start research on
the now-approved cell lines with the help of his federal grant money.

The director of the health institutes, Dr.~Francis S.~Collins, said he believed that most
researchers would be satisfied with the outcome, even though they were still barred from deriving
the cells themselves.

``I'm not sure everyone is interested in deriving their own cell lines as long as they can get lines
from others,'' Dr.~Collins said.

Researchers' interest in human embryonic stem cells has abated since the discovery in 2007 by the
Japanese biologist Dr.~Shinya Yamanaka that the mature cells of the body can be reprogrammed to the
embryonic state.

These induced embryonic cells are highly similar to the real thing but may not be exactly the same.
One reason is that the mature cell may perceive the forced walk-back to embryonic state as
unauthorized and switch on its anticancer defenses.

Because the reprogrammed cells and those derived from leftover human embryos may not be identical,
researchers need to work with both kinds, Dr.~Collins said.

\section{North Korea Revalues Its Currency}

\lettrine{S}{eoul}\mycalendar{Dec.'09}{03}, South Korea -- North Korea has taken confiscatory steps
to curb unofficial economic activity and suppress inflation, issuing a revalued currency and sharply
limiting the amount of old money people can exchange for new bills, South Korean news reports said
Tuesday.

The measures appeared to be a crude form of shock therapy intended to severely punish people
involved in black market trading, which had begun to take root in North Korea in recent years as the
country's central planning system broke down.

If carried out as described by South Korean and Chinese news agencies, the adjustment would wipe two
zeroes off the value of North Korea's currency, the won, and allow people to exchange only about
100,000 to 150,000 old won for the new notes. That would effectively decimate private stores of cash
wealth in local currency. At the official exchange rate, 100,000 won is the equivalent of \$690, and
at the black market rate it is worth only \$35.

Unconfirmed reports on South Korean news Web sites said that the revaluation and exchange limits
briefly set off street protests in North Korea's capital, Pyongyang, forcing the authorities to
revise slightly upward the amount of currency people could exchange.

The value of the dollar and the Chinese currency skyrocketed as people rushed to underground
money-exchangers, said the Yonhap news agency of South Korea, citing anonymous North Korean traders
in China.

Rail stations were crammed with ``uncontrollable crowds'' of people rushing home to swap their
money, said Daily NK, a Seoul-based Web site that specializes in gathering North Korean news, citing
informants in Haesan, a North Korean town near China, and other cities.

``It's all chaos in markets and railway stations,'' Daily NK quoted a North Korean as saying.

But Good Neighbors, a Seoul-based Buddhist relief agency that collects information on North Korea
through people inside the country, said that Pyongyang and other major cities had returned to
``utter quietness'' by late Tuesday as the North's secret police, the national security agency and
the military increased patrols, and most shop owners closed their stores.

North Koreans were shocked by the currency reform, Good Neighbors said. Many had stored cash to help
tide them over during the lean winter and spring months, but now they have suddenly seen much of
their savings wiped out, the group said.

There was no official announcement of the measure in North Korea's official media, and the accounts
of protests in the North, which would be rare in the tightly policed country, could not be
independently confirmed.

``As people rushed to swap money, commercial activities have virtually come to a standstill,'' Good
Neighbors said in a statement. ``The purpose of the reform is to kill private market activities that
stoked anti-socialism.''

South Korea has received intelligence about the reported currency move but has been unable to
confirm it, said Chun Hae-sung, a government spokesman.

Choi Soo-young, an expert on the North Korean economy at the government-financed Korea Institute for
National Unification in Seoul, said: ``The primary aim is to control inflation. But the North also
wants to stop the market from prospering too fast.''

Since a famine that killed many North Koreans and shook the rationing system in the mid-1990s, the
North's centrally planned economy, with state-run stores that sell goods at government-set prices,
has coexisted with an unofficial economy where people sell homegrown food or goods smuggled in from
China.

For a time, those nascent signs of private-sector activity were cited as evidence that the North
Korean leader, Kim Jong-il, might be inclined to follow the example of China, the North's neighbor
and most important backer, which began to move away from its dysfunctional central planning system
in the late 1970s. At that time, Beijing first allowed entrepreneurs to supplement the country's
state-run production and distribution systems.

But after years of tolerating unofficial markets, North Korea declared in 2005 that it was
reinforcing the ration system. Since then, it has restricted the kinds of goods that can be sold in
the markets, how often they should open and the age of people who can sell there.

Earlier this year, Vitit Muntarbhorn, the United Nations' special representative on the human rights
situation in North Korea, said that the government had singled out women in a crackdown on marketing
activities.

He said that by suppressing the private markets, the government was demonstrating its intention to
keep North Koreans dependent on the state.

Women have been especially active in the private markets, and there have been reports in the South
Korean news media about young North Korean women protesting against the government's renewal of
antimarket measures while food shortages have persisted.

In June, the authorities shut down the largest unofficial market for food and goods outside
Pyongyang, dispersing traders to smaller markets in other districts, according to South Korean news
reports at the time.

One reason is that the mix of state and private activity has fueled inflation, partly by drawing
goods that would otherwise be sold at low, state-subsidized prices into more lucrative private
markets.

``With the government stores unable to provide enough supplies, unofficial markets have fueled
inflation,'' said Dong Yong-sueng, an economist at the private Samsung Economic Research Institute
who monitors the North Korean economy.

Food prices have risen so sharply in recent years that workers can buy only two kilograms of rice,
about four and a half pounds, with their monthly wages, according to Open Radio for North Korea, a
Seoul radio station and Web site that specializes in collecting news from informants in the North.

The North Korean officials also saw the increasingly vibrant markets as a conduit of capitalist
ideas and outside influence on the tightly controlled populace.

By wiping out much of the wealth that traders have accumulated, the currency measure could anger
many people who lose money. But from the government's perspective it could have the effect of
reducing a growing income gap among citizens, Mr.~Choi said.

The measure will have an immediate but short-lived impact in fighting inflation, Mr.~Choi and
Mr.~Dong said. But it could also worsen the North's already acute shortages by discouraging the
markets, which have emerged as an important source of food for North Koreans.

``As long as the North's official economy cannot generate enough goods, and unless it can invite
foreign investors by improving external relations, inflation will recur,'' Mr.~Dong said. ``The
vicious circle will continue.''

\section{H.I.V. Tests Turn Blood Into Cash in China}

\lettrine{A}{}\mycalendar{Dec.'09}{03} young, boisterous\cite{boisterous} crowd gathered in front of
the Purple Tribe nightclub on a recent Friday night, but hardly anyone was interested in going
inside.

Instead the men, most of them gay, waited their turn to duck into a dingy storage space next to the
club. A needle prick and a wince\cite{wince} later, they emerged with a triumphant\cite{triumphant}
grin, having exchanged a test tube of blood for a pocketful of cash. ``This is my third time in two
weeks,'' Zhang Haoyun, an 18-year-old store clerk, boasted as he walked away holding a cotton swab
to the bend of his arm.

On any given night, in 14 cities around the country, hundreds of people flock to makeshift blood
collection centers in bars, bathhouses and apartments where workers test for syphilis and H.I.V.,
the virus that causes AIDS. The ambitious testing initiative, started in 2007, is financed by the
Bill \& Melinda Gates Foundation, which will spend \$50 million over five years in an effort to slow
the spread of AIDS in China. So far, more than 110,000 people have been tested.

But the Gates H.I.V. prevention program in China is unusual because it offers a financial incentive
to those drawing the blood -- about \$9 per sample and an additional \$44 for those that come back
positive -- which is shared with donors. The program has provoked a flurry\cite{flurry} of criticism
from some established AIDS organizations that say the money has given rise to a network of
fly-by-night groups whose only interest is collecting money.

Here in Tianjin, a northern city of 11 million people, two dozen organizations have sprung up in the
past year, many of them run by bar owners or bureaucrats affiliated with the government. Some of the
groups do not provide counseling to those giving blood and make little effort to help those who test
positive get medical treatment.

``Gates has created a huge blood-buying operation that only cares about money, not about people,''
said Ma Tiecheng, who runs a seven-year-old AIDS organization in the northwest city of Shenyang.
``I've seen people getting four H.I.V. tests a day.''

With more than \$3 billion in annual giving, the Gates Foundation tends to make waves wherever it
directs its largess. It is widely lauded for taking on malaria, tuberculosis and AIDS, among the
biggest killers in the developing world, but it has aroused criticism for what some describe as the
foundation's resistance to outside viewpoints when charting its programs. Others complain that the
huge grants alter the local health care landscape by creating a realm of haves and have-nots among
grass-roots organizations.

Here in China, the foundation's stated mission is to identify H.I.V.-infected people as the first
step in getting them treated. Those who know their status, architects of the program say, are also
more likely to modify behavior that puts others at risk. The payments, they say, are crucial for
bringing people face to face with outreach workers.

Although not trumpeted in its promotional materials, the foundation's other goal is more
far-reaching: to empower the small but growing crop of nongovernmental groups that stand a better
chance of addressing the AIDS epidemic than China's lumbering bureaucracy does.

To carry out its mission, the foundation has linked up with the Ministry of Health, which funnels
\$20 million to about 200 nonprofits, many of which exist in a bureaucratic gray zone and are viewed
suspiciously by China's authoritarian government.

The distrust flows both ways.

By compelling the government to work with privately run organizations, the foundation is hoping to
foster a lasting relationship between them -- and over time contribute to creating more profound
changes in Chinese society.

Dr.~Ray Yip, who runs the foundation's China effort, acknowledges problems with the program but
likens them to growing pains.

``We are experiencing some of the hiccups of a less-than-perfect arrangement, but we expected
that,'' he said. ``If you look historically at arranged marriages, some of them last.''

Dr.~Yip, who is the former director of the Centers for Disease Control and Prevention office in
China, embraces the Gates Foundation's philosophy of bold initiatives and risk taking -- traits
often lacking in government-run global health behemoths. He said that if some of the money ended up
in the pockets of corrupt officials running fake organizations, it was the cost of doing business in
China, where government malfeasance is endemic.

``We don't expect every grant in every city to be spectacularly successful,'' he said. ``That's like
buying 30 stocks and expecting them all to go up.''

Sun Jiangping, deputy director of China's National Center for AIDS/STD Control and Prevention, said
the program had already had a positive impact on government attitudes toward private AIDS
organizations, whose numbers have increased to more than 400 from just a few dozen when the
initiative started. He said his agency was working to weed out illegitimate groups.

Compared with the rates in other developing nations, the prevalence of H.I.V. in China is relatively
low, with fewer than a million people thought to be infected, according to government figures
released last week.

But public health experts are alarmed by an infection rate among gay men that has been doubling
annually. By the end of 2008, nearly 5 percent of gay men in China's largest cities were thought to
be H.I.V. positive; in some cities, that figure exceeds 10 percent. Health officials say gay men now
account for a third of all new transmissions, up from 12 percent in 2007.

Advocates for people with AIDS say the government has been ham-handed in its efforts to prevent the
spread of H.I.V., in some cases banning condoms in bars or hounding activists who become too vocal.

In recent years, organizations have sprung up to help those with AIDS who are refused care by
hospitals. Many, like Deep Blue, a group that operates from an apartment on the outskirts of
Tianjin, are largely financed by grants from abroad.

Deep Blue's two counselors meet with the 50 people who come each week for an H.I.V. test.

``If you have any questions, contact the volunteers with the red armbands,'' says one poster. ``You
can stop the testing at any time,'' says another. About 65 percent of those who test positive come
back for counseling, said the group's director, Yang Jie.

Tong Ge, a veteran AIDS activist who has advised the Gates Foundation on its China program, said he
was pleased with the foundation's work, though he said he wished there was more of an emphasis on
training government workers and less money spent on testing. His biggest regret, he said, is that
the foundation chose to funnel the money through the government.

``So much of the Gates money has ended up nurturing corruption in a place it didn't exist before,''
he said. Then, after a pause, he added: ``But the truth is we can't blame them. The real problem is
with China.''

\section{In NBC Interview, a Failure to Note Network Ties}

\lettrine{T}{o}\mycalendar{Dec.'09}{03} the untrained eye, it might appear that NBC Universal has
given a lesson in corporate synergy\cite{synergy}.

When two television cast members sneaked into a White House state dinner last week, Bravo cameras
were filming them all the way to the east gate. Days of publicity followed. When they finally
decided to explain themselves, they gave their first -- and so far only -- interview to Bravo's
corporate sibling, NBC.

NBC News says the appearance of the party-crashers, Tareq and Michaele Salahi, was merely the result
of hard journalism and old-fashioned hustle. But NBC's ``Today Show'' faced sharp criticism on
Tuesday and Wednesday when it did not disclose that the couple was being filmed for ``The Real
Housewives of D.C.,'' the latest extension of Bravo's glamorous and lucrative franchise.

NBC and Bravo are both units of NBC Universal, and are both based at Rockefeller Center. By
neglecting to acknowledge the corporate relationship, ``NBC is expanding the field of quicksand
they're standing in,'' said Bob Steele, a journalism ethics professor at DePauw University.

NBC's news division and Bravo have denied assisting each other, and there is no evidence that they
have done so.

``I would be perfectly comfortable admitting that this interview was a result of corporate synergy
if that were the truth,'' said Lauren Kapp, an NBC News spokeswoman. ``But in this case I can
honestly say that it was good old-fashioned booking.''

Critics have not questioned the appearance as much as how it was handled. When the ``Today'' co-host
Matt Lauer interviewed the Salahis on Tuesday, he did not mention Bravo or the fact that the couple
is under contract to the network (Bravo said that it was still considering whether to feature them
on the show). Mr.~Lauer also omitted the information in a follow-up ``Today'' segment with
commentators on Wednesday.

Such disclosures are common on ``Today'' and other news programs.

``Clearly, there's a competing loyalties issue and, I would say, a conflict of interest that should
be revealed,'' said Mr.~Steele, who added that the perception of unethical entanglements could erode
the credibility of journalists.

The White House has said that the Salahis had not been invited to the state dinner, prompting an
investigation by the Secret Service. By coincidence, Jeffrey R.~Immelt, the chief executive of
General Electric, which owns NBC, attended the dinner last week, as did Brian Williams, the ``NBC
Nightly News'' anchor. (A deal giving control of NBC Universal to Comcast is expected to be
announced Thursday.)

The Salahis have been hotly pursued by bookers for all the major television networks. They initially
scheduled an interview with Larry King of CNN, only to postpone it a day later and, according to a
network executive who declined to be identified discussing negotiations, the couple requested
appearance fees in the six figures.

Both NBC and the couple say they received no money for appearing on the ``Today Show.'' An NBC staff
member suggested Wednesday that the couple selected Mr.~Lauer in a good-will gesture to NBC and, by
extension, Bravo. The staff member spoke on condition of anonymity because the network declined
interview requests about the subject on Wednesday.

Some television critics knocked Mr.~Lauer for not asking the Salahis any questions about the
``Housewives'' production, which is still under way. The Salahis told Mr.~Lauer that they hoped to
appear on the show again soon to explain why they believed they were invited to President Obama's
first state dinner.

The evening counterpart to ``Today,'' the ``NBC Nightly News,'' did note the corporate connection on
Tuesday when it reported on e-mail messages between the Salahis and a Pentagon official who had
tried and failed to gain a dinner invitation for the couple.

As it deliberates about whether to feature the Salahis on the ``Housewives'' series -- a decision
that is months away, according to Bravo -- NBC Universal may also have to consider its corporate
interests. Frazier Moore, the television critic for The Associated Press, asked in a column this
week, ``Would a cable network whose hydra-headed corporate parent counts on Washington for such
things as favorable regulation and fat defense contracts want to tie itself to a 'D.C. housewife'
who royally ticked off the feds?''

\section{Recession Elsewhere, but It's Booming in China }

\lettrine{F}{or}\mycalendar{Dec.'09}{10} the first time, Chinese will buy more cars this year than
Americans. Demand is so high that drivers put their names on long waiting lists for the most popular
models.

``I'm disappointed, but what can I do?'' asked Zhang Ge Lu, a 28-year-old interior designer. He came
recently with two friends to a row of dealerships here in southeastern China to buy a black Toyota
RAV4, only to be told that he would have to wait two months for delivery.

And it is not just cars. For more and more consumer goods, China is surpassing the United States as
the world's biggest market -- from cars to refrigerators to washing machines, even desktop
computers.

The Chinese market is ``on full tilt -- booming is an understatement these days,'' said John
Bonnell, the director of Asia vehicle forecasting at J.D. Power \& Associates.

China is pulling ahead at this particular moment partly because Americans, debt-laden and worried
about their jobs, are pulling back. After decades of gorging on consumption, Americans are saving.
And the Chinese, whom economists thought were addicted to saving, are spending more.

Among China's 1.3 billion people, rising incomes are finally making large numbers of Chinese
prosperous enough to make big-ticket purchases.

The question is: will they keep spending? The Beijing government is increasing consumption with
rebates, subsidies and heavy bank lending. Whether China can turn the spending spree into the seeds
of a true consumer society matters not just to China, but to the world.

For years, the West has pushed China to increase domestic consumption and reduce its dependence on
exports -- that's because its overdependence on exports has distorted global trade.

To keep its export machine humming, China kept its currency undervalued to make its goods more
competitive in foreign markets. The county beggared its own citizens, keeping salaries and bank
deposit interest rates artificially low to support exporters.

China's trade surpluses and extensive intervention in currency markets have led it to amass \$2.27
trillion in reserves, mainly in United States Treasuries, mortgage-backed securities and other
dollar-denominated investments, helping to keep interest rates low and finance Americans' borrowing.
Chinese parsimony enabled American profligacy.

If the Chinese buy more and Americans save more, a more stable global economic exchange can take
shape. In the meantime, China's rapid consumption growth is good news for the whole world. For the
first time, China, not the United States, is a locomotive helping to pull the global economy out of
a slump. But China's tiny appetite for American exports means that the main benefit has gone to
commodity exporters and to businesses in China.

Automakers are on track to sell 12.8 million cars and light trucks in China this year, virtually all
of them made in China (although many are foreign brands), compared with 10.3 million in the United
States. Appliance manufacturers expect to sell 185 million refrigerators, washing machines and other
pieces of kitchen and laundry equipment in China this year, compared with 137 million in the
American market.

In desktop computers, China moved solidly ahead of the United States in the third quarter, buying
7.2 million compared with 6.6 million in the United States.

Retail sales are growing 17 percent a year in China after adjusting for inflation, almost twice as
fast as the overall economy.

Americans have been cutting back on purchases of everything from shoes to furniture to jewelry. But
Chinese households are crossing a series of income thresholds at which cars and other big-ticket
purchases become affordable.

At the same time, Chinese banks are stepping up consumer lending. The proportion of car sales
financed with loans has doubled this year, to nearly 25 percent, although most Chinese still head
for dealerships with bricks of 100-renminbi notes, each note worth about \$14.62. Credit card
spending rose 40 percent in the first nine months of the year compared with the same period last
year, yet China still has just one credit card for every eight people, compared to two credit cards
for each American man, woman and child.

While it is spreading creature comforts, China's lending-based prosperity may also be sowing the
seeds of future economic problems. China's Banking Regulatory Commission recently told banks to show
restraint in lending for the rest of the year, fearful that some of this year's loans could become
bad debts in the next several years, as happened with the mortgage lending spree in the United
States.

The regulator threatened to block banks' overseas investments and branch openings unless they can
demonstrate adequate capital to cover risks.

The size of China's consumer market, notwithstanding its growth, will make it hard for China to
rescue the world economy by itself. Total consumer spending in China is still less than a sixth of
American consumer spending at current prices and exchange rates. That is mainly because China has
relatively few restaurants, hotels and other service businesses, even as sales of manufactured goods
have risen.

The average price tags on most Chinese products are much lower than in Western markets. For many
products, including some in which China leads in the sheer number of goods, the total dollar value
of sales in China is still smaller than in the United States.

The average new car sells for \$17,000 in China compared with almost \$30,000 in the United States,
according to J.D. Power. This is because Chinese consumers buy more subcompacts and fewer sport
utility vehicles. While the Chinese market is one-quarter larger in the number of cars sold, the
American market is still about two-thirds larger in dollar terms.

Similarly, the United States market for household appliances is a third larger in dollars, even
though the Chinese market is a third larger in the number of appliances. Cooking ranges in China are
sold for countertop installation without a lot of other equipment, for example.

``You don't have the cook-a-turkey-in-the-oven type of product in China, because we don't have that
kind of cooking,'' said Philip S.~Carmichael, the president of Asian operations at Haier, China's
biggest appliance manufacturer.

But in some sectors, Chinese buyers are already proving more lavish than Americans. The average
flat-panel television sold in China is bigger than in the United States, according to AU Optronics
of Taiwan, the world's third-largest manufacturer of flat-panel televisions.

When car sales began surging early this year, many auto executives attributed the boom to government
incentives. To stimulate the economy, the government has offered rebates for rural families to buy
cars and household appliances, and has cut sales taxes on cars with small engines.

But the boom has broadened to categories that barely qualify for incentives.

S.U.V. sales rose 72 percent in October from a year earlier. At Nissan, sales of cars with larger
engines that do not qualify for the sales tax reduction are growing even faster than sales of
small-engine cars.

Auto sales jumped 42 percent in the first 11 months of this year compared with sales in the same
period last year. And sales are still accelerating, soaring 96 percent in November compared with the
same month a year ago. Auto sales in the United States plunged 37 percent last month on the same
basis.

China's consumers have the potential to buy even more in the years ahead. The savings rate is close
to 40 percent -- and will remain high unless and until Beijing creates a social safety net for
things like health care or retirement, which would encourage Chinese to spend more today.

And though annual incomes still average just \$2,775 a person in cities and \$840 in rural areas,
Western economists predict the economy will grow almost 12 percent in each of the next two years and
the renminbi is widely expected to appreciate someday, further increasing consumers' buying power.

\section{Ready to Exercise? Check Your Watch }

\lettrine{M}{y}\mycalendar{Dec.'09}{11} friend Jen Davis and I often run together in the morning
because it can be easier to fit in a run before work than after. But we always thought we ran better
in the evening.

Then I accidentally discovered something weird. I took a spinning class one Thursday night, and my
heart rate, measured by a monitor strapped around my chest, soared. I don't usually use a heart-rate
monitor, but with stationary bikes, heart rate is pretty much the only way to know how hard you are
working. And that night, my high heart rate told me it really was a tough workout.

The next morning I did a workout in my garage on a trainer -- a device that holds a road bike,
turning it into a stationary bike and yet allowing you to use its gears. My heart rate was about 15
beats a minute lower than it had been the night before. It seemed like a pitiful workout.

So the next night I got on the trainer again. I had the same playlist (I use music to set my
cadence). I used the same gears for each song. And during the hourlong workout, my average heart
rate and my maximum heart rate were about 15 beats a minute higher than they'd been the morning
before.

I tried again the next morning. My heart rate was low. Intrigued, I tried my experiment for a week,
alternating between early morning and early evening workouts. I got really sick of that playlist,
but I wanted to control every variable.

And the pattern persisted: high heart rate at night, low in the morning for the identical workout.
Once I even tried the workout in midday -- that time, my heart rate was in between.

Could it be that I actually was a more efficient athlete in the morning, doing the same work but
with less effort, as measured by a lower heart rate?

Jen reminded me that we'd seen the heart-rate effect last year but had not appreciated it. I had a
stress fracture and was confined to pool running, which involves sprinting in the deep end of a
pool. Your feet never touch the bottom. It was hard to gauge how hard we were working, so Jen and I
wore heart rate monitors, just as we do in spinning classes.

We did the pool workouts together, and neither of us got our heart rates as high as we wanted in the
morning. Evenings were fine, though. We thought we were just sluggish in the morning.

I also asked some friends who use heart rate monitors if they'd noticed anything like what I'd
experienced.

Tara Martin, a triathlete, said she could never get her heart rate up in the morning.

Richard Friedman, a swimmer, said his heart rate was always lower in the morning. His swim team does
the same workout in the morning as in the evening, and he swims it just as fast. He had assumed that
somehow he was just not putting in the same effort early in the day. ``Still,'' he said, ``I'm
pretty energetic all the time.''

I asked Dr.~William Haskell, an exercise researcher and emeritus professor of medicine at Stanford,
if I'd stumbled on a known fact about heart rates. But he was baffled. Maybe I didn't have caffeine
in the morning? So I tried taking NoDoz before the next morning workout. It made no difference.

Dr.~William Roberts, a former president of the American College of Sports Medicine and a family
physician at the University of Minnesota, said it was a ``tough question.'' He added, ``I do not
have a good physiologic explanation for the phenomenon you are describing.''

But, it turns out, a small group of researchers has studied the question of exercise performance and
time of day, even doing studies of heart rates. And not only are performances better in the late
afternoon and early evening, but, contrary to what exercise physiologists would predict, heart rates
are also higher for the same effort.

One recent study, by the late Thomas Reilly and his colleagues at the Research Institute for Sport
and Exercise Sciences at Liverpool John Moores University in England, found that people's maximum
heart rates and sub-maximal heart rates were lower in the morning but that their perception of how
hard they were working was the same in the morning as it was later in the day.

Dr.~Reilly and his colleague Jim Waterhouse, in a review published this year, also noted that
athletes' best performances, including world records, were typically set in the late afternoon or
early evening.

Greg Atkinson, also at Liverpool John Moores University, said that some researchers, noticing that
heart rates during exercise were lower in the morning, reasoned the way I did -- that people must be
more efficient in the morning. It would mean that exercise was easier in the morning. Of course, it
seemed harder to me, but I could have been deluding myself. Not really, Dr.~Atkinson said. It
actually is harder to exercise in the morning.

``Most components (strength, power, speed) of athletic performance are worst in the early hours of
the morning,'' he wrote in an e-mail message. ``Ratings of perceived exertion during exercise have
generally been found to be highest in the early morning.''

If you exercise later in the day, your muscles are more flexible and stronger and your heart and
lungs are more efficient, said Michael H.~Smolensky, an expert in chronobiology, the study of the
body clock.

``Is a heart rate of 140 in the morning indicative of the same level of workout cost as in the
afternoon?'' asked Dr.~Smolensky, a visiting professor at the University of Texas Health Sciences
Center in Houston.

``I would say no,'' he added. ``Exercise physiologists say you should be able to perform at the same
level with a heart rate of 140 in the morning as in the afternoon or early evening. But
chronobiologists say your capacity to generate and tolerate a higher heart rate is better later in
the day.''

``In the afternoon and evening,'' Dr.~Smolensky said, ``you are in a different biological state.''

But, he added, all this applies to people who are regular exercisers, who work out vigorously three
or more times a week. People who are not regular exercisers, Dr.~Smolensky said, put much more
strain on their hearts in the morning, making their heart rates higher then.

In fact, Dr.~Smolensky added, people at risk for a heart attack should plan their workouts for late
afternoon or early evening.

But if you are used to regular exercise, is it better to train in the early evening?

``I really don't know the answer,'' Dr.~Smolensky said.

``My personal approach is to train when your biological efficiency is greatest, which means late
afternoon or early evening for most people,'' he said. ``Others say if you train when your
biological efficiency is least you will get a harder workout.''

Some elite athletes prefer morning workouts for reasons that have nothing to do with research
studies.

Deena Kastor, who holds the American marathon record, said her former coach and mentor, Joe Vigil,
insisted on morning workouts. He told her that there was more fluid between the vertebrae of the
spine after a night in bed, Ms.~Kastor said. And, she said, ``fluid made your spine more forgiving
and more able to absorb the pounding of running.'' She noted that she had been running in the
morning for the last 13 years ``with very little injury.''

But when people compete, if, for example, they want a personal best time, they might want to seek
out one of the few events that start late in the day. Or, even better, it might make sense for
endurance events, like marathons, to start in the afternoon instead of the morning, when they almost
always are held. Maybe they could be held later in the year, to avoid afternoon heat.

Dr.~Smolensky agreed.

``Most marathons start early under the guise that it's cooler then,'' he said. ``That needs to be
looked at.''

\section{China Stampede Kills 8 Students }

\lettrine{E}{ight}\mycalendar{Dec.'09}{11} students died and 26 were injured when hundreds of
students surged down a narrow staircase at a middle school in south-central China on Monday night,
the Chinese media reported.

The accident occurred as more than 400 students at Yucai Middle School in Hunan Province headed back
to their dormitories after evening classes. Most of them funneled down a single staircase about four
feet wide because it was close to their dorms and it was raining.

One provincial newspaper quoted students saying a few ``naughty'' boys tried to block the stairwell.
Xinhua, the state-run news agency, said a girl tripped.

Seven boys and one girl, aged 11 to 14, were trampled to death. An administrative aide at one
hospital said the injured students suffered from broken bones, concussions and punctured lungs.

Such accidents are strikingly common in Chinese primary and middle schools. In the past seven years,
more than 50 students have been killed and more than 200 have been injured in stampedes at nine
different schools, usually down staircases.

After six middle school students were crushed to death in a stairwell in eastern China in 2006,
education officials in Beijing vowed to improve school safety. Local officials were ordered to
inspect school staircases, corridors, cafeterias and bathrooms.

Many of the accidents have occurred in the evening, when Chinese students typically return to the
classroom to prepare for arduous exams that determine their academic futures.

At Yucai Middle School, a privately run 12-year-old school with 3,500 students and a reputation for
academic rigor, eight classrooms were filled with ``self-study'' students on Monday night. It was
unclear whether any teachers monitored their departure at 9:10 p.m.

CCTV, the state-controlled television broadcaster, broadcast video of a bloodstained landing between
the first and second floors. ``It's like a nightmare,'' one school official, Chen Xinwei, told
Xinhua. ``It's hard to believe those children are gone forever.''

Xinhua reported that the local education chief had been removed and that officials were
investigating how the accident occurred.

The injured were taken to three hospitals in Xiangxiang, a city of about 900,000. ``They are so
young,'' said an aide at Hospital No.~2. ``They looked so pitiful.''

The accident appeared to be the worst of its kind since 2002, when 21 middle school students were
crushed to death in northern China. The students were rushing down a pitch-dark staircase after
evening classes when the railing collapsed.

\section{Illegal Immigrant Students Publicly Take Up a Cause}

\lettrine{I}{t}\mycalendar{Dec.'09}{11} has not been easy for the Obama administration to deport
Rigoberto Padilla, a Mexican-born college student in Chicago who has been an illegal immigrant in
this country since he was 6.

On Thursday, Immigration and Customs Enforcement officials said they would delay Mr.~Padilla's
deportation for one year.

Mr.~Padilla's case had seemed straightforward to immigration agents who detained him for deportation
in January after he was arrested by the Chicago police for running a stop sign and charged with
driving under the influence.

But since then, students held two street rallies on his behalf and sent thousands of e-mail messages
and faxes to Congress. The Chicago City Council passed a resolution calling for a stay of his
deportation and five members of Congress from Illinois came out in support of his cause. One of them
was Representative Jan Schakowsky, a Democrat, who offered a private bill to cancel his removal.

Obama administration officials said they would review cases like Mr.~Padilla's as they arose. They
said the situation of Mr.~Padilla, 21, pointed to the need for an immigration overhaul that would
include a path to legal status for people in the United States illegally.

``We are committed to confronting these problems in practical, effective ways, using the current
tools at our disposal while we work with Congress to enact comprehensive reform,'' said Matthew
Chandler, a spokesman for the Department of Homeland Security.

Behind Mr.~Padilla's case -- and others in Florida of students who fought off deportation -- is
activism by young immigrants, many of them illegal, which has become increasingly public and
coordinated across the country, linked by Web sites, text messages and a network of advocacy groups.
Spurred by President Obama's promises of legislation to grant them legal status, and frustration
that their lives have stalled without it, young illegal immigrants are joining street protests
despite the risk of being identified by immigration agents.

With many illegal immigrants lying low to avoid a continuing crackdown, immigrant students have
become the most visible supporters of a legislative overhaul, which Mr.~Obama has pledged to take up
early next year. In the meantime, their protests are awkward for the administration, with young,
often high-achieving illegal immigrants asking defiantly why the authorities continue to detain and
deport them.

``Maybe our parents feel like immigrants, but we feel like Americans because we have been raised
here on American values,'' said Carlos Saavedra, national coordinator of a network of current and
former students called United We Dream.

``Then we go to college and we find out we are rejected by the American system. But we are not
willing to accept that answer,'' said Mr.~Saavedra, 23, a Peruvian who lived here illegally until he
gained legal status two years ago.

Young people who were brought to the United States by illegal immigrant parents draw a certain
degree of sympathy even from some opponents of broader legalization programs. Roy Beck, the
executive director of NumbersUSA, a group that has staunchly opposed a legal path for the estimated
12 million illegal immigrants, said in an interview that he could support legal status for some
young immigrant students. Mr.~Beck said he would do so, however, only if Congress eliminated the
current immigration system based on family ties and imposed mandatory electronic verification of
immigration status for all workers -- conditions that Democrats in Congress are not likely to
accept.

The students' goal is to gain passage of legislation that would give permanent resident status to
illegal immigrants who had been brought to the United States before they were 15, if they have been
here for at least five years, have graduated from high school and attend college or serve in the
military for two years.

Known to its supporters as the Dream Act, it has been offered in the Senate by Richard J.~Durbin,
Democrat of Illinois, and Richard G.~Lugar, Republican of Indiana. An effort to bring it to the
Senate floor was defeated in 2007, and proponents now consider it part of a package that includes a
path to legal status for illegal immigrants in general, an estimated 12 million people. Mr.~Beck
said he continued to oppose that proposal.

Many illegal immigrant students who were brought to the United States as children receive a shock
when they get ready to go to college. They are generally not eligible for lower in-state tuition
rates or government financial aid. In most states they cannot get drivers' licenses.

In recent years, student groups joined battles in several states for in-state tuition for illegal
immigrants, some successful and some not. This year, student organizers said, they worked to tie
those state efforts into a national network, hoping to match the mobilization networks of opponents
of the immigration overhaul, which proved far superior in the past.

The troubles for Mr.~Padilla began when he drove home after watching a football game and drinking
beer with friends. He ran the stop sign, and the traffic police arrested him because he did not have
a driver's license and had been drinking. Eventually, he pleaded guilty to a misdemeanor.
Immigration agents found him in the county jail.

Mr.~Padilla, now enrolled at the University of Illinois at Chicago, had no prior record and had been
an honors student and president of the Latino student organization at Harold Washington College,
which he attended for two years. Friends from both schools mobilized after his arrest.

Similar rallies took place in November in Miami, when immigration agents detained two brothers from
Venezuela who were illegal immigrants -- Jesús Reyes Mendoza, 21, a former student government
president at Miami Dade College, and his brother Guillermo, 25. Students from the college held a
protest in front of the immigrant detention center where the brothers were held.

``The undocumented youth are losing our fear of being undocumented,'' said Carlos Roa, an illegal
immigrant student from Venezuela who joined that rally. ``I'm public with this. I'm not hiding
anymore.''

Miami Dade College, with 170,000 students, has become a center for immigrant activism. After the
protests, and letters from Eduardo Padron, the college president, the immigration authorities on
Nov. 8 deferred the deportation of the Reyes brothers for one year.

\section{For Judges on Facebook, Friendship Has Limits}

\lettrine{J}{udges}\mycalendar{Dec.'09}{11} and lawyers in Florida can no longer be Facebook
friends.

In a recent opinion, the state's Judicial Ethics Advisory Committee decided it was time to set
limits on judicial behavior online. When judges ``friend'' lawyers who may appear before them, the
committee said, it creates the appearance of a conflict of interest, since it ``reasonably conveys
to others the impression that these lawyer 'friends' are in a special position to influence the
judge.''

In practice, of course, actual friends and Facebook friends can be as different as leather and
pleather, and the committee did recognize that online friends were not the same as friends in the
traditional sense. A minority of the panel would have allowed Facebook friendship, which it
characterized as more like ``a contact or acquaintance'' without conveying the notion of ``feelings
of affection or personal regard.''

But the committee's majority concluded that the possibility of the appearance of impropriety
required that they recommend against friending, said Judge T.~Michael Jones of the First Judicial
Circuit Court, a committee member. He emphasized that the committee's role was advisory, and that
the opinion ``does not have the force of a Supreme Court opinion'' in Florida.

The decision was issued last month, and was first reported Tuesday by the Legal Profession Blog.

Stephen Gillers, a legal ethics expert at New York University, said the Florida rule went too far.
``In my view, they are being hypersensitive,'' Professor Gillers said. He noted that the differences
within the committee probably indicated a generational gap, which Judge Jones said was not the case.
In the case of a truly close friendship between a judge and a lawyer involved in a case, the other
side can simply seek to disqualify the judge, Professor Gillers said.

Judges do not ``drop out of society when they become judges,'' he said. ``The people who were their
friends before they went on the bench remained their friends, and many of them were lawyers.''

Nonetheless, Professor Gillers said, many judges are careful not to socialize with friends during
cases that involve them.

One Florida county judge, Nina Ashenafi Richardson of Tallahassee, said the rule was ``probably a
good idea, just to avoid any perceptions of impropriety.'' Although the judge has a Facebook page
that a friend put together for her political campaign -- ``it was an amazing tool to get my message
out'' -- she said she had not used it since.

Still, legal sycophants can take heart: lawyers can declare themselves Facebook ``fans'' of judges,
the committee said, ``as long as the judge or committee controlling the site cannot accept or reject
the lawyer's listing of himself or herself on the site.''

\section{Threatened G.M. and Chrysler Dealers Win a Round in the House}

\lettrine{T}{he}\mycalendar{Dec.'09}{11} House approved legislation on Thursday that would grant
Chrysler and General Motors dealerships the right to challenge the companies' decisions to close
them in third-party arbitration\cite{arbitration}.

The provision was part of a spending bill that passed 221 to 202 and is now headed to the Senate. No
Republicans in the House voted for the bill.

Under the legislation, 2,000 or so car dealers that have already lost their franchises or are
scheduled to lose them would be allowed to enter binding arbitration. A neutral arbitrator would
decide whether to restore the dealer's franchise agreement. Among factors to be considered in the
decision are the profitability of the franchise from 2006 to 2009, how long the dealer has been in
business and the dealer's current economic viability.

The provision was announced this week by Representative Steny H.~Hoyer, the House majority leader,
and Senator Richard J.~Durbin of Illinois, the No.~2 Democrat in the Senate.

Lawmakers in both parties, especially the House, have expressed their support for car dealers since
Chrysler and G.M. announced plans earlier this year to shut many operations.

G.M. and Chrysler announced last week that they would re-evaluate the closings, offering either
reviews that would result either in the dealers having their franchise agreements restored or being
offered arbitration. At a conference call on Thursday joined by other Democratic supporters of the
dealers, Mr.~Hoyer declared that the companies' plans did not go far enough, suggesting that
Chrysler and G.M. would have decided whether they had fairly used their own criteria in closing
dealerships.

The arbitration provision in the legislation, said Mr.~Hoyer, Democrat of Maryland, was ``fair to
the dealers and fair to the companies.''

This week, after the compromise provision was announced, both auto companies released statements
indicating they were interested in working with lawmakers and dealers to find a resolution that
balances the interests of both sides.

On the Republicans' side on Thursday, the House minority leader, John A.~Boehner, gave qualified
support to the auto dealer provision, even as he was critical of the spending measure as a whole and
voted against the bill.

``I don't think it's quite as much as I'd have hoped, but it's something.'' Mr.~Boehner, an Ohio
Republican, said in a news conference. ``And I hope it does resolve some of the unfair closings that
I and many of my colleagues on both sides of the aisle have witnessed.''

Groups backing the auto dealers, like the Committee to Restore Dealer Rights and the National
Automobile Dealers Association, praised the House vote.

The House had voted in July to restore the franchise agreements for the Chrysler and G.M. dealers
facing closure. The White House opposed that provision, saying the companies' move to a leaner
dealer network was a crucial part of their reorganization. The federal government provided bailout
aid to both companies, giving it ownership stakes in each.

In June, Chrysler, now allied with Fiat, closed almost 800 dealerships, or roughly a quarter of its
dealer network.

G.M. informed more than 1,000 of its dealers earlier this year that their franchise agreements would
not be renewed when they expired in October 2010.

\section{Washington Fuss Over White House Hanukkah Party}

\lettrine{A}{t}\mycalendar{Dec.'09}{11} the first Hanukkah party in the Obama White House, a Jewish
student choir will sing in sweet harmony, the two young children of a soldier deployed in Iraq will
light a 19th-century silver menorah from Prague and President Obama and his wife, Michelle, will
greet more than 500 guests in a celebration that is expected to spill from the State Room to the
East Room.

But to the dismay of some administration officials, the plans for next week's party -- one of the
hottest holiday events for the nation's Jewish elite -- have been overtaken by feverish debate over
the size of the guest list, the language on the invitations and what this says (or does not say)
about Mr.~Obama's relationship with Jews.

President George W.~Bush, who began the tradition of White House Hanukkah parties, invited 600
people to his last party, administration officials say. But rumors spread wildly, first in the
Israeli press and then locally, that President Bush had invited 800 people and that the Obamas were
planning to invite only 400. (Administration officials say they have invited 550 people.)

The invitations have also caused some consternation because they make no mention of Hanukkah,
inviting guests to ``a holiday reception'' on Dec.~16.

In an opinion article published by JTA, the Jewish news agency, Tevi Troy, a former Bush
administration liaison to Jewish groups, warned that the Obama White House had given Jewish
Americans ``a number of reasons to fear that it takes its votes for granted.'' Mr.~Troy cited as
examples the administration's call for a freeze on Jewish settlements in the West Bank and the
decision to honor Mary Robinson, the former president of Ireland, who has been accused by some
Democratic lawmakers of anti-Israel bias.

Mr.~Troy said the reduced guest list created ``a nagging sense that there may be a studied
callousness at work here.''

His commentary, published on Nov. 23, and an article a week earlier in The Jerusalem Post, touched
off a flurry of news articles, blog postings and kitchen table discussions. This week, the Israeli
newspaper Yediot Achronot published photographs of President Bush lighting a menorah and Mr.~Obama
standing with Santa Claus alongside an article headlined, ``Obama Downsizes Hanukkah in the White
House.''

Rabbi Levi Shemtov, who is overseeing the process of making the White House kitchen kosher for the
party, said he was besieged with questions about the issue on a recent trip to Israel.

``I usually get asked when I came and how long I'm staying; this time, all anyone wanted to know was
whether I was getting invited to the White House Hanukkah party,'' said Rabbi Shemtov, who heads the
Washington office of the American Friends of Lubavitch, which lobbies for the Lubavitch movement.

Rabbi Shemtov, who has attended Hanukkah parties at the White House, said he raised an eyebrow when
he received his invitation, but noted that the Bush administration once sent invitations out with
Christmas trees on them.

``This is all one big overblown latke,'' the rabbi said.

``I feel that we need to save our communal kvetching in reserve for when it's more called for and
really matters,'' he continued.

Jewish Democrats accused Republicans of using the party for political ammunition. Advisers to
Mr.~Obama described the focus on the guest list as disappointing.

``Hanukkah is a wonderful holiday to celebrate, but that's not the whole ballgame, by any means, in
terms of outreach to the Jewish community,'' said Susan Sher, one of the president's two liaisons to
Jewish groups.

Ms.~Sher noted that Mr.~Obama held the first White House seder, invited the leaders of more than a
dozen Jewish organizations for a wide-ranging discussion at the White House in July, held a
conference call with 900 rabbis in August and videotaped a message to Jews for the High Holy Days in
the fall.

Administration officials also noted that White House records showed that Mr.~Bush never had more
than 584 guests at his Hanukkah parties. Most years there were fewer than 500, they said.

Josh Block, a spokesman for the American Israel Public Affairs Committee, the pro-Israel lobby,
which has disagreed with aspects of Mr.~Obama's policy toward Israel, praised his outreach and said
what mattered was that he would continue to hold the party.

Mr.~Troy, who said he was astonished by the reaction to his article, agreed. He said the Obamas were
``doing what they need to do in terms of outreach.''

As for the party, he is not expecting an invitation.

``The people who are invited will have a great time,'' he said. ``And a lot of people who didn't get
in will grumble. But you won't hear any grumbling from me.''

\section{Lessons From the War Zone}

\lettrine{O}{ne}\mycalendar{Dec.'09}{11} morning as a medical student on the surgery service, I
learned about a patient who had been hemorrhaging on the operating table the night before. The
intern who had assisted during the operation took great pains to describe every detail of the failed
efforts of several senior surgeons and the final, ultimately lifesaving, maneuvers of the department
chairman. ``He came in and just got control of the bleeding,'' the intern concluded, waving his
hands as if the chairman's work had involved magic.

``How did he manage that?'' one of my classmates asked.

``He's one of the best,'' the intern answered matter-of-factly. ``He was a surgeon in Vietnam.''

From then on, the idea that our chairman had served as a war doctor loomed large in our collective
imaginations. We were in awe of him, dazzled by his skill. He was someone who could handle even the
most challenging injuries on the fly, and without the luxury of resources that other doctors staked
their reputations upon, and we dreamed one day of becoming as good as he was.

But unlike him, we had been lucky enough to have come of age at a time and in a place when no war
was being fought. All that our minds could muster when we tried to imagine what he might have faced
were the technical challenges of caring for wartime injuries. And after any of his remarkable
operative feats, we would repeat to one another: ``Of course. He was a surgeon in Vietnam!''

Not long ago I saw my former chairman again for the first time in nearly 20 years. He was older and
slightly grayer, and still possessed of the stunning carriage that made him stand out in a room. But
as I stood before him, what I remembered was not his having served in Vietnam nor even his
remarkable skill, but his profound respect for the humanity of those he cared for, whoever they
might have been.

For what he had taught me, first as a doctor-in-training and later as a fully trained surgeon and
teacher, and what I recalled were the moments he urged us not to be tardy with our patients, to
change out of our surgical scrubs into neat street clothes when leaving the operating rooms to see
patients, and to be mindful, always, of even the smallest details of our patients' experiences.

War, we are now learning, can have wide-ranging, complex and not always positive effects on doctors
who serve. But one thing is certain: seeing the casualties of combat does more than produce war
doctors capable of caring for any injury. It has the paradoxical power to create doctors with an
extraordinary appreciation for all humanity.

Dr.~Chris Coppola illustrates this power in his memoir, ``A Pediatric Surgeon in Iraq'' (NTI
Upstream, 2009). Now practicing surgery at Geisinger Medical Center in Danville, Pa., Dr.~Coppola
served two medical tours of duty in Iraq as a lieutenant colonel in the Air Force. His memoir
chronicles his 10 months as one of two dozen surgeons among 300 military medical personnel stationed
in Balad, Iraq.

Dr.~Coppola pulls few punches in describing the war he saw. While a defter editing hand might have
resulted in a more polished and cohesive book, it is the unflinching honesty -- and poignancy -- of
Dr.~Coppola's voice that ultimately makes this a powerfully thought-provoking and occasionally
disturbing read.

He writes, for example, about a soldier so badly injured that even gender is mistaken. ``The soldier
is encased in a thick vinyl body bag. At the head of the stretcher, his puffy charred face juts out
of the opening of a silver survival blanket.''

But the soldier turns out to be a woman. ``I feel no pulse in her carotid artery, and above the
neck, her head is a swollen, burned, androgynous mess. I feel the wind sucked out of me, and a
pressure in my chest.''

Iraqis from both sides of the conflict are also among Dr.~Coppola's patients, some of them injured
during the daily late afternoon ambush or I.E.D. detonation. One victim is an insurgent who panicked
and prematurely detonated his explosive when he saw American soldiers coming.

``I look for signs in his swollen face of malice,'' Dr.~Coppola writes, ``but all I see is a dying
man who needs our treatment.''

The patient soon bleeds to death, but the event leaves Dr.~Coppola ``cold.'' The insurgent ``was on
a mission to kill my countrymen and may have succeeded if he were still alive. How am I supposed to
comprehend his death when my duty as a doctor to heal contradicts my duty as an officer to defend?''

Dr.~Coppola never offers pat answers. Over and over again he simply shows what he and his fellow
soldiers have chosen to do in unimaginably trying circumstances. Dr.~Coppola goes out of his way to
use his specialized training to treat Iraqi children who would otherwise die in understaffed and
underequipped Iraqi hospitals. On the military base where blood reserves are severely limited,
soldiers line up repeatedly to donate their own blood, for fellow servicemen and for Iraqis from
both sides of the conflict.

Even amid the horrors of war, these soldiers and doctors hold fast to a concern for others that
transcends not only politics but also their own mortality.

``What I learned just seeing war firsthand,'' Dr.~Coppola told me recently, ``was that people die
and it makes no sense whatsoever. I'm a doctor because of the reward of helping people; and that was
there every day in Balad. But civilians might just be going to the market or to a wedding and be
wiped out. Or I'd see the troops, these 20-year-old kids, who would be laughing, sending home e-mail
and opening packages one day, and the next day we would be calling family and doing everything we
could to keep them alive so their families could see them in Germany.''

It has been two years since Dr.~Coppola's last tour of duty. I asked him how he and his practice
have changed. ``I'm older and grayer,'' he said, ``and what I see now doesn't shock or bother me.''

He continued: ``Any time you interact with patients, there are two agendas. There is something you
as the physician want to achieve medically, and there is something that is important for the
patient. Since coming back, I've made it a more important part of my practice to understand my
patient's agenda.''

``I don't think it takes going to war to realize how important it is to appreciate our patients as
people, not as lines on a to-do list.''

What it takes is a book or a teacher, or for those of us who are fortunate enough, both.

\section{Hamstrung by Delays, Fitbit Explains and Tries to Deliver}

\lettrine{N}{icholas}\mycalendar{Dec.'09}{11} Cole was determined to shed a few pounds after the
holiday season. So when he heard about a nifty new fitness device called the Fitbit Tracker, a
wearable sensor that tracks movement, distance traveled and calories burned, he decided to take the
plunge and order one.

That was 13 months ago. Mr.~Cole is still waiting for the \$99 gadget to arrive, as are legions of
other eager customers.

``I haven't heard anything yet, but I'm hoping I'll be able to get it in a month or two,'' said
Mr.~Cole, a 20-year-old student living in New Brighton, Pa.

Others, unwilling to wait any longer, have forked over as much as \$350 to buy a Fitbit on eBay from
other consumers.

Fitbit's tale of expectation and delay is a classic start-up story: a couple of entrepreneurs with a
hot idea generate excitement, then run into a range of real-world problems in actually trying to
make their product and get it to customers. With bigger companies like Nike and Philips Electronics
making similar fitness devices, Fitbit runs the risk of getting stomped by competitors before it can
really get going.

But the company's unusual frankness about its problems may also help it survive its growing pains.

``They've been really open and transparent about where they are in the process, and that's made it
easier to tolerate the wait,'' Mr.~Cole said.

A prototype of the Fitbit Tracker was introduced in September 2008 by the co-founders, James Park
and Eric Friedman, at TechCrunch 50, an annual showcase of innovative products and Web services.

The thumb-size device uses an accelerometer to sense a user's movement, then translates that into
calories burned. In addition, users can wear the device at night to track the quality of their sleep
and can manually input their food consumption to get a better grasp of their overall health and well
being. A home base station collects information each time the user passes by and uploads it to
Fitbit.com.

The concept appealed to fitness enthusiasts, and Fitbit began taking preorders right away, expecting
to be able to ship them within a few months.

Instead, the company, which is based in San Francisco, found that it took eight months to refine its
prototype into something that was ready to manufacture, Mr.~Park said in a recent interview.

Mr.~Park and Mr.~Friedman are experienced entrepreneurs, having started two previous tech companies
together.

But this was their first foray into hardware. The Fitbit has more than 100 electronic components and
22 plastic and metal parts. Its complexity resulted in unexpected problems in making everything work
together.

For example, ``we would discover the product used more power than we'd originally thought and have
to decide if we wanted to include a larger battery,'' Mr.~Park said, ``which meant going back to the
drawing board to figure out how that impacted the size and form factor.''

The company also encountered setbacks during the testing phase, including equipment that at one
point got stuck in customs in Indonesia.

Even now, once the products arrive in California from the manufacturer in Singapore, Fitbit
employees must spend several minutes updating the software on each device before sending it out to
fulfill orders.

To assuage the growing restlessness of its customers, Fitbit has been e-mailing them and posting
updates on the company blog about the progress of the Fitbit Tracker, including photos and videos
showing production and testing of the devices.

``It's definitely a way to help people stay interested,'' Mr.~Park said. ``Most products on back
order are in a black hole of information.''

Some customers who have received their Fitbit, like Andrew Chen, a Silicon Valley entrepreneur, say
it is worth the wait. ``I love how simple and convenient it is,'' he said. ``I just put it on and
don't think about it.''

However, Mr.~Chen managed to get his device only after complaining on Twitter about his mounting
impatience over his back order. One of the investors in the company then shipped him a unit
directly.

Now, Fitbit says it is planning to ship the remainder of its back orders by the end of January, when
the stand-alone devices will also be available in retail stores.

Jon Callaghan, founder and managing partner of True Ventures, which led a \$2 million round of seed
financing to Fitbit, acknowledged that it would have been better not to miss the holiday shopping
season.

But it was more important to get the product right. ``I would not like a warehouse full of products
customers didn't like,'' he said. ``That would be disastrous and short-term thinking. We're going to
be around for next Christmas.''

Bernie Tenenbaum, an expert on small businesses who has no connection to Fitbit, said the delay was
not likely to do permanent damage to the company's reputation.

``Yes, they've missed some sales and some revenue opportunity'' and created some excitement
prematurely, ``but that's not fatal'' compared to shipping a flawed product, said Mr.~Tenenbaum,
managing partner of the investment firm China Cat Capital.

Still, Fitbit may have lost valuable momentum, he said. ``It's not as if there aren't enough other
shiny new tech things floating around at this time of year. That's not to say the passionate
aficionados won't be there with bated breath, but for everyone else, there are too many other
choices in the world.''

While the Fitbit struggles to fill old orders, the competition has moved in.

Many Web sites allow users to track various aspects of their health, similar to what users can do on
Fitbit.com.

And other companies are getting into the gadget side of the business. Philips, a consumer
electronics heavyweight, has introduced a fitness tracking device called DirectLife. And other
rivals have come out with the BodyBugg, a calorie-tracking monitor worn on the arm, and WakeMate,
which monitors sleep cycles and selects the ideal time to rouse its wearer within a given time slot.

So far, none of this appears to have diminished demand for the Fitbit -- which does not surprise Jim
Silver, an industry analyst and editor in chief of TimetoPlayMag.com.

In the age of Twitter and Facebook, word of mouth operates on hyperdrive, whether it is for the
Fitbit or this holiday season's hot toy, the battery-powered hamsters called Zhu Zhu Pets.

``Even people who aren't parents or associated with toys know what they are, which opens the door to
more sales, just because it's a hot item,'' he said.

\section{President Obama in Oslo}

\lettrine{A}{ccepting}\mycalendar{Dec.'09}{11} the Nobel Peace Prize on Thursday, President Obama
gave the speech he needed to give, but we suspect not precisely the one the Nobel committee wanted
to hear.

Mr.~Obama was appropriately humble. He said that ``compared to some of the giants of history who
have received this prize,'' his accomplishments ``are slight'' and suggested that he had been chosen
not so much for what he had done but for what he is expected to do.

He then acknowledged that most of what he called ``the considerable controversy'' surrounding his
selection came from the fact that he is ``the commander in chief of the military of a nation that is
in the midst of two wars.'' He made no apologies for that.

In a speech that was both somber and soaring, he returned again and again to Afghanistan, arguing
that the war was morally just and strategically necessary to defend the United States and others
from more terrorist attacks.

In a moving passage, he invoked the memories of Mahatma Gandhi and the Rev. Dr.~Martin Luther King
Jr., saying that without Dr.~King's vision, leadership and sacrifice, he never would have been
standing at that lectern in Oslo.

But he said he could not be guided by their examples alone. ``For make no mistake: Evil does exist
in the world. A nonviolent movement could not have halted Hitler's armies. Negotiations cannot
convince Al Qaeda's leaders to lay down their arms.''

In his introduction, the chairman of the Nobel committee, Thorbjorn Jagland, made only a brief,
forbearing reference to Afghanistan. He made clear that Mr.~Obama was chosen because of his
commitment, and early steps, to unwind the worst policies and abuses of George W.~Bush's presidency.
He pointed to Mr.~Obama's embrace of ``multilateral diplomacy,'' his offer to negotiate with Iran,
his decision to ban torture, his efforts to revive arms control negotiations and address global
warming. ``President Obama is a political leader who understands that even the mightiest are
vulnerable when they stand alone,'' Mr.~Jagland said.

It is a great relief to hear an American president described with such hope and respect. In his
speech, Mr.~Obama recommitted himself to those policies and principles, warning that ``we lose
ourselves when we compromise the very ideals that we fight to defend.''

What struck us most is how often Mr.~Obama used the war in Afghanistan to make his points. He said
that even as the United States confronts ``a vicious adversary that abides by no rules,'' this
country must remain ``a standard-bearer in the conduct of war.''

While he reserved the right to act unilaterally in a world where threats are ``more diffuse and
missions more complex,'' he said, ``America alone cannot secure the peace. This is true in
Afghanistan.'' And he directly challenged the widespread ambivalence and aversion toward the war in
the United States and in Europe. ``The belief that peace is desirable is rarely enough to achieve
it,'' he said.

When he announced his plan to send an additional 30,000 troops last week, Mr.~Obama's speech was
well argued but sounded more like a legal brief than an exemplar of presidential oratory. At the
time, he was coming out of months of difficult internal debates and girding himself for the
skepticism and disappointment of many members of his own party.

On Thursday in Oslo, Mr.~Obama argued his case far more eloquently.

We'll leave it to the philosophers to debate what is and what is not a just war. But we agree that
this war is a very difficult but necessary one.

We also know that there is no chance at all of winning it, and the broader fight against terrorism,
unless the United States hews to international standards and upholds its own ideals. That is Mr.
Obama's promise and his challenge going forward.

\section{N. Korea Sees Progress in Envoy's Visit}

\lettrine{N}{orth}\mycalendar{Dec.'09}{11} Korea said Friday that its differences with the United
States had been diminished during the visit this week of President Obama's special envoy, but that
it needed more talks before deciding whether to return to six-nation nuclear disarmament talks.

``Through businesslike and candid talks, both sides deepened mutual understandings, narrowed
differences in their respective views and identified not a small number of things in common,'' an
unidentified spokesman for North Korea's Foreign Ministry said, according to the state-run Korean
Central News Agency. ``Both sides agreed to continue to cooperate to narrow the remaining
differences.''

Stephen W.~Bosworth, Mr.~Obama's envoy, released a similar assessment of his three-day visit to
North Korea, which ended Thursday. He said that North Korea and the United States shared an
understanding of the necessity but not the timing of reconvening six-nation talks and carrying out a
2005 agreement to dismantle the North's nuclear weapons programs.

``We identified some common understandings on the need for and a role of six-party talks and the
importance of implementation of the 2005 joint statement,'' Mr.~Bosworth said at a news conference
in Seoul on Thursday. But, he said, ``it remains to be seen when and how'' North Korea will return
to talks.

He said that the talks -- the Obama administration's first senior-level negotiations with the North
Korean government -- were ``exploratory'' and that it would take more consultation with the other
nations in the six-nation talks before it would be clear when the North would rejoin the
discussions.

Secretary of State Hillary Rodham Clinton told reporters in Washington that ``for a preliminary
meeting it was quite positive.''

Although the North has indicated that it may return to the six-nation meetings, it has also insisted
on a new track of bilateral talks with the United States to negotiate the normalization of
diplomatic relations and a peace treaty to replace the truce that ended the Korean War in 1953. The
North said it would not give up its nuclear arms until it believed that the United States was no
longer ``hostile.''

Mr.~Bosworth said that once the six-nation talks reconvened and gained ``a significant traction'' on
denuclearizing the Korean Peninsula, the United States and its allies would be ready to discuss the
incentives offered in the 2005 agreement: economic aid, security guarantees, normalized relations
with Washington and a peace treaty.

After a stop in South Korea, Mr.~Bosworth was heading to China, Japan and Russia to brief officials.
Those four countries have been the other participants in the six-nation talks.

Mr.~Bosworth updated South Korean officials on the talks with North Korea's first vice foreign
minister, Kang Sok-ju, who is considered the main architect of his country's nuclear policy and one
of the closest aides to North Korea's leader, Kim Jong-il. Mr.~Bosworth also met with the North's
chief negotiator, Kim Kye-gwan. Mr.~Bosworth's visit was meant to help Washington decide whether it
should maintain sanctions or switch to engagement to prod North Korea toward giving up its nuclear
weapons.

American and South Korean officials remain unconvinced that the North would give up its nuclear
weapons, fearing that it wants to use a new round of talks to escape sanctions and obtain aid. North
Korea has struggled to maintain control of its economy, most recently by revaluing its currency.

In the past year, North Korea has tested ballistic missiles and a long-range rocket, detonated a
nuclear device and reactivated its plutonium production plant. It has also declared that it began
enriching uranium for possible use in nuclear weapons. But recently, it has begun to reach out.

On Thursday, in a rare admission of a domestic problem and a sign of new openness, North Korea
accepted a South Korean offer to send it swine flu medication, South Korean officials said. A day
earlier, the North acknowledged nine cases of H1N1 influenza in Pyongyang, the capital, and Sinuiju.

\section{In Holiday Crush, a Fatal Shootout in Times Sq.}

\lettrine{I}{n}\mycalendar{Dec.'09}{11} the preholiday swirl of Broadway in Times Square, where
wide-eyed tourists and time-starved New Yorkers pass through an urban playground vastly different
from the one Damon Runyon inhabited, two people came face to face on Thursday morning: a police
sergeant and a street vendor.

From inside his unmarked patrol car, the sergeant -- in plain clothes -- recognized the vendor. He
left the car and asked a question. The vendor answered by darting off through the crowd, the
sergeant at his heels, running along Broadway.

A quick turn down a side street, another fast turn into a covered outdoor promenade -- a gift shop
on one side, the revolving doors of a hotel on the other. The vendor pulled something from under his
coat: a semiautomatic gun, its 10-inch clip filled with 30 bullets, the police said. Then there were
shots, two from the vendor and four from the sergeant.

The vendor fell to the pavement, dying. The sergeant was not wounded. And suddenly, the crossroads
of the world was a very public crime scene, and the tourists had a story they could tell when they
got back home.

A visitor from Australia, Suzanne Davis, 42, stopped to take images with a video recorder. ``It's my
first day in New York, so it makes very real what you see in the movies,'' she said.

The drama unfolded around 11:15 a.m., when the sergeant, identified by police officials as
Christopher Newsom, 41, confronted two men outside 1515 Broadway, the skyscraper between West 44th
and 45th Streets that has the studio most famous as the setting of the old MTV show ``Total Request
Live.''

Sergeant Newsom is normally assigned to Midtown South Precinct's peddler unit, but was on loan for
the day to an anticrime unit that patrols Times Square. He recognized the two men, who the police
said were taking part in a street ruse using compact discs to hoodwink the gullible. Vendors would
approach someone on the street, ask his or her name, write the name on a CD and demand \$10.

Trailed by a couple of other anticrime officers in plain clothes, Sergeant Newsom asked for the
vendors' tax stamp, something the two men would need to sell merchandise on the street legally.

One of the men, Raymond Martinez, 25, of the Bronx, took off, dashing north and making a left on
West 45th Street, the police said. The sergeant ran after him, following as he turned into the
promenade leading to the main entrance of the Marriott Marquis Hotel. The police said they had
pieced together their account of what happened next from interviews with witnesses and from video
from a surveillance camera.

A security guard at the Marquis reported hearing the sergeant yell, ``Stop! Stop! Show me your
hands!'' The guard told the police he saw the vendor pull something from his coat -- the submachine
gun, which the police said later had been stolen in Virginia in October.

The video showed Mr.~Martinez turning, the police said, but he moved out of camera range. Police
officials, who did not immediately release the video, said it also showed the sergeant reaching for
his weapon and raising it.

It also showed Sergeant Newsom, who has been on the force for 17 years, raising his left arm over
his chest in hopes of protecting his heart. It is a defensive move rookies are taught in the Police
Academy. Police officials were astonished that the sergeant, less than 15 feet from the stubby
barrel of a semiautomatic weapon with no hope of taking cover, was cool-headed enough to remember to
do so.

Police Commissioner Raymond W.~Kelly said Mr.~Martinez fired first, getting off two shots. Then his
gun jammed.

Mr.~Kelly said the sergeant fired four shots. All four hit Mr.~Martinez: in the chest, below the
neck and in the left arm; he also suffered a graze wound to the right arm.

Of the two bullets that Mr.~Martinez fired, one smashed into the plate-glass window of a souvenir
shop in the promenade; the glass shattered, leaving a bullet hole to the right of a ``Wicked'' snow
globe souvenir. The other bullet hit a glass panel near the marquee for the Marquis Theater, in the
same building as the hotel.

The police said they found 27 more bullets in Mr.~Martinez's gun.

Dave Kinahan of Boston was standing about 60 feet away, having just pulled up to the hotel's valet
stand in his sport-utility vehicle. His wife, Emma, who is pregnant, was inside, along with their
two young children and Mr.~Kinahan's parents. Mr.~Kinahan, who was unloading luggage, described
hearing at least three gunshots.

He said he turned around in time to see Sergeant Newsom, standing over the vendor, shoot once at a
downward angle. Mr.~Kinahan said the sergeant leaned down on the man's body to subdue him and
yelled, ``Call 911! Call 911!'' Mr.~Kinahan also said the sergeant pulled a badge from around his
neck and held it aloft.

``At first I wasn't sure: Was it real, or was it a movie?'' Mr.~Kinahan said, adding that when he
realized it was real, he was stunned.

Emergency workers took Mr.~Martinez to St.~Luke's-Roosevelt Hospital Center, where he was pronounced
dead.

Police officials said he had been issued a summons earlier this year -- for selling compact discs
and DVDs without a tax stamp -- at the same spot where the chase began. They also said he was wanted
for questioning in connection with an assault in the Bronx, and on a warrant for failing to appear
in court on a disorderly conduct charge, also in the Bronx.

The police said they found business cards for several Virginia gun dealers in Mr.~Martinez's pocket.
One was from Gary A.~Lewis, who runs Gary's Guns \& Transfers in Manakin-Sabot, a pair of villages
northwest of Richmond.

Mr.~Lewis said he could not explain how Mr.~Martinez came to be carrying his card.

``I don't know who the guy is,'' Mr.~Lewis said by telephone on Thursday. ``It's not my gun. I don't
know anything about the story. I don't sell those types of guns.''

The police said a cryptic comment that referenced a 1985 martial-arts movie was written on the back
of the card: ``I just finished watching 'The Last Dragon.' I feel sorry for a cop if he think I'm
getting into his paddy wagon.''

The police said they did not know whether the handwriting was Mr.~Martinez's.

What was more certain was that the scene on Broadway was surreal to tourists and New Yorkers alike.

Emer Rooney, 33, on the last day of a visit from Ireland, took pictures of the scene. She said that
despite the shooting, she had never felt unsafe in New York, noting the presence of ``all the police
officers.'' She cited the shooting, in fact, as one of the more exciting moments of her trip --
right up there, she said, with recovering lost luggage at the Port Authority Bus Terminal and
getting tickets to ``Wicked.''

\section{Stephanopoulos to Join 'Good Morning'}

\lettrine{A}{bc}\mycalendar{Dec.'09}{11} News has made official the long-rumored appointment of
George Stephanopoulos to succeed Diane Sawyer as the anchor of ``Good Morning America'' -- and its
top executive promised on Thursday that the show would move ``in a new direction,'' toward more news
and substantive information.

David Westin, the president of ABC News, who made the announcement on Thursday, conducted a
three-month search for a successor for Ms.~Sawyer, who was named in September as the next anchor of
ABC's evening news program, ``World News,'' succeeding Charles Gibson.

Throughout the process, Mr.~Stephanopoulos, the anchor of ABC's Sunday morning discussion program
``This Week'' and the chief political correspondent for the network, was considered the
front-runner, although ``there were many candidates who could have done the job,'' Mr.~Westin said
in an interview by phone.

With his background in politics -- he worked in the Clinton White House -- and political reporting,
Mr.~Stephanopoulos had, according to some accounts, insisted that he not lead some of the lighter
segments, like cooking features and celebrity interviews, that are often a staple of morning
television.

But ``that was not as much a point of discussion as has been reported,'' Mr.~Stephanopoulos said in
an interview by phone. ``David told me he was already thinking of reorienting the show.''

Mr.~Westin said that even before Mr.~Gibson made clear his intention to retire from the evening
anchor position -- leading to the appointment of Ms.~Sawyer on ``World News'' -- he had decided that
``we needed to be providing more information about things that matter to our audience.''

Mr.~Stephanopoulos said: ``I think I have a pretty good sense of what makes sense to do and what
doesn't make sense. If there's something that makes you uncomfortable, maybe it shouldn't be on the
show.''

Jim Murphy, the show's executive producer, pointed out that Mr.~Stephanopoulos had substituted as
anchor on many occasions. ``George's No.~1 attribute is his interviewing skill,'' Mr.~Murphy said.
``But he has interests. It's not like we're getting some wonk who was born on a Sunday morning
show.''

Mr.~Stephanopoulos, 48, is married to the actress Ali Wentworth; they have two young daughters. He
acknowledged that the decision to join the show was a difficult one for him. ``I was definitely
torn,'' he said. ``I love the job I have. And Ali and I love our life in Washington. That said,
'G.M.A.' is an incredible franchise.''

The show has long been in second place in the morning behind NBC's ``Today,'' and the gap between
the shows is wide at the moment. Last week ``Today'' beat ``Good Morning America'' by its widest
margin in three years -- more than 1.5 million viewers -- and continued the longest winning streak
in television history when it recorded a 14th consecutive year of finishing first every week.

Mr.~Westin said the pursuit of the top-rated spot in the morning remains a priority and that the
decision to pair Mr.~Stephanopoulos with the holdover anchor Robin Roberts ``is going to help us to
be No.~1.''

The pursuit begins immediately: Mr.~Stephanopoulos starts Monday.

In another adjustment, Juju Chang will succeed Chris Cuomo as the show's news reader.

Mr.~Cuomo, who had sought the show's anchor job, agreed on Thursday to stay with ABC News as the
co-anchor of the network's prime-time magazine show ``20/20.''

ABC also announced that Mr.~Stephanopoulos would remain in his old position as anchor on ``This
Week'' until a successor was chosen. Mr.~Stephanopoulos, who said he was proud of how much progress
that show had made in closing the ratings gap in recent months with another perennial NBC leader,
``Meet the Press,'' expressed regret that he would not be able to keep both jobs permanently.

``I just can't do both jobs for long and also help raise a family,'' he said. He will, however,
retain his job as the chief political reporter for ABC News.

\section{Accepting Peace Prize, Obama Offers 'Hard Truth'}

\lettrine{P}{resident}\mycalendar{Dec.'09}{11} Obama used his acceptance of the Nobel Peace Prize on
Thursday to defend the idea that some wars were necessary and just, remind the world of the burden
the United States had borne in the fight against oppression and appeal for greater international
efforts for peace.

``We must begin by acknowledging the hard truth: we will not eradicate violent conflicts in our
lifetimes,'' Mr.~Obama said, addressing the paradox of receiving an award for peace as commander in
chief of a nation that is escalating the war in Afghanistan as it continues to fight in Iraq.
``There will be times when nations -- acting individually or in concert -- will find the use of
force not only necessary but morally justified.''

He delivered a mix of realism and idealism, implicitly criticizing both the Rev. Dr.~Martin Luther
King Jr.~as inadequately appreciating the dangers of the world, and President George W.~Bush as too
quick to set aside fundamental American values in pursuit of security. And he embraced the concept
of American exceptionalism, the idea that the United States has a special role as a defender of
liberty, even as he promoted multilateralism.

In that way, he continued a pattern evident throughout his public career of favoring pragmatism over
absolutes.

The address -- delivered at once to a European audience that has grown skeptical about American
power and to a domestic audience watching closely to see how he would handle the acceptance of an
award that even he acknowledged he did not yet deserve -- represented one of the broadest
declarations of his foreign policy doctrine. He said that others deserved the award more, noting
that his ``accomplishments are slight,'' but he accepted the prize with a strong endorsement of
America's place in the world.

``Whatever mistakes we have made, the plain fact is this,'' Mr.~Obama said. ``The United States of
America has helped underwrite global security for more than six decades with the blood of our
citizens and the strength of our arms.''

The Nobel lecture, a 36-minute address that the president and his aides completed on an overnight
flight from Washington, carried echoes of several American presidents, from Jimmy Carter to
Mr.~Bush, but Mr.~Obama singled out one above all: John F.~Kennedy.

Mr.~Obama cited Mr.~Kennedy's focus on ``not a sudden revolution in human nature but on a gradual
evolution in human institutions.''

Mr.~Obama called for more robust international sanctions against nations like Iran and North Korea
that defy demands for them to curtail their nuclear programs.

Weeks after being criticized for not speaking out more publicly in defense of human rights while in
China, he suggested that quiet diplomacy was sometimes the most productive path, even if it ``lacks
the satisfying purity of indignation.''

The ceremony was the focal point of a series of events celebrating Mr.~Obama's entry into the ranks
of Nobel laureates. On Thursday night, the president and his wife, Michelle, appeared in a window of
the Grand Hotel, waving to thousands of people below who had gathered for a torch-light parade.

Trumpets sounded when Mr.~Obama walked down the long aisle of a soaring auditorium to deliver his
address. He escorted his wife, who took her seat in the front row, before he assumed his position on
the stage and faced the king and queen of Norway.

The Nobel chairman, Thorbjorn Jagland, opened the ceremony by explaining how the committee came to
its decision two months ago. He said Mr.~Obama's leadership had been a ``call to action for all of
us.'' As he invoked the story of Dr.~King, the winner of the prize in 1964, he turned to Mr.~Obama,
saying, ``Dr.~King's dream has come true.''

Mr.~Obama pursed his lips and nodded gently as the audience applauded loudly. When he was presented
his gold medal and Nobel diploma, he received a standing ovation that stretched for more than a
minute. The crowd did not rise again until the conclusion of his remarks.

Mr.~Obama's speech was sober, with his remarks only sparingly interrupted by applause. He was
applauded when he renewed his pledge to ban torture and close the prison at the American base at
Guant\'anamo Bay, Cuba.

``We lose ourselves when we compromise the very ideals that we fight to defend,'' Mr.~Obama said.
``And we honor those ideals by upholding them not when it is easy, but when it is hard.''

To a European audience of academics, diplomats and Nobel laureates, he said there was ``a deep
ambivalence about military action today,'' which he said he suspected was rooted in ``a reflexive
suspicion of America.'' But he offered a forceful defense of the United States, saying the lessons
of history should ease those suspicions. And he urged his audience to envision a hopeful future.

``Let us reach for the world that ought to be,'' he said, ``that spark of the divine that still
stirs within each of our souls.''

He did not dwell on the specifics of his announcement last week that he would send 30,000 more
American troops to Afghanistan. But that decision, which attracted scores of peaceful demonstrators
here, set the framework that Mr.~Obama returned to again and again as he sought to explain his
policy as an extension of the post-World War II system that contained the cold war.

``A decade into a new century, this old architecture is buckling under the weight of new threats,''
Mr.~Obama said. ``The world may no longer shudder at the prospect of war between two nuclear
superpowers, but proliferation may increase the risk of catastrophe. Terrorism has long been a
tactic, but modern technology allows a few small men with outsize rage to murder innocents on a
horrific scale.''

Mr.~Obama, who is scheduled to stay in Oslo for about 26 hours, miffed some Norwegians by not
participating in some of the traditional events surrounding the peace prize ceremony, including a
luncheon and a concert.

Mr.~Obama, sensitive to the criticism, explained the brevity of his visit. ``I only wish that my
family could stay longer in this wonderful country,'' he told reporters, ``but I still have a lot of
work to do back in Washington, D.C., before the year is done.''

The president is scheduled to return to Washington on Friday.

\section{Crew of Detained Plane Denies Knowledge of Arms}

\lettrine{A}{}\mycalendar{Dec.'09}{14} shipment of arms and apparently sophisticated missiles from
North Korea seized here on a tip from American intelligence agencies has set off a series of
investigations, as officials try to determine whether the cargo was headed to South Asia or the
Middle East.

The Obama administration welcomed the interception\cite{interception} by the Thai authorities as
evidence that it had scored a success in its effort to enforce a United Nations Security Council
resolution banning weapons exports by the North Korean government, an attempt to cut off the
country's most profitable export.

A senior White House official in Washington, declining to speak on the record because the
interception on Friday had required high-level coordination among intelligence agencies that watched
the shipment as it was loaded in Pyongyang, said the seizure was evidence that the administration
would ``aggressively and rigorously implement'' the Security Council resolution, passed in June
after North Korea set off its second nuclear test.

The resolution, No.~1874, authorizes any country to inspect and seize North Korean weapons shipments
that pass through its territory, regardless of the cargo's ultimate destination.

The Obama administration is trying to show that it will choke off the North's illicit exports even
while attempting to reopen nuclear talks. In interviews in recent weeks, several officials said they
believed that the Bush administration, which came to office seeking to topple the North Korean
government of Kim Jong-il, mistakenly eased pressure on him while pursuing disarmament accords,
almost all of them since renounced by the North.

But on Sunday, United States officials were largely silent on the destination of the shipment,
saying they were still trying to sort through conflicting reports and misleading flight plans.

One senior official said he believed that the shipment was headed to Iran, a major buyer of North
Korean missiles and arms, some of them passed on later to Hezbollah or Hamas.

Others said they thought the more likely destination was Pakistan, which also has a long history of
purchases from North Korea, but which has been under heavy American pressure to cut its ties to
Mr.~Kim's government.

Thai officials arrested the five-man crew on Friday after seizing the large cargo craft at Bangkok's
military airport, where the crew had landed to refuel. It was unclear why the crew would have chosen
a close American ally for a refueling stop, rather than neighboring Myanmar, which has deep business
and political ties to North Korea.

Thai officials said they intended to charge the crew members with possession of weapons of war. On
Monday, a court here approved a 12-day detention for the five men.

But in their first interview since they were arrested, the crew members insisted that they did not
know they had been transporting an arsenal of at least 30 tons of missiles, grenade launchers and
other weapons, some of which the Thai authorities were still trying to identify.

``They said it was oil drilling equipment,'' said Viktor Abdullayev, the plane's co-pilot.

``That's what the manager told us,'' he said, referring to the crew's employer, a civilian cargo
company from the former Soviet republic of Georgia.

Mr.~Abdullayev and his colleagues said they started their trip in Ukraine, picked up cargo in North
Korea and were traveling back to Ukraine via Thailand, Sri Lanka and the United Arab Emirates. They
declined to say in which of those locations the cargo was meant to be delivered. The emirates,
officials note, are often used as a transit point to Iran.

Mr.~Abdullayev said it never occurred to him to ask about the cargo. ``I have no interest in what I
carry,'' he said. ``Like a truck driver: just keep driving.''

An intriguing hint about his cargo came from a photograph published in Thailand. It showed a series
of rockets stacked in the cargo hold, with the crates marked ``K 100.'' The rockets were visible,
though partly shrouded. Numbers on individual crates that ran in sequence -- like 78, 83, 86 and 87
-- presumably denoted the individual weapons in the shipment.

Charles P.~Vick, a missile expert at GlobalSecurity.org, a research group in Alexandria, Va., said
that if the markings were correct, the rockets might be K-100s, a type of Russian missile designed
to destroy sophisticated radar planes. They are advertised as ``Awacs killers,'' a reference to the
Airborne Warning and Control System planes used by the United States, Israel and soon India, which
can orchestrate combat plans.

Mr.~Vick noted that the diameter and length of the packaged missiles in the photograph appeared to
match the specifications of the K-100: 16 inches wide and 20 feet long.

``It's just a guess,'' he cautioned.

But he said it was also possible that the tubular weapons might be smaller artillery rockets packed
end to end to fit in the large metallic crates. Those types of rockets are ``the kind of thing''
that Hamas and Hezbollah use against Israel, he said.

Panitan Wattanayagorn, the Thai government spokesman, said in an interview that the aircraft, a
Russian-made Ilyushin 76 registered in Georgia, passed through Bangkok twice -- on its way to North
Korea and on its return.

Four of the crew members are from Kazakhstan: Mr.~Abdullayev, Ilyas Issakov, Aleksandr Zrybney and
Vtaliy Shurmnov. The fifth, Mikhail Prtkhou, is from Belarus.

The aircraft was searched on the return journey after the Thai authorities were tipped off by
American officials that it might be carrying weapons. The crew members were detained and the cargo
was confiscated\cite{confiscate}, but not immediately. The crew members had enough time to buy six
large bottles of beer at a duty-free shop, but those were confiscated in the detention center where
they have been kept since.

``There's much more to this story than what has been made public so far,'' said Bertil Lintner, an
author who has written extensively on North Korea. ``Why would you refuel in Thailand?''

Mr.~Lintner said that if the aircraft had landed in neighboring Myanmar, which is ruled by generals
friendly to the North Korean government, ``there would have been no problem.''

No major seizures of weapons have been made public since the passage of Resolution 1874. This summer
for about three weeks, the United States Navy tracked a North Korean freighter suspected of carrying
banned cargo, and the ship eventually turned back to its home port without incident. The resolution
does not permit the boarding of ships on the high seas.

\section{Obama Says Afghan Buildup Must Show Results}

\lettrine{P}{resident}\mycalendar{Dec.'09}{14} Obama said in a taped interview that military
officials should know by the end of December 2010 whether a strategy to secure population centers in
Afghanistan is meeting its objectives.

``If the approach that's been recommended doesn't work, we're going to be changing approaches,''
Mr.~Obama said in the interview that aired on CBS's ``60 Minutes'' Sunday night.

The interview was conducted last week at the White House, before the president traveled to Europe to
accept the Nobel Peace Prize on Thursday. Much of the interview focused on Mr.~Obama's decision to
send an additional 30,000 troops to Afghanistan.

Under Mr.~Obama's strategy, the additional 30,000 troops he added to forces in Afghanistan would
begin withdrawing in July 2011 as part of a transition phase.

Gen.~Stanley A.~McChrystal, the top American military commander in Afghanistan, has suggested that
the United States wants the new forces to blunt the Taliban's recent momentum and to buy time to
train Afghan soldiers and police officers to take over security duties.

Mr.~Obama said that he had set the 2011 deadline to put Afghan officials on notice that the United
States does not intend to carry the entire burden of securing the nation.

``That's not what the American people signed off for when they went into Afghanistan in 2001,''
Mr.~Obama said. ``They signed up to go after Al Qaeda.''

Mr.~Obama also said that the United States would need more cooperation from Pakistan in pursuing Al
Qaeda because tribal territories that straddle the border along Afghanistan and Pakistan harbor
enemy fighters.

The interview came as Mr.~Obama approached his first anniversary in the White House and some of the
public's enthusiasm for his agenda has waned. Ever since Mr.~Obama took office, critics of his
leadership style have accused him of tackling too many initiatives at once. Recently, his approval
rating in some polls has dwindled to 50 percent or below.

The interviewer, Steve Kroft also asked the president his thoughts about some Wall Street banks that
had recently recovered enough to repay government loans, but once again are giving large bonuses to
employees. At three of the largest banks, bonuses are expected to total about \$30 billion.

``I did not run for office to be helping out a bunch of fat-cat bankers on Wall Street,'' Mr.~Obama
said. He said the only firms paying out such large bonuses were the ones that had paid back money
used from the Troubled Assets Relief Program, or TARP.

``What's most frustrating me right now is you've got these same banks who benefited from taxpayers'
assistance who are fighting tooth and nail with their lobbyists up on Capitol Hill fighting against
financial regulatory reform,'' Mr.~Obama said. He added that he thought that giving the bonuses may
have been a motivation in some cases for banks to repay their TARP loans.

Mr.~Obama was also asked if he thought the latest health care bill would pass.

``Yes,'' he replied, then added that he thought the bill would pass before Christmas Day. Health
care reform was a signature issue during Mr.~Obama's presidential campaign. This week, the
Congressional Budget Office was expected to complete the latest cost projections of the latest
version of the health care bill being promoted by Senator Harrry Reid, the majority leader.

The interview also touched on the Nov.~24 incident in which a Virginia couple, Michaele and Tareq
Salahi, slipped past multiple layers of high-level White House security to attend Mr.~Obama's first
state dinner without being on the guest list. The breach put a spotlight on Desiree Rogers, the
White House social secretary, who has acknowledged that no one from her office was at the checkpoint
to help identify guests.

Mr.~Obama said the incident angered him, but he did not directly address Ms.~Rogers, who was invited
to testify before the House Homeland Security Committee, which held a hearing into the security
lapse. Ms.~Rogers did not appear and the White House press secretary, Robert Gibbs, cited the
separation of powers and a history of White House staff not testifying before Congress.

``I was unhappy with everybody who was involved in the process,'' Mr.~Obama said. ``It was a
screw-up.''

\section{Chinese Company Seeks to Make Old Saab Models}

\lettrine{A}{s}\mycalendar{Dec.'09}{14} some of the world's best-known auto brands stumble in the
global downturn, the Chinese auto industry is again on the prowl\cite{prowl} for bargains.

This time, Beijing Automotive Industry Holding is moving to get a piece of Saab, striking a
tentative accord on Sunday for the right to produce versions of the beleaguered Swedish
manufacturer's older models in China.

If completed, people close to the negotiations said on Sunday, the deal would not prevent a separate
buyout of all of Saab, a unit of General Motors that was left flailing after negotiations last month
between G.M. and the high-end Swedish automaker Koenigsegg to buy it failed.

The fate of Saab's 3,500 workers in Trollhattan, Sweden, remains uncertain. But a Saab official
briefed on the discussions said that a pact with the state-owned Beijing Automotive ``would be good
for Sweden, good for China and good for Saab.''

As American automakers have contended with huge losses -- and bankruptcy in the case of G.M. and
Chrysler -- Asian buyers have stepped in to scoop up several of their luxury brands. In March 2008,
Tata Motors of India bought Jaguar and Land Rover from Ford for \$2.3 billion, and Ford is now in
final negotiations to sell Volvo to Zhejiang Geely Holding of China.

A spokesman for G.M. declined to comment on Sunday about Saab's talks with Beijing Automotive, which
is based in Beijing. Several buyers are reportedly still negotiating to buy Saab. These include
Renco, owned by the American financier Ira L.~Rennert, and Spyker Cars, a specialty automaker in the
Netherlands.

Spyker has confirmed its interest in acquiring Saab, but a spokesman declined to comment further.
Spyker sells 30 to 50 high-performance sports cars a year, which are made to individual order and
cost just under a quarter-million dollars each.

As was the case for Koenigsegg, taking over Saab would mean a large increase in production as well
as a formidable business challenge for Spyker, especially given G.M.'s inability to succeed in
making Saab profitable.

It would also expose Spyker and its Russian backers to more public scrutiny, some of it potentially
unwelcome.

The main investor in Spyker is the Russian bank Convers Group, which is controlled by Alexander
Antonov, a Russian tycoon who was shot seven times and reportedly lost a finger in a failed
assassination attempt in Moscow in March. His son Vladimir Antonov, a 34-year-old banker who is a
top executive at Convers, is chairman of Spyker.

Like China, Russia is trying to increase its profile in the global auto industry, and Sberbank of
Russia was a major investor in the failed effort by Magna International to acquire G.M.'s European
operations earlier this year.

Beijing Automotive was one of Koenigsegg's backers in the earlier Saab acquisition plan, potentially
contributing \$200 million to \$300 million toward the deal. The Chinese automaker has long coveted
Saab's engineering expertise, but the leaders of Saab are eager to prevent a breakup of the company,
which has been making cars for 60 years.

The new agreement would not cover Saab's new version of the 9-5 luxury sedan, scheduled to arrive in
showrooms in April. It is the first update of Saab's top-end car in 12 years.

The new 9-5 would compete with other European luxury cars like the Audi A6 and BMW 5 series and be
priced at about \$40,000.

In 1990, G.M. paid \$600 million for half of Saab, and anted up \$125 million in 2000 for the rest
of the company, which has long commanded a small, loyal following in the United States for its
idiosyncratic but stylish models.

But sales suffered in recent years as drivers believed Saab had lost its distinctive identity under
the sway of Detroit.

The new 9-5, company officials said, is an effort to revive Saab's traditional appeal in Europe and
win back loyal customers in the United States while increasing production in Trollhattan.

If Saab production in Sweden survives, it will be a boost to the country's industrial base. The area
around Trollhattan and Goteborg to the south is home to both Saab and Volvo factories, and a network
of auto-parts makers and other suppliers in the region is dependent on the two companies.

``The Saab organization is still intact,'' said Eric Geers, a spokesman for Saab in Sweden. ``We're
very excited about the new 9-5, which has already been received positively by enthusiasts around the
world.''

While the survival of Volvo and Saab would be welcome news for their workers and the Swedish
economy, it could worsen the broader problem of automobile industry overcapacity in Europe.

Many European automakers have had steep losses during the economic downturn, but a combination of
strong unions and political pressure from governments has prevented factory closures or layoffs
across the Continent.

What's more, under their new owners, Saab and Volvo would be likely to try to raise production,
threatening profit margins in the middle to upper range of the auto market.

Under G.M.'s control, Saab was hobbled by a small range of models as well as a sense that it was
less connected to its Scandinavian base in terms of engineering and design.

``The whole point of the new business plan,'' Mr.~Geers said, ``is everything comes back to Sweden
and we emphasize the true Swedishness of the brand.''

``We have a loyal base of customers and dealers and employees around the world,'' he added.
``Everybody is extremely eager to start up the new Saab.''

\section{China's Economic Power Unsettles the Neighbors}

\lettrine{I}{n}\mycalendar{Dec.'09}{14} the Dickensian depths of the Dunia Metal Works here, all is
cacophony\cite{cacophony}: the bam bam bam of grease-drenched punches; the rhythmic clank of
unspooling steel wire; the storm and stress of glinting, freshly minted nails cascading onto a broad
metal table for boxing.

But for all the industrial din, Dunia is undergoing a painful slump\cite{slump}. Today it runs at 40
percent of its capacity, its domestic nail business imperiled -- and its exports wiped out -- by
cheaper Chinese alternatives.

``We have been competing with the Japanese and the Koreans,'' said Juniarto Suhandinata, the
factory's director. ``But the Chinese -- no chance.''

The Chinese are tough competitors, and Dunia is hardly the first to find out. But Mr.~Suhandinata's
lament speaks to something different: a sense of disquiet, even in developing Asian nations in
Beijing's orbit, over the implications of China's swift, seemingly boundless economic growth.

China has long claimed to be just another developing nation, even as its economic power far
outstripped that of any other emerging country.

Now, it is finding it harder to cast itself as a friendly alternative to an imperious American
superpower. For many in Asia, it is the new colossus.

``China 10 years ago is totally different with China now,'' said Ansari Bukhari, who oversees
metals, machinery and other crucial sectors for Indonesia's Ministry of Industry. ``They are
stronger and bigger than other countries. Why do we have to give them preference?''

To varying degrees, others are voicing the same complaint. Take the 10 Southeast Asian nations in
the Association of Southeast Asian Nations, known as Asean, a regional economic bloc representing
about 600 million people. After a decade of trade surpluses with Asean nations that ran as high as
\$20 billion, the surplus through October totaled a bare \$535 million, according to Chinese customs
figures, and appears headed toward a 10-year low. That is prompting some rethinking of the
conventional wisdom that China's rise is a windfall for the whole neighborhood.

Vietnam just devalued its currency by 5 percent, to keep it competitive with China. In Thailand,
manufacturers are grousing openly about their inability to match Chinese prices. India has filed a
sheaf of unfair-trade complaints against China this year covering everything from I-beams to coated
paper.

The Asia-Pacific Economic Cooperation forum, the biggest regional group, last month urged the
adoption of ``market-oriented exchange rates'' for Asian currencies without mentioning -- or needing
to mention -- China's currency, which many economists say China keeps artificially undervalued to
promote its own exports.

In Southeast Asia, Indonesia is having second thoughts about a free-trade pact China negotiated with
the six core Asean nations.

Under strong pressure from industries as varied as steel and motorbike makers, the Trade Ministry
said last week that it would seek to renegotiate some of the 350-odd tariff reductions that were
envisioned in the first year of the accord, set to take effect in January.

Jong-Wha Lee, the chief economist for the Asian Development Bank, noted that Japan and South Korea
were also seen as juggernauts -- and were criticized -- when their state-backed industries rapidly
increased exports. But the challenge from China seems different.

``Not just the size, but the speed of China's emerging power is really unprecedented in the
region,'' Mr.~Lee said. ``So it creates a lot of issues -- not just trade and exchange-rate
policies. But in the future, what will be the role of China?''

China has taken some steps to mollify complainers. In April, it proposed a \$10 billion investment
fund to help build badly needed roads, railways and ports in Southeast Asia, and a \$15 billion fund
to give Asian nations low-interest development loans.

But it has so far done little to address regional and global unease over the value of its currency,
the renminbi. Because the currency is lashed by effective government fiat to the sinking American
dollar, China's exports have become significantly cheaper in countries whose own currencies have not
compensated for the dollar's recent fall.

In Asia, the renminbi is doubly significant. During the 1997 Asian economic crisis, the values of
many regional currencies collapsed, making their goods cheap to foreign buyers. The Chinese then won
the gratitude of their neighbors -- and cast their country as a responsible power -- by keeping the
renminbi's value fixed. That prevented a competitive spiral of devaluations that many economists
feared might make the crisis much worse.

The latest financial crisis tells a different story: China's exchange rate controls are cited as a
leading cause of huge global imbalances that contributed to the collapse of 2008.

This time, China has resisted pressure to untie the renminbi from the dollar and let it rise. And
its neighbors' exports have suffered as a result.

Michael Pettis, an economist and scholar with the China program of the Carnegie Endowment for
International Peace, argues that China can no longer pursue the same export-driven development model
at a time when Western consumers no longer are able to gobble up whatever it and other Asian
manufacturers produce.

Until 2008, Mr.~Pettis said, ``most of these countries ran trade surpluses, and the U.S.~was the
countervailing trade deficit.''

``The entire model depended on the ability of an external agent -- the United States -- to absorb
trade deficits,'' he added.

Indonesia is especially vulnerable to the shift. It is the most populous and arguably the least
economically advanced nation among the onetime Asian Tigers, and perhaps the least able to
accommodate itself to a new regional order dominated by China.

Didik J.~Rachbini, a professor and the founder of an economic research institute here, said that in
the past four years, Indonesia had swung from more or less parity in bilateral trade to a deficit
equal to one-third of its annual exports to China -- and rising.

The lowly nail is one focus of tension. Making nails is not complicated: start with a bale of steel
wire, shave it down to the proper diameter, then feed it into a punch that shapes the nail, cuts it
and spits it into a bin. Labor and machinery account for 10 or 15 percent of the cost of a nail. The
rest is the cost of the wire.

And that is Indonesia's problem.

``Many Chinese steel factories have overcapacity, so they sell their wire very cheap,'' said Ario
N.~Setiantoro, who leads the Indonesia Nail and Wire Factory Association. ``Chinese nails enter the
market here at about the same price as our wire.''

He is right. Most analysts say China has too many steel mills. Its excess steelmaking capacity
equals the entire annual production of the world's No.~2 steelmaker, Japan. Every Chinese province
wants a steel industry, because it conveys prestige, creates jobs and attracts other business.

Beyond supply, Chinese state-run banks support industry with construction loans so cheap that credit
can be almost free, holding down operating costs. China's vast purchases of iron ore lock in volume
discounts that Indonesia's small steelmakers cannot match.

Export markets have dried up.

Like Dunia Metal, Surabaya Wire, a nail maker in east Java, has given up on exports altogether. ``I
used to have 450 workers,'' said, Sindu Prawira, the chief executive of Surabaya. ``Now, we have
170. Almost everybody is like that.''

Industries everywhere tend to accuse competitors of dirty tricks when they lose market share, of
course, and Indonesia's anemic steel industry shoulders its own share of blame for the nation's
competitive problems.

But as layoffs mount, the Indonesian government has been forced to try to shore up ailing producers.

In October, Indonesia's Trade Ministry invoked World Trade Organization rules and slapped a 145
percent safeguard tariff on Chinese nail imports, pending negotiations to settle complaints that the
Chinese are competing unfairly.

Irvan K.~Hakim, a co-chairman of the Indonesian Iron and Steel Industry Association, said he had
aired those sorts of complaints to Chinese officials for years. He did not appear optimistic about a
meeting of the minds.

``China is China, you know?'' he said, shrugging. ``Even the U.S.~cannot talk to China.''

\section{Church Works With U.S.~to Spare Detention}

\lettrine{W}{hen}\mycalendar{Dec.'09}{14} the young pastor started his ministry here at the
century-old Reformed Church in 2001, he gave little thought to the separate congregation of
Indonesian Christians who shared the sanctuary. They worshiped quietly in their own language on
Sunday afternoons, at the end of a hard week's work in the factories and warehouses of central New
Jersey.

But by May 2006, when they began pleading to sleep at the church, the pastor, the Rev. Seth
Kaper-Dale, had to pay attention. At the apartment complex where many Indonesians lived, armed
federal immigration agents in a single night had rounded up 35 men with expired visas and
outstanding deportation orders, as their wives and children cried and other families hid.

Suddenly a prosperous suburban congregation was confronted with the labyrinthine world of
immigration law and detention. This year, when one of its own leaders, an Indonesian, was detained
for months, only the pastor's passionate, last-ditch efforts saved him from deportation. And the
church reached a new level of activism -- with extraordinary results.

Under an unusual compact between the pastor and Immigration and Customs Enforcement officials in
Newark, four Indonesians have been released from detention in recent weeks, and 41 others living as
fugitives from deportation have turned themselves in under church auspices. Instead of being jailed
-- as hundreds of thousands of immigrants without criminal records have been in recent years -- they
have been released on orders of supervision, eligible for work permits while their lawyers consider
how their cases might be reopened.

Though agency officials say the arrangement is simply an example of the case-by-case discretion they
often use, the outcome has astonished advocates and experts in immigration enforcement, and raised
hopes that it signals some broader use of humanitarian release as the Obama administration vows to
overhaul the immigration system.

Still, for those who turn themselves in, the leap of faith carries big risks. For now, they can
check in at a federal office every three months and, if granted a work permit, can secure a driver's
license. But they are also vulnerable to immediate deportation. Just this fall, nine Indonesian
Christians in Seattle who had been on supervised release for years were abruptly detained, and some
were deported.

The immigration agency issues about 10,000 orders of supervision annually, but they typically
involve people who cannot be deported for practical reasons, like a homeland that will not take them
back. The agency detains roughly 380,000 people a year.

``I'm totally on uncharted waters,'' Mr.~Kaper-Dale, 34, a Vermont native who shares the pulpit with
his wife, Stephanie, said in October as he began seeking volunteers willing to place themselves in
the government's hands, from about 200 candidates not only at his church, but at several other New
Jersey congregations.

The first ones to step up had to overcome fear born of experience.

``Very, very scary,'' said Augus Alex Assa, 46, who fought tears as his 5-year-old daughter,
Christia Celine, clung to him in the van from the church, in Middlesex County, to an immigration
enforcement unit in Newark. ``In my heart, I hope I will stay in the United States.''

Like most of the Indonesians, Mr.~Assa and his wife, Grace, came on tourist visas that were suddenly
easy for poor people to get in the 1990s, when a booming economy welcomed foreign labor with a wink
and a nod. Everything changed after 9/11, when a government directive required the ``special
registration'' of men ages 16 to 65 who had entered the country on temporary visas from a list of
predominantly Muslim countries, including Indonesia. If they did not register, it was understood,
they would be considered terrorist fugitives.

Most of the Indonesian Christians complied, on the advice of pastors. They hoped that honesty would
open a path to legal status rather than deportation to their homeland, where many had faced
discrimination and sectarian violence.

Instead, their appeals for asylum were denied in most cases, some through inattention by inept or
overburdened lawyers. And those who registered became easy targets when national immigration
politics demanded a crackdown.

During the 2006 raid, Mr.~Assa hid in a closet when immigration agents came to the door, as his wife
covered their daughter's mouth. For two weeks afterward, they and others slept at the church.

About 50 men were eventually deported, typically after lengthy stays in immigration jails, leaving
wives struggling to support American-born children. ``We were shocked, but we were kind of
paralyzed,'' the pastor said.

On Jan.~12, the detention of one of their own spurred the congregation to action. Harry Pangemanan,
a popular Bible study leader, was picked up by immigration agents as he left for work as a warehouse
supervisor. He and his wife, Mariyana, parents of two American-born daughters, were the only
Indonesians among the 300 people in the main congregation.

Church members organized daily visits to the detention center, a 40-minute drive away in Elizabeth,
N.J., while the pastor appealed to Congressional and immigration offices. When Mr.~Pangemanan
reached out with his Bible to fellow detainees, the congregation visited them, too. Appalled to find
asylum-seekers behind barbed wire and plexiglass, they began holding vigils outside the center, run
for profit by the Corrections Corporation of America.

Some church members resisted. ``As a construction worker who is directly affected by immigration,
it's very hard,'' said Rich Lord, 39. ``I felt like, they're taking my jobs away.''

But his union and his faith changed his mind, he said: ``There's pregnant women so desperate in
Mexico that they're willing to cross the desert so their child will be born in the United States.
And as a Christian, I have to remember that Mary, the mother of Jesus, had to flee their homeland.''

Then, at 5 a.m. on March 31, came bad news: Mr.~Pangemanan was being put on a plane to Indonesia.
The pastor threw on his clerical collar and ran through Newark Liberty International Airport in a
frantic search for the right gate, determined to pray with his friend before he was sent away.

By the time the pastor found the flight, the passengers had already boarded. As he tells the story,
he prayed at the gate, so visibly upset that an airline worker let him on the plane.

Mr.~Pangemanan was in the last row between two immigration agents -- bound not for Jakarta but for a
detention center in Tacoma, Wash. -- when he saw his pastor coming down the aisle. An astonished
agent asked, ``How did this guy get in here?''

``And I just put my finger up,'' Mr.~Pangemanan recalled, pointing heavenward.

The agents let them pray briefly; the pastor said goodbye but vowed to keep trying. Back at the
church, he phoned every number on the immigration agency's Web site.

He still cherishes the recording of the only message that came back, from Dora B.~Schriro, who has
since left the agency but was then special detention adviser to Janet Napolitano, secretary of
homeland security. Within a week of their conversation, Mr.~Pangemanan was back in New Jersey with
his family, his case under reconsideration by the Board of Immigration Appeals.

When immigration agents arrested several more Indonesian men in late September, church leaders took
their effort to a new level, meeting with Scott Weber, director of the detention and removal field
office in New Jersey, and agency envoys from Washington.

David J.~Venturella, acting director of the agency's national detention and removal operations, said
he approved the discussions. ``We encourage all of our field office directors to exercise
prosecutorial discretion on a case-by-case basis,'' he said. ``This is a perfect example.''

Mr.~Weber rejected the ministers' proposal for a church-run alternative to detention, but offered
his own: In groups of 5 or 10, twice a week, the church could bring in the Indonesians they vouched
for, and lawyers committed to the lengthy process of seeking their full case files.

Unless something was amiss -- a hidden criminal conviction, a false address -- the former fugitives
could walk out the same day. Even before the details were arranged, Mr.~Weber released four recent
Indonesian detainees, one a Muslim.

Amy Gottlieb, immigrant rights director for the American Friends Service Committee in New Jersey,
who has been dealing with the field office since 1996, called it ``an amazing moment.''

``One, you just never believe that ICE is going to work with you on anything, given the history,''
she said. ``And given the intensive arrest efforts for the last two or three years, it's hard to
believe that people are ready to recognize that every single case has a human angle.''

Rex Chen, the supervising lawyer at Catholic Charities of the Archdiocese of Newark, remains more
pessimistic, likening himself to a financial adviser who warns, ``This mutual fund could collapse.''

While the arrangement may buy the Indonesians a year or two, he said, unless grounds are found to
reopen their cases, or Congress changes immigration law, they could find ``they just moved up from
not known, to on the list, to you're taking the steps up to the airplane.''

There are no guarantees, acknowledged Melinda Basaran, another participating lawyer and chairwoman
of the New Jersey chapter of the American Immigration Lawyers Association. But many of the
Indonesian wives, who did not have to register after 9/11, will soon have been here 10 years without
drawing official attention, making them eligible to apply for green cards.

The more pressing question is who is included in the supervised release, said Joan Pinnock, another
lawyer involved. Word of mouth has brought calls from Washington State, Pennsylvania and New
Hampshire, where many Indonesians fled after the New Jersey raid -- and where their detention and
deportation continues unabated. But Newark immigration authorities have ruled out their return to
New Jersey.

``I would love to get this for my Jamaican clients,'' Ms.~Pinnock said, echoing others who pointed
to different groups, like the many Muslims affected by special registration.

On a recent Wednesday night, in a church meeting room hung with the quilts of four generations of
grandmothers, fathers restored to their families thanked God and the congregation.

``I'm proud of my church,'' Mr.~Pangemanan said. ``Not just the pastor, the whole church.''

\section{China Indicts Prominent Dissident}

\lettrine{L}{iu}\mycalendar{Dec.'09}{14} Xiaobo, one of China's best-known dissidents and a
principal author of a pro-democracy manifesto that has attracted more than 10,000 signatures from
Chinese supporters, was indicted Thursday on charges of trying to subvert the state, his lawyer
said.

Mr.~Liu was expected to be tried in four to six weeks, the lawyer, Shang Baojun, said Friday.

The authorities disclosed the decision to prosecute Mr.~Liu -- a step that almost invariably ends in
imprisonment -- exactly one year and a day after the manifesto, Charter 08, was published. Other
Charter 08 signers said in interviews that the government was using Mr.~Liu's case to send a strong
message to Chinese intellectuals that it would not tolerate organized, independent efforts to foster
democracy.

``The government is trying to tell us to stop trying to push for human rights and democracy in
China,'' said Xu Youyu, a Charter 08 signer and a philosophy professor who recently retired from the
Chinese Academy of Social Sciences. ``Secondly, he has been the biggest threat inside of China, and
they want to get rid of him.''

Mr.~Liu's supporters had hoped that Chinese leaders would be persuaded to release Mr.~Liu, who has
been detained for more than a year, when President Obama visited China last month. During the visit,
United States officials gave Chinese leaders a list of ``cases of concern'' that included Mr.~Liu
and 11 other political and religious activists who are imprisoned or facing charges, according to
Nicholas Bequelin, an Asia researcher for Human Rights Watch, which is based in New York.

Mr.~Xu said: ``I think the message to the outside world is, it doesn't really matter to the
government how this case is viewed by the international community. It can do whatever it wants.''

Many activists viewed Charter 08 as the most important pro-democracy effort in China since the 1989
Tiananmen protests. Although censors swiftly deleted the document from Internet pages and chat
rooms, more than 10,000 people managed to sign it. It stated that Chinese citizens should be able to
elect their own government, that power should be divided among different branches of government and
that the military should come under government, not party, control.

Many of the initial signers were interrogated\cite{interrogate}, but only Mr.~Liu was arrested, two
days before the manifesto was published. He was confined for more than six months to a windowless
room in Beijing, then charged by the police with subversion and trying to overthrow the government.

Charter 08 signers have tried to show their solidarity by issuing letters of protest over his case.
The most recent one, signed by 165 people, states, ``If Liu Xiaobo is found 'guilty' that means each
one of us is guilty, and we have to shoulder the punishment together with Liu Xiaobo.''

Ran Yunfei, a well-known blogger in Sichuan Province who signed the manifesto and the most recent
letter, said the government wanted to silence the movement. ``But this is the historic trend of
democratization, and it cannot be stopped by anyone,'' he said in an interview.

Mr.~Liu spent two years in prison after the 1989 pro-democracy protests and three years in a labor
camp starting in 1996 for challenging single-party rule and advocating negotiations with the Dalai
Lama over Tibet. He was prepared to return to prison over Charter 08, Mr.~Bequelin said.

Chinese defendants in subversion cases are typically sentenced to three to eight years in prison, he
said. They are almost never acquitted.

Mr.~Bequelin said Mr.~Liu's fate depended partly on the reaction of the international community to
the indictment. ``Whatever they have been doing, it has not been enough,'' he said.

\section{Long-Term Care Is Latest Issue in Health Care Debate}

\lettrine{E}{mbedded}\mycalendar{Dec.'09}{14} in sweeping health legislation passed by the House and
being debated on the Senate floor is a major new federal insurance program for long-term care
intended to help people like Anne M.~Rader.

Ms.~Rader, 45, works at Booz Allen Hamilton as a consultant to federal agencies on emergency
preparedness. Even though she has cerebral palsy and multiple sclerosis, she leads a full, active
life. But she worries that she will lose her independence if her conditions grow worse.

``Having two disabilities, two disabling conditions, I can't predict what will happen in the
future,'' said Ms.~Rader, who lives alone in a condominium in Arlington, Va.

Advocates for older Americans and people with disabilities see the program as a long-overdue effort
to address needs that will explode as baby boomers age. It is meant for people with severe
disabilities who want to live in the community, though the benefits could also be used to help pay
for nursing home care or assisted living.

But critics say that the program is unsustainable and that it could ultimately create serious fiscal
problems for the government.

Senator Edward M.~Kennedy, Democrat of Massachusetts, drafted the proposal several years before he
died. Barack Obama, as a senator from Illinois, was a co-sponsor. Now, as president, Mr.~Obama wants
Congress to include it in the health care bill.

Senator Tom Harkin, Democrat of Iowa and chairman of the health committee, said: ``This is the next
logical step after the Americans With Disabilities Act. It will provide people with security and
peace of mind. They won't have to go to a nursing home or an institution if they become disabled
through an accident or an illness.''

The bill would provide cash benefits if a person had a substantial cognitive impairment or was
unable to perform two or three ``activities of daily living,'' like eating, bathing or dressing. The
program would be financed with premiums paid by participants, through voluntary payroll deductions,
with no federal subsidy. People could qualify for lifetime benefits if they became disabled after
paying premiums for at least five years and working for three of those years.

The Congressional Budget Office assumes that premiums would be \$123 a month for benefits expected
to average \$75 a day, or about \$27,000 a year. The amount of benefits would vary, depending on the
degree of a person's disability. The secretary of health and human services could increase premiums
to ensure ``the financial solvency'' of the program over 75 years.

The Senate bill says, ``No taxpayer funds shall be used for payment of benefits.''

Senator Christopher J.~Dodd, Democrat of Connecticut, said the benefits would allow people with
disabilities to ``live out their lives with decency and dignity.''

``What's the alternative?'' Mr.~Dodd asked. ``Getting rid of all your assets, impoverishing
yourself, relying on your family or friends to take care of you in order to try to survive.''

But Republicans and some fiscally conservative Democrats said they feared that the government would
eventually have to bail out the program because it would prove unsustainable.

``It would create a huge new liability down the road,'' said Senator John Thune, Republican of South
Dakota.

Paradoxically, the proposed new program accounts for more than half of the bills' deficit reduction
in the first 10 years -- because the government would pay out far less in benefits than it would
collect in premiums. But costs would grow later.

The Congressional Budget Office estimates that the long-term-care insurance program in the Senate
bill would reduce federal budget deficits by \$72 billion from 2010 to 2019. For the House bill, the
comparable figure is \$102 billion.

On Dec.~4, the Senate voted 51 to 47 to strip the program from the bill. The Senate had previously
agreed to set a 60-vote threshold, so the effort failed, and the program remains in the bill. But it
could become a bargaining chip in negotiations over the measure.

The effort to eliminate the program won support from 11 Democrats, including the chairman of the
Finance Committee, Max Baucus of Montana, and the chairman of the Budget Committee, Kent Conrad of
North Dakota.

Six of those Democrats said the program ``would not be fiscally responsible.'' The bill would
``create a new federal entitlement program with large, long-term spending increases that far exceed
revenues,'' they said in a letter to the Senate majority leader, Harry Reid, Democrat of Nevada.

Richard S.~Foster, chief actuary at the federal Centers for Medicare and Medicaid Services, has
expressed a similar concern.

``There is a very serious risk that the program would become unsustainable,'' Mr.~Foster said,
because people who have or anticipate health problems would be more likely to sign up than people in
better-than-average health.

Some companies that sell long-term-care insurance are lobbying against the proposal, known as the
Community Living Assistance Services and Supports Act, or Class Act.

One big insurer, Genworth Financial, said the program ``would give individuals a false sense of
security'' and could discourage them from buying private insurance to cover the costs of long-term
care.

``A government-run program that covers only a small fraction of Americans' total long-term-care
needs will mislead the general public and make it even harder for agents and advisers to encourage
their clients to plan for this important retirement protection,'' Genworth says on a Web site for
its employees and agents.

Supporters of the program say it will help not only people with disabilities, but also those who
care for them as well.

Carolyn A.~Martin, 85, has a touch of dementia, kidney problems and severe arthritis. She cannot
prepare meals, wash her clothes or bathe herself, and she often has trouble getting out of bed.

She lives in Columbia, Md., with her daughter, Alma M.~Gill, who cares for her while holding a
full-time job at a nonprofit organization.

``If I had someone to care for my mother four hours a day, it would change her life, and mine,'' Ms.
Gill said. ``I feel guilty about leaving her alone when I go to work. If something were to happen to
her, it would be my fault.''

\section{Lieberman Rules Out Voting for Health Bill}

\lettrine{I}{n}\mycalendar{Dec.'09}{14} a surprise setback for Democratic leaders, Senator Joseph
I.~Lieberman, independent of Connecticut, said on Sunday that he would vote against the health care
legislation in its current form.

The bill's supporters had said earlier that they thought they had secured Mr.~Lieberman's agreement
to go along with a compromise they worked out to overcome an impasse within the Democratic Party.

But on Sunday, Mr.~Lieberman told the Senate majority leader, Harry Reid, to scrap the idea of
expanding Medicare and abandon any new government insurance plan or lose his vote.

On a separate issue, Mr.~Reid tried over the weekend to concoct a compromise on abortion that would
induce Senator Ben Nelson, Democrat of Nebraska, to vote for the bill. Mr.~Nelson opposes abortion.
Any provision that satisfies him risks alienating supporters of abortion rights.

In interviews on the CBS News program ``Face the Nation,'' Mr.~Lieberman and Mr.~Nelson said the
bill did not have the 60 votes it would need in the Senate.

Senate Democratic leaders, including Mr.~Reid and Senator Charles E.~Schumer of New York, said they
had been mindful of Mr.~Lieberman's concerns in the last 10 days and were surprised when he assailed
major provisions of the bill on television Sunday. He reiterated his objections in a private meeting
with Mr.~Reid.

A Senate Democratic aide, perplexed by Mr.~Lieberman's stance, said, ``It was a total flip-flop, and
leaves us in a predicament as to what to do.''

Democrats are desperately trying to round up 60 votes and conclude Senate debate on the health care
bill before Christmas.

Mr.~Reid could not immediately figure out how to achieve that goal at a meeting he held Sunday with
senior Democratic senators and White House officials, including Mr.~Obama's chief of staff, Rahm
Emanuel, according to Senate Democratic aides.

Marshall H.~Wittmann, a spokesman for Mr.~Lieberman, said the Connecticut senator ``notified Senator
Reid on Friday that he had severe misgivings about the Medicare buy-in proposal, so his comments on
'Face the Nation' should not have come as a surprise to the leadership.''

The Senate Republican leader, Mitch McConnell of Kentucky, said that passage of the bill was looking
less and less inevitable. The Democrats ``are in serious trouble on this,'' he said, ``and the core
problem is the American people do not want us to pass it.''

On television Sunday, Mr.~Lieberman said: ``We've got to stop adding to the bill. We've got to start
subtracting some controversial things. I think the only way to get this done before Christmas is to
bring in some Republicans who are open-minded on this, like Olympia Snowe.''

Senator Snowe, of Maine, has tried to find common ground with Democrats, but has rejected Mr.~Reid's
proposal to let uninsured people 55 to 64 years old purchase coverage under Medicare.

``You've got to take out the Medicare buy-in,'' Mr.~Lieberman said. ``You've got to forget about the
public option. You probably have to take out the Class Act, which was a whole new entitlement
program that will, in future years, put us further into deficit.''

Class Act refers to a federal insurance program for long-term care, known as the Community Living
Assistance Services and Supports Act.

Mr.~Lieberman said he would have ``a hard time'' voting for a bill with the Medicare buy-in.

``It has some of the same infirmities that the public option did,'' he said. ``It will add taxpayer
costs. It will add to the deficit. It's unnecessary. The basic bill, which has a lot of good things
in it, provides a generous new system of subsidies for people between ages 55 and 65, and choice and
competition.''

Mr.~Nelson said he wanted to know the cost of the Medicare buy-in. ``I am concerned that it's the
forerunner of single payer, the ultimate single-payer plan, maybe even more directly than the public
option,'' he said.

Mr.~Lieberman said: ``The bill itself does a lot to bring 30 million people into the system. We
don't need to keep adding onto the back of this horse, or we're going to break the horse's back and
get nothing done.''

Even if Senate Democratic leaders were prepared to meet Mr.~Lieberman's demands, they would still
need to resolve intraparty disputes over insurance coverage for abortion.

Aides to Mr.~Reid met Saturday with advocates of abortion rights to explore ideas for a compromise.

Details were sketchy. Under one idea, some health plans receiving federal subsidies could offer
optional coverage for abortion, but they could not use federal money to pay for the procedure. They
would have to use money taken from premiums paid by subscribers and would have to keep it separate
from federal money.

Critics of abortion say such requirements for the segregation of funds are an accounting gimmick.

In hopes of placating opponents of abortion, Mr.~Reid is also considering an increase in the federal
tax credit for adoption of children and a new program to provide services to pregnant high school
and college students.

Both ideas were proposed by Senator Bob Casey, Democrat of Pennsylvania, who opposes abortion but
generally supports the overall bill.

``Many teens and women who face an unplanned pregnancy do so with little or no support,'' Mr.~Casey
said.

\section{`Smart' Electric Utility Meters, Intended to Create Savings, Instead Prompt Revolt}

\lettrine{M}{illions}\mycalendar{Dec.'09}{14} of households across America are taking a first step
into the world of the ``smart grid,'' as their power companies install meters that can tell them how
much electricity they are using hour by hour -- and sometimes, appliance by appliance. But not
everyone is happy about it.

Customers in California are in open revolt, and officials in Connecticut and Texas are questioning
whether the rush to install meters benefits the public.

Some consumers argue that the meters are logging far more kilowatt hours than they believe they are
using. And many find it unfair that they will begin to pay immediately for the new meters through
higher rates, when the promised savings could be years away.

Power companies say the meters will allow utilities to vary the price charged to their customers by
the hour to correspond to what those utilities are paying for energy in the wholesale market. This
can help consumers save money, they say.

They also say the meters will be crucial to remaking the electric system to handle intermittent
power sources like wind turbines and solar cells while continuously meeting customers' needs.

Pacific Gas \& Electric, which distributes power to Northern and Central California, has so far
installed four million meters in households and businesses and plans to install six million more
within the next three years. The meters cost the utility roughly \$220 apiece, including
installation.

Elizabeth Keogh, a retired social worker in Bakersfield, Calif., who describes herself as ``a bit
chintzy,'' has created a spreadsheet with 26 years of electric bills for her modest house. She
decided that her new meter was running too fast.

Ms.~Keogh reported to the utility that the meter recorded 646 kilowatt-hours in July, for which she
paid \$66.50; last year it was 474 kilowatt-hours, or \$43.37.

At a hearing in October organized by her state senator, Ms.~Keogh took out two rolls of toilet paper
-- one new, one half used up -- and rolled them down the aisle, showing how one turned faster than
the other. ``Something is wrong here,'' she said.

Scores of electric customers with similar complaints have turned out at similar hearings. At one in
Fresno, Calif., Leo Margosian, a retired investigator, testified that the new meter logged the
consumption of his two-bedroom townhouse at 791 kilowatt-hours in July, up from 236 a year earlier.
And he had recently insulated his attic and installed new windows, Mr.~Margosian said.

At the urging of the state senator, Dean Florez, Democrat of Fresno and the chamber's majority
leader, and others, the California Public Utilities Commission is moving to bring in an outside
auditor to determine whether the meters count usage properly.

In response to a wave of complaints from the Bakersfield area in the Central Valley, Pacific Gas \&
Electric has been placing full-page advertisements in newspapers in the area promising benefits from
the new meters. It says customers will save money not only by paying rates based on hourly
fluctuations in the wholesale market, but also eventually by displaying real-time rates.

To reduce their bills, customers could cut back at pricey peak times and shift some activities, like
running a clothes dryer or a vacuum cleaner, to off-peak periods. Utilities will then have lower
costs, the argument goes, because the grid will need fewer power plants as demand levels out.

Customers will become ``structural winners,'' said Andy Tang, senior director of the company's Smart
Energy Web program.

Someday utilities hope to use the meter to control consumption by major appliances like air
conditioners. But experts are still debating what technical standards the meters and appliances
should use to communicate.

While the costs of installing the meters is substantial -- \$2.2 billion in the case of PG\&E -- the
utility reaps some immediate benefits that eventually will be passed along to the consumer.

The most obvious one is that utilities can eliminate their meter readers, along with an expenditure
of 50 cents to a dollar to read each meter every month. And with smart meters, utilities are alerted
immediately if a customer's power is out.

If a utility decides to shut off a customer for nonpayment, it can do so by remote control; if the
customer pays enough money to allow resumption of service, the utility can also do that from a
central office without sending out a representative.

PG\&E attributes the higher bills that some consumers complain about to recent rate increases and to
quirks in California's pricing system. Electricity in the state is priced in so-called tiers:
consumers get the first few hundred kilowatt-hours at a low rate, but the next few units of
consumption are billed at a high rate. A small increase in use can therefore result in a big
increase in the bill, the utility says. It adds that an unusually high number of very hot days were
recorded last summer.

But people in other parts of the country are also wary of the meter switch. Attorney General Richard
Blumenthal of Connecticut frets that consumers could be shouldering the costs of the transition long
before they realized the savings.

``The major benefits come in the second and third stages,'' Mr.~Blumenthal said, referring to
instantaneous rate information, the ability to adjust use and the prospect of developing appliances
that can be set to switch on or off when the meter announces that prices have reached a certain
level.

So Mr.~Blumenthal has helped persuade regulators in his state and Connecticut Light \& Power to
scale back a plan for widespread installation of smart meters, and to run a pilot program first.

In Texas, where state law encourages installation of smart meters, the public utility counsel, a
consumer advocate, got the utilities to agree to pay tens of millions of dollars for public
education and to subsidize the cost of an in-home display to give an instantaneous price in
low-income households.

(While meters in Texas will bill customers based on time of day, the utilities do not routinely
provide the indoor hardware that furnishes such information.)

Complicating the transition, the widespread introduction of smart meters comes amid a recession and
a decline in electricity demand.

Two years ago, experts agree, it was cheaper to persuade 100,000 customers to each use two kilowatts
less energy at any given moment than for a utility to build and run a 200-megawatt power plant.

But today, reining in energy consumption is less of a corporate priority: generating capacity is in
surplus in almost all parts of the United States because the recession has shuttered so many
factories. And in swaths of the eastern United States, the wholesale price difference between peak
and off-peak demand is far smaller lately.

The long-term impact of the smart meters is uncertain. Some studies show that people use less
electricity when they can see the numbers ticking higher on the meter.

Among residential customers who volunteered for a test program in California last summer, 70 percent
saved money and 97 percent said they were satisfied with the program and wanted to stay in it,
Mr.~Tang of PG\&E said.

Yet of about a million customers with smart meters at the time, only 24,000, or 2.4 percent, chose
to take part in the test program.

\section{McCain, Distanced From Race, Raises Senate Voice}

\lettrine{M}{uch}\mycalendar{Dec.'09}{14} of the old presidential campaign gang has moved on. The
governor he made famous -- Sarah Palin, his vice-presidential pick in 2008 -- is the one exciting
crowds these days. He is facing the possibility of a primary challenge at home, one more reminder of
his uneasy relationship with his own party.

Yet at the age of 73, one year after his defeat by President Obama, Senator John McCain of Arizona
is trying to make the most of the platform where he has always been most comfortable, the United
States Senate.

The Republican Party's leadership vacuum has given Mr.~McCain an opening, and he is charging through
it, tacking right on some issues and loudly embroiling himself in battles with the White House and
Democratic leaders over health care, stimulus spending, foreign policy and the style of the Obama
presidency.

He is more visible now than at any time since the end of his presidential campaign.

``Let's do what the president said last October a year ago,'' Mr.~McCain said the other day at one
of what has become a geyser of appearances on the Senate floor, in Capitol hallways and at news
conferences. ``Let's all sit down together, Republicans and Democrats, with C-Span in the room, and
negotiate so that the American people can see what's going on here.''

Mr.~McCain has cut back his dealings with many of the people who were at his side while he was
running for president. Steve Schmidt, his campaign manager, has returned to California, while Mike
Murphy, a longtime adviser, has not talked to him since last summer, according to associates of both
men. Mark Salter, his closest aide, alter ego and book collaborator, has left Mr.~McCain's Senate
staff and gone into private business, and now speaks to him about once a week.

Mr.~McCain has even been a bit at odds with his closest friends on the campaign trail, Senator
Joseph I.~Lieberman, independent of Connecticut, and Senator Lindsey Graham, Republican of South
Carolina, attacking a global warming bill that they proposed along with Senator John Kerry, Democrat
of Massachusetts, though Mr.~Graham said he and Mr.~McCain remained as close as ever.

``His presidential aspirations are over -- he knows he's never going to be president,'' Mr.~Graham
said. ``Most people in that position have a hard time re-engaging, but he's really engaged. I've
never seen him like this before. He's really going down there, he's really making it tough for
them.''

Mr.~McCain, who did not respond to a request for an interview, told advisers the day after he lost
to Barack Obama that he intended to seek a fifth term in the Senate. Yet he is facing a threatened
primary challenge from J.D. Hayworth, a conservative talk show host and former Republican member of
Congress, who said he was encouraged by polls that showed Mr.~McCain was vulnerable.

``The question that people are asking is this,'' Mr.~Hayworth said in an interview. ``Do we want to
send John McCain back to the United States Senate again, or is it time to change to a clear,
consistent, common-sense Republican?''

Don Bivens, the chairman of the Arizona Democratic Party, said he believed that Mr.~McCain would be
vulnerable in a general election, though as of now, the Democrats' best candidate is a little-known
city councilman from Tucson.

``I think he's taken a hard turn to the right,'' Mr.~Bivens said of Mr.~McCain. ``There's no
maverick left. There's no bipartisan guy left. He's just sort of a Republican attack dog, and as a
result I think he's made himself vulnerable here because we're pretty much a middle-of-the-road kind
of state.''

Mr.~McCain's friends said that in raising his profile, he was motivated not by concern at home, but
by philosophical differences over the scope of Mr.~Obama's health care proposals and spending
measures.

``Had they reached out to him in a more genuine way, and not tried to pursue a pretty leftist
agenda, I think they might have had a potential ally in John on certain things,'' said Senator Jon
Kyl, Mr.~McCain's fellow Republican from Arizona.

Mr.~Salter said: ``I think he's doing what anybody who knew him knew he would do. He deals with
loss, he deals with adversity, he deals with everything by being active when it's over. Whether he
won or lost. He wasn't going to go back and stare out the window or anything.''

But Mr.~McCain has also distanced himself from some long-held positions: he once backed measures to
deal with climate change, which made his criticism of Mr.~Graham's bill so striking. He denounced
efforts to curb Medicare costs by Congressional Democrats a year after he said such cuts were
critical.

Mr.~McCain pointedly praised Ms.~Palin after she recently released a book attacking his former
campaign aides, putting him in the position of embracing someone who has become the symbol of the
wing of the party he has long warned against.

And these days he is regularly at the side the Senate minority leader, Mitch McConnell, a
conservative Republican from Kentucky who almost did not endorse Mr.~McCain's presidential campaign
because of Mr.~McCain's history of breaking with conservatives on issues like immigration and tax
cuts.

``I'll leave others to interpret what some might call a campaign conversion,'' Mr.~Hayworth said.

If Mr.~McCain has had a history of being a happy warrior, that is not the phrase used by many of his
friends to describe his demeanor these days. There are few glimpses of the winks, wry smiles and
one-liners that were once an integral part of his character. More typically, his remarks are tinged
with sarcasm or anger, delivered with a wave of the arm or both hands chopping through the air, like
those he delivered Friday, when he declared that a report on Medicare benefits ``should put a dagger
in the heart of the Reid bill.''

Some of his associates noted similarities between Mr.~McCain's current demeanor and the way he
handled himself in 2001, when he returned to the Senate after losing the 2000 Republican
presidential nomination to George W.~Bush. He voted against Mr.~Bush's first round of tax cuts and
was one of his toughest critics, though Mr.~McCain grew closer to Mr.~Bush as he began preparing for
his own second race for the presidency.

``I think there's always going to be some maverick to McCain which makes him unpredictable and hard
to pin down,'' said Mark McKinnon, a longtime adviser. ''Which is what makes him so interesting. The
book on McCain is far from over.''

\section{Viruses That Leave Victims Red in the Facebook}

\lettrine{I}{t}\mycalendar{Dec.'09}{14} used to be that computer viruses attacked only your hard
drive. Now they attack your dignity.

Malicious programs are rampaging through Web sites like Facebook and Twitter, spreading themselves
by taking over people's accounts and sending out messages to all of their friends and followers. The
result is that people are inadvertently telling their co-workers and loved ones how to raise their
I.Q.'s or make money instantly, or urging them to watch an awesome new video in which they star.

``I wonder what people are thinking of me right now?'' said Matt Marquess, an employee at a public
relations firm in San Francisco whose Twitter account was recently hijacked, showering his followers
with messages that appeared to offer a \$500 gift card to Victoria's Secret.

Mr.~Marquess was clueless about the offers until a professional acquaintance asked him about them
via e-mail. Confused, he logged in to his account and noticed he had been promoting lingerie for
five days.

``No one had said anything to me,'' he said. ``I thought, how long have I been Twittering about
underwear?''

The humiliation sown by these attacks is just collateral damage. In most cases, the perpetrators are
hoping to profit from the referral fees they get for directing people to sketchy e-commerce sites.

In other words, even the crooks are on social networks now -- because millions of tightly connected
potential victims are just waiting for them there.

Often the victims lose control of their accounts after clicking on a link ``sent'' by a friend. In
other cases, the bad guys apparently scan for accounts with easily guessable passwords.
(Mr.~Marquess gamely concedes that his password at the time was ``abc123.'')

After discovering their accounts have been seized, victims typically renounce the unauthorized
messages publicly, apologizing for inadvertently bombarding their friends. These messages -- one
might call them Tweets of shame -- convey a distinct mix of guilt, regret and embarrassment.

``I have been hacked; taking evasive maneuvers. Much apology, my friends,'' wrote Rocky Barbanica, a
producer for Rackspace Hosting, an Internet storage firm, in one such note.

Mr.~Barbanica sent that out last month after realizing he had sent messages to 250 Twitter followers
with a link and the sentence, ``Are you in this picture?'' If they clicked, their Twitter accounts
were similarly commandeered.

``I took it personally, which I shouldn't have, but that's the natural feeling. It's insulting,'' he
said.

Earlier malicious programs could also cause a similar measure of embarrassment if they spread
themselves through a person's e-mail address book.

But those messages, traveling from computer to computer, were more likely to be stopped by antivirus
or firewall software. On the Web, such measures offer little protection. (Although they are
popularly referred to as viruses or worms, the new forms of Web-based malicious programs do not
technically fall into those categories, as they are not self-contained programs.)

Getting tangled up in a virus on a social network is also more painfully, and instantaneously,
public. ``Once it's delivered to everyone in three seconds, the cat is out of the bag,'' said Chet
Wisniewski of Sophos, a Web security firm. ``When people got viruses on their computers, or fell for
scams at home, they were generally the only ones that knew about it and they cleaned it up
themselves. It wasn't broadcast to the whole world.''

Social networks have become prime targets of such programs' creators for good reason, security
experts say. People implicitly trust the messages they receive from friends, and are inclined to
overlook the fact that, say, their cousin from Ohio is extremely unlikely to have caught them on a
hidden webcam.

Sophos says that 21 percent of Web users report that they have been a target of malicious programs
on social networks. Kaspersky Labs, a Russian security firm, says that on some days, one in 500
links on Twitter point to bad sites that can infect an inadequately protected computer with typical
viruses that jam hard drives. Kaspersky says many more links are purely spam, frequently leading to
dating sites that pay referral fees for traffic.

A worm that spread around Facebook recently featured a photo of a sparsely dressed woman and offered
a link to ``see more.'' Adi Av, a computer developer in Ashkelon, Israel, encountered the image on
the Facebook page of a friend he considered to be a reliable source of amusing Internet content.

A couple of clicks later, the image was posted on Mr.~Av's Facebook profile and sent to the ``news
feed'' of his 350 friends.

``It's an honest mistake,'' he said. ``The main embarrassment was from the possibility of other
people getting into the same trouble from my profile page.''

Others confess to experiencing a more serious discomfiture.

``You feel like a total idiot,'' said Jodi Chapman, who last month unwisely clicked on a Twitter
message from a fellow vegan, suggesting that she take an online intelligence test.

Ms.~Chapman, who sells environmentally friendly gifts with her husband, uses her Twitter account to
communicate with thousands of her company's customers. The hijacking ``filled me with a sense of
panic,'' she said. ``I was so worried that I had somehow tainted our company name by asking people
to check their I.Q. scores.''

Social networking attacks do not spare the experts. Two weeks ago, Lee Rainey, director of the Pew
Internet and American Life Project, a nonprofit research group, accidentally sent messages to dozens
of his Twitter followers with a link and the line, ``Hi, is this you? LOL.'' He said a few people
actually clicked.

``I'm worried that people will think I communicate this way,'' Mr.~Rainey said. `` 'LOL,' as my
children would tell you, is not the style that I want to engage the world with.''

\section{Paul A.~Samuelson, Economist, Dies at 94}

\lettrine{P}{aul}\mycalendar{Dec.'09}{14} A.~Samuelson, the first American Nobel laureate in
economics and the foremost academic economist of the 20th century, died Sunday at his home in
Belmont, Mass. He was 94.

His death was announced by the Massachusetts Institute of Technology, which Mr.~Samuelson helped
build into one of the world's great centers of graduate education in economics.

In receiving the Nobel Prize in 1970, Mr.~Samuelson was credited with transforming his discipline
from one that ruminates about economic issues to one that solves problems, answering questions about
cause and effect with mathematical rigor and clarity.

When economists ``sit down with a piece of paper to calculate or analyze something, you would have
to say that no one was more important in providing the tools they use and the ideas that they employ
than Paul Samuelson,'' said Robert M.~Solow, a fellow Nobel laureate and colleague of
Mr.~Samuelson's at M.I.T.

Mr.~Samuelson attracted a brilliant roster of economists to teach or study at the university, among
them Mr.~Solow as well as others who would go on to become Nobel laureates like George A.~Akerlof,
Robert F.~Engle III, Lawrence R.~Klein, Paul Krugman, Franco Modigliani, Robert C.~Merton and Joseph
E.~Stiglitz.

Mr.~Samuelson wrote one of the most widely used college textbooks in the history of American
education. The book, ``Economics,'' first published in 1948, was the nation's best-selling textbook
for nearly 30 years. Translated into 20 languages, it was selling 50,000 copies a year a half
century after it first appeared.

``I don't care who writes a nation's laws -- or crafts its advanced treatises -- if I can write its
economics textbooks,'' Mr.~Samuelson said.

His textbook taught college students how to think about economics. His technical work -- especially
his discipline-shattering Ph.D. thesis, immodestly titled ``The Foundations of Economic Analysis''
-- taught professional economists how to ply their trade. Between the two books, Mr.~Samuelson
redefined modern economics.

The textbook introduced generations of students to the revolutionary ideas of John Maynard Keynes,
the British economist who in the 1930s developed the theory that modern market economies could
become trapped in depression and would then need a strong push from government spending or tax cuts,
in addition to lenient monetary policy, to restore them. Many economics students would never again
rest comfortably with the 19th-century view that private markets would cure unemployment without
need of government intervention.

That lesson was reinforced in 2008, when the international economy slipped into the steepest
downturn since the Great Depression, when Keynesian economics was born. When the Depression began,
governments stood pat or made matters worse by trying to balance fiscal budgets and erecting trade
barriers. But 80 years later, having absorbed the Keynesian teaching of Mr.~Samuelson and his
followers, most industrialized countries took corrective action, raising government spending,
cutting taxes, keeping exports and imports flowing and driving short-term interest rates to near
zero.

\emph{Lessons for Kennedy}

Mr.~Samuelson explained Keynesian economics to American presidents, world leaders, members of
Congress and the Federal Reserve Board, not to mention other economists. He was a consultant to the
United States Treasury, the Bureau of the Budget and the President's Council of Economic Advisers.

His most influential student was John F.~Kennedy, whose first 40-minute class with Mr.~Samuelson,
after the 1960 election, was conducted on a rock by the beach at the family compound at Hyannis
Port, Mass. Before class, there was lunch with politicians and Cambridge intellectuals aboard a
yacht offshore. ``I had expected a scrumptious meal,'' Mr.~Samuelson said. ``We had franks and
beans.''

As a member of the Kennedy campaign brain trust, Mr.~Samuelson headed an economic task force for the
candidate and held several private sessions on economics with him. Many would have a bearing on
decisions made during the Kennedy administration.

Though Mr.~Samuelson was President Kennedy's first choice to become chairman of the Council of
Economic Advisers, he refused, on principle, to take any government office because, he said, he did
not want to put himself in a position in which he could not say and write what he believed.

After the 1960 election, he told the young president-elect that the nation was heading into a
recession and that Kennedy should push through a tax cut to head it off. Kennedy was shocked.

``I've just campaigned on a platform of fiscal responsibility and balanced budgets and here you are
telling me that the first thing I should do in office is to cut taxes?'' Mr.~Samuelson recalled,
quoting the president.

Kennedy eventually accepted the professor's advice and signaled his willingness to cut taxes, but he
was assassinated before he could take action. His successor, Lyndon B.~Johnson, carried out the
plan, however, and the economy bounced back.

\emph{Adding Bite to Academia}

In the classroom, Mr.~Samuelson was a lively, funny, articulate teacher. On theories that he and
others had developed to show links between the performance of the stock market and the general
economy, he famously said: ``It is indeed true that the stock market can forecast the business
cycle. The stock market has called nine of the last five recessions.''

His speeches and his voluminous writing had a lucidity and bite not usually found in academic
technicians. He tried to give his economic pronouncements a ``snap at the end,'' he said, ``like
Mark Twain.'' When women began complaining about career and salary inequities, for example, he said
in their defense, ``Women are men without money.''

Remarkably versatile, Mr.~Samuelson reshaped academic thinking about nearly every economic subject,
from what Marx could have meant by a labor theory of value to whether stock prices fluctuate
randomly. Mathematics had already been employed by social scientists, but Mr.~Samuelson brought the
discipline into the mainstream of economic thinking, showing how to derive strong theoretical
predictions from simple mathematical assumptions.

His early work, for example, presented a unified mathematical structure for predicting how
businesses and households alike would respond to changes in economic forces, how changes in wage
rates would affect employment, and how tax rate changes would affect tax collections.

His relentless application of mathematical analysis gave rise to an astonishing number of
groundbreaking theorems, resolving debates that had raged among theorists for decades, if not
centuries.

\emph{An Economic Theorem}

Early in his career, Mr.~Samuelson developed the rudimentary mathematics of business cycles with a
model, called the multiplier-accelerator, that captured the inherent tendency of market economies to
fluctuate.

The model showed how markets magnify the impact of outside shocks and turn, say, an initial one
dollar increase in foreign investment into a several dollar increase in total domestic income, to be
followed by a decline.

Mr.~Samuelson provided a mathematical structure to study the impact of trade on different groups of
consumers and workers. In a famous theorem, known as Stolper-Samuelson, he and a co-author showed
that competition from imports of clothes and similar goods from underdeveloped countries, where
producers rely on unskilled workers, could drive down the wages of low-paid workers in
industrialized countries.

The theorem provided the intellectual scaffold for opponents of free trade. And late in his career,
Mr.~Samuelson set off an intellectual commotion by pointing out that the economy of a country like
the United States could be hurt if productivity rose among the economies with which it traded.

Yet Mr.~Samuelson, like most academic economists, remained an advocate of open trade. Trade, he
taught, raises average living standards enough to allow the workers and consumers who benefit to
compensate those who suffer, and still have some extra income left over. Protectionism would not
help, but higher productivity would.

Mr.~Samuelson also formulated a theory of public goods -- that is, goods that can be provided
effectively only through collective, or government, action. National defense is one such public
good. It is nonexclusive; the Navy, for example, exists to protect every citizen. It also eliminates
rivalry among its many consumers; that is, the amount of security that any one citizen derives from
the Navy subtracts nothing from the amount of security that any other citizen derives.

The features of public goods, Mr.~Samuelson taught, stand in direct contrast to those of ordinary
goods, like apples. An apple eaten by one consumer is not available to any other. Public goods, he
concluded, cannot be sold in private markets because individuals have no incentive to pay for them
voluntarily. Instead they hope to get a free ride from the decisions of others to make the public
goods available.

\emph{A Predictive Principle}

Mr.~Samuelson pushed mathematical analysis to new levels of sophistication. For example, economists
routinely write mathematical models of market economies that assume consumers and producers make
choices to maximize their well-being. The question arises when such economies are stable: if
disturbed by, say, droughts or wars or technological change, will the economy return to appropriate
levels of prices and output or, instead, fly out of control? What Mr.~Samuelson's ``correspondence
principle'' shows is the theoretical link between the behavior of individuals and the aggregate
stability of the entire economic system. Information about individual responses, Mr.~Samuelson's
theorem holds, shapes predictions about overall economic stability.

He analyzed the evolution of economies with a mathematical model, called an overlapping generations
model, that scholars have since used to study, for example, the functioning over time of the Social
Security system and the management of public debt.

He also helped develop linear programming, a mathematical tool used by corporations and central
planners in socialist countries to calculate how to produce pre-set levels of various goods and
services at the least cost.

Late in his career, Mr.~Samuelson laid out the mathematics of stock price movements, an analysis
that became the basis for Nobel Prize-winning research by his student Mr.~Merton and Myron
S.~Scholes. They designed formulas that Wall Street analysts use to trade options and other
complicated securities known as derivatives.

But beyond his astonishing array of scientific theorems and conclusions, Mr.~Samuelson wedded
Keynesian thought to conventional economics. He developed what he called the Neoclassical Synthesis.
The neoclassical economists in the late 19th century showed how forces of supply and demand generate
equilibrium in the market for apples, shoes and all other consumer goods and services. The standard
analysis had held that market economies, left to their own devices, gravitated naturally toward full
employment.

Economists clung to this theory even in the wake of the Depression of the 1930s. But the need to
explain the market collapse, as well as unemployment rates that soared to 25 percent, gave rise to a
contrary strain of thought associated with Keynes.

Mr.~Samuelson's resulting ``synthesis'' amounted to the notion that economists could use the
neoclassical apparatus to analyze economies operating near full employment, but switch over to
Keynesian analysis when the economy turned sour.

\emph{Midwestern Roots}

Paul Anthony Samuelson was born on May 15, 1915, in Gary, Ind., the son of Frank Samuelson, a
pharmacist, and the former Ella Lipton. His family, he said, was ``made up of upwardly mobile Jewish
immigrants from Poland who had prospered considerably in World War I, because Gary was a brand new
steel town when my family went there.''

But after his father lost much of his money in the years after the war, the family moved to Chicago.
Young Paul attended Hyde Park High School, where as a freshman he began studying the stock market.
At one point he helped his algebra teacher select stocks to buy in the boom of the 1920s.

``Hupp Motors and other losers,'' he remembered in an interview in 1996. ``Proof of the fallibility
of systems,'' he said.

He left high school at age 16 to enter the University of Chicago. ``I was born as an economist on
Jan.~2, 1932,'' he said. That was the day he heard his first college lecture, on Thomas Malthus, the
18th-century British economist who studied the relation between poverty and population growth.
Hooked, he began taking economics courses.

The University of Chicago developed the century's leading conservative economic theorists, under the
later guidance of Milton Friedman. But Mr.~Samuelson regarded the teaching at Chicago as
``schizophrenic.'' This was at the height of the Depression, and courses about the business cycle
naturally talked about unemployment, he said. But in economic-theory classes, joblessness was not
mentioned.

``The niceties of existence were not a matter of concern,'' he recalled, ``yet everything around was
closed down most of the time. If you lived in a middle-class community in Chicago, children and
adults came daily to the door saying, 'We are starving, how about a potato?' I speak from poignant
memory.''

After receiving his bachelor's degree from Chicago in 1935, he went to Harvard, where he was
attracted to the ideas of the Harvard professor Alvin Hansen, the leading exponent of Keynesian
theory in America.

As a student at Chicago and later at Cambridge, Mr.~Samuelson had at first reacted negatively to
Keynes. ``What I resisted most was the notion that there could be equilibrium unemployment'' -- that
some level of unemployment would be impossible to eliminate and have to be tolerated. ``I spent four
summers of my college career on the beach at Lake Michigan,'' he said. ``It was pointless to look
for work. I didn't even have to test the market because I had friends who would go to 350 potential
employers and not be able to get any job at all.''

Eventually he was converted. ``Why do I want to refuse a paradigm that enables me to understand the
Roosevelt upturn from 1933 to 1937?'' he asked himself.

Mr.~Samuelson was perceived at the outset of his career as a brilliant mathematical economist. He
shot to academic fame as a 22-year-old prodigy at Harvard when he began a boldly sweeping and highly
technical doctoral dissertation, published as a book in 1947 by Harvard University Press.

At Harvard, as at Chicago, he was not shy about criticizing his professors -- ``respecting neither
age nor rank,'' according to James Tobin, a Nobel laureate of Yale University. The young
Mr.~Samuelson's chief complaint against economists was that they preoccupied themselves with finer
economic principles while all around them people were being thrown into bread lines.

\emph{A Bold Dissertation}

His attitudes did not endear him to the austere chairman of the economics department at Harvard,
Harold Hitchings Burbank, with whom he had a rocky relationship.

But the publication of his dissertation was an immediate success. It won him the John Bates Clark
Medal awarded by the American Economic Association to the economist showing the most scholarly
promise before the age of 40; it would eventually help him win his Nobel, and it was frequently
reprinted despite the heavy resistance of Professor Burbank, selling to economists around the world
for more than 20 years. (``Sweet revenge,'' Mr.~Samuelson said.)

Among Mr.~Samuelson's fellow students was Marion Crawford. They married in 1938. Mr.~Samuelson
earned his master's degree from Harvard in 1936 and a Ph.D. in 1941. He wrote his thesis from 1937
to 1940 as a member of the prestigious Harvard Society of Junior Fellows. In 1940, Harvard offered
him an instructorship, which he accepted, but a month later M.I.T. invited him to become an
assistant professor.

Harvard made no attempt to keep him, even though he had by then developed an international
following. Mr.~Solow said of the Harvard economics department at the time: ``You could be
disqualified for a job if you were either smart or Jewish or Keynesian. So what chance did this
smart, Jewish, Keynesian have?''

During World War II, Mr.~Samuelson worked in M.I.T.'s Radiation Laboratory, developing computers for
tracking aircraft, and was a consultant for the War Production Board. After the war, having resumed
teaching, he and his wife started a family. When she became pregnant the fourth time, she gave birth
to triplets, all boys.

Marion Samuelson died in 1978. Mr.~Samuelson is survived by his second wife, Risha Clay Samuelson;
six children from his first marriage: Jane Raybould, Margaret Crawford-Samuelson, William and the
triplet sons, Robert, John and Paul; and 15 grandchildren. Mr.~Samuelson is also survived by a
brother, Robert Summers, a professor emeritus of economics at the University of Pennsylvania and
father of Lawrence H.~Summers, director of President Obama's National Economic Council and former
secretary of the Treasury under President Clinton and former president of Harvard.

\emph{A Keynesian Textbook}

The birth of the triplets doubled the number of children in the Samuelson household, which soon
found itself sending 350 diapers to the laundry each week. His friends suggested that Mr.~Samuelson
needed to write a book to earn more money.

He decided to write an economics textbook, but one that would not only be compelling for students
but also sophisticated and comprehensive. And he wanted to center it on the still poorly understood
Keynesian revolution. President Herbert Hoover, he noted, had never referred to Keynes other than as
``the Marxist Keynes.''

``I never quite understood that venom,'' Mr.~Samuelson said.

He said he ``sweated blood'' writing his book, employing detailed charts, color graphics and humor.
He wrote: ``Economists are said to disagree too much but in ways that are too much alike: If eight
sleep in the same bed, you can be sure that, like Eskimos, when they turn over, they'll all turn
over together.''

It would be difficult to overestimate the influence of ``Economics.'' Business Week, taking note of
the textbook's publication in Greek, Punjabi, Hebrew, Russian, Serbo-Croatian and other languages,
once said that it had ``gone a long way in giving the world a common economic language.'' Students
were attracted to its lively prose and relevance to their everyday lives. Many textbook authors
began to copy its presentation.

``Economics,'' together with shrewd investing, made Mr.~Samuelson a millionaire many times over.

\emph{Friendship With a Rival}

A historian could well tell the story of 20th-century public debate over economic policy in America
through the jousting between Mr.~Samuelson and Milton Friedman, who won the Nobel in 1976.
Mr.~Samuelson said the two had almost always disagreed with each other but had remained friends.
They met in 1933 at the University of Chicago, when Mr.~Samuelson was an undergraduate and
Mr.~Friedman a graduate student.

Unlike the liberal Mr.~Samuelson, the conservative Mr.~Friedman opposed active government
participation in most areas of the economy except national defense and law enforcement. He thought
private enterprise and competition could do better and that government controls posed risks to
individual freedoms.

Both men were fluid speakers as well as writers, and they debated often in public forums, in
testimony before Congressional committees, in op-ed articles and in columns each of them wrote for
Newsweek magazine. But Mr.~Samuelson said he always had fear in his heart when he prepared for
combat with Mr.~Friedman, a formidably engaging debater.

``If you looked at a transcript afterward, it might seem clear that you had won the debate on
points,'' he said. ``But somehow, with members of the audience, you always seemed to come off as
elite, and Milton seemed to have won the day.''

Mr.~Samuelson said he had never regarded Keynesianism as a religion, and he criticized some of his
liberal colleagues for seeming to do so, earning himself, late in life, the label ``l'enfant
terrible emeritus.'' The experience of nations in the second half of the century, he said, had
diminished his optimism about the ability of government to perform miracles.

If government gets too big, and too great a portion of the nation's income passes through it, he
said, government becomes inefficient and unresponsive to the human needs ``we do-gooders extol,''
and thus risks infringing on freedoms.

But, he said, no serious political or economic thinker would reject the fundamental Keynesian idea
that a benevolent democratic government must do what it can to avert economic trouble in areas the
free markets cannot. Neither government alone nor the markets alone, he said, could serve the public
welfare without help from the other.

As nations became locked in global competition, and as the computerization of the workplace created
daunting employment problems, he agreed with the economic conservatives in advocating that American
corporations must stay lean and efficient and follow the general dictates of the free market.

But he warned that the harshness of the marketplace had to be tempered and that corporate downsizing
and the reduction of government programs ``must be done with a heart.''

Despite his celebrated accomplishments, Mr.~Samuelson preached and practiced humility. The M.I.T.
economics department became famous for collegiality, in no small part because no one else could play
prima donna if Mr.~Samuelson refused the role, and, of course, he did. Economists, he told his
students, as Churchill said of political colleagues, ``have much to be humble about.''

\section{With a New Phone, Google May Challenge Apple}

\lettrine{T}{wo}\mycalendar{Dec.'09}{14} titans of the tech world, Google and Apple, may soon be
engaged in hand-to-hand combat. Or, more precisely, handset-to-handset combat.

Google plans to begin selling its own smartphone early next year, company employees say, a move that
could challenge Apple's leadership in one of the fastest-growing and most important technologies in
decades.

Google's new touch-screen Android phone, which it began giving to many employees to test last week,
could also shake up the fundamentals of the cellphone market in the United States, where most phones
work only on the networks of the wireless carriers that sold them.

The company, using the power of its brand, plans to market and sell the new phone directly to
consumers over the Internet, and buyers would be able to sign up for service from any compatible
provider, the employees say.

The introduction of a Google phone -- manufactured to its hardware and software specifications by an
Asian maker of handsets -- would be an important and risky departure for Google. Until now, it has
made software to power cellphones that are built and marketed by partners, and it has largely
avoided selling hardware.

The apparent shift underscores the fact that mobile phones are quickly becoming the biggest
technology battleground of the future, as consumers increasingly rely on their phones to browse the
Internet and perform other computing tasks.

It also indicates Google's determination to make its mark on yet another industry, as it has done
previously in advertising, books and online videos.

But analysts say it is not clear that Google's success on the Internet will carry over into the
design, marketing and distribution of hardware. Many companies have tried to make similar shifts and
stumbled. Microsoft turned the Xbox into a hit, but when it pushed aside its partners in the
music-player business in favor of its own Zune, it failed to gain traction against Apple and its
iPods.

The phone's success could also depend on how Google chooses to price it. Most Americans buy phones
that are subsidized by wireless carriers, which recoup that cost by locking customers into
contracts. IPhones that cost consumers \$199 actually cost AT\&T about \$550, analysts say.

Katie Watson, a Google spokeswoman, declined to comment on the company's plans. She referred
reporters to a blog post published Saturday in which Google said that the new device was a ``mobile
lab'' that would let employees test out new technologies.

Google employees who spoke on the condition of anonymity because plans for the phone were supposed
to be confidential said that the device, manufactured by the Taiwanese company HTC, was thinner than
Apple's iPhone, with a slightly larger touch screen. It could be available as early as January, they
said.

Employees say the phone will be sold unlocked, meaning that buyers will be able to choose a service
provider, and will be based on G.S.M. technology, which is used by AT\&T and T-Mobile in the United
States and by most other carriers around the world. It is named Nexus One, according to various
reports and to digital traces that the phones have been leaving on Web sites.

Google had long insisted that it was not interested in building and selling phones, saying it
preferred to rely on hardware partners and wireless carriers to market a wide variety of phones
powered by Android, the operating system that it offers free.

In October, Andy Rubin, vice president of engineering for Android at Google, scoffed\cite{scoff} at
the idea that the company would ``compete with its customers'' by releasing its own phone, according
to the technology news service CNet.

Analysts say that the apparent shift signals a recognition by the company's executives that Google
needs to take more control of its destiny in the wireless world.

``They perceive mobile as the next major opportunity,'' said Jeffrey Lindsay, an analyst at Sanford
C.~Bernstein. ``It is too big a risk to drive the strategy through their partners. They want more
say and more control.''

In addition, analysts say that the iPhone, despite prominently featuring some Google services, makes
the company nervous.

``They don't want to have access to Google being controlled or influenced by one player like
Apple,'' said Ben Schachter, an analyst with Broadpoint AmTech, a research firm.

Google wants to get more people using Web-friendly phones in part because it depends on the growth
of search advertising, which is slowing on PCs. On cellphones, however, use of Google's mobile
search engine grew 30 percent in six months this year, Mr.~Schachter said. ``That's huge, and a
majority of that growth is coming from the iPhone,'' he said. ``When that happens, Apple has a lot
of power over influencing users' behavior.''

Until recently, Google and Apple were considered close allies with a common enemy: Microsoft. They
shared two board members, Eric E.~Schmidt, the chairman and chief executive of Google, and Arthur
Levinson, the former chief executive of Genentech.

But the rise of mobile phones has increasingly pitted the two companies against each other. More
than a year ago, Mr.~Schmidt began recusing himself from discussions of the iPhone in Apple board
meetings.

Google and Apple compete in other areas, like online distribution of music and video, Web browsers,
and soon, computer operating systems. But it is in mobile phones that the two companies appear to be
battling most aggressively. In recent weeks, Google has unveiled a string of cutting-edge services
for smartphones that run Android. They included an application that gives turn-by-turn driving
directions, which Google is giving away free. A comparable application from the satellite navigation
company TomTom for the iPhone costs \$99.

``Finally, Apple will meet a worthy competitor,'' said Ashok Kumar, an analyst with Northeast
Securities, who predicted two months ago that Google would develop and market its own phone.

Natalie Kerris, an Apple spokeswoman, declined to comment on Google's plans.

Analysts say that Google's gamble is a risky one, in particular because it could alienate Google's
Android partners, which include handset makers like Motorola and Samsung, and carriers like Verizon
Wireless, Sprint and T-Mobile.

Those partners are largely responsible for the recent momentum of Android. They have introduced a
handful of new devices in recent months, including the Motorola Droid, a well-reviewed and powerful
device. But Android phones continue to be outsold by the iPhone, which is also the phone of choice
for developers of tens of thousands of cellphone applications.

Analysts say Google may believe that only its own backing could make Android into a serious
competitor to the iPhone, which for now in the United States is available only through AT\&T. That
exclusivity is expected to end next year if the deal is not renewed.

It is not clear whether Google will team up with carriers to sell service plans for the phone, or
whether Google will subsidize the devices in the hope of making up the difference with profits from
mobile advertising.

``If it is not subsidized, then I suspect the impact will be small,'' said Charles S.~Golvin, a
wireless industry analyst at Forrester Research. Mr.~Golvin said that Nokia had tried to sell
unsubsidized high-end phones in the United States without much success. ``If they choose to sell it
subsidized, then the impact could be very, very significant,'' he added.

Analysts also said that Google's gambit might be aimed at important overseas markets, where
unlocked, unsubsidized phones tend to be the norm.

Mr.~Kumar said that while Google's plans may upset carriers, the company could find ways to
placate\cite{placate} them and to continue to work with them. ``As long as Google provides a
platform that the carriers can profit from, they will be happy participants,'' he said. ``Google
could share profits from advertising with the carriers.''

Much will depend, of course, on how appealing the phone is. Some people who say they have had an
early look have gushed about it on Twitter, with one calling it ``a sexy beast.''

\section{As Prices Fall, Blu-ray Players Are Invited Home}

\lettrine{B}{lu}\mycalendar{Dec.'09}{14}-ray, a high-definition variation of the DVD format
introduced three years ago, was initially met by a collective shrug from most consumers. Who needed
another black box to connect to the TV, the thinking went, even if it did promise to play movie
discs in clear, crisp high-definition?

But this year, even as the country moves tentatively out of a recession, consumers are buying the
devices at a faster pace than they bought previous generations of movie players like the VCR and the
DVD player. Analysts predict that sales of Blu-ray machines will be up 112 percent over last year,
one of the true bright spots in retailing this holiday season.

Blu-ray's household penetration is higher than that of DVD for the same period after introduction,
according to Shawn DuBravac, chief economist for the Consumer Electronics Association. Three years
after introduction, Blu-ray stand-alone players, which excludes Sony's PlayStation 3 game machines
that also play Blu-rays discs, are in 7 percent of American homes.

Amazon.com reports that sales of Blu-ray players are outnumbering those of standard DVD units,
according to Paul Ryder, the company's vice president for consumer electronics. Among the top 10
disc players sold, eight were Blu-ray, and five of the top 10 movie titles sold were in the Blu-ray
format.

At a Best Buy in Thousand Oaks, Calif., small mountains of Blu-ray players are stacked seven units
high wherever flat-panel TVs are on display, while the handful of standard DVD players are in a
distant aisle.

The main reason for the shift in thinking is right there on the stacks of Blu-ray players that Best
Buy has piled at its video wall. Prices for high-quality models have dropped below \$150, a steep
drop from the \$300-plus that retailers were charging when Blu-ray had its debut.

``The price for an impulse buy is under \$100, and we're getting there,'' said Andy Parsons,
president of the Blu-ray Disc Association, a trade group.

Amazon's best-selling Blu-ray model, Panasonic's DMP-BD60, is available for \$129, while a no-frills
Magnavox Blu-ray player is on sale at Wal-Mart Stores for \$78.

``We're sanguine about Blu-ray taking over as the physical disc format of choice,'' said Tom Adams,
the president of Adams Media Research. ``The differences in cost between DVD and Blu-ray players and
software is going away. So there will be a natural evolution from standard DVD to Blu-ray.''

It also helps that more Americans now have 1080p HDTV displays, sets with sufficiently high screen
resolution to fully display the clarity of a Blu-ray disc. Forty percent of all TVs sold are 1080p,
according to Riddhi Patel, an analyst at iSuppli. She says that number should rise to 46 percent
next year.

In addition to superior picture quality, Blu-ray also offers features not available on standard DVD
players, like pop-up menus during play. Some Blu-ray players can also connect to the Internet; the
player can then receive additional information about a movie, offer movie-related games, provide
updated trailers, and allow friends to simultaneously watch a film while writing comments on screen.

The newest generation of Blu-ray players lets viewers receive streaming movies and TV shows from
online providers like Amazon Video on Demand, CinemaNow, Netflix and YouTube.

As a result, Blu-ray manufacturers have placed themselves in a seemingly awkward position: They are
selling a device that relies on people to continue to buy discs, but the same device gives them a
way to download videos -- bypassing the discs the machines were built to play.

Consumer electronics companies believe that this strategy makes sense. The additional programming
will add to the Blu-ray experience, executives say, which still offers superior picture and sound
quality to what can be currently streamed over the Internet.

And since most HDTVs do not yet come with their own built-in Internet capability, a Blu-ray player
can become an inexpensive but essential all-in-one content source, according to Todd Richardson,
senior vice president for P\&F USA, the marketing arm for Philips.

``These services are supplemental. They fill out the consumers' demand for more and more content,''
added Tim Alessi, LG's director of new product development.

Yet, as high-speed broadband becomes ubiquitous, the ability to quickly download Blu-ray-quality
content will become a reality. That day is probably 10 years off, according to Ross Rubin, an NPD
Group analyst, as physical discs continue to provide a higher-quality image and an easier way to
move programming throughout the house than by trying to create an in-home network.

But the consumer electronics industry is no stranger to product life cycles and planned
obsolescence. Already, manufacturers are readying a new line of Blu-ray players and TVs that can
display video in 3-D. They should be in place by next year's holiday shopping season.

\section{Big Risk in a One-Man Brand Like Tiger Woods}

\lettrine{A}{ccenture}\mycalendar{Dec.'09}{14}, the giant consulting firm, ended its six-year
marketing relationship with Tiger Woods on Sunday, showing once again that in advertising as in
sports, there are no sure things.

Mr.~Woods had been featured in Accenture campaigns with the tagline, ``Go On. Be a Tiger,'' splashed
in business magazines and airport waiting rooms since 2003. Since most consumers have no idea what a
company like Accenture does, Mr.~Woods became the human face of the corporation and a means to
extol\cite{extol} the corporate virtues of performance and risk-taking.

In a statement Sunday, the company said Mr.~Woods, 34, was ``no longer the best representative'' for
its advertising. The action came more than two weeks after a bizarre car accident involving
Mr.~Woods, who is married, led to a series of embarrassing revelations about affairs he has carried
on with other women. Then on Friday he announced that he was taking a hiatus\cite{hiatus} from golf.

Accenture is the first of Mr.~Woods's many sponsors to withdraw its support completely. Gillette has
said that it would remove its Woods-related advertising for now and other sponsors, like AT\&T,
PepsiCo and Nike, have said they will wait to see when and how Mr.~Woods returns to golf before
making any permanent decisions. EA Sports, which has a line of best-selling Tiger Woods video games,
is also standing by the golf star.

In response to Accenture's announcement, Mark Steinberg, Mr.~Woods's agent, said, ``We are
disappointed but respect their decision.''

Accenture's decision is another example of how risky it can be for corporations to base their
marketing efforts on individual celebrities in an age when blogs, camera phones and the digital
footprint left by the celebrities themselves can quickly expose their foibles\cite{foible} to the
world.

Earlier this year, Kellogg dropped its sponsorship of Michael Phelps, the gold-medal swimmer, after
he was photographed using a bong at a party, saying, in language similar to Accenture's, that his
acts were not ``consistent with the image of Kellogg.'' In 2003, Nutella and McDonald's ended their
deals with Kobe Bryant after he was accused of sexual assault in Colorado; the charges were
eventually dropped.

Among corporate pitchmen, however, Mr.~Woods is unique. He dominates a sport for individuals that is
marketed around the world while still being hugely popular in the United States, in contrast to
sports like tennis.

More than any other athlete, Mr.~Woods has built a billion-dollar brand around himself that has not
only benefited him but has also made millions of dollars for his sponsors. Perhaps the best analogy
for Mr.~Woods's predicament is not another athlete but Martha Stewart -- another example of a
billion-dollar brand based on the cultivated image of one person.

Ms.~Stewart put her public company and television contracts at risk when she was convicted of lying
to prosecutors about a minor stock deal. She emerged from her five-month sentence in prison to
retake the reins at Martha Stewart Living Omnimedia, but the company has not returned to the
financial heights it once enjoyed.

``Tiger is the best example of a walking, individual corporation,'' said Ben Porritt, a public
relations executive who advised Alex Rodriguez, the Yankees third baseman, last spring after his
onetime steroid use was disclosed.

``Tiger is going to come out of this as somewhat of a bankrupt brand,'' said Mr.~Porritt. ``He will
have to restructure and go forward.'' But, he said, ``It's going to be an ugly few months.''

Several companies that measure consumer reaction believe the ugliness has already started. Zeta
Interactive, a digital ad agency that monitors message boards, blogs and social media posts, said
that positive sentiment toward Mr.~Woods had already plummeted. Before the accident, buzz about the
golfer was 91 percent positive; by Friday, that figure had sunk to 43 percent.

The turnabout ``is the quickest fall from positive to negative we've ever seen,'' said Al DiGuido,
chief executive at Zeta Interactive.

The Nielsen IAG unit of the Nielsen Company found that the accident quickly affected the brands
Mr.~Woods peddles. There were more than 20 instances through Monday, Dec.~7, of jokes being made on
late-night talk shows that paired Mr.~Woods with one of his sponsors by name, according to Nielsen
IAG.

The recall among viewers of the brand mentioned in the joke was 55 percent, according to Nielsen
IAG, compared with a norm in late-night shows of 39 percent.

About 11 percent of those viewers who recalled a brand with ties to Mr.~Woods said they had a
negative opinion of the brand, Nielsen IAG reported, compared with the average negative opinion of a
brand mentioned in a late-night show of 6 percent.

David Arluck, an independent sports marketing executive, said that Accenture was in a different spot
than other sponsors because it pitched a service to businesses, not consumers, and so it had
centered its branding on Mr.~Woods's persona.

``In business, credibility means a lot and honesty means a lot, and Accenture wants to be a company
that can be trusted,'' he said. ``Obviously Tiger Woods did some things that betrayed the trust of
people who were close to him, and I can see how they would want to distance themselves from that.''

Even the most upright athletes seem to be aware of the tenuousness of image in the Internet age. On
a recent evening in Manhattan, the Yankees shortstop Derek Jeter -- a close friend of Mr.~Woods --
stood on a dais after receiving Sports Illustrated's Sportsman of the Year award and was asked,
essentially, how he was able to enjoy the perks of athletic superstardom for so many years and
maintain his pitch-perfect image.

``Don't jinx me,'' Mr.~Jeter said, to nervous laughter.

\section{Remodeling ABC News, Without Too Much Dust}

\lettrine{F}{or}\mycalendar{Dec.'09}{14} changes that seem so sweeping, the reordering of the most
important players at ABC News this month is being played more like a couple of casual early moves in
a game of checkers.

Diane Sawyer is off ``Good Morning America'' and into the anchor chair of the network's evening
newscast, ``World News,'' starting Dec.~21; George Stephanopoulos is off the Sunday show ``This
Week'' and into Ms.~Sawyer's old slot, starting Monday.

For ABC News, the moves have broad implications. Both programs rank No.~2 to NBC's offerings, and
the competitive situation is only getting worse.

But the changes have been accompanied by a scrupulous\cite{scrupulous} and clearly conscious lack of
fanfare. Not only are both new anchors sliding onto the air with a minimum of promotion, they are
doing it in December, when viewing levels are lower -- and most first-stringers in all forms of
television depart for extended holiday leaves.

David Westin, the president of ABC News, said that although they might be fewer in number in
December, ``viewers are still there, and there is still news.''

But there is little doubt that ABC is deliberately trying to avoid the kind of media frenzy that
surrounded the arrival of Katie Couric in the anchor chair at CBS. One ABC executive said that
Ms.~Sawyer's first newscast is unlikely to be subjected to a raft of critical pieces in the press or
on Internet blogs -- as Ms.~Couric's was.

Though he did not mention Ms.~Couric's example specifically, Mr.~Westin said that ``some of these
blowouts can be a little self-referential,'' adding that ``I always prefer to overperform rather
than overpromise.''

``World News'' with Charles Gibson as anchor broke through briefly to claim the top-rated spot more
than a year ago. But for 59 of the last 60 weeks, NBC and its anchor, Brian Williams, have had the
most-watched evening newscast. And this year, Mr.~Williams has widened the gap to an average of
880,000 viewers, up from 318,000 last year. (CBS has been stuck in third place.)

Nor will it help that after September 2011, ABC stations will be without one of the most valuable
lead-ins in television, ``The Oprah Winfrey Show.''

Still, it is the morning -- where most of the money is made in network news -- that remains the
daunting challenge for ABC. NBC's ``Today'' show hasn't lost so much as a single week in 14 years,
and the trend line is hardly encouraging for its competitors. In the past year, ``Today'' has
increased its average weekly advantage over ``Good Morning America'' to 1.3 million viewers, from
880,000.

And ABC is now changing the leading player on both its shows at the same time.

``That does add an element of drama, now doesn't it?'' said Andrew Lack, the former president of NBC
News, who is now the chief executive of Bloomberg's multimedia group.

Mr.~Lack agreed with the network's low-key approach, but another executive with experience working
with network newscasts and anchors in the midst of transitions suggested another reason for the soft
sell.

``They are clearly using this low-key launch as a strategic cover for their declining ratings,''
said the executive, who asked not to be identified because of professional connections to people
affected by the changes at ABC.

Mr.~Westin said that both shows also would undergo change over that period. For the morning, that is
likely to mean an approach that is ``a little harder, a little more urgent,'' said Jim Murphy, the
executive producer of ``Good Morning America.'' In part, that is because Mr.~Stephanopoulos, like
all anchors, will bring with him his preferences -- especially political reporting.

``You always play to the strengths of your anchor,'' Mr.~Murphy said.

The morning has so many elements -- not just newsmakers, but also latte makers -- that adjustments
may be more noticeable than on the evening newscast, where the format has resisted changes.
Ms.~Couric learned that when she tried a number of innovations early on that have since been
dropped.

Ms.~Sawyer, who is likely to benefit from not having to blaze the trail as the first woman to anchor
solo, is not setting out to change the model.

Jon Banner, the executive producer of ``World News,'' said there might be some changes in the show's
look, but even the set Mr.~Gibson has used will remain the same. Mr.~Westin said Ms.~Sawyer might
engage in more cross-talk with correspondents (something Ms.~Couric also does) because ``Diane's
style is more conversational.''

But Ms.~Sawyer is also known for pushing herself -- and her staff -- extremely hard. One producer
who has worked with her in the past said, ``They have no idea what they are in for; she will shake
up that show.''

\section{Can Nokia Recapture Its Glory Days?}

\lettrine{I}{f}\mycalendar{Dec.'09}{14} there's anywhere left in the world where it's still impolite
to flash a BlackBerry or an iPhone, it's Nokia's annual analyst meeting.

But earlier this month, as executives talked up the company's plans for 2010, the optimistic message
from the stage was belied by the behavior of the audience. In the back of the room, one money
manager after another distractedly toyed with a competing device, typically a BlackBerry, even as
cheery PowerPoint slides promoted Nokia's latest offerings.

Francois Meunier, an analyst with Cazenove in London, whispered doubts about the presentation as he
tried to catch the eye of one of the floor managers handing out microphones for the
question-and-answer session. Finally, it was Mr.~Meunier's turn, but before he could ask an actual
question, he couldn't resist declaring publicly what he'd been muttering all afternoon.

``I don't think anyone in this room is expecting an improvement in earnings next year,'' he told the
assembled executives, before asking whether Nokia's 4 percent dividend is sustainable.

Mr.~Meunier's downbeat assessment of the once-mighty mobile phone maker's prospects in 2010 comes
after an equally gloomy 2009, a year the company would just as soon forget.

Although Nokia, based near Helsinki in Espoo, still commands 37 percent of the world's handset
market, it's facing bruising competition in the lucrative high end of the industry, where Apple's
iPhone and Research in Motion's BlackBerry have grabbed the cool factor in smartphones that can surf
the Web and handle e-mail.

``The whole user experience is a nightmare,'' moans Nick Jones, a senior analyst with Gartner, which
tracks the technology sector. ``It's just not in any sense a competitive experience with iPhone.''

Olli-Pekka Kallasvuo, the company's taciturn chief executive, admits the mood out there is gloomy,
especially on Wall Street. ``We are not getting the benefit of the doubt,'' he said in an interview
the day after the analysts' meeting. ``We need to change that.''

Nokia's problems are especially acute in North America, where its hold on smartphones equals a
barely visible 3.9 percent, compared with 51 percent for Research in Motion and 29.5 percent for
Apple, according to Gartner. As if to underscore its problems in the United States, Nokia announced
Thursday that it would shutter its flagship stores in New York and Chicago.

``We made wrong decisions in the American market,'' says Kai Oistamo, executive vice president for
devices. For example, Nokia was slow to make the change to so-called clamshell phones, sticking with
``monoblock'' models even as consumers abandoned them.

And while Nokia first offered touch-screen technology in 2004 -- three years before the debut of the
iPhone -- Apple's models quickly made Nokia's competing products look stodgy. Most of Nokia's
touch-screen phones can't quickly transform their screen with the jab of a finger, which is among
the factors that make the iPhone seem so much more slick.

Until recently, according to both Nokia executives and industry experts, the company didn't want to
produce phones specifically tailored for American consumer tastes, and it resisted demands from the
major carriers to come up with phones based around their brands and individual specifications.

``The market in the U.S.~has always been dominated by the carriers, so they call the shots,'' says
Carolina Milanesi of Gartner. ``And Nokia has had a difficult relationship with the carriers.''

Nokia has also been hobbled by its traditional weakness in phones employing C.D.M.A., the wireless
technology offered by Sprint and Verizon Wireless that's used by about 50 percent of American
consumers. (Sprint's current lineup does not include any Nokia models.) Nokia focuses instead on
G.S.M. phones for AT\&T and T-Mobile. However, AT\&T's exclusive deal with Apple has hurt Nokia in
the high-end smartphone market.

And though Nokia sells a lot of smartphones elsewhere in the world, its share of the global
smartphone market has fallen to 39.3 percent, down from 42.3 percent a year ago. Even in Nokia's
home base of Europe, the iPhone is rapidly gaining in popularity.

Nokia is finally responding -- its lithe, BlackBerry-like E72 appeared in the United States on
Tuesday -- but it is facing looming threats in other segments.

Google is offering Android, a rival to Nokia's own operating system, which has been picked up by
competitors like HTC, Motorola and Dell, while Asian manufacturers are turning up the heat with
low-priced handsets in emerging economies where Nokia has long enjoyed outsize market share.
Meanwhile, Apple and Nokia are locked in a legal battle over patents.

``Nokia faces competition everywhere,'' says Sherief Bakr, a Citigroup analyst. ``At the high end
from Apple, in the midrange by Research in Motion, and by the Koreans and the Chinese in the low
end.''

ALL in all, it's enough to make the mood as grim as a December day in Helsinki, where the sun
struggles to get above the horizon by 9 a.m and night falls at 4 p.m.

Once a stock market darling, Nokia shares have fallen 20 percent since September even as the broader
market has rallied. The company reported its first quarterly loss in more than a decade in October
after a \$1.3 billion write-down in its equipment business.

Here in Finland, Nokia's problems are felt especially keenly. Nokia accounts for 25 percent of the
Helsinki stock exchange's capitalization and one-third of Finland's total research and development
spending, according to Jyrki Ali-Yrkko, of the private Research Institute of the Finnish Economy.

Deeper than the numbers, however, has been the damage to Nokia's role as a wellspring of pride in a
country historically known for exporting wood and paper products, not high tech.

Nokia's roots go back to 1865, and as recently as the 1980s, its products included not only cable
and telecom equipment but also rubber boots and toilet paper. But in the early 1990s, many
businesses were spun off in favor of the growing cellphone sector. By the mid-1990s, under its
former chief executive and current chairman, Jorma Ollila, the profits were rolling in.

Nokia quickly became one of Europe's rare technology success stories, an exception in an industry
dominated by American and Japanese giants. And in a traditional, social-democratic Nordic country
where ostentatious displays of wealth are frowned upon, hundreds of long-time employees became Nokia
millionaires, says Mr.~Ali-Yrkko.

``Nokia has been the flagship of Finland in terms of a company succeeding on a global scale,'' he
says. ``But that sense of glory we had has disappeared, or at least diminished.''

The problems have reached all the way to Finland's national coffers. In 2007, Nokia paid 18 percent
of Finland's overall corporate taxes, but that dropped to 9 percent last year, and the contribution
is expected to be even lower in 2009. The Finnish government may have to increase borrowing to make
up for the shortfall, warns Mr.~Ali-Yrkko.

A lawyer by training, the C.E.O. Mr.~Kallasvuo is a much more cautious leader than his predecessor,
the charismatic Mr.~Ollila, who some Finns thought might go into politics after he stepped down as
chief executive in 2006. At times, Mr.~Kallasvuo seems uneasy when pressed for his vision of Nokia's
future, and repeated earnings disappointments have led many analysts to question whether his dour
style is what's needed as new competitors circle.

``The market believes this is a management team that can't and won't execute,'' says Mr.~Bakr.
``There is a large element of investors who are not convinced that Kallasvuo is the man who can make
this transition and compete with the likes of Steve Jobs.''

Despite the pessimism outside, Mr.~Kallasvuo insists spirits are still high inside the company.
``Competition is nothing new; we've been attacked by many players,'' he says. And while last
quarter's performance ``was difficult for me and the C.F.O., it wasn't a difficult moment for an
excited Nokia engineer who wants to change the world.''

IF Olli-Pekka Kallasvuo seems to have stepped off the set of an Ingmar Bergman movie, then Anssi
Vanjoki's charisma and sculpted features recall Michael Douglas. An 18-year veteran of Nokia, he is
the executive vice president for markets as well as something of a standout in Nokia's geeky
culture. In a country where speeding tickets are directly tied to income on a sliding scale, he
racked up a 116,000 euro (\$170,000) fine racing his motorcycle through Helsinki, although he was
able to negotiate that figure down somewhat.

Much more feisty than Mr.~Kallasvuo, he is unwilling to admit Nokia has lost any of its competitive
edge. ``We have not lost our ability to innovate; we have not lost our ability to truly understand
the consumer and make intuitive solutions for them,'' Mr.~Vanjoki says.

Indeed, for all the new competition in smartphones, Nokia remains the dominant player in
conventional handsets, selling roughly 15 phones a second worldwide, according to the company,
including the Nokia 1201, a basic model that is its best seller. Analysts project revenue in 2010
will top \$60 billion, while profit is expected to equal \$3.5 billion next year as the overall
phone market grows 10 percent.

And while market share might be minuscule in North America, the company commands a whopping 62.3
percent of the market in the Middle East and Africa, as well as 48.5 percent in Eastern Europe and
41.8 percent in Asia. ``We are the incumbent behemoth of the mobile arena,'' Mr.~Vanjoki boasts.

What's more, Nokia has been written off before.

Citing past crises in 1998 (the advent of smaller phones), 2001 (the bursting of the tech bubble)
and 2004 (the sudden popularity of flip phones), Mr.~Vanjoki says. ``we've always had points where
technology hit a plateau and had to be reconfigured.''

SO why didn't Nokia move more quickly to counter Apple and Research in Motion in smartphones? ``We
didn't execute; we were aiming at too geeky a community,'' he says. ``Apple is made for the common
man. It's more for Joe Six-Pack than techno-geeks. But we understand Joe Six-Pack too.''

The coming 12 months will show whether Mr.~Vanjoki's confidence is warranted, and he better be right
as far as shareholders are concerned, since smartphones are where the growth is.

By 2013, Gartner predicts smart device sales will represent 82.5 percent of the mobile phone
business in Western Europe, and 58.2 percent of sales in North America and 18.2 percent in Asia.
Nokia generated \$5.6 billion in sales from conventional phones in the third quarter worldwide,
compared to \$4.6 billion for smartphones.

Nokia executives say new offerings like the N900, which is as much a mobile computer as it is a
phone, or the N97 Mini, which combines touch-screen technology with a qwerty keyboard, will win back
buzz from Apple and BlackBerry while appealing to the company's 1.1 billion customers. Then there's
the X6, out this month, which includes Nokia's Comes With Music plan, allowing users to choose from
millions of songs they can download free from Nokia's Music Store.

Another crucial development in 2010, according to Mr.~Kallasvuo, will be a bigger push for North
American market share, as Nokia works more closely with carriers and brings out more smartphones.
``We have not invested enough there,'' he says. ``It's a necessity for us.''

Although it's still secret, Nokia executives are also promising a smartphone for next year that will
update the company's aging Symbian operating system, combining the touch-screen coolness of the
iPhone with a BlackBerry-like e-mail solution. ``We intend to give R.I.M. a run for their money,''
says Mr.~Kallasvuo.

And though Nokia's flagship outlets in the United States may be folding, the Finnish giant is still
trying to compete directly with Apple online, opening Ovi (``door'' in Finnish) in May to compete
with Apple's hugely successful Apps Store.

Looking out further, Nokia's engineers are promising nifty new features like the ability to simply
point your phone at a friend to connect to the person's Facebook page.

For all of Nokia's mistakes, Citigroup's Mr.~Bakr says the company can bounce back yet again.
``Sitting in London or New York, you don't appreciate the dynamics of Nokia's huge market share,
especially in emerging markets,'' he says. ``I think they know what they've done wrong and what they
need to do right. It's just a question of whether they can execute in time.''

\section{Menopause, as Brought to You by Big Pharma}

\lettrine{M}{illions}\mycalendar{Dec.'09}{14} of American women in the 1990s were told they could
help their bodies ward off major illness by taking menopausal hormone drugs. Some medical
associations said so. Many gynecologists and physicians said so. Respected medical journals said so,
too.

Along the way, television commercials positioned hormone drugs as treatments for more than hot
flashes and night sweats -- just two of the better-known symptoms of menopause, which is technically
defined as commencing one year after a woman's last menstrual cycle.

One commercial about estrogen loss by the drug maker Wyeth featured a character named Dr.~Heartman
in a white coat discussing research into connections between menopause and heart disease,
Alzheimer's disease and blindness.

``When considering menopause, consider the entire body of evidence,'' Dr.~Heartman said. ``Speak to
your doctor about what you can do to help protect your health during and after menopause.''

Connie Barton, then a medical office assistant in Peoria, Ill., was one woman who responded to such
messages. She says she took Prempro, a hormone drug made by Wyeth, from 1997, when she was 53, until
2002, when she received a diagnosis of breast cancer. As part of her cancer treatment, she had a
mastectomy to remove her left breast.

Now Ms.~Barton, who said in an interview that she used Prempro in part because her doctor told her
it could help prevent heart disease and dementia, is one of more than 13,000 people who have sued
Wyeth over the last seven years, claiming in courts across the country that its menopause drugs
caused breast cancer and other problems.

The suits also assert, based on recently unsealed court documents, that Wyeth oversold the benefits
of menopausal hormones and failed to properly warn of the risks.

In October, a jury in a Pennsylvania state court awarded Ms.~Barton \$75 million in punitive damages
from Wyeth on top of compensatory damages of \$3.75 million.

The drug giant Pfizer, which absorbed Wyeth and its hormone drugs in a merger this year, says that
Prempro is a safe, federally approved drug that did not cause Ms.~Barton's breast cancer. Chris
Loder, a Pfizer spokesman, says Wyeth acted responsibly by including a clear warning about a breast
cancer risk on Prempro labels and by updating the warning as new evidence emerged.

Mr.~Loder also notes that Pfizer plans to appeal every product-liability case on menopausal drugs it
loses, including Ms.~Barton's.

While Wyeth has faced periodic complaints about its blockbuster menopause drugs, the latest lawsuits
have turned the company's menopausal hormone franchise into the kind of case study dissected at Ivy
League business schools. Lawyers have made some documents public in the suits, and The New York
Times and the nonprofit Public Library of Science filed successful motions to unseal thousands of
documents in July.

To be sure, even some doctors who think hormone therapy has risks say it is the most effective
treatment for symptoms directly associated with menopause.

The documents that have surfaced in the Wyeth cases offer a rare glimpse inside the file cabinets
and hard drives of a major drug company. And the cases demonstrate the importance of litigation in
detailing exactly how drug makers operate their businesses, says Dr.~Jerome L.~Avorn, a professor of
medicine at Harvard Medical School who has written about the subject in The Journal of the American
Medical Association.

``The information coming out in litigation helps us understand how a belief in a 'protective
benefit' of estrogens on the heart was able to spread like wildfire through the medical community,''
says Dr.~Avorn, who is not involved in the Wyeth litigation.

``Thousands of doctors prescribed the drugs for millions of women on that basis,'' he says, adding
that studies later contradicted the belief. ``It will be very interesting to see whether the courts
are able to connect the dots and make it clear whether this was a kind of medical ventriloquism on
Wyeth's part.''

PREMPRO is a combination of Premarin, an estrogen drug derived from the urine of pregnant mares and
first approved by the Food and Drug Administration in 1942, with an additional hormone, progestin.

Part of the Premarin saga shows how a drug maker successfully and cannily expanded a franchise whose
central ingredient is horse estrogens into a billion-dollar panacea for aging women. Yet several
hundred pages of court documents also raise questions about another aspect of Premarin's trajectory:
how Wyeth worked over decades to maintain the image and credibility of its hormone drugs even as the
products were repeatedly under siege.

Pfizer representatives say court documents paint an unfair picture of Wyeth's practices and that
plaintiffs' lawyers have cherry-picked documents for out-of-context comments to sway juries.

Still, the documents offer a snapshot of Wyeth's efforts. Taken together, they depict a company that
over several decades spent tens of millions of dollars on influential physicians, professional
medical societies, scientific publications, courses and celebrity ads, inundating doctors and
patients with a sea of positive preventive health messages that plaintiffs' lawyers say deflected
users' attention from cancer concerns.

Even as evidence mounted of an association of the drugs with cancer -- first in the 1970s with
Premarin and endometrial cancer, then in the 1990s with Prempro and breast cancer -- Wyeth tried to
contain the concerns, the court documents show. (A note handwritten in 1996 by a Wyeth employee
responding to a new report of breast cancer risks associated with hormone therapy said:
``Dismiss/distract.'')

In 2002, researchers halted the largest clinical trial ever conducted of women's health because
participants who took certain combined hormones had an increased risk of breast cancer -- as well as
a higher risk of heart attack, stroke and blood clots in the lungs -- compared with those taking a
placebo.

Other parts of the same federal study, called the Women's Health Initiative, later found that
hormone drugs increased the risk of dementia in a subset of participants, those age 65 and older.

Sales of Wyeth's hormone drugs peaked at about \$2 billion in 2001, but after results of the 2002
study came out sales plummeted.

Pfizer is now fighting the Prempro litigation along with lawsuits over its progestin drug, Provera.
Mr.~Loder, the Pfizer spokesman, says Pfizer and Wyeth had fully informed patients, doctors and
regulators of the risks of their menopause drugs, based on the best available science at the time of
the disclosures.

``They provided accurate warnings, performed studies on benefits and risks, and kept the F.D.A.
fully informed,'' he says.

But last month, a federal appellate court in St.~Louis ruled in the case of a plaintiff named Donna
Scroggin that Wyeth's inaction over accumulating evidence -- and the company's attempts to mitigate
cancer concerns by trying to undermine unfavorable scientific reports -- could allow a jury to find
Wyeth guilty of malicious conduct and award punitive damages.

For its part, Pfizer contends that two state judges in Pennsylvania have reached the opposite
conclusion: that juries should not be allowed to award punitive damages because there was
insufficient evidence of corporate misconduct.

Whichever direction the different cases ultimately follow, the court papers associated with them
illustrate a pattern in the history of hormone therapy. First, many doctors enthusiastically
prescribe hormone therapy drugs. Then a few researchers publish studies cautioning about risks,
causing sales to fall. And finally, some doctors start prescribing a new iteration of hormone drugs.

For example, Prempro now comes in lower doses. Prempro labels say the drug should be prescribed for
the shortest duration appropriate for the treatment goals and risks of the individual woman;
previous labels on Wyeth's hormone drugs for decades gave the same advice. The current label also
says that using estrogens, with or without progestins, may increase a woman's chance of heart
attack, stroke, breast cancer and blood clots.

MENOPAUSAL hormone therapy has long been pitched as a way to stave off what some doctors viewed as
the undesirable aspects of female aging.

In the popular 1966 book ``Feminine Forever,'' Dr.~Robert A.~Wilson, a gynecologist, used
disparaging descriptions of aging women (``flabby,'' ``shrunken,'' ``dull-minded,'' ``desexed'') to
upend the prevailing idea of menopause as a normal stage of life. Women and their physicians,
Dr.~Wilson wrote, should regard menopause as a degenerative disease that could be prevented or cured
with the use of hormone drugs.

``No woman can be sure of escaping the horror of this living decay,'' Dr.~Wilson wrote. ``There is
no need for either valor or pretense. The need is for hormones.''

Premarin had been available for decades, but Dr.~Wilson's book propelled the idea of hormone
``replacement'' into the popular consciousness and onto physicians' prescription pads. The
revivifying drugs promised to inhibit the ravages of time on the appearance and the psyche,
Dr.~Wilson wrote.

As the popularity of estrogen grew, an increasing number of women developed cancer of the uterine
lining, the endometrium. In 1975, an F.D.A. panel concluded there was a link between Premarin and
endometrial cancer. The company then sent a letter to doctors trying to mitigate such concerns,
documents show.

``Dear Doctor,'' wrote Dr.~John B.~Jewell, at the time the medical director of Ayerst, the Wyeth
predecessor. ``It would be simplistic indeed to attribute an apparent increase in the diagnosis of
endometrial carcinoma solely to estrogen therapy.'' Women may still receive ``proven benefits,'' he
wrote, by using ``the lowest maintenance dose needed to control the menopausal symptoms.'' He added
that the company planned to study the issue further.

F.D.A. officials then met with company officials, saying they were ``incensed'' that the letter was
``intended to obfuscate the issues,'' according to a 1976 memo signed by the F.D.A. and the company.
The F.D.A. said it would issue a bulletin saying there was a clear link between estrogen therapy and
endometrial cancer. In 1976, the maker of Premarin added a warning to the label about the risk of
endometrial cancer.

But the company never conducted further studies on the risk of developing endometrial cancer,
according to the St.~Louis appeals court decision.

The company instead focused its risk research on the possibility of breast cancers associated with
hormone replacement therapy. But two studies published in the mid-1970s in The New England Journal
of Medicine reported that taking estrogen therapy had increased the risk of endometrial cancer by at
least five times.

Reports in 1975 about endometrial cancer ``resulted in a precipitous decrease in estrogen use,''
according to a history of hormone therapy in The American Journal of Medicine in 2005.

In 1980, researchers at Boston University Medical Center estimated that the use of hormone therapy
had caused more than 15,000 cases of endometrial cancer in the United States between 1971 and 1975
alone.

``This represents one of the largest epidemics of serious iatrogenic disease'' -- meaning disease
caused by physician-administered treatments -- ``that has ever occurred in this country,'' the
authors wrote. (Mr.~Loder said Pfizer was not familiar with that report.)

Today, physicians prescribe Premarin to women who have had hysterectomies and therefore are not at
risk for endometrial cancer.

BY the mid-1990s, after a few studies had reported a protective effect of hormone drugs on the
heart, Wyeth had begun to reposition menopausal hormone therapy as a preventive health choice that
could help inhibit heart disease and other maladies, according to court documents.

That marketing strategy coincided with the introduction of Wyeth's new combination hormone drug
Prempro, which included a progestin hormone to keep estrogen from causing excessive cell growth in
the uterine lining.

In one commercial from a Wyeth research institute, the model Lauren Hutton runs down a beach and
warns of the health risks of estrogen loss.

``My doctor said if you don't replace estrogen that you lose at menopause, your risk for certain
age-related diseases could increase,'' Ms.~Hutton said in the commercial. In a voice-over, a
narrator told viewers about studies looking into the connections between menopause and heart
disease, memory loss and sight loss.

``Believe me,'' Ms.~Hutton said, ``the time to protect your future is now.''

Sally Beatty, a spokeswoman for Pfizer, said this was a ``help seeking'' ad, of the type encouraged
by the F.D.A. She added that the promotion did not mention any specific drugs, not did it suggest
that drugs could cure the ailments described.

The labels for Premarin and Prempro at the time said the drugs were approved to treat moderate to
severe symptoms of menopause like hot flashes, night sweats and vaginal dryness and to prevent
osteoporosis.

But Wyeth also positioned its menopausal hormone drugs as having larger protective benefits, court
documents show.

Wyeth used proxies to promote a wide range of health benefits from hormone therapy, paying millions
of dollars to influential doctors and medical groups and helping them develop abstracts for medical
conferences and articles for medical journals, according to court documents.

The company also paid \$12 million to sponsor continuing medical education programs from 2002
through 2006 at the University of Wisconsin, Madison. The programs, including an assertion that the
Women's Health Initiative and another heart-risk study ``miss the mark on quality of life,'' reached
thousands of doctors.

Doctors were aware in the 1990s that hormone therapy could increase a woman's risk of breast cancer,
says Dr.~Carol Bates, the director of the primary care program at Beth Israel Deaconess Medical
Center in Boston.

But based on the results of observational studies that had been published, many physicians, herself
included, believed that the drugs' ostensible ability to reduce heart attacks and perhaps
Alzheimer's would outweigh a breast cancer risk, she says.

``In the 1990s, there was actually tremendous pressure to put women on hormone therapy, and it came
from a good place,'' Dr.~Bates says. ``But it was taken a bit to the extreme.''

HORMONE therapy -- aimed at the symptoms it effectively treats and with full disclosure about its
possible risks -- has many advocates. Dr.~Lynne T.~Shuster, the director of the women's health
clinic at the Mayo Clinic in Rochester, Minn., says such regimens can be very worthwhile for
treating hot flashes, night sweats and vaginal dryness associated with menopause.

And some users agree.

Irene Fisher, a kitchen and bath designer in Baldwin, N.Y., says she has been taking Prempro for 16
years to control hot flashes and night sweats.

``I always feel good when I take it,'' she says. The benefits are worth a small risk, Ms.~Fisher
says, adding that she has an annual mammogram to check for breast cancer and that ``I think you have
to know your own body.''

But many women were not so fully informed in the 1990s.

In 1996, for example, a federal study reported that breast cancer risk may have been ``substantially
underestimated.'' Wyeth reacted with plans to dismiss it as ``just one more paper,'' and try to
``overshadow'' it by directing journalists to other studies, according to documents cited in the
court of appeals decision in Missouri.

In 1997, Wyeth began working with DesignWrite, a company in Princeton, N.J., that is paid by drug
makers to develop manuscripts for publication in medical journals. The specific objective of a
publication plan for Premarin was to ``increase physician awareness on the multitude of benefits
that hormone replacement therapy provides'' and ``diminish the negative perceptions associated with
estrogens and cancer,'' according to a 1997 DesignWrite proposal prepared for Wyeth.

Over the next decade, Wyeth paid DesignWrite to prepare at least 60 articles for publication in
medical journals on the potential benefits of hormone therapy for cardiovascular disease,
Alzheimer's disease, diabetes, colon cancer, vision loss and other health problems, the court
documents show.

In response to an e-mail query, Michael Platt, president of DesignWrite, wrote that the articles
were all medically and scientifically accurate and valid and peer reviewed.

But Wyeth's and DesignWrite's effort hit an obstacle in 2002 when researchers reported the results
of the Women's Health Initiative.

The National Institutes of Health ultimately decided to start using the term ``menopausal hormone
therapy'' instead of ``hormone replacement therapy,'' says Marcia L.~Stefanick, a professor of
medicine at the Stanford University medical school who was principal investigator on the Women's
Health Initiative study at her institution.

While the drugs are effective at treating symptoms like hot flashes, she says, the word
``replacement'' implies that women need hormone drugs after menopause. ``But there is no good
evidence that women need this after menopause,'' Professor Stefanick says.

In 2003, Wyeth added a ``black box'' warning to the label saying Prempro should not be prescribed to
prevent cardiovascular disease.

The same year, the F.D.A. approved a lower dose version of Prempro so women would have more options.

Today, many doctors who once offered hormone therapy to women without symptoms like hot flashes
limit the use of the drug to those with symptoms, prescribing low doses for a short time.

``Right now, the big difference is we do not recommend hormone therapy for good health or health
promotion or anti-aging,'' says Dr.~Shuster of the Mayo Clinic.

And even with lower-dose hormones, doctors do not have a uniform view on what constitutes the
optimal duration.

Dr.~Adriane Fugh-Berman, an associate professor at the medical school of Georgetown University,
considers both Premarin and Prempro examples of drugs that gained widespread popularity before
science had established the full extent of their risks.

``Where there has always been a push is where there isn't data,'' says Dr.~Fugh-Berman, who has been
a paid expert witness for plaintiffs in the hormone litigation. ``Now, low-dose hormones are being
pushed.''

LIKE Dr.~Wilson, the gynecologist in the 1960s who identified the evils of menopause, contemporary
voices are advocating hormones as an anti-aging treatment.

The actress Suzanne Somers, for example, has identified her own lineup of maladies, which she calls
the Seven Dwarves of Menopause: ``Itchy. Bitchy. Sweaty. Sleepy. Bloated. Forgetful. All Dried Up.''

In books with titles like ``The Sexy Years'' and ``Ageless,'' Ms.~Somers has promoted the use of
``bio-identical'' hormones, which can be prescribed by doctors in customized doses and can be
prepared individually by pharmacies.

But Dr.~Shuster of the Mayo Clinic says the hormones have not been extensively studied for safety
and efficacy. And, unlike branded hormone therapy, she says, they have not been approved by the
F.D.A.

Women, Dr.~Shuster says, should not assume that compounded hormones are safer than F.D.A.-approved
menopausal hormone drugs. Nevertheless, with sales of more than two million books, Ms.~Somers has
become a menopause guru to millions.

``I think I had a lot to do with making the word 'menopause' mentionable,'' Ms.~Somers, 63, said in
a phone interview last week. She said the compounded hormones were safe, and she sent some articles
from medical journals to back up her point.

In fact, much of Ms.~Somers's description of menopause as a deficiency that can be rebalanced with
hormones sounds like a modern take on ``Feminine Forever.''

``Hormones,'' Ms.~Somers said last week, ``are the juice of life.''

\section{Building a Baby, With Few Ground Rules}

\lettrine{U}{nable}\mycalendar{Dec.'09}{14} to have a baby of her own, Amy Kehoe became her own
general contractor to manufacture one. For Ms.~Kehoe and her husband, Scott, the idea seemed like
their best hope after years of infertility.

Working mostly over the Internet, Ms.~Kehoe handpicked the egg donor, a pre-med student at the
University of Michigan. From the Web site of California Cryobank, she chose the anonymous sperm
donor, an athletic man with a 4.0 high school grade-point average.

On another Web site, surromomsonline.com, Ms.~Kehoe found a gestational carrier who would deliver
her baby.

Finally, she hired the fertility clinic, IVF Michigan, which put together her creation last
December.

``We paid for the egg, the sperm, the in vitro fertilization,'' Ms.~Kehoe said as she showed off
baby pictures at her home near Grand Rapids, Mich. ``They wouldn't be here if it weren't for us.''

On July 28, the Kehoes announced the arrival of twins, Ethan and Bridget, at University Hospital in
Ann Arbor. Overjoyed, they took the babies home on Aug.~3 and prepared for a welcoming by their
large extended family.

A month later, a police officer supervised as the Kehoes relinquished the swaddled infants in the
driveway.

Bridget and Ethan are now in the custody of the surrogate who gave birth to them, Laschell Baker of
Ypsilanti, Mich. Ms.~Baker had obtained a court order to retrieve them after learning that Ms.~Kehoe
was being treated for mental illness.

``I couldn't see living the rest of my life worrying and wondering what had happened, or what if she
hadn't taken her medicine, or what if she relapsed,'' said Ms.~Baker, who has four children of her
own.

Now, she and her husband, Paul, plan to raise the twins.

The creation of Ethan and Bridget tested the boundaries of the field known as third-party
reproduction, in which more than two people collaborate to have a baby. Five parties were involved:
the egg donor, the sperm donor, Ms.~Baker and the Kehoes. And two separate middlemen brokered the
egg and sperm.

About 750 babies are born each year in this country through gestational surrogacy, and twice that
many surrogacies are attempted. Most are less complicated than the arrangement that resulted in the
birth of Ethan and Bridget.

But as the dispute over the Michigan twins reveals, surrogacy arrangements that go badly can have
profound implications, particularly for the children. Surrogacy is largely without regulation, with
no authority deciding who may obtain babies through surrogacy or who may serve as a surrogate,
according to interviews and court records.

Instead, surrogacy is controlled mainly by fertility doctors, who determine which arrangements are
carried out and also earn money by performing the procedures. And while some agencies that
coordinate surrogacies and some clinics that carry them out strictly adhere to guidelines, others do
not, the interviews and records show.

The lax atmosphere means that it is now essentially possible to order up a baby, creating an
emerging commercial market for surrogate babies that raises vexing ethical questions.

In some cases, parents must go through adoption proceedings to gain legal custody of the children.
But even in those situations, the normal adoption review process is upended. In surrogacy,
prospective parents with no genetic link often create their own baby first, then ask for legal
approval, potentially leaving judges with little alternative. Some states allow prebirth orders that
place the parents' names on the birth certificates without any screening.

When disputes arise after the babies are born, the outcome can vary from state to state. In
California, considered a friendly state for surrogacy, courts have upheld the validity of surrogacy
contracts, meaning that the people who hire surrogates are very likely to keep the babies if a
dispute arises.

But a statute in Michigan, where Ethan and Bridget were born, holds that surrogacy is contrary to
public policy and that surrogacy agreements are unenforceable, giving the woman who gives birth a
strong case if she decides to keep the babies.

A handful of other states have similar laws, according to an analysis by the Center for American
Progress, a liberal research group.

About 10 states have laws that allow for surrogacy but impose restrictions; several of those states
require at least one parent to have a genetic relationship to the baby. But the majority of states
are silent on surrogacy, according to the analysis. Legal uncertainty in some states means that
babies are sometimes left in limbo, their parentage left up to courts.

``When they go bad, it's so sad,'' said Mitzi Heineman, the Michigan broker who supplied Ms.~Kehoe's
donor eggs. ``You feel sorry for the baby. Who are the baby's parents?''

Four-year-old twin girls in Union City, N.J., have lived under such uncertainty. Their short lives
have included two tours in the foster care system.

New Jersey child welfare officials alleged earlier this year that the girls were neglected by
Stephen Melinger, 62, who arranged their birth almost five years ago. In July, a New Jersey judge
exonerated Mr.~Melinger of those accusations. But the Supreme Court in Indiana, where the girls were
born, recently ruled that his adoption of the twins was improperly executed and must be redone.

Fewer problems arise when the prospective parents have a genetic link to the offspring, lawyers who
have handled such cases say. Gestational surrogacy frequently involves couples who can produce their
own eggs and sperm, but in which the female partner cannot carry a baby.

Doctors say that when surrogacy arrangements go smoothly, they are very rewarding.

``It's been unbelievably satisfying seeing these families grow that otherwise wouldn't have,'' said
Dr.~James Goldfarb, director of fertility services at the Cleveland Clinic and president of the
Society for Assisted Reproductive Technology. Dr.~Goldfarb was involved in one of the world's first
gestational surrogacies in 1986. Today, the Cleveland Clinic takes part in 8 to 10 such arrangements
a year, he said.

But even less complex arrangements can lead to legal challenges. In another case in New Jersey, a
woman agreed to be a surrogate for her brother and his male partner, who donated sperm. But the
three are now playing tug of war over 3-year-old twin girls. The woman is seeking custody and a
declaration that she is the mother, even though she did not supply the eggs. Lawyers in both New
Jersey cases asked that the children's names be withheld for reasons of privacy.

The New Jersey physician who performed the procedure, Dr.~Susan Treiser, did not require
psychological screening and waived what is commonly a prerequisite for being a surrogate -- that she
must have given birth to her own child.

Ms.~Baker, the surrogate in Michigan, says the fertility clinic where she was impregnated failed to
perform psychological screening of the Kehoes, which is recommended by professional societies. Such
screening, she believes, might have prevented her from going through with the pregnancy.

Partly in hopes of standardizing the disparate laws governing surrogacy, the American Bar
Association has developed a model act for state legislatures. Judges across the country have said
they need guidance to sort out complex legal issues posed by reproductive technology. One section of
the proposal says that when prospective parents have no genetic link to the babies, surrogacies
would require preapproval by a court in a process that would include a home study.

Lawyers who handle surrogacy arrangements say those cases represent only 5 percent of surrogacy
cases, but they are the riskiest.

George J.~Annas, a bioethicist who is chairman of the health law program at Boston University, said,
``This is the main problem with commercialization, seeing children as a consumer product.''

``This is especially true when there is no genetic connection with the child,'' he said. ``It really
does treat children like commodities. Like pets.''


Brokers and Fees

It was a pet -- a pet bird -- that transformed the birth of the Melinger twins from a private
transaction into a public controversy.

Employees at Methodist Hospital in Indianapolis became alarmed when the man who had commissioned
their creation, Stephen Melinger, took his pet bird to the neonatal intensive care unit where they
were hospitalized. It was among several things that raised concerns about Mr.~Melinger's ability to
care for the two girls, according to court documents.

When Mr.~Melinger, a single man who taught elementary school in Union City, decided he wanted a
child, he enlisted the help of an agency called Surrogate Mothers in Monrovia, Ind.

The American College of Obstetricians and Gynecologists has adopted a set of guidelines for
surrogacy arrangements. Among its recommendations are that surrogacy be handled by nonprofit
agencies. Currently it is largely for-profit and can be very lucrative.

Between brokers, legal and medical expenses and surrogate fees, a successful surrogacy can cost
prospective parents \$80,000 to \$120,000. About an estimated 100 agencies advertise themselves as
surrogacy brokers.

``People can get into this business easily,'' said Charles P.~Kindregan Jr., a professor at Suffolk
University Law School who was co-chairman of the American Bar Association committee that drafted the
model legislation.

Surrogate Mothers, one of the older agencies, advertises on its Web site that it can arrange
surrogacies for under \$50,000.

On April 8, 2005, the twins intended for Mr.~Melinger were born in Indianapolis to a surrogate
mother from South Carolina. The girls were 9 weeks premature and weighed only 3 pounds each.

Steve Litz, a lawyer who runs Surrogate Mothers, filed a petition on behalf of Mr.~Melinger seeking
to adopt the twins, identified in court papers as the ``infants H.'' According to court documents,
the petition identified Mr.~Melinger as an Indiana resident, born in Indiana and employed as a
teacher.

The woman who carried the children was giving up her rights to them. She had listed Mr.~Melinger as
``father'' on the birth certificate.

The private adoption was on track to be granted, almost as a perfunctory matter. But hospital
employees became concerned by Mr.~Melinger's eccentric behavior.

On one day, he arrived at the intensive care unit carrying his pet bird, which posed a risk of
infection. Mr.~Melinger testified that his bird was not near the babies because he had stayed in the
office area. Yet on a separate visit, hospital workers said he had gone into the intensive care unit
with bird feces on his clothing.

The hospital staff was also worried about what they considered Mr.~Melinger's unrealistic plan for
taking the babies home. He hoped to make the 12-hour drive from Indianapolis to Union City in his
car, alone, with the two premature infants strapped in car seats.

Hospital workers asked Indiana's child welfare agency to investigate.

After learning of the investigation and the possibilities that there would be difficulties in
getting Mr.~Melinger's adoption approved, Mr.~Litz filed another motion on his client's behalf,
calling the children ``hard to place'' because their mother was African-American, he said, making
the girls biracial.

It is easier for out-of-state residents, like Mr.~Melinger, to adopt ``hard to place'' children. But
it was not true that the girls were biracial. The surrogate mother was African-American, but the
babies she had carried grew from eggs from a white donor. The twins were white.

It was merely one of the assertions in papers filed on Mr.~Melinger's behalf by Mr.~Litz that turned
out to be false, according to the Indiana Supreme Court.

``An earlier representation that Mr.~Melinger was a sperm donor likewise turned out not to be
true,'' the court said, so the children were not his biological offspring.

In addition, Mr.~Melinger was not born in Indiana, but New York. The Indiana residence he claimed
was a hotel room.

A lower court had criticized the submissions for ``lack of candor.''

In an e-mail message, Mr.~Litz denied that he had misled the courts. ``I have never knowingly made a
misrepresentation to a judge in my life,'' he said. Mr.~Melinger declined to be interviewed.

As the case moved forward amid a swirl of Indiana news reports, the girls were placed in foster
care. Mr.~Melinger continued his fight to adopt the girls, finally prevailing in 2006. Mr.~Melinger
returned to New Jersey with the girls, but the Indiana Department of Child Services appealed the
adoption to the Indiana Supreme Court.

The appeal was still pending in January when the girls had another encounter with the child welfare
system, according to records disclosed by Mr.~Melinger's lawyer, Anthony Carbone of Jersey City.

It began as a simple family outing to a park in January.

A woman who saw Mr.~Melinger with his children complained to the police that the girls were dirty
and inappropriately dressed for cold weather, the records show.

One was wearing a pink coat, a skirt, ankle socks and black shoes. The other was wearing pajamas, a
yellow coat and sneakers with no socks. Mr.~Melinger later said that it had been a particularly warm
winter day and that he had taken extra clothes for the girls.

The complaint prompted a review by the New Jersey Department of Youth and Family Services, which
sent a worker the next day to Mr.~Melinger's apartment in Union City.

When a caseworker arrived ``she noticed a strong smell of urine in the apartment,'' according to a
court document. Mr.~Melinger later said that the girls were not completely toilet trained and had
accidents, and that he tried to clean up after them as best he could.

The home was ``particularly dirty,'' the caseworker said, with inadequate clean clothes for the
twins. Department workers also said the children's pediatrician, Dr.~Pearl Cenon, had concerns about
their care and had considered contacting the agency. The girls were removed from Mr.~Melinger's
custody.

But in a hearing last summer, a parade of witnesses came to Mr.~Melinger's defense. They included
Dr.~Cenon, who denied being concerned about the girls and testified that Mr.~Melinger was an
excellent father.

In July, Judge Bernadette N.~DeCastro of New Jersey Superior Court ruled that the Department of
Youth and Family Services had failed to prove its claim against Mr.~Melinger. The girls had already
been returned to his custody in April.

Meanwhile, the Indiana Supreme Court had also ruled.

In a decision issued in April and reaffirmed in October, the court said the adoption of the
4-year-old twins must be repeated. Among missing elements in the original adoption, the Supreme
Court said, was a letter from New Jersey authorities stating that the placement was in the twins'
best interest.

The Indiana court said that as the case continued, the girls would be allowed to remain with
Mr.~Melinger. Frances Watson, a professor at the Indiana University School of Law in Indianapolis
who briefly served as the appointed legal representative of the children, said the case provided a
stark example of what the state's adoption laws attempted to prevent.

``You should not be able to come from out of state on some contract and order up some babies and
then go about your business,'' Ms.~Watson said.

Barriers to Adoption

On July 11, 2004, Donald W.~Robinson, a Manhattan accountant, boarded the Norwegian Dawn cruise ship
with his partner, Sean Hollingsworth.

The departure from New York was the maiden voyage of the ``Rosie'' cruises, named after the
entertainer Rosie O'Donnell, which cater to gay men and lesbians and their families.

The cruise was also the genesis of a major surrogacy dispute.

As the 1,600 passengers sailed from New York to Florida and the Bahamas, one of the speakers was
Dr.~Treiser, the fertility specialist.

The weeklong cruise was an excellent way to promote her clinic, IVF New Jersey, to an important
surrogacy niche market: gay partners who might want to become parents.

The shift from traditional surrogacy, in which women carry their own biological children after
artificial insemination, to gestational surrogacy, as well as the wide availability of donor eggs,
has opened the possibility of parenthood to a variety of people who cannot have children of their
own.

In Manhattan, the Lesbian, Gay Bisexual \& Transgender Community Centersponsors monthly seminars on
having families through surrogacy. The well-attended sessions often feature speakers with children
born through surrogacy arrangements.

In many of those cases, one of the male partners donates sperm that is used, along with a donor egg,
to impregnate a surrogate.

Many of the people who have children through surrogates would have had difficulty adopting because
of sexual orientation, marital status or age. Some foreign countries place upper age limits on
adoptive parents. And birth mothers giving up their children in the United States often hand-pick
the adoptive parents of their children.

``The default position for young birth moms tends to be a mother and a father in a stable
relationship and a white picket fence around the yard,'' said David C.~Cole, a Dallas lawyer with
Little Flower Adoptions, which also handles surrogacy arrangements.

After Dr.~Treiser's speech, Mr.~Robinson and Mr.~Hollingsworth approached her to discuss their plans
for having a child through surrogacy.

As Dr.~Treiser remembered during a deposition, they wanted to use Mr.~Hollingsworth's sperm and an
egg from Mr.~Robinson's sister, Angelia Gail Robinson, a resident of Texas. That way, both
Mr.~Robinson and Mr.~Hollingsworth would have genetic ties to the child.

But testing later revealed that Ms.~Robinson, already in her 40s, could not produce viable eggs.
Instead, the couple decided to use another egg donor. Ms.~Robinson agreed to serve as the
gestational carrier and intended to play a role in the life of the baby.

``She was going to be the doting aunt and live close by,'' Dr.~Treiser testified in a deposition.

Ms.~Robinson sold her home in Texas and went to work in her brother's Manhattan accounting office.

As the agreement proceeded, there were several things that should have waved cautionary flags.
Foremost among them was that Ms.~Robinson did not have her own children.

A previous birth provides proof that a surrogate can deliver a baby without medical complications,
fertility doctors said. And it gives assurance that the surrogate understands the biological and
emotional implications of pregnancy and childbirth.

``If a surrogate has not had a baby before, we won't use her,'' said Dr.~Goldfarb of the Cleveland
Clinic.

In an interview, Dr.~Treiser said she made an exception because Ms.~Robinson was carrying a child
for her sibling and expressed no interest in children of her own.

In a deposition, Dr.~Treiser said that she offered Ms.~Robinson psychological screening, but that it
was declined.

As her pregnancy progressed, Ms.~Robinson now says, fissures developed in her relationship with her
brother. At the same time, she says in court papers, she began to bond with the twins she was
carrying.

It turned out to be an extremely difficult pregnancy that ended on Oct.~4, 2006, when Ms.~Robinson
was rushed to the hospital suffering from pre-eclampsia, a pregnancy-induced condition that includes
extremely high blood pressure.

In March 2007, after Mr.~Robinson and Mr.~Hollingsworth had custody of the children for five months,
she filed papers for custody of the children in family court in Jersey City, where the men live. The
two were married in California in September 2008, and Mr.~Robinson has taken his spouse's name.

Ms.~Robinson has also asked to be declared the legal mother of the children. Her lawyer, Harold
J.~Cassidy of Shrewsbury, N.J., has cited the Baby M decision two decades ago, in which the New
Jersey Supreme Court upheld the maternal rights of Mary Beth Whitehead, who delivered her own
biological child for another couple after artificial insemination with the man's sperm. Mr.~Cassidy
also represented Ms.~Whitehead.

The court ruled that even though Ms.~Whitehead had agreed to a payment of \$10,000 for the service,
``There are, in a civilized society, some things that money cannot buy.''

Unlike Ms.~Whitehead, Ms.~Robinson has no genetic relationship to the girls. But as the case
continues, the family court has temporarily awarded Ms.~Robinson three days a week of parenting
time, according to records. The girls are shuttled back and forth between Ms.~Robinson's frame house
in Keansburg, N.J., and the home of their father.

Their fate may be determined by a trial as early as April.

Charges of Betrayal

Ethan and Bridget, the babies born in Michigan, are propped in their car seats in a booth at Stoney
Creek Koney Island, a diner in a strip mall in Ypsilanti. They are out for breakfast with the woman
who gave birth to them, Laschell Baker, and her husband, Paul.

The Bakers have picked out new names for them. They are calling the boy Peyton and the girl Dani. As
soon as they can spare \$320, they will file papers for legal name changes. It is a way to leave the
past behind.

Someday, though, the twins will know all about what happened. ``I'll tell them the truth,''
Ms.~Baker said.

For Ms.~Baker, 35, they are babies No.~8 and 9. In addition to her four children, she has delivered
three other surrogate babies, including another set of twins. Her previous surrogacy arrangements
went smoothly, and the children are with the family who requested them.

Ms.~Baker said she had been vilified by the national community of professional surrogates, who chat
regularly online. The Internet community is mad at her, siding mostly with the Kehoes. They even
collected money for legal funds for the Kehoes, who say they were betrayed by Ms.~Baker.

``They don't want anything to do with me,'' Ms.~Baker said. ``I'm the bad apple that ruins the name
of surrogacy.''

But, she says, this is not a story about a surrogate who changed her mind.

``My husband and I would not do something like this unless we thought it was given to us to do,'' Ms
Baker said. ``My belief is that God placed this on my heart for a reason.''

In the fall of 2007, Ms.~Baker advertised in surromomsonline.com saying she would carry a baby for a
Christian couple.

Amy Kehoe saw it and was delighted to find that Ms.~Baker lived only two hours away.

Ms.~Baker said she chose the Kehoes for the same reason. ``I picked them because I wanted a couple
that was local so they could enjoy the whole pregnancy with me,'' she said.

They traded e-mail messages and phone calls and met for dinner before agreeing to go forward with
the surrogacy.

Under Michigan's law, commercial surrogacy is punishable by five years in prison and a \$50,000
fine. Ms.~Baker said she did not carry the children for money and was reimbursed only for actual
expenses like doctor's appointments. Neither she nor the Kehoes have disclosed exactly how much that
was.

Ms.~Baker said she was the one who recommended Dr.~Jonathan Ayers of IVF Michigan for fertility
services. Dr.~Ayers was involved in her two previous surrogate pregnancies.

She has generally praised Dr.~Ayers, but says the failed arrangement might have been avoided if IVF
Michigan had required psychological screening.

A nurse at IVF Michigan said Dr.~Ayers would not comment on his clinic's policies because of patient
privacy laws.

On Tuesday, July 28, the babies were born by Caesarean section. The following Monday, in court in
Ann Arbor, Ms.~Baker said she first learned of Ms.~Kehoe's psychiatric history.

During a hearing to transfer guardianship to the Kehoes, Scott Kehoe said his wife had paranoid
schizophrenia. Ms.~Kehoe's psychiatrist listed the diagnosis as a ``psychotic disorder not otherwise
specified.'' Ms.~Kehoe takes an antipsychotic to control her symptoms.

Before her diagnosis in 2001, Ms.~Kehoe told the judge, she had self-medicated, and that was the
reason for her arrest on charges of cocaine use and driving under the influence.

Adoption experts said that mental illness was not a bar to adoption if the illness was under control
and the patient went to doctor's appointments and took medications. And Ms.~Kehoe's psychiatrist
wrote a letter saying she would be a good mother because her disease had been fully controlled for
eight years and she currently had no symptoms.

Ms.~Baker, however, said she was stunned at the disclosure of Ms.~Kehoe's mental illness, which she
believes she should have known in advance. And she became concerned that Ms.~Kehoe might relapse and
be unable to take care of the twins.

``I'm not going to be the one that's going to feel guilty if something happens,'' Ms.~Baker said.

Ms.~Kehoe said Ms.~Baker's decision made no sense in light of her doctor's statement and other
letters of strong support. ``Does she really think she knows better than a psychiatrist who has
known me for nine years?'' Ms.~Kehoe said.

Instead, she says, Ms.~Baker ``legally stole our babies from us.''

Because Michigan law states that surrogacy contracts are void and unenforceable, it was an easy
matter for Ms.~Baker to go to court and have the Kehoes' guardianship rescinded.

Last month, Amy and Scott Kehoe made a decision.

``We are stopping the fight to get our babies back,'' Ms.~Kehoe wrote in an e-mail message. ``The
reason is because of the slow court system, and because of the terrible Michigan laws. JUSTICE DOES
NOT PREVAIL in this case due to Michigan laws.''

Ms.~Kehoe still has hope, though. It is stored in a tank of liquid nitrogen at IVF Michigan. The
tank contains 20 frozen embryos made from the eggs and sperm she bought.

\section{When Woods Goes Missing, So Does the Audience}

\lettrine{A}{s}\mycalendar{Dec.'09}{14} Tiger Woods's major sponsors plot their own courses in the
enlarging wake of multiple accusations of Woods's marital infidelities, the PGA Tour and its players
must also deal with the prospect of yet another protracted absence by its biggest rainmaker.

``We need him out here because of sponsorships and just the awareness in our tour in general,''
Steve Stricker told PGATour.com.

Stricker, the No.~3-ranked golfer in the world, was describing what came to be known as the Tiger
Effect from the first day Woods joined the tour in 1996 -- no Woods, no ratings.

Sponsors have reacted quickly after Woods announced Friday that he would take an ``indefinite
break'' from the PGA Tour to work on his marriage. The global consulting giant Accenture ended its
endorsement contract with Woods on Sunday, a day after Gillette announced it would reduce his
presence in its advertising.

His presence on the tour is in question at least for the start of the 2010 season, which creates
different concerns for the PGA Tour and golf in general. Most of his fellow competitors are
prefacing or punctuating their remarks with expressions of concern for Woods's well-being, but all
are aware of the effect he has on everything from the number of people who attend tournaments to the
number of zeroes on the checks they deposit if they manage to beat him.

Although Woods is not solely responsible for the economic growth of the tour, he is given much of
the credit for the quadrupling of prize money since he joined it -- from \$70 million in 1996 to
\$278 million in 2009. Most of the larger purses directly result from higher revenue from title
sponsors, and the PGA Tour is in the midst of negotiating new deals with the sponsors of a dozen
events that will expire by the end of 2010. Therefore, uncertainty about his availability will have
a negative impact on the negotiations.

``We all know what kind of effect there is when he doesn't play in tournaments, let alone taking
time off,'' the Hall of Famer Nick Price said. ``The golf world is going to miss him.''

The main uncertainties facing the tour in 2010 depend on what the definition of ``indefinite'' turns
out to be. In 2008, Woods played six events, including the United States Open, before he had surgery
to repair his knee and missed about 10 tournaments, including two major championships.

Television ratings for golf during his absence were down by an average of 50 percent from 2007, and
down a similar level for the two events he missed from the 2009 schedule before returning to the
W.G.C.-Accenture World Match Play.

Was there an upside during the eight-month absence? There was the emergence of some new young stars.
Anthony Kim, 24, and Camilo Villegas, 27, won twice -- though neither has won again. And the veteran
Steve Stricker rediscovered his game and rose to No.~2 in the world golf rankings. And Phil
Mickelson, who regained the No.~2 spot with his win at the Tour Championship, re-established himself
as clearly the second-best player in the world and the only player besides Woods capable of moving
the needle.

Could Rory McIlroy, the 21-year-old sensation from Northern Ireland who will join the PGA Tour in
2010, become the new world star? Could Adam Scott, 29, who recently broke his slump with a win in
Australia, ditch the matinee-idol work ethic and finally emerge as the golf star he was touted to
be?

Those are large question marks attached to any upside possibility for golf without Tiger Woods.
Since no one, probably including Woods, knows how many events he will miss in 2010, any forecast is
guesswork. But should he miss six months the tournaments that would suffer most based on his playing
schedules from 2005-09 include the San Diego Open, where he started four of five seasons since 2006
at Torrey Pines Golf Course; the W.G.C.-Accenture in Marana, Ariz.; the W.G.C.-CA Championship in
Miami; the Arnold Palmer Invitational in Orlando, Fla.; the Masters in Augusta, Ga.; the Quail
Hollow Championship in Charlotte, N.C.; the Players Championship in Ponte Vedra, Fla.; and the
Memorial Tournament in Dublin, Ohio.

It is doubtful Woods would play in the United States Open in Pebble Beach, Calif., without a prior
competitive start. Should he choose to play the Open, he will probably play at the Memorial two
weeks before Pebble Beach, where he set the United States Open record with a 15-stroke victory in
2000.

If Woods skips the entire season, he will also miss the remaining two major championships -- the
British Open at St.~Andrews, and the P.G.A. Championship at Whistling Straits in Kohler, Wis. -- and
the AT\&T, for which Woods is the host, at Aronimink Country Club near Philadelphia.

``Contrary to what everyone believes, the tour will go on,'' the golfer Chris DiMarco said in an
interview with PGATour.com. ``Anytime you take the greatest competitor of all time out of the game,
though, it will miss him.

``Sometimes you think he is superhuman, and he is with some of the things he's done on the golf
course. This is about Tiger and his family, though. Golf is secondary right now. We were without him
and it was awful, but this is way worse. This is life. This is reality.''

\section{Obama Has Goal to Wrest a Deal in Climate Talks}

\lettrine{P}{resident}\mycalendar{Dec.'09}{18} Obama arrives here on Friday morning bent on applying
a combination of muscle and personal charm to secure a climate change agreement involving nearly 200
countries.

He injects himself into a multilayered negotiation that has been far more chaotic and contentious
than anticipated -- frozen by longstanding divisions between rich and poor nations and a legacy of
mistrust of the United States, which has long refused to accept any binding limits on its greenhouse
gas emissions.

The world is looking to Mr.~Obama to wrest some credible success from this process. And on Thursday,
with almost 120 heads of state and government in attendance, there were some signs that a meaningful
political deal might be at hand, including a slight shift in China's position and a pledge by the
United States to help the poorest nations cope financially with global warming.

But top negotiators here said that the talks could also prove a humiliating failure, because China
and the United States, the world's two largest emitters, remain deeply divided over a number of
difficult problems.

Mr.~Obama is putting a measure of his and the nation's prestige on the line by entering a debate
with so much still unresolved. It was only 11 weeks ago that he left this same city empty-handed
after pleading for Chicago to be selected as the site of the 2016 Olympics.

But the maneuvering and brinksmanship that have characterized the final week of the talks are also a
sign of their seriousness; never before have global leaders come so close to a meaningful agreement
to reduce the greenhouse gases linked to warming the planet.

The administration provided the talks with a palpable boost on Thursday when Secretary of State
Hillary Rodham Clinton declared that the United States would contribute its share of \$100 billion a
year in long-term financing to help poor nations adapt to climate change.

The administration had remained silent for months as other major economic powers came forward with
similar or even more generous financial proposals.

A senior Obama administration official said the announcement was carefully timed to resuscitate the
talks before Mr.~Obama's arrival.

``It's a negotiation,'' he explained.

But Mrs.~Clinton's offer came with two significant conditions. First, the 192 nations involved in
the talks here must reach a comprehensive political agreement that takes effect immediately. Second,
and more critically, all nations must agree to some form of verification -- she repeatedly used the
term ``transparency'' -- to ensure they are meeting their environmental promises.

China, the world's largest producer of greenhouse gases, has brought the talks to a virtual
standstill all week over this issue, which its leaders claim to be an affront to national
sovereignty.

But the Chinese balkiness on the issue is matched in large measure by Mr.~Obama's own constraints.
The Senate has not yet acted on a climate bill that the president needs to make good on his promises
of emissions reductions and on the financial support that he has now promised the rest of the world.

``The president and his team have been doing everything possible to create a deal that is fair to
the U.S.~and facilitates international agreement,'' said Paul Bledsoe of the National Commission on
Energy Policy, a bipartisan advisory group. ``But if the Chinese will not accept monitoring of
emissions, then a deal is not worth doing.''

China appeared to crack the door a bit toward a system of reporting its emissions and its actions to
reduce them on Thursday. He Yafei, the vice foreign minister, repeated China's opposition to any
intrusive international monitoring regime in a news conference on Thursday. But he said his country
would consider voluntary ``international exchanges'' of information on its climate programs.

Administration officials here were not ready to publicly declare any breakthroughs in their talks
with China and other nations on verification measures.

``We're making progress on all the outstanding issues with the Chinese,'' a senior administration
official said in a conference call with reporters. ``But it's still a very challenging task. It's
impossible to anticipate where this will end.''

Detailed talks were still continuing Thursday afternoon among at least 16 different subgroups on a
range of issues, including finance for developing countries and mechanisms to preserve forests --
even on whether to keep the Kyoto Protocol or commit to rolling that treaty into a new agreement.

A contingent of Democratic leaders in Congress flew in Thursday to highlight support for the
administration's position at the summit meeting. House Speaker Nancy Pelosi and a number of House
committee chairmen declared that the climate change bill they passed in June would provide a large
share of the money needed to redeem the administration's promises. They met privately with a number
of delegations to urge them to encourage China and a number of other nations in the so-called Group
of 77 developing countries to drop their objections to the proposed agreement.

``I believe the leaders appreciate the magnitude of the challenge they're facing and will come to an
agreement by the end of the week,'' said Representative Edward J.~Markey, Democrat of Massachusetts.
``And the \$100 billion really does help with the concerns of the G-77. And now we're left with
China.''

Republicans mounted a counteroffensive. Senator James M.~Inhofe, Republican of Oklahoma and
Congress's most vocal climate change skeptic, showed up in the press area of the Bella Center early
on Thursday to deliver what he called a reality check to the proceedings here.

``There is going to be no cap and trade or binding legislation in the United States,'' he said.
``It's dead. It's not going to happen.''

Somewhat more substantively, in Washington, a group of House Republicans said they planned to
introduce a resolution formally disapproving of the Environmental Protection Agency's finding that
greenhouse gases endangered public health and safety, a step that could lead to economy-wide
regulation of such emissions.

The Republicans said the finding would lead to job losses and take money out of the pockets of
consumers ``so that radical environmentalists can wage a war against nature.''

The resolution, if it passes, will not have the force of law.

The \$100 billion figure proposed by Mrs.~Clinton is similar to estimates by the European Union of
the needed contributions, although the amount is below the \$150 billion or so that experts at the
European Union have pushed for.

Mrs.~Clinton said the money would be a mix of public and private funds, including ``alternative
sources of finance,'' but declined to explain what that might mean. Nor did she say what the
American share of the fund would be, although typically in such multilateral financial efforts the
United States contributes about 20 percent. She said the money should contain billions of dollars to
slow deforestation, which contributes to concentrations of carbon dioxide in the atmosphere.

Some delegates and observers said that Mrs.~Clinton's announcement might weaken the solidarity of
the G-77 developing countries with China.

Apparently attempting to appeal to these divisions, the United States Congressional delegation
immediately went into a meeting with Indian officials, focusing on verification of emissions
reductions programs, Ms.~Pelosi said.

Prime Minister Wen Jiabao, meanwhile, appeared focused on consolidating China's G-77 ties, meeting
with representatives of Brazil, Bangladesh, Trinidad, Ethiopia and Sudan since his arrival on
Wednesday night.

The head of Brazil's delegation, Dilma Rousseff, said her country, too, objected to ``intrusive
verification,'' and felt strongly that the Kyoto track should be maintained, for now anyway.

``No one is going to give up a bird in the hand for a bird that hasn't been introduced to us yet,''
she said.

Fander Falconí, Ecuador's foreign minister, was skeptical of the American financial offer.

``What we really need are firm mechanisms to reduce emissions from industrialized countries,'' he
said. ``Financial mechanisms are useful, but not central, not a solution.''

\section{At Japanese Cliffs, a Campaign to Combat Suicide}

\lettrine{T}{he}\mycalendar{Dec.'09}{18} towering cliffs of Tojimbo, with their sheer drops into the
raging, green Sea of Japan, are a top tourist destination, but Yukio Shige had no interest in the
rugged scenery. Instead, he walked along the rocky crags searching for something else: a lone human
figure, usually sitting hunched at the edge of the precipice.

That is one of the telltale signs in people drawn here by Tojimbo's other, less glorious,
distinction as one of the best known places to kill oneself in Japan, one of the world's most
suicide-prone nations. Mr.~Shige, a 65-year-old former policeman, has spent his five years since
retirement on a mission to stop those who come here from jumping.

His efforts have helped draw attention to the grim fact that Japan's suicide rate is again on the
rise. Police figures show that the number of suicides this year could approach the country's record
high of 34,427, reached in 2003, almost 95 suicides a day. The World Health Organization says that
people in Japan are now almost three times as likely to kill themselves as are Americans.

Mr.~Shige and a group of volunteers he put together have saved 222 people so far, a tally that has
made Mr.~Shige a national figure in a country that often seems apathetic about its high rate of
self-destruction. But he has also met with criticism from a conformist society that can look dimly
on people who draw attention by engaging in activism, even of the most humanitarian kind.

``In Japan, we say the nail that sticks up gets hammered down,'' said Mr.~Shige, who says he started
the patrols after he grew angry at inaction by local authorities. ``But I'll keep sticking up. I
tell them, hit me if you can!''

In part, public health experts blame Japan's romanticized image of suicide as an honorable escape,
going back to ritual self-disembowelment by medieval samurai, for the high suicide rate. But the
main cause, they say, is the nation's long economic decline. Suicides first surged to their recent
high levels in 1998, when traditional lifetime employment guarantees began to vanish, and they have
remained high as salaries and job security continued to erode.

The situation has worsened during the recent global financial crisis, which is driving this year's
increase, experts say. While Prime Minister Yukio Hatoyama, in his first policy speech in October,
referred to Japan's suicide rate in calling for ``mutual support'' among Japanese, experts say the
government's limited steps to deal with suicide have made little difference.

While preventing suicides is a universally difficult task, it is particularly challenging in Japan.
Depression remains a taboo topic here, making it hard for those most at risk to seek the help of
family and friends. Many Japanese view suicide as an issue of private choice rather than public
health, and there are few efforts to highlight the problem.

``Americans raise awareness with grass-roots action, but Japanese just wait for the government to
take care of them,'' said Yoshitomo Takahashi, a professor of behavioral science who researches
suicide at the National Defense Medical College in Tokorozawa, Japan.

Officials in Sakai, the small city in Fukui Prefecture, where Tojimbo is located, have installed
outdoor lighting at the cliffs along with two pay phones and plenty of the 10-yen coins needed to
dial up the national suicide hot line.

Nevertheless, city officials call this the grimmest year on record, with the police saying they know
of more than 140 people who came here intending to commit suicide, twice the average in recent
years. Most of them were stopped by the police or nearby tourists, or decided not to jump for other
reasons, the police say.

The police figure does not include the 54 people this year whom Mr.~Shige says he and his group have
stopped. City officials credit Mr.~Shige with helping keep the number of deaths here down to 13 so
far this year, about the same as the 15 suicides last year.

Mr.~Shige says his approach to stopping suicides is quite simple: when he finds a likely person, he
walks up and gently begins a conversation. The person, usually a man, quickly breaks down in tears,
happy to find someone to listen to his problems.

``They are just sitting there, alone, hoping someone will talk to them,'' Mr.~Shige said.

As an officer stationed at Tojimbo at the end of his 42-year career, he said he was appalled by all
the bodies he had to pluck out of the sea. He said he once stopped an elderly couple from Tokyo from
jumping and turned them over to city officials who he said gave them money and told them to buy a
ticket to the next town. Days later he received a letter from the couple, mailed just before they
committed suicide in a neighboring prefecture.

``The authorities' coldness outraged me,'' said Mr.~Shige, whose cellphone rings to the tune of
``Amazing Grace,'' though he is not religious. He now has 77 volunteers patrolling the cliffs and
providing food, lodging and assistance in finding work to those it helps. He said they tried to
patrol two or three times a day.

On a recent afternoon, Mr.~Shige checked three of the most popular sites for jumpers -- all with
drops of at least 70 feet. He said the loners were easy to spot because most visitors moved in
groups behind flag-waving guides. Speaking through bullhorns, the guides loudly describe the morbid
fame of the cliffs, which were named for an evil Buddhist monk who was said to have fallen to his
death there.

One of those whom Mr.~Shige stopped was Yutaka Yamaoka, 29, a factory worker who tried to commit
suicide last year after being laid off. Mr.~Yamaoka visited Mr.~Shige's tiny office by the cliffs on
a recent day to thank him and tell him that he had found a job.

When Mr.~Shige found him last year, Mr.~Yamaoka said, he was sitting silently near the cliffs
clutching his knees. He said Mr.~Shige spoke with him for two hours, then allowed him to stay in an
apartment rent free for a month until he felt better.

``I felt saved. I felt I could live,'' recalled Mr.~Yamaoka, who spoke haltingly in a barely audible
voice. ``My feelings of panic and unease just built up. I had no one to talk with.''

Mr.~Shige's efforts have stirred local resentment, particularly from a local tourist association
that says his activities are bad for business. But Mr.~Shige is not easily deterred.

``I will continue until the government finally gets its act together and takes over,'' he said. ``I
can't let their inaction cost another precious life.''

\section{Brain Damage Found in Hockey Player}

\lettrine{A}{}\mycalendar{Dec.'09}{18} deceased professional hockey player has been found to have
had brain damage associated with repeated head trauma, connecting hockey for the first time to
health risks linked to boxers and, most recently, football players.

Reggie Fleming, a defenseman and left wing known for fighting as much as scoring in a long career
from 1959 to 1974, was found by Boston University researchers to have chronic traumatic
encephalopathy, a neurodegenerative disease known to cause cognitive decline, behavioral
abnormalities and ultimately dementia. Fleming died in July at age 73 and was the first hockey
player known to have been tested for the disease, known as C.T.E.

Fleming's having had C.T.E. will stoke further debate in the National Hockey League this season over
rules to decrease player concussions. Eleven former National Football League players have been found
to have the same disease, catalyzing questions of football's long-term health risks and miring the
N.F.L. in three years of controversy over its handling of brain injuries.

``Boxing we've known for a long time, football we've recently become aware of -- now hockey,'' said
Dr.~Ann McKee, a neuropathologist at Boston University and the Bedford Veterans Administration
Medical Center, who similarly diagnosed C.T.E. in several former N.F.L. players. ``Repetitive head
injuries can have very serious long-term consequences, regardless of how you get them.''

Bill Daly, the deputy commissioner of the N.H.L., said in a statement that the league would have no
comment on McKee's diagnosis until it ``had a chance to read and digest it.''

``This very preliminary finding will set forth a hypothesis for further scientific research,'' said
Jonathan Weatherdon, a spokesman for the N.H.L. Players Association.

The N.H.L. and the N.H.L. Players Association jointly administer the league's protocol for players
returning from concussions, which has been in place since 1997-98. The program was the first in pro
sports to mandate independent baseline and postconcussion neuropsychological testing, as well as
clearance from independent doctors before a player can return to game action. A similar protocol was
adopted by the N.F.L. only last month.

In recent years the N.H.L. has experienced several incidents of head injuries to its players. But
unlike in Fleming's day, most of today's concussions are caused by checks to the head delivered at
high speed with players' shoulder pads. The league's general managers have consistently declined to
take steps curbing such hits, which are legal under N.H.L. rules.

However, last month they agreed to form a committee to study the issue, and a recommendation for new
rules designed to cut down on hits that resulted in concussions was expected before the end of the
season.

Checks to the head are barred in the Ontario Hockey League, the top developmental league for
aspiring professional players, and at international tournaments like the Olympics. Any blow to the
head in those games results in an automatic penalty.

At the N.H.L. Board of Governors meeting at Pebble Beach, Calif., this week, the Anaheim Ducks'
general manager, Bob Murray, distilled the cautious stance being taken by league executives.

``You've got to be careful what you do when you talk about rule changes,'' Murray said. ``Hitting is
part of our game, and you don't want to change the fundamental nature of the game.''

Fleming relished his persona as one of the most bruising players of his era, in which fights
featured bare fists punching helmetless skulls. (During his seasons with the Rangers from 1966
through 1970, he spent so much time in the Madison Square Garden penalty box that he once joked, ``I
got my mail delivered there.'') Fleming spent 12 full seasons in the N.H.L. from 1960 to 1971 --
playing 749 games, scoring 108 goals and compiling 1,468 penalty minutes -- before finishing his
career with two seasons in the World Hockey Association.

Fleming's son, Chris, said in an interview this week that his father never knew how many concussions
he had sustained playing hockey. That number, as well as if or how they were treated, will most
likely never be known because such records were not kept in his era, if they were considered
injuries at all.

Chris Fleming said that his father went through decades of emotional problems after retiring. He was
found to be manic depressive in his early 40s, drank excessively during that period, and exhibited
striking short-term memory problems in his late 50s. Chris Fleming said that his father had trouble
controlling his temper his entire life -- that was one of the reasons for his hockey success -- but
that it worsened post-retirement.

``He'd get in fistfights with people on the street, and kicked out of the racetrack,'' Chris Fleming
said. ``It just didn't make sense, someone snapping so quickly and violently. Other hockey players
didn't stay like that. But he didn't know how to react.''

In the 1990s, Chris Fleming said his father would ask him a question, hear the answer, and then five
minutes later ask it again. This could happen three or four times as tension mounted.

``I just told you -- you didn't listen!'' Chris recalled saying.

``You didn't listen!'' the father snapped back, having already forgotten the previous exchanges.

Chris Fleming posted several videos on YouTube in which his father reminisced about his life from
his hospital bed.

C.T.E. can be diagnosed only post-mortem, with brain tissue being subjected to special staining
techniques for neurofibrillary tangles and protein deposits. Experts believe those structural
abnormalities bring on significant cognitive and behavioral disturbances for years and sometimes
decades at the end of the patient's life.

Some former N.H.L. players have expressed concern about the repeated blows to the head they took
during their careers.

``My memory has gotten worse the last 10 years or so,'' said Ron Duguay, who played helmetless for
the Rangers and three other N.H.L. teams from 1977-78 through 1988-89 and who is taking a series of
neurological tests as a result of his concerns. He agreed to share the results of his tests in an
interview last month.

``I fail a lot of the memory tests,'' said Duguay, 52. ``I took a lot of hits to the head with no
helmet, and if you've taken hits to the head you've suffered damage. Now I'm seeing what I can do to
keep my health.

``I had fun as a New York Ranger,'' said Duguay, who was known as a bon vivant during his playing
days. ``People say you should write a book, and I would, but I can't remember.''

McKee said that because C.T.E. symptoms resembled those of Alzheimer's disease -- although they
appear sooner, as early as the person's 30s, and last longer -- many athletes currently being
treated for Alzheimer's might have been misdiagnosed. She added that patients with C.T.E. appeared
to show considerably more aggression and anger-management problems than patients with Alzheimer's
did, and could therefore be misunderstood as psychiatric.

``This is not a psychiatric disorder or a postcareer adjustment issue -- the individual is
struggling with a disease that is short-circuiting his nerve connections inside the brain,'' McKee
said. ``That is compromising his ability to deal with the world as he used to. I can't imagine the
chaos that these individuals are suffering.''

The Boston University group is collaborating with the Sports Legacy Institute to collect brain
tissue of athletes and nonathletes to explore and better understand the effects of sport-related
concussion. A dozen hockey players are among 250 current and retired athletes who have pledged to
donate their brains to the study.

\section{Panel Backs Bernanke for Second Term}

\lettrine{T}{he}\mycalendar{Dec.'09}{18} Senate Banking Committee voted Thursday to approve the
nomination of Ben S.~Bernanke for a second four-year term as chairman of the Federal Reserve.

The 16-to-7 vote, which came after flashes of populist anger erupted amid a sluggish and uneven
economic recovery, was not unexpected.

President Obama has said the nomination would provide crucial continuity as the country works
through its gravest economic challenges in decades.

But it represented a retreat from Mr.~Bernanke's near-unanimous approval by the committee four years
ago. Neither party's members voted as a bloc, and even some of Mr.~Bernanke's supporters said they
harbored reservations and might reconsider when the vote went to the Senate.

Even as Mr.~Bernanke is widely credited for helping stabilize markets and averting economic
calamity, outspoken critics, not all of them Republicans, have blamed him for enormously costly
initiatives that have bolstered some Wall Street financial firms while leaving ordinary Americans
staring at persistent double-digit jobless rates.

Mr.~Bernanke, 56, has also been criticized for not doing more to prevent the financial crisis. He
was originally named by President Bush in 2005 to succeed Alan Greenspan as the Fed chairman.

Senator Richard C.~Shelby of Alabama, the committee's ranking Republican, said that Mr.~Bernanke,
first as a member of the board of governors and then as chairman, was complicit in creating an
atmosphere that encouraged risk-taking and contributed to the disastrous housing price bubble --
signs of which, Mr.~Shelby said, Mr.~Bernanke had missed.

``Our trust and confidence were misplaced,'' Mr.~Shelby said, shortly before voting against the
nominee.

Senator Jeff Merkley, Democrat of Oregon, also opposed the nomination, saying that even if
Mr.~Bernanke ``has shown himself to be quite adroit with the fire hose'' since the crisis erupted,
he bore some responsibility for its outbreak, and thus for the loss of jobs and depleted savings
accounts.

Some of the critics' complaints had merit, said the committee chairman, Senator Christopher J.~Dodd
of Connecticut. Still, Mr.~Dodd added, ``Chairman Bernanke must also receive credit for the critical
role he played in the events of last fall.'' Had he and others not acted, Mr.~Dodd said, the nation
probably would have faced ``utter economic catastrophe -- and I believe nothing short of that was at
risk.''

Instead, Mr.~Dodd added, ``I believe better days do lie ahead.''

A Republican, Senator Judd Gregg of New Hampshire, offered a strong endorsement of the nominee.
While ``mistakes were made'' under his purview, Mr.~Gregg said, Mr.~Bernanke's swift reaction to the
financial crisis had proved crucial. ``I tell you, it worked,'' he said. ``It's that simple.''

But Senator Bernard Sanders, independent of Vermont, strongly opposes Mr.~Bernanke's nomination
because of his role in the bailouts and has put a hold on the nomination. That means it will require
60 votes, not a simple majority.

Although his confirmation appears likely, a number of no votes are expected from members of both
parties when the full Senate votes and if Mr.~Bernanke is not confirmed by Jan.~31, when his term
ends, then the Fed's vice chairman, Donald L.~Kohn, would temporarily take over.

President Obama, in announcing his nomination of Mr.~Bernanke in August, said he had led the Fed
through ``one of the worst financial crises that this nation and this world have ever faced'' and,
because of his background, his expertise on the Depression and his unflappable personality, was
ideally suited to help lead the recovery.

During his four years, Mr.~Bernanke has emerged as a favorite target of liberals and conservatives
alike. Critics have tried to portray him as cozy with Wall Street elites and slow to react to signs
that the financial system was on the brink of collapse.

Mr.~Bernanke, a former professor at Princeton, was initially criticized as too academic, but he has
since spearheaded some of the most interventionist actions by the Fed in its history.

At a hearing earlier this month, Mr.~Bernanke acknowledged mistakes by the Fed in the runup to the
financial crisis. But Mr.~Bernanke adamantly defended the Fed's response to the crisis, particularly
its efforts to crack down on banks, subprime mortgages, credit card fees and executive compensation.

Mr.~Bernanke has often been forced to navigate the line between doing too much and too little, and
he has faced the constant danger of pushing the economy into a deeper downturn if the central bank's
extraordinary measures were withdrawn too quickly.

At the hearing earlier this month, Mr.~Bernanke bristled at proposals to strip the Fed of its power,
saying the central bank had unique expertise and knew the inner workings of major financial
institutions better than any entity.

Adding to the intensity surrounding Mr.~Bernanke's appointment, Time Magazine named him this week as
its Person of the Year, citing his role in leading the economic recovery. Mr.~Bernanke's visage
appears on the magazine's cover, superimposed on a dollar.

When the Senate Banking Committee took a voice vote on Mr.~Bernanke's first nomination as Fed
chairman, only one senator opposed the nominee: Senator Jim Bunning, Republican of Kentucky.

On Thursday, Mr.~Bunning, who has one of the most conservative voting records in the Senate,
remained a thorn in the chairman's side.

In delivering his criticism, he cited a blogger's remark that giving Mr.~Bernanke the Person of the
Year citation was like rewarding the captain of the Titanic.

\section{Liberal Revolt on Health Care Stings White House}

\lettrine{I}{n}\mycalendar{Dec.'09}{18} the great health care debate of 2009, President Obama has
cast himself as a cold-eyed pragmatist, willing to compromise in exchange for votes. Now ideology --
an uprising on the Democratic left -- is smacking the pragmatic president in the face.

Stung by the intense White House effort to court the votes of moderate holdouts like Senator Joseph
I.~Lieberman, independent of Connecticut, and Senator Ben Nelson, Democrat of Nebraska, liberals are
signaling that they have compromised enough. Grass-roots groups are balking, liberal commentators
are becoming more critical of the president, some unions are threatening to withhold support and
Howard Dean, the former Democratic Party chief, is urging the Senate to kill its health bill.

The White House scrambled Thursday to tamp down the revolt, which has been simmering for weeks but
boiled over when the Senate Democratic leadership, bowing to Mr.~Lieberman, scrapped language
allowing people as young as 55 to buy into Medicare.

David Axelrod, Mr.~Obama's senior adviser, began the day by calling in to MSNBC to urge the party to
hold together, warning of a ``tragic outcome'' if Democrats failed to pass a bill that the White
House says would expand health coverage while reining in costs.

``I don't think you want this moment to pass,'' Mr.~Axelrod said. ``It will not come back again.''

Former President Bill Clinton echoed the theme, issuing a statement in which he urged fellow
Democrats to ``take it from someone who knows: these chances don't come around every day.''
Mr.~Clinton's former chief of staff, John Podesta, a close Obama ally, pushed back against Mr.~Dean
on the Internet and on television.

The backlash is building as Mr.~Obama is under increasing pressure from his party's left flank on
issues ranging from the economy to Afghanistan. This week alone, the House speaker, Nancy Pelosi,
declared that she would not insist that her caucus vote to finance Mr.~Obama's troop buildup, while
Senator Bernard Sanders, the Vermont independent and self-avowed democratic socialist, chastised the
president for reappointing Ben S.~Bernanke as chairman of the Federal Reserve.

The left's disenchantment with Mr.~Obama on health care harks back to his decision, before he became
president, not to try to replace the current private insurance system with a single-payer,
government-run ``Medicare for all'' system. Throughout the year, at town-hall-style meetings and
other public appearances, he has been dogged by advocates who have complained that he has sold out.

Instead, the president proposed what he called the public option, a government-run insurance plan
that would coexist with private plans on a new health insurance exchange. For many on the left who
are also dissatisfied with a deal the White House cut with the pharmaceutical industry, the public
option was itself a compromise.

Now that the Senate Democratic leadership has stripped the last vestige of the public option -- the
Medicare buy-in provision -- from its bill, progressives are feeling doubly betrayed.

``It's time for the president to get his hands dirty,'' Representative Anthony Weiner, Democrat of
New York, said in a statement this week. ``Some of us have compromised our compromised compromise.
We need the president to stand up for the values our party shares.''

The passion is most intense in the House, but it appears to be spilling over to the Senate. In an
interview on Thursday, Mr.~Sanders, another advocate of a single-payer system, said he was not
certain how he would vote on the bill, though Democratic leaders have been assuming he would back it
in the end.

``I don't sleep well,'' Mr.~Sanders said. ``I am struggling with this issue very hard, trying to
sort out what is positive in this bill, what is negative in the bill, what it means for our country
if there is no health insurance legislation, when we will come back to it.''

The senator added, ``And I have to combine that with the fact that I absolutely know that the
insurance companies and the drug companies will be laughing all the way to the bank the day after
this is passed.''

The White House insists that is not the case. Aides to Mr.~Obama argue that the bill includes a
variety of consumer protections -- like ending lifetime coverage caps and the practice of refusing
patients with pre-existing conditions -- that will make insurers unhappy.

``People need to put this in perspective,'' Dan Pfeiffer, Mr.~Obama's communications director, said
in an interview on Thursday. ``Two years ago, the Democratic Party would have done anything for the
opportunity to pass a health care reform like this. Let's realize how far we've come, and how close
we are to making history.''

\section{At a Mine's Bottom, Hints of Dark Matter}

\lettrine{A}{n}\mycalendar{Dec.'09}{18} international team of physicists working in the bottom of an
old iron mine in Minnesota said Thursday that they might have registered the first faint hints of a
ghostly sea of subatomic particles known as dark matter long thought to permeate the cosmos.

The particles showed as two tiny pulses of heat deposited over the course of two years in chunks of
germanium and silicon that had been cooled to a temperature near absolute zero. But, the scientists
said, there was more than a 20 percent chance that the pulses were caused by fluctuations in the
background radioactivity of their cavern, so the results were tantalizing, but not definitive.

Gordon Kane, a physicist from the University of Michigan, called the results ``inconclusive,
sadly,'' adding, ``It seems likely it is dark matter detection, but no proof.''

Dr.~Kane said results from bigger and thus more sensitive experiments would be available in a couple
of months.

The team, known as the Cryogenic Dark Matter Search, announced its results in a pair of simultaneous
talks by Jodi Cooley from Southern Methodist University at the SLAC National Acceleratory Laboratory
in California and by Lauren Hsu of the Fermi National Accelerator Laboratory in Illinois at
Fermilab, and they say they plan to post a paper on the Internet.

The stakes for astronomy and physics could hardly be greater. If the particles are confirmed by
tests at other detectors, it would mean that, after more than half a century of speculation,
astronomers are zeroing in on the identity of the invisible material that accounts for 25 percent of
the universe and determines the architecture of the visible universe.

Confirmation of the particles would also constitute the first evidence for a new feature of nature,
called supersymmetry, that physicists have been seeking as avidly as the astronomers have been
seeking dark matter. It is central to theoretical efforts like string theory, which unify all of the
forces of nature into one mathematical expression.

The report ended weeks of speculation on physics blogs and in laboratory cafeterias around the
world. At the Kavli Institute for Theoretical Physics in Santa Barbara, Calif., where dark matter
experts who had gathered for a two-week workshop watched the talks on the Web, Dr.~Kane, who was
present, described the mood at the workshop as ``a high level of serious hysteria.''

Dark matter became a serious issue in the 1970s, when Vera Rubin of the Carnegie Institution of
Washington and her colleagues charted the rotation speeds of galaxies and found that they seemed to
be enveloped in halos of dark matter, then called missing mass.

A wide range of astrophysical and cosmological measurements have subsequently converged on an
intimidating recipe for the cosmos of 4 percent atoms, 25 percent dark matter and 70 percent a
mysterious energy that has been called dark energy and has nothing to do with the news on Thursday.

The cryogenic experiment is nearly half a mile underground in an old iron mine in Soudan, Minn., to
shield it from cosmic rays. It consists of a stack of germanium and silicon detectors, cooled to
one-hundredth of a degree Kelvin. When a particle hits one of the detectors, it produces an
electrical charge and deposits a small bit of energy in the form of heat, each of which are
independently measured.

By comparing the amounts of charge and heat left behind, the collaboration's physicists can tell
so-called wimps from more mundane particles like neutrons, which are expected to flood the
underground chamber from radioactivity in the rocks around it.

The team is planning a larger detector, called SuperCDMS. In the meantime, Elena Aprile of Columbia,
who was also present in Santa Barbara, said the results would be tested soon by her own detector,
called Xenon, filled with liquid xenon, which just began working this fall under the Alps in Italy.

``All eyes will be on Xenon,'' she said in an interview a few days before, explaining that her
detector, which is bigger, should see more events, adding, ``Otherwise there will be a big clash.''

\section{Mobile Phones Become Essential Tool for Holiday Shopping}

\lettrine{T}{he}\mycalendar{Dec.'09}{18} mobile phone is quickly becoming Santa's biggest helper.

Powerful software applications for devices like the Apple iPhone are making it easy for
bargain-hunting consumers to see if another retailer is offering a better deal on a big-screen HDTV
or pair of shoes and to use it to haggle at the cash register.

Online retailers are revamping the mobile versions of their sites so consumers can make purchases
without tedious typing. And offline retailers, battling for every last dollar, are sending cellphone
users electronic coupons\cite{coupon} to lure them away from competitors.

One in five shoppers said they intended to use their cellphones to shop this holiday season,
according to an annual survey by Deloitte, the accounting and consulting firm. Of those, 45 percent
said they would use their phone to research prices, 32 percent said they would use it to find
coupons or read reviews and 25 percent said they would make purchases from their phones.

``We are at the cusp of this technology really driving a lot of activity during the shopping
season,'' said Stacy Janiak, United States retail practice leader at Deloitte. ``It is both an
opportunity and a challenge for a retailer, because you can have a consumer who can cross-shop your
store with other bricks-and-mortar stores or online, all from the convenience of your aisle.''

Heather Reed, a mother in Cypress, Tex., is one of those mobile power-shoppers. She uses several
apps on her Samsung Moment phone to whittle down her spending. She was recently considering a
\$29.99 Hot Wheels video game for her son at Wal-Mart. With a quick scan of the bar code, an
application called ShopSavvy found it at Target, just across the freeway, for \$19.99. Another app
from MyCoupons.com provided a Target coupon that sliced off \$10 more.

``It went from \$29.99 to \$9.99, all in five minutes, no searching the Internet or spending hours
trying to find a deal or a coupon,'' she said. ``It's all right there in your hand.''

Of course, mobile shopping technology is still somewhat clunky, between erratic Internet
connections, outdated pricing data and balky product scanners.

But smarter phones and a heightened bargain-consciousness among consumers are spurring a level of
innovation in e-commerce last seen during the height of the dot-com boom a decade ago.

In addition to ShopSavvy, mobile apps from RedLaser, TheFind, ShopStyle and PriceGrabber.com allow
customers to compare prices across a range of retailers. Retrevo, an electronics review site, has a
service called RetrevoQ that lets users send a text or Twitter message with the name of the product
they are considering and get an immediate response with a recommendation of whether to buy it and a
range of online prices.

While searching for prices is easy, buying from the screen of a cellphone is more difficult. It
generally involves clicking to the retailer's Web site, which is often not customized for a mobile
phone's tiny screen, and then entering shipping and billing information using the tiny keyboard.

To fix that, some retailers are building sites and applications specifically for cellphones. The
iPhone app for the Tommy Hilfiger online store, for instance, shows select products based on what
shoppers are looking for so they do not have to scroll through pages of clothes. Those who are
registered on the Web site need only enter their e-mail address and password to check out.

``Retailers need to realize that if you give people a way to make it easy, people will shop on their
phones,'' said Kelly O'Neill, product marketing director for ATG, which provides e-commerce
technology to retailers and built Tommy Hilfiger's app.

EBay's iPhone app sends people notifications if they are outbid in an auction and lets people check
out with just a few clicks if they have a PayPal account. Mobile shoppers will spend \$500 million
on eBay this year, the company said.

By improving ease of use, savvy online retailers are snatching sales from bricks-and-mortar ones.
Matthew Tractenberg, for example, was recently shopping in a Silicon Valley bookstore, where he
picked out five books for a total of \$80. Before taking them to the counter, he typed the titles
into the Amazon app on his BlackBerry Curve. Amazon had the books for \$50 and would not charge
sales tax or shipping. He placed the order on the spot and left his small pile of books in the
store.

``It's almost easier than doing it on a computer,'' Mr.~Tractenberg said.

Offline retailers are feeling the pain. Armed with competitive price information, shoppers are
haggling as never before.

Although most stores refuse to match prices, especially from Web retailers, it is difficult to
simply allow a customer brandishing a lower price to walk out the door.

Best Buy, for example, officially says it will not match prices of online electronics retailers and
will match offline prices only if the customer brings in an ad or receipt. But several ShopSavvy
users report having luck getting individual stores to match prices they find using the app.

Pacific Sunwear, a clothing and accessories retailer, said it would match lower prices found in
stores or online. According to Chad Petrillo, a clerk at the chain's San Francisco store, more
people have been showing him competing prices on their phones, most often for shoes. The store will
honor them after calling the other store to verify the price, he said.

For most shoppers, price is only one factor, to be weighed against the time it takes to drive to
another store or wait for a Web site to ship an item. That could be a boon for offline stores,
according to Ron Levi, vice president of products at TheFind, a shopping comparison Web site. ``Your
proximity to that retailer gives them an advantage,'' he said. ``It's theirs to lose.''

Michael Robison, a Coast Guard petty officer from Guernewood Park, Calif., routinely uses ShopSavvy
to check prices, but that doesn't mean he always goes with the lowest one. He just bought a
Victorinox laptop case for \$45 at Macy's, even though it was \$30 at eBags. For that amount of
money, ``I would much rather walk out with it than wait,'' he said.

Another problem with the mobile apps is accuracy. When Mr.~Robison scanned a Nintendo hand-held
gaming device at Radio Shack recently, ShopSavvy told him he could get it for \$110 online instead
of paying \$170 at the store. When he got home, he discovered that online bargain was for a used
machine.

Aware of the power of mobile phones, some offline retailers are using the technology to fight back.

If someone standing in one store scans a product with ShopSavvy, for example, a retailer down the
street could deliver the shopper a coupon for the same item. A major retailer is already doing that
in a few test cities, including Seattle, said Alexander Muse, co-founder of Big in Japan, the
start-up that created ShopSavvy.

Other applications, including Yowza, use the GPS location information in cellphones to send shoppers
coupons for stores within walking distance of where they're standing.

``This empowers consumers to make a smart decision,'' Mr.~Muse said. ``Already, retailers are
starting to figure out, 'I need to be in this game.' ''

\section{China Imposes New Internet Controls}

\lettrine{C}{hina}\mycalendar{Dec.'09}{18}'s government censors have taken fresh aim at the
Internet, rolling out new measures that limit its citizens' ability to set up personal Web sites and
to view hundreds of Web sites offering films, video games and other forms of entertainment.

The authorities say the stricter controls are intended to protect children from pornography; to
limit the piracy of films, music, and television shows; and to make it hard to perpetuate Internet
scams. But the measures also appear devised to enhance the government's already strict control of
any political opposition.

In various pronouncements, top propaganda and security officials have stressed anew the need to
police the Internet on ideological and security grounds.

The ``Internet has become an important avenue through which anti-China forces infiltrate, sabotage
and magnify their capabilities for destruction,'' wrote the public security minister, Meng Jianzhu,
in the Dec.~1 issue of Qiushi, a magazine published by the Communist Party's Central Committee.

``Therefore it represents a new challenge to the public security authority in maintaining national
security and social stability,'' he said.

The newly announced restrictions are the government's broadest effort to control the Internet since
June, when it tried to require manufacturers to install Internet filtering software on all new
computers, experts said. Officials scaled back that program, known as the Green Dam-Youth Escort,
after an outcry by both individual Internet users and corporations.

Under the new controls, more than 700 Web sites have been shut down, including many that offered
free movies, television dramas and music downloads. BT China, which recorded at least 250,000 visits
daily, was among them. China's largest file-sharing site, Very CD, must obtain a new license or face
possible shutdown as well, according to news media reports.

In addition, individuals have been banned from registering Web sites ending in .cn, China's country
code domain name. That domain is now limited to registered businesses. Although individuals can
still register Web sites in other domains, like .com and .net, the new rule ``will have a negative
impact on the vibrancy of the Chinese Internet,'' Kenneth Jarrett, vice chairman of the
communications company APCO Worldwide's China region, said in an e-mail message.

``Local e-mail e-commerce start-ups and individuals will find it difficult to apply,'' he wrote.

Huang Xiwei, the founder of BT China, criticized the move in an interview posted on Sina.com, a
popular Chinese Internet portal. ``Not just film and video sites are affected,'' Mr.~Huang said.
``All Web sites owned by individuals will gradually exit the arena. All paths leading to a future
have been blocked.''

The government has also intensified pressure on cellphone companies to prevent transmissions of
online pornography. In response, China Mobile, the nation's largest cellphone operator, has
suspended its practice of allowing third-party providers to sell content over its cellular network,
according to the Chinese media. The disruption has dealt a blow to an industry that serves an
estimated nearly 200 million mobile Internet users, industry specialists said.

Experts say the latest measures are a continuation of the state's increasingly sophisticated effort
to control the Internet's influence on more than 300 million Chinese users. This year, China blocked
Facebook, Twitter, YouTube and thousands of other Web sites.

Some analysts had predicted that those restrictions would be lifted after a spate of potentially
troublesome anniversaries, including the 20th anniversary of the crackdown on pro-democracy
demonstrators in Tiananmen Square. But they have remained in place.

``The trend in China is toward tighter and tighter control,'' said Rebecca MacKinnon, an assistant
professor of journalism at the University of Hong Kong who specializes in Chinese Internet issues.
``They are basically improving their censorship mechanisms.''

Reaction to the government's latest crackdown runs the gamut from enthusiastic support from Chinese
parents who want to shield their children from pornography to harsh criticism from those who view
the Internet as the best antidote to government propaganda and state-controlled media.

Among the critics are college students who are accustomed to downloading music, films and other
material easily and cheaply. Some students at Beijing universities predicted that Internet users
would find ways around the new online obstacles.

Still, Wang Shuang, 20, a student at Beijing Foreign Studies University, complained, ``Since the BT
was closed, I cannot find all the American television series that I have been watching, like 'The
Mentalist.' ''

\section{China Is Disputing Status of Uighurs in Cambodia}

\lettrine{T}{he}\mycalendar{Dec.'09}{18} Chinese Foreign Ministry has hinted that it is seeking or
will seek the return of 22 Uighurs who fled to Cambodia after the eruption of deadly ethnic riots in
July in western China and a subsequent government crackdown.

The Uighurs are a Turkic-speaking ethnic minority group concentrated in western China whose members
often say the Chinese government, dominated by ethnic Han, discriminates against them. The 22
Uighurs in Cambodia entered the country about a month ago with the aid of an underground network of
Christian missionaries in China that usually helps North Koreans reach nations where they can seek
refugee status.

The Uighurs made their way to the Cambodian capital, Phnom Penh, where they applied for refugee
status at a United Nations refugee office.

This week, a Chinese Foreign Ministry spokeswoman said at a news conference that the Uighurs were
suspected of criminal activities and that the ``relevant departments'' were investigating them. She
said at the news conference and in a written statement sent to The New York Times on Thursday that
criminals should not be allowed to take advantage of the United Nations' refugee system.

``China's stance is very clear: the international refugee protection system shouldn't become a
shelter where criminals stay to escape legal punishment,'' she said.

The Foreign Ministry sent the statement in response to a question from a reporter asking whether
Chinese officials had contacted the Cambodian government about the Uighurs. The Phnom Penh Post, an
English-language newspaper, reported Tuesday that the Chinese Embassy had sent a note to the
Cambodian government in early December about the Uighurs.

Koy Kuong, a Cambodian Foreign Ministry spokesman, said he had no additional details about the note,
the newspaper reported.

This week, Amnesty International sent a letter to Sar Kheng, deputy prime minister and interior
minister of Cambodia, saying that Cambodia was bound by the terms of a 1951 convention related to
refugees that prohibited signatory countries from forcibly returning refugees to nations where they
could face torture or other ill treatment.

``Since September 2001, Amnesty International has documented cases in which Uighur asylum seekers
who were forcibly returned to China were detained, reportedly tortured and in some cases sentenced
to death and executed,'' Sam Zarifi, Asia-Pacific director of Amnesty International, wrote in the
open letter.

Three of the Uighurs who made it to Cambodia are children. Two Uighurs were detained in Vietnam en
route to Cambodia, and five others who fled China have disappeared, according to Uighur advocacy
groups in the West.

The United States has declined to send Uighurs detained in Guant\'anamo Bay back to China and has
also refused to grant them refugee status. Some of them have ended up in Albania, Palau and Bermuda.

The Uighurs in Cambodia fled their homeland after the deadliest ethnic rioting in decades in China.
Uighurs clashed on July 5 with riot police officers sent to put down a protest in Urumqi, the
capital of the western region of Xinjiang, and then went on a rampage through neighborhoods, killing
scores of people. In all, at least 197 people were killed, most of them ethnic Han, the Chinese
government said; in the next few days, vengeance-seeking Han attacked Uighurs across the city.

The government detained hundreds of Uighurs, and at least 43 Uighur men have disappeared, according
to Human Rights Watch. The Xinjiang government has executed several people for involvement in the
rioting.

\section{American Engineer Detained for Year in China}

\lettrine{A}{n}\mycalendar{Dec.'09}{18} American automotive engineer has been detained by the
Chinese police for more than a year on charges of violating trade secrets, according to United
States officials.

Hu Zhicheng, a Chinese-born American citizen who won awards in the United States for his work
developing catalysts that control auto emissions, is being held in a case that involves a dispute
with a former business partner in the northern port city of Tianjin.

Mr.~Hu's detention is the latest trade-secrets case to shed light on China's harsh legal system and
the complexities that can emerge in its dynamic but often ruthless business system.

The government has been holding another American, an oil geologist named Xue Feng, for nearly two
years on suspicion of stealing state secrets. President Obama raised Dr.~Feng's detention during his
visit to China last month.

Another ongoing detention case -- that of Stern Hu, a Chinese-born Australian executive working for
Rio Tinto, the Australian mining giant -- has received widespread coverage in part because it caused
diplomatic tensions between Australia and China.

United States officials say that they have been allowed to visit Hu Zhicheng at least 10 times but
that he is not allowed to receive any written communication from his wife.

Mr.~Hu's lawyer could not be reached for comment. A spokeswoman for the Tianjin police, who would
give only her surname, Liu, said by telephone Thursday that it was inappropriate to discuss the
case.

United States officials said he was initially arrested in November 2008, but formally arrested and
charged in January 2009, while working for a company called Hysci Specialty Materials, based in
Tianjin.

The news of the detention was first reported by The Associated Press, and it is unclear why the
detention was not disclosed sooner.

\section{Nuclear Power Expansion in China Stirs Concerns}

\lettrine{C}{hina}\mycalendar{Dec.'09}{18} is preparing to build three times as many nuclear power
plants in the coming decade as the rest of the world combined, a breakneck pace with the potential
to help slow global warming.

China's civilian nuclear power industry -- with 11 reactors operating and construction starting on
as many as an additional 10 each year -- is not known to have had a serious accident in 15 years of
large-scale electricity production.

And with China already the largest emitter of gases blamed for global warming, the expansion of
nuclear power would at least slow the increase in emissions.

Yet inside and outside the country, the speed of the construction program has raised safety
concerns. China has asked for international help in training a force of nuclear inspectors.

The last country to carry out such a rapid nuclear expansion was the United States in the 1970s, in
a binge of reactor construction that ended with the Three Mile Island accident in Pennsylvania in
1979. And China is placing many of its nuclear plants near large cities, potentially exposing tens
of millions of people to radiation in the event of an accident.

In addition, China must maintain nuclear safeguards in a national business culture where quality and
safety sometimes take a back seat to cost-cutting, profits and outright corruption -- as shown by
scandals in the food, pharmaceutical and toy industries and by the shoddy construction of schools
that collapsed in the Sichuan Province earthquake last year.

``At the current stage, if we are not fully aware of the sector's over-rapid expansions, it will
threaten construction quality and operation safety of nuclear power plants,'' Li Ganjie, the
director of China's National Nuclear Safety Administration, said in a speech this year.

A top-level corruption scandal is already unfolding in the nuclear industry.

In August, the Chinese government dismissed and detained the powerful president of the China
National Nuclear Corporation, Kang Rixin, in a \$260 million corruption case involving allegations
of bid-rigging in nuclear power plant construction, according to official media reports. No charges
have been reported against Mr.~Kang, who is being held incommunicado for interrogation.

While none of Mr.~Kang's decisions publicly documented would have created hazardous conditions at
nuclear plants, the case is a worrisome sign that nuclear executives in China may not always put
safety first in their decision-making.

In contrast with its performance in industries like toys, China has a strong safety record in
industries like aviation, which receive top-level government attention.

The challenge for the government and for nuclear companies as they increase construction is to keep
an eye on a growing army of contractors and subcontractors who may be tempted to cut corners.

``It's a concern, and that's why we're all working together because we hear about these things going
on in other industries,'' said William P.~Poirier, a vice president for Westinghouse Electric, which
is building four nuclear reactors in China.

Philippe Jamet, the director of the division of nuclear installation safety at the International
Atomic Energy Agency in Vienna, said that China had welcomed foreign inspectors at its reactors and
that ``they show pretty good operations safety.''

But he added that the international agency was concerned about whether China would have enough
nuclear inspectors with adequate training to handle the rapid expansion.

``They don't have very much staff, when you compare their staff with how many they will need,''
Mr.~Jamet said. The agency accepted a Chinese request to send a team of international experts to the
country next year to assess staffing and training, he added.

In late October, Prime Minister Wen Jiabao ordered a quintupling of the safety agency's staff by the
end of next year, to 1,000, according to United States regulators. Chinese officials did not respond
to requests for confirmation.

China has two rival state-owned nuclear power giants: the China National Nuclear Corporation, mainly
in northeastern China, and the China Guangdong Nuclear Power Group, mainly in southeastern China.

Western experts regard the Daya Bay nuclear power plant in Shenzhen, which mainly uses French
designs and is run by China Guangdong Nuclear, as evidence that China can run reactors safely. A
display case holds trophies the power plant won in global safety competitions.

China National Nuclear likewise cooperates with international inspectors and has had no reported
mishaps. But its roots are in a government ministry with close ties to the former Soviet Union,
making it more of an enigma to most Western experts, and the corruption case has added to their
concern. China National Nuclear was on track to grow faster than China Guangdong over the next
decade.

China National Nuclear has sought to hush up the case involving the arrest of its president,
deleting from its Chinese-language Web site even the most minor news releases that mentioned
Mr.~Kang. In a faxed response to questions, China National Nuclear made no mention of Mr.~Kang, but
emphasized that its plants met international standards.

The arrest of Mr.~Kang, a member of the Chinese Communist Party's powerful Central Committee, can be
seen as evidence of China's seriousness about safety.

Today, China's nuclear plants can produce about nine gigawatts of power when operating at full
capacity, supplying about 2.7 percent of the country's electricity. Three years ago, the government
set a goal of increasing that capacity more than fourfold by 2020.

The government will soon announce a further increase in its targets, to 70 gigawatts of capacity by
2020 and 400 gigawatts by 2050, said Jiang Kejun, an energy policy director at the National
Development and Reform Commission, the main planning agency.

Electrical demand is growing so rapidly in China that even if the industry manages to meet the
ambitious 2020 target, nuclear stations will still generate only 9.7 percent of the country's power,
by the government's projections.

Bringing so much nuclear power online over the next decade would reduce the country's energy-related
emissions of global warming gases by about 5 percent, compared with the emissions that would be
produced by burning coal to generate the power.

``For anyone concerned about carbon dioxide emissions, it's heartening, but it's only a piece of the
puzzle,'' said Jonathan Sinton, a China specialist at the International Energy Agency in Paris.

China, which by most estimates overtook the United States in 2006 to become the largest emitter of
greenhouse gases, is seeking sharp improvements in the energy efficiency of its economy.

But the economy is growing so fast that even if the country can meet its goals, total emissions will
rise 72 to 88 percent by 2020, Mr.~Sinton said.

The challenge for China is to build and operate its nuclear reactors without the equivalent of the
Three Mile Island accident, in which a reactor core partly melted and released radioactivity, or the
Chernobyl disaster in the former Soviet Union in 1986, the world's worst civilian nuclear accident.

China does not use the kind of reactor that exploded at Chernobyl. And engineers in China study the
mistakes that poorly trained operators made at Three Mile Island.

Liu Yanhua, a vice minister of science and technology, said China believed that its nuclear industry
would continue to grow safely.

``So far,'' Mr.~Liu said, ``there is no damage.''

\section{Obama Interviewer in China Is Said to Be Demoted}

\lettrine{T}{he}\mycalendar{Dec.'09}{18} top editor of Southern Weekly, one of China's most
influential newspapers, apparently has been demoted, just weeks after his interview with President
Obama annoyed Communist Party propaganda officials, according to people familiar with the decision.

Although the interview with Mr.~Obama at the end of his visit to China last month did not yield any
significant news, or even include unusually tough questions, it seemed to bother propaganda
officials by putting the spotlight on the feisty, liberal-leaning paper from southern China's
Guangdong Province, rather than China's official news agency, Xinhua, or its main broadcaster, China
Central Television, according to several people close to the paper who asked not to be named for
fear of being punished.

Xiang Xi, who had been the paper's editor in chief, was named its executive editor last week. But
several people familiar with the paper's decision said his role at the paper was diminished.

Reached by telephone Tuesday afternoon, Mr.~Xiang said he was not allowed to comment on the move.

Huang Can, who was named the new editor in chief, said in a telephone interview Tuesday that the
move was ``not really'' a demotion for Mr.~Xiang, and added, ``His specific job has not been clearly
decided. We'll know soon.''

The White House chose Southern Weekly for Mr.~Obama's interview, rather than the more official
Chinese media outlets, in what seemed to be a move intended to appeal to a publication with a more
independent voice. But the Nov.~19 interview was cut out of many copies of Southern Weekly the
following day, particularly in Beijing, and was not featured prominently on the paper's Web site.

The White House later published a full transcript of the interview on its own Web site.

\section{At White House, Holidays Are Both Treat and Test}

\lettrine{F}{or}\mycalendar{Dec.'09}{18} many in this town, December is a magical season of
flickering menorahs and sparkling pine trees, spicy eggnog and savory latkes. But at the White House
it is also a grueling whirlwind of nonstop parties and receptions -- the biggest test yet of the
Obamas' entertaining endurance and party-hosting prowess.

The White House is holding 17 parties and 11 open houses, all over the course of about three weeks.
President Obama and his wife, Michelle, are throwing open the doors to more than 50,000 visitors --
supporters, lawmakers, rabbis, soldiers, generals, journalists, White House staff members and
ordinary people -- who will ogle the decorations and angle for coveted close encounters with the
president and first lady.

While this public display of power entertaining is an annual tradition, the Obamas have put their
stamp on the festivities by highlighting their eco-friendly decorations, an understated d\'ecor to
reflect the difficult economic times, and a buffet menu that showcases healthful fare like salmon,
fresh vegetables and sushi.

Unlike their predecessors, the Obamas face heightened scrutiny as the White House puts in place
tougher measures to prevent security breaches like the one in November, when a couple seeking a spot
on a reality television show crashed the state dinner for the prime minister of India.

``It's a massive, massive undertaking,'' said Kim Merlin, the Chicago floral consultant who helped
coordinate the White House's holiday d\'ecor with Simon Doonan, the creative director of Barneys New
York.

Ms.~Merlin said the planning required ``an army of people,'' including aides from the social
secretary's office, florists, chefs, electricians, butlers, welders, carpenters and volunteers.
There are 27 Christmas trees, 129 wreaths, 2,700 yards of ribbon and a 390-pound gingerbread White
House (with Bo, the first dog, depicted in marzipan), not to mention some 48,000 sweets for guests.

So far, the Obamas have played down the organizational weightlifting required.

``It is just absolutely magical,'' Mrs.~Obama told Oprah Winfrey in an interview about Christmas in
the White House that was broadcast Sunday on ABC.

And during a seemingly endless receiving line at one of the two holiday parties for journalists this
week, Mr.~Obama chuckled and dismissed questions about his party stamina. ``We don't get tired,'' he
said.

But Anita McBride, who was Laura Bush's chief of staff during the second Bush term, said White House
staff members, particularly those in the East Wing, had always viewed the holiday season with a
sense of eager anticipation and some dread.

``It's beautiful; it's a special time,'' Ms.~McBride said. ``But at the beginning of the season, you
would look at the calendar and say, 'How are we going to get through this?' There was a part of you
that really was bracing yourself for being on your feet so long.''

The Obamas have opted not to stand as much, abandoning the practice of shaking hands and posing for
photographs with every guest at every party, a custom embraced (though perhaps a bit wearily) by
other presidents and first ladies. Their aides say the Obamas prefer smaller gatherings where they
can mix and mingle, but some guests have expressed dismay.

``People who didn't know were disappointed,'' said Dee Dee Myers, a former press secretary to
President Bill Clinton, who attended the party that included Democratic donors and party officials
and described the glum chatter there. ``I totally understand that doing receiving lines is
exhausting, yet people have come to expect it. It means a lot to them.''

There are other changes, too. The Obamas are having fewer parties than the Bushes did and none on
weekends, aides say, perhaps reflecting that the White House is also home to their daughters, Malia
and Sasha.

As for the d\'ecor, Mrs.~Obama decided to display fewer trees and wreaths. There is no artificial
snow and no Santa Claus, who was displayed in the East Garden Room last year. Aides said the choices
reflected the Obamas' personal taste and their empathy for the hardships facing many Americans.

``It should be beautiful and refined, but it also should be restrained, because it's a difficult
time,'' Ms.~Merlin said, describing the principles that guided her mission.

So the more than 800 ornaments adorning the Christmas tree in the Blue Room were recycled from
earlier administrations but updated with new designs by community groups across the country.

The cranberry Christmas tree in the Red Room -- a holiday staple since the Reagan administration --
is gone, but cranberry garlands are on the mantle.

As for security, visitors' names were being checked twice before they entered the White House this
week. Guests were also being asked to show photo identification at two separate checkpoints, and
Desir\'ee Rogers, the social secretary, was on hand to welcome the invitees.

Meanwhile, like other first ladies, Mrs.~Obama is also making time in her hectic social schedule to
help the needy, organizing a toy drive for poor children and urging Americans to help the hungry by
volunteering in food banks and delivering meals to the homebound.

``That's what America is all about, people already sacrificing, stepping up and doing a little bit
more,'' Mrs.~Obama said this week as she delivered toys to the Toys for Tots warehouse in Virginia.

\section{'American Idol' Creator Plans Web Show}

\lettrine{T}{he}\mycalendar{Dec.'09}{18} creator of ``American Idol,'' Simon Fuller, is developing a
new reality show that will unfold on the Internet in real time, largely sidestepping television.

The venture, called ``If I Can Dream,'' is expected to start in February on Hulu and other Web
sites.

Mr.~Fuller's company, 19 Entertainment, said the venture would follow the lives of five aspiring
actors and musicians who move to Hollywood in search of stardom. Their daily activities will be
shown live on a Web site, IfICanDream.com, and they will communicate with viewers via blog posts and
social networks.

A half-hour recap show will be produced once a week and streamed exclusively on Hulu, a popular Web
site for watching shows from NBC, Fox and other networks. The project represents the first time that
Hulu has distributed a full-length original program. ``Everyone has a dream. Watch them follow
theirs,'' says a promotional video for the show on Hulu.

Notably, ``If I Can Dream'' will not be shown on television, at least not at first. But 19
Entertainment would like it to earn a place on TV eventually.

The company announced marketing and advertising deals on Wednesday with MySpace, Clear Channel
Radio, Pepsi and Ford Motor Company for the project.

Several companies have tried to create hit series for the Internet in recent years, with little to
show for it. But ``If I Can Dream'' stands out because it has the backing of Mr.~Fuller, a respected
music and media producer who created ``Pop Idol,'' a British show that spawned ``American Idol.''
About to begin its eighth season, ``Idol'' is the most popular television show in the United States.

In a statement, Mr.~Fuller said that the program would push the boundaries of mainstream
entertainment. ``The next frontier is the video world of authentic real-time interaction,'' he said.

People will be able to audition for ``If I Can Dream'' on MySpace; radio hosts on Clear Channel
stations will help to promote it.

Media companies and producers are becoming increasingly interested in distributing shows on the Web.
In a separate development, the actor and producer Ashton Kutcher announced Wednesday that he was
giving his failed TV show ``The Beautiful Life'' an encore on YouTube.

Only two episodes of ``The Beautiful Life'' were shown on the CW network before the show was
canceled because of poor ratings. (Only 1.1 million people watched the show in its second week.)
Those two episodes and several that were not broadcast will stream on YouTube free through a
sponsorship with Hewlett-Packard.

Sounding ambitious, Mr.~Kutcher said in a YouTube video that ``I want this to be the first show ever
that gets more viewers on the Web than what it got on terrestrial television.''

\section{BlackBerry's Maker Posts a 59\% Increase in Profit}

\lettrine{T}{he}\mycalendar{Dec.'09}{18} BlackBerry maker Research In Motion, while dealing with a
service outage on Thursday, posted quarterly results that surprised analysts and set off a 12
percent gain in its share price.

The company said that profit in its most recent quarter rose by 59 percent as the company attracted
new subscribers and sold a record 10 million phones.

The company said it earned \$628.4 million, or \$1.10 a share, compared with \$396.3 million, or 69
cents a share, in the period a year earlier.

Revenue in the period, which ended Nov.~28 and was the third quarter of the company's fiscal year,
was up 41 percent, to \$3.92 billion, from \$2.78 billion.

The company's performance surpassed the expectation of analysts polled by Thomson Reuters, who were
expecting net income of \$1.04 a share and revenue of \$3.78 billion.

``We are pleased to report record shipments of more than 10 million BlackBerry smartphones during
the third quarter with higher-than-expected revenue, earnings and subscriber growth,'' the co-chief
executive, Jim Balsillie, said in a statement.

During a conference call, Mr.~Balsillie said that more than 80 percent of new subscribers in the
quarter were nonbusiness customers.

After initially focusing on corporate customers, Research In Motion has expanded its reach into the
consumer market in recent years. The devices face competition from devices like the iPhone from
Apple, the Palm Pre and the Motorola Droid.

Research In Motion said fourth-quarter revenue was expected to be in the range of \$4.2 billion to
\$4.4 billion, beating analysts' expectation of \$4.1 billion. The company also said that gross
margin was expected to be at 43.5 percent and earnings per share in the range of \$1.23 to \$1.31.

Research in Motion's report came as BlackBerry users in North America faced delays in receiving
e-mail. The company said that technicians isolated and resolved the issue and that it was
investigating the cause.

Users were still able to make phone calls, browse the Internet and send and receive text messages.
The company did not say how many users were affected or how long the outage lasted. It said some
customers might still experience delays as queues are processed.

The company, which is based in Waterloo, Ontario, released its results after the close of market
trading. In after-hours trading, shares of Research In Motion surged \$7.76 to \$71.22.

\section{Waking Up Serious Instead of Silly}

\lettrine{T}{he}\mycalendar{Dec.'09}{18} best way to assess George Stephanopoulos's first day as
Diane Sawyer's replacement on ``Good Morning America'' was to watch the competition on NBC.

The ``Today'' show, as is its wont, started on Monday with a melodramatic report on missing hikers
on Mount Hood and then lingered on Tiger Woods's misfortunes by way of a long, lurid interview with
Cori Rist, who claims to be one of his mistresses.

Mr.~Stephanopoulos didn't adjust to the tabloid tone of morning news so much as bring his sober
Washington talk show with him to the ABC studio in Times Square. Monday's show led with President
Obama's attack on Wall Street bankers, illustrated with a clip from a ``60 Minutes'' interview on
CBS from the night before. Mr.~Stephanopoulos then chatted with David Axelrod, a senior adviser to
Mr.~Obama.

``Good Morning America'' did report the latest on Tiger Woods -- the consulting firm Accenture
dropped its sponsorship of him -- but interviewed only advertising and media experts, with no
detours to mistresses or Vegas blondes. Even the musical guest of the day was high-minded: the
cellist Yo-Yo Ma playing ``Dona nobis pacem.''

Mr.~Stephanopoulos is a policy-minded, even wonky addition to a morning show that in the past mostly
relied on Everyman bathos and Hollywood glitz. His debut marks the start of a busy season of
turnover in the ABC News division: Ms.~Sawyer, who takes over the evening news from Charles Gibson
later this month, bade a tearful goodbye to her morning viewers last week. Mr.~Gibson's tenure as
the anchor of ``World News'' ends on Friday.

ABC News hopes that its anchor shakeup will help close its ratings gap with NBC News, but new faces
are not necessarily a guarantee of success: the anointment of Katie Couric as the first solo female
network anchor has not budged CBS's evening news out of a distant third place.

Despite ever-dwindling audiences, network news programs cling to established formula and
cookie-cutter formats. Changes in personnel don't radically alter the program, but personality does
have an effect on its content.

While it takes time for anchors to adjust to new desks, and first days do not always reflect
subsequent performance, it was notable that Mr.~Stephanopoulos did not try to morph into Morning
George right out of the gate. This anchor -- who went to ABC News after working as a campaign
operative and White House adviser to President Bill Clinton, and will continue to host ``This Week
With George Stephanopoulos'' until a successor is named -- led with his strong suit. He questioned
Mr.~Axelrod about the health care bill; the president's plans to badger ``fat cat'' bankers about
bonuses; and the economy. Mr.~Axelrod presented him with an alarm clock permanently set to 3:30
a.m., a gift from, as he put it, ``your friends in the White House.''

He was welcomed with fanfare by his co-host, Robin Roberts, but it wasn't clear that the two have
natural chemistry. Ms.~Roberts, who was of a height with Ms.~Sawyer, towers over the diminutive
Mr.~Stephanopoulos; their senses of humor seem equally unaligned. Ms.~Roberts is more literal than
funny and is given to outsize expressions of feeling. Mr.~Stephanopoulos is reserved and permits
himself small, sarcastic asides.

After noting that surveys suggest that women view Mr.~Woods less favorably than men do, the new
anchor said dryly, ``Well that's a shocker, huh?''

Ms.~Roberts, using a ``Romper Room'' tone that suggested that she didn't realize he was kidding,
replied, ``We're learning something new every day.''

Mr.~Stephanopoulos wasn't the only new face on the ``Good Morning America'' set. Sam Champion stayed
on as weatherman, but JuJu Chang took over the news desk from Chris Cuomo, who will be an anchor on
``20/20.''

Mr.~Stephanopoulos wasn't spared the more frivolous consumer tips that are the backbone of morning
talk shows. He sat, gamely but rigidly, through a lesson on Internet coupon shopping. He chatted
with Dr.~Mehmet Oz about holiday heart attacks: Dr.~Oz reported that booze, fatty food and stressful
relatives are a ``hurricane of factors that make this holiday so dangerous.''

Mr.~Stephanopoulos will almost certainly have more silly features in his future, but it looks as if
he won't entirely shun his political past. On a network morning show, that's a refreshing change.

\section{In Senate Vote on Health Care Overhaul, New Partisan Vitriol}

\lettrine{T}{he}\mycalendar{Dec.'09}{24} vote on Monday, in the dead of night, was 60 to 40. The
vote on Tuesday, just after daybreak, was 60 to 39. And the vote on Wednesday afternoon, at a civil
hour but after less-than-civil debate, was 60 to 39 again -- an immutable tally that showed
Democrats unwavering in the march to adopt a far-reaching overhaul of the health care system over
united Republican opposition.

The votes also marked something else: the culmination of more than a generation of partisan
polarization of the American political system, and a precipitous decline in collegiality and
collaboration in governing that seemed to move in inverse proportion to a rising influence of
lobbying, money, the 24-hour news cycle and hostilities on talk shows and in the blogosphere.

The health care legislation is likely to be approved Thursday morning, with the Senate divided on
party lines -- something that has not happened in modern times on so important a shift in domestic
policy, or on major legislation of any kind, lawmakers and Congressional historians said.

The Democrats flaunted their unity on Wednesday at a news conference with nearly their entire caucus
in attendance.

Many senators said the current vitriol, which continued on the floor on Wednesday with a fight over
when to cast the final health care vote, was unlike anything they had seen. ``It has gotten so much
more partisan,'' said Senator John D.~Rockefeller IV, Democrat of West Virginia. ``This was so
wicked. This was so venal.''

Even in a bitter fight over President Bill Clinton's budget in 1993, decided 51 to 50 with a rare
tie-breaker vote by Vice President Al Gore, the partisanship was not as stark as it is today.
Although no Republicans voted for Mr.~Clinton's budget, six Democrats joined them in what amounted
to bipartisan opposition.

Mr.~Rockefeller said the health bill had created an almost perfect storm of political and policy
disagreements, so that some of the bitterness reflected basic philosophical disputes crystallized by
President Obama's agenda. ``If there was ever a time for that kind of partisanship to come out, this
was the bill to do it,'' he said.

Ross K.~Baker, a political scientist at Rutgers University and an expert on the history of the
Senate, said that in earlier eras, senators would routinely cross party lines to vote in favor of
major legislation on issues like civil rights and social welfare policy.

In 1965, the Senate created the Medicare program by a vote of 68 to 21, with 13 Republicans joining
55 Democrats in favor, and 7 Democrats joining 14 Republicans in opposition. In 2003, some Democrats
in both the House and the Senate voted with most Republicans to add a prescription drug benefit to
Medicare.

``It certainly is a culmination of a long period of intensifying political polarization,'' Mr.~Baker
said of this year's showdown over health care. ``It has gotten so bad now that Republicans don't
want to be seen publicly in the presence of Democrats or have a Democrat profess friendship for them
or vice versa.''

With Democrats nominally controlling 60 seats, the precise number needed to overcome Republican
filibusters, there is no room for wavering Democrats to break ranks. If they held one less seat,
there would be no choice but to win over a Republican; one or two more, and one or two senators with
apprehensions could be released to vote no.

Some lawmakers predicted that the Senate would eventually rediscover its genteel equilibrium.

``There's a tolerance level here for what we have just been through, and I think we have hit the
tipping point,'' said Senator Christopher J.~Dodd, Democrat of Connecticut. ``It got rougher than it
should. We are getting precariously close to fracturing an institution where no one wins, so I think
we are going to be back on track.''

But some experts said that the divide in the Senate reflected a broader political shift that
lawmakers cannot easily reverse. ``In the 1970s, for instance, there was a much wider political
spectrum in both parties,'' said Donald A.~Ritchie, the Senate historian. ``You had conservative and
liberal wings in both parties.''

Mr.~Ritchie and many senators said they had witnessed the change in the last 30 years.

``You have got this divide, this polarization in America,'' said Senator Olympia J.~Snowe of Maine,
the only Republican in recent weeks to seriously consider supporting the health bill. ``People
become risk-averse, politically risk-averse. There is no incentive to reach across the divide and
appeal to a broader inclination. It looks like pragmatism is a political cop-out; compromise is
certainly viewed that way.''

But even as senators complained about the rancor and expressed nostalgia for a kinder era, they
conceded that the hyper-partisanship was likely to continue, potentially coloring coming debates on
other major issues including financial regulation, climate change and, perhaps, immigration.

Senator Max Baucus, Democrat of Montana, and chairman of the Finance Committee, said the political
-- and often personal -- divisions that now characterize the Senate were epitomized by the empty
tables in the senators' private dining room, a place where members of both parties used to break
bread.

``Nobody goes there anymore,'' Mr.~Baucus said. ``When I was here 10, 15, 30 years ago, that was the
place you would go to talk to senators, let your hair down, just kind of compare notes, no spouses
allowed, no staff, nobody. It is now empty.''

For more than 30 years, the major parties -- Democrats and Republicans -- worked every angle to
transform politics into a zero-sum numbers game. State legislatures redrew Congressional districts
to take advantage of party affiliation in the local population. The two-year campaign cycle became a
never-ending one.

Senator Orrin G.~Hatch, Republican of Utah, who worked on many bipartisan health care bills over the
years, often with a close friend, the late Senator Edward M.~Kennedy, Democrat of Massachusetts,
said that the both parties were to blame but that external factors including ethics rules also
discouraged senators from fraternizing.

``Both parties have become very polarized,'' Mr.~Hatch said. ``A lot of that is because of the
stupid ethics rules. We can't get together at various events. A lot of people complain about taking
foreign trips, which are really critical for us to understand foreign policy. The Internet is
constantly badgering everybody. In the process, it's gotten pretty doggone partisan, both ways. It's
bad.''

Mr.~Hatch and Republican leaders said the lack of any support on their side showed that the health
bill was mortally flawed.

The majority leader, Harry Reid of Nevada, at a news conference on Wednesday with most of his caucus
standing behind him, offered a different take.

``I don't see this as 60 Democrats versus 40 Republicans,'' he said. ``I see it as 60 leaders who
stood up to insurance companies and stood up for working families all across America.''

\section{Taking Hold in Silicon Valley, a Ping-Pong Boom}

\lettrine{Y}{oung}\mycalendar{Dec.'09}{24} people who were serious about table tennis used to have
to make the trip to Beijing, Stockholm or Moscow to train with world-class coaches.

Now they go no farther than this Silicon Valley suburb.

``I'm trying to become one of the greatest players in the nation,'' Srivatsav Tangirala, 14, said
matter-of-factly between drills at the huge new table tennis facility here. He and three dozen
players, some as young as 5, sprinted sideways along the edge of the tables, 45 times in a row,
perfecting their footwork.

``Lean forward, lean, lean, lean, lean!'' their coach implored.

This is the largest training program for youths in the country, run by the India Community Center in
an area that is 60 percent Asian. Here, Ping-Pong parents who grew up with the sport in Sichuan
Province or Hyderabad are the new soccer moms and Little League dads.

One of 12 table tennis clubs in the area, up from 5 clubs in 1990, the India Community Center's
Ping-Pong facility was started last year with seed money from two Indian entrepreneurs and has
already become an influential hatchery for Olympic hopefuls, most of whom banter in Hindi or
Mandarin at home.

Ariel Hsing, 14, the top-ranked United States junior, from San Jose, Calif., and Lily Zhang, 13 and
ranked No.~2, from Palo Alto, Calif., are a fearsome twosome, with matching teal braces, bulging
calf muscles and a dream of playing in the 2012 Olympics. Ariel cradles the ball in her palm like a
baby chick -- before she lets go and smashes it.

They and over 100 other teenagers, many the daughters and sons of technology professionals, are
being coached by talent from around the world: Gaolin Tang from Sichuan Province; Stellan Bengtsson,
the Swedish champion; and Rajul Sheth, the center's executive director, a veteran of the Indian
national team.

In the past, top players grew up in China and became American citizens in order to play for the
United States Olympic team. Today, 80 percent of players age 14 and younger are Asian-Americans,
according to USA Table Tennis, the sport's national governing body.

``Hyphenated kids who are born and raised here and have a foot in both worlds are the ones taking
the lead, '' said David Del Vecchio, a board member of the National Collegiate Table Tennis
Association.

In Milpitas, the hollow knocking sounds of Ping-Pong balls reverberate off walls lined with
triumphal newspaper clippings in The Sing Tao Daily and The India Express featuring India Community
Center offspring.

Ariel's mother, Xian Hua Jiang, a 46-year-old hardware engineer, was weaned on two-volley games on
concrete tables in the schoolyard in Henan Province. Growing up poor, she had to borrow white shoes
to participate in a tournament. ``Rain or shine, during school breaks everyone rushes to the
tables,'' she recalled.

Today, she and her husband, Michael Hsing, a software engineer from Taiwan, spend at least \$40,000
a year fostering their daughter's talent, and have added an indoor table tennis practice room to
their house. Their plastic Home Depot window shades are pocked with holes from the velocity of
Ariel's balls.

Named for Disney's ``Little Mermaid,'' Ariel juggles school and international tournaments in Tokyo,
Chile and elsewhere. Her ferocious close-to-the-table backhand, honed from age 6, has generated so
much buzz that two years ago she was invited to play with ``Mr.~Bill'' -- as in Bill Gates, the
chairman of Microsoft -- and ``Mr.~Warren'' -- Warren E.~Buffet, the investor. She overwhelmed the
billionaire Ping-Pongers at Mr.~Buffet's 75th birthday party in San Francisco. ``They're very humble
and down to earth,'' she said.

Although the sport's visibility is growing nationally -- thanks in part to celebrities like Susan
Sarandon, who recently opened SPiN New York, a table tennis social club in Manhattan -- it does not
yet have Little League-style cultural clout in the United States, which parents say makes it
difficult to compete. In Germany, table tennis engenders a Los Angeles Lakers-like fever, with
televised games and some professionals earning \$1 million a year in endorsements.

In this country, the sport is still considered a hobby; only three colleges -- Texas Wesleyan
University in Fort Worth; Lindenwood University in St.~Charles, Mo.; and the University of Puerto
Rico in San Juan -- offer table tennis scholarships.

Randy Capps, an analyst at the Migration Policy Institute in Washington, said Silicon Valley's
rapidly growing Ping-Pong scene reflected the state's demographics, in which half its school-age
children are the offspring of immigrants. ``The parents have competed hard to get where they are,''
Mr.~Capps said. ``They expect their children to do the same.''

Min Zhou, a professor of sociology at the University of California, Los Angeles, said that among
Asian-Americans, there was a perception that the children would not excel at football or basketball,
and that table tennis ``is a sport where they have an advantage because of cultural affinity.''

``Being too academically oriented has become a stereotype of the Asian-American kid,'' she said.
``Parents are grounding them in sports so the kid does not appear as nerdy.''

Now junior teams from India and Hungary come to Milpitas to work with coaches like Sean O'Neill, an
Olympian from McLean, Va. ``I had to go to Sweden'' to train, Mr.~O'Neill said. ``But they've
brought the world's borders inside'' the India Community Center.

It is a global game, with the Chinese coaches specializing in a fast ``flat game'' in which the ball
is hit hard close to the table, while their European and Indian counterparts are ``topspinners'' who
move away from the table to put more spin on the ball.

The sprawling new complex is a satellite of the India Community Center, an ambitious,
one-stop-shopping center for Indian culture modeled on Jewish community centers. It includes a free
medical clinic, a program for retirees and Bollywood aerobics classes.

``In India, you walk out of the house and there is a community all around you,'' said Anil Godhwani,
a co-founder of the center. ``In the U.S., we felt we were missing something.''

The program started small in 2005 with five Indian players. ``The Chinese people didn't want to
learn table tennis from some Indian,'' as Mr.~Sheth put it. Winning 16 medals the following year at
the Junior Olympics helped persuade the Chinese of the India Community Center's serious intentions.
Today parents have nicknamed it ``the India-China center.''

Last week, at the United States table tennis national championships in Las Vegas, 21 players from
Milpitas competed against 653 athletes from around the country, garnering 15 awards -- the most of
any club.

Mr.~Sheth's team included Krish Avvari, a bespectacled fourth grader known for his mean loop
forehand, and Ariel Hsing, known as a powerful ``two-winged attacker,'' comfortable with backhand or
forehand.

``I want to control who wins and who loses,'' Ariel said of her style -- an exercise perhaps in the
most subtle art of spin.

\section{Senate Clears Final Hurdle to Vote on Health Care Bill}

\lettrine{T}{he}\mycalendar{Dec.'09}{24} Senate trudged Wednesday toward passage of sweeping health
legislation after disposing of Republican claims that it would be unconstitutional to require
Americans to have health insurance, as the bill does.

The Senate was poised to take a final vote on the legislation, President Obama's top priority, on
Thursday morning.

Senator Orrin G.~Hatch, Republican of Utah, said Democrats had secured the 60 votes they needed with
``a grab bag of backroom Chicago-style buyoffs'' for specific states and favored constituencies.

But Republicans could not crack the 60-member Democratic caucus, which hung together Wednesday to
advance the legislation in the face of solid Republican opposition.

The bill cleared a last procedural hurdle on Wednesday when the Senate voted 60 to 39 to wind up
debate on the measure, which would extend coverage to more than 30 million uninsured Americans.
Sixty was the precise number needed to limit debate.

The Senate on Wednesday rejected a Republican effort to require disclosure of all items inserted in
the bill for the benefit of specific states and entities.

The Senate also rejected a constitutional point of order against the bill by Senator John Ensign,
Republican of Nevada.

``I don't think Congress has ever required Americans to buy a product or service, such as health
insurance, under penalty of law,'' Mr.~Ensign said. ``I doubt Congress has the power to do that.
Under this bill, if an American chooses not to buy minimal essential health coverage, he or she will
face rapidly increasing taxes -- up to \$750 or 2 percent of taxable income, whichever is greater,
by the year 2016.''

But Senator Max Baucus, Democrat of Montana and an architect of the bill, said Congress had
authority to impose a health insurance mandate under its constitutional power to regulate interstate
commerce and to tax and spend for the general welfare.

``Health care and health insurance affect and are distributed through interstate commerce,''
Mr.~Baucus said. ``Claims payments flow through interstate commerce. That gives Congress the power
to legislate a coverage requirement using its commerce powers.''

The issue, like many other provisions of the bill, is almost sure to be tested in court if the bill
becomes law.

The Senate session on Wednesday felt a bit like the last day of school before a vacation, as
lawmakers joked and laughed on the floor and exchanged holiday greetings.

But the camaraderie did nothing to dissolve Republican objections to the bill. Members of the two
parties skirmished over when to hold the final vote.

Given the bad weather threatening holiday travel across much of the country, Democrats asked
unanimous consent to hold the vote on Wednesday night. Republicans balked.

Senator Jim Risch, Republican of Idaho, jumped to his feet to make his own request: that the vote be
held at 2 p.m. on Jan.~20, when senators are scheduled to return from recess. Numerous Democrats
objected.

Since 1850, the Senate has held only eight sessions on Christmas Eve. Vice President Joseph R.~Biden
Jr.~plans to preside over the Senate Thursday when it votes on final passage of the health bill.

In an interview with PBS, President Obama said Wednesday that he was pleased with the Senate bill,
as with a companion bill passed by the House last month.

``I'm getting 95 percent of what I want,'' Mr.~Obama said.

He vowed to sign the final legislation even if it did not include a government-run insurance plan.
The House bill, but not the Senate measure, includes such a public option.

``I've been in favor of the public option,'' Mr.~Obama said. ``I think the more choice, the more
competition we have, the better.'' But, he said, the legislation includes many consumer protections
and ``extraordinary reforms,'' so it would make no sense to reject it just because it had no public
insurance plan.

Mr.~Obama said he would be deeply involved in negotiations to work out differences between the
Senate and the House. And he dismissed the idea that he had already compromised too much.

``This notion that somehow the health care bill that is emerging should be grudgingly accepted by
Democrats as half a loaf is simply incorrect,'' Mr.~Obama said. ``This is nine-tenths of a loaf. And
for a family out there that right now doesn't have health insurance, it is a great deal. It's a full
loaf for a lot of families who have nothing to fall back on if they get into a medical emergency.''

Republicans said Democrats were trying to jam the bill through the Senate before the public could
register opposition.

Senator Charles E.~Grassley, Republican of Iowa, who spent weeks last summer in unsuccessful
negotiations with Democrats, said Wednesday, ``This bill slid rapidly down the slippery slope to
more and more government control of health care.''

But after the last procedural hurdle, Democrats rejoiced.

``We stand a few short steps from the most significant finish line we've had in Congress for many
decades,'' said the Senate majority leader, Harry Reid, Democrat of Nevada.

\section{Trial in China Signals New Limits on Dissent}

\lettrine{L}{iu}\mycalendar{Dec.'09}{24} Xiaobo, one of China's most prominent advocates of
democratic change, was tried Wednesday on charges of subversion, a sign that Chinese leaders are
reducing their already limited tolerance for peaceful political dissent.

Mr.~Liu, an essayist and social critic who has spent more than a year in detention, faces up to 15
years in prison if convicted of charges legally defined as ``incitement to subvert state power.''
His lawyer said a verdict in the closed trial was likely to be announced on Friday. Charging one of
the best-known dissidents with subversion is a disheartening milestone in the eyes of some Chinese
legal experts and human rights advocates.

China has rarely brought political charges against people advocating peaceful dissent in recent
years, though it often accuses those who offend the authorities of other crimes, like tax evasion,
leaking state secrets or violating business regulations.

Now, flush with record foreign exchange reserves and buoyed by the world's most resilient major
economy, Chinese officials seem less hesitant to call a crackdown by its own name: the charge of
subversion has now been brought against not only Mr.~Liu but also Hu Jia, an AIDS activist and
environmentalist, who was convicted of that crime last year and sentenced to three and a half years
in prison.

``Many people see this trial as a tipping point,'' said John Kamm, the founder of the Dui Hua
Foundation, a group that advocates for human rights and works behind the scenes to free Chinese
political prisoners. ``The government seems to be getting tougher and more unyielding.''

During the past year, the government has tightened restrictions on access to the Internet,
suppressed the country's small band of public advocacy lawyers and jailed activists who blamed poor
school construction for the deaths of thousands of children during the 2008 Sichuan earthquake.

Legal scholars say they worry that top party leaders seem less amenable to building an impartial
legal system and allowing people to exercise the political rights in China's Constitution, which
could mean that intellectuals and civic groups have less room to operate.

Although Mr.~Liu is charged with writing six articles recently published on overseas Web sites, the
main accusation seems to be a role in creating Charter 08, a political reformer's wish list that
attracted 10,000 signatures this year during its brief life on China's heavily censored Internet.

The petition called for rule of law, expanded human rights and an end to the Communist Party's
monopoly on power. Most poignantly, it sought to guarantee the right to unfettered speech and the
abolition of the very law under which Mr.~Liu is now being tried. ``We should end the practice of
viewing words as crimes,'' the document said.

A literature professor and prolific writer, Mr.~Liu, 53, was a visiting scholar at Columbia
University when student protesters occupied Tiananmen Square in the spring of 1989. He quickly
returned home, joined hunger strikers and, as the military closed in, encouraged students to leave
the square peacefully.

After soldiers regained control of the capital, Mr.~Liu was detained and held for 21 months without
trial. In 1996, after publicly demanding the release of those still imprisoned for their roles in
the protests, he was sent to a labor camp for three years.

But although he was fired from his teaching job, Mr.~Liu became a symbol of the small but real space
for expressing opinions in a changing China. He continued to write provocative essays with titles
like ``The Chinese Communist Party's Dictatorial Patriotism'' and meet with like-minded
intellectuals who urge China to embrace democracy without violent upheaval.

``We couldn't have a repeat of June 4, where all sides lose, so we came up with a constructive way
forward,'' said Zhang Zuhua, a lead author of Charter 08, referring to the date of the 1989
Tiananmen crackdown.

In an interview earlier this year, Mr.~Zhang described how he and about 100 other people laid out
their aspirations and batted the document back and forth for a few months. Posted online last
December, the manifesto, whose name is a reference to Charter 77, a Soviet-era petition by Czech
dissidents, quickly gained thousands of signatures from workers, teachers and retired party members.

Party leaders and security officials saw the charter as a step too far.

The police questioned most of the original signers, about 300 people, and detained two of them:
Mr.~Liu, who has been held in a secret location, and Mr.~Zhang, who was later released but is now
subjected to round-the-clock surveillance.

In his defense, Mr.~Liu's lawyers argued that he had written more than 490 articles since 2005 but
that the authorities chose only six as evidence that he sought to subvert the state, according to
his younger brother, Liu Xiaoxuan, who attended the trial. ``He rejected their argument that Charter
08 brought about a 'malevolent social impact' and told the court that his remarks are within the
realm of free speech, which is protected by the Constitution,'' the brother said.

Mr.~Liu's lawyers have expressed frustration with the judicial process, saying for months that they
had no access to their client and that they received the indictment less than two weeks ago, leaving
little time to prepare a defense.

Officials in the West have also taken up Mr.~Liu's case, although it is unclear whether the
diplomatic pressure would have any impact. On Wednesday, a contingent of two dozen consular
officials from European nations, Canada and the United States lingered outside the courthouse,
having been barred from entry at the last minute despite applying for permits a week earlier.

``We were simply told there were no more seats,'' said Nicholas Weeks, the first secretary of the
Swedish Embassy.

The defendant's wife, Liu Xia, was also kept from the trial by security officials who watch her
apartment building day and night. ``I'm already prepared for the worst,'' she said by telephone.

Despite the heavy security, about two dozen supporters of Mr.~Liu gathered at the courthouse on
Wednesday morning stamping their feet against the cold.

Lei Ji, 48, an unemployed meat plant worker and self-described social critic, made an 18-hour train
ride to show his solidarity outside the courthouse. He said that he had never met Mr.~Liu but that
they had exchanged e-mail messages in years past. And yes, he said proudly, he had signed Charter
08.

Then he reached into his fraying leather satchel to reveal a bundle of photocopied manifestoes.
``I'm not afraid,'' he said after handing one to a court official who had stepped outside to look at
the scene. ``I love China. I just want my country to have freedom and human rights.''

\section{BlackBerry Breakdown Puts Users in Uproar}

\lettrine{G}{iven}\mycalendar{Dec.'09}{24} their dependence on the device, it is no surprise that
many BlackBerry owners are quick to complain when their smartphone fails to deliver on its promise
of offering e-mail anywhere, at any time.

But will a major disruption of BlackBerry e-mail and Internet service from Tuesday night to
Wednesday morning -- the second failure in a less than a week -- prompt any of them to make the
switch to another smartphone, threatening the BlackBerry's market-leading position?

The shutdown, which affected users in North and South America, appears to have started about 6:30
p.m. Tuesday. Service was generally restored by 2:45 a.m. Wednesday. BlackBerry users turned to
another communication platform, Twitter, to express outrage over the latest interruption. ``Never
Again BlackBerry,'' one user post read.

Like other shutdowns in the past, including a major disruption 18 months ago, this week's failure
appeared to be an unintended consequence of a software upgrade. But Research In Motion, the company
in Ontario that makes the BlackBerry, again provided relatively little information to users,
wireless carriers and the news media.

In a brief statement, R.I.M. said its analysis showed that recent versions of its proprietary
instant-messaging software ``caused an unanticipated database issue within the BlackBerry
infrastructure.''

A spokeswoman for the company said in an e-mail message that no one was available for comment.

Eric Miller, an e-commerce consultant in Philadelphia, said he had spent three hours trying to
``debug'' his wife's BlackBerry on Tuesday night. He learned of the network problem on Wednesday
morning.

``As a consultant, my time is money, and those three hours could have bought her a new iPhone and
plan,'' Mr.~Miller wrote in an e-mail message.

For businesses and governments, the traditional customer base for the BlackBerry, the service
interruption was a distraction and a warning sign.

Sparkplug, a wireless-broadband Internet service provider based in Scottsdale, Ariz., lost an
automated alert service for its service technicians when BlackBerry's system failed. Since the
technicians still had phone service through their BlackBerrys, a manager worked through the night
calling each of them to monitor the progress of repairs.

``We were O.K., but it was problematic,'' said Joel Payne, Sparkplug's vice president for
engineering. He added that the service interruption caused him to re-evaluate contingency plans.
``People are trying to figure out a way to diversify for this,'' he said.

But several wireless-industry analysts say that people looking for a more reliable system than the
one offered by the BlackBerry are headed for disappointment.

Consumer Reports magazine surveyed 54,332 readers about their wireless services for a feature in its
January issue. Michael Gikas, the senior editor for electronics, said that BlackBerry service is the
front-runner for wireless-data services regardless of the carrier providing the service.

That, he said, leaves dissatisfied BlackBerry users with an unappealing choice: ``Is it worse to
have an occasional complete outage or just generally poor service all the time?''

The strength and weakness of the BlackBerry system is its highly centralized data network. All
BlackBerry e-mail passes through one of three data centers operated by the company.

That gives a high degree of security to messages traveling through the system and allows R.I.M. to
use technologies to minimize the system overloads that Apple's iPhone has sometimes brought to
AT\&T's network. The downside is that failures of the centralized BlackBerry system can immediately
cut off millions of users.

Srinivasan Keshav, a professor of computer science the University of Waterloo, which is in R.I.M.'s
hometown, said that the immense complexity of R.I.M.'s network and services made it impossible to
ensure reliable and error-free upgrades to the BlackBerry system.

``The upgrades are very much like trying to change the wheels on a car while it's going down the
highway at 60 miles an hour,'' Professor Keshav said.

Unless the BlackBerry shutdowns become a regular feature of the service, most analysts foresee few
users switching, particularly at the corporate and government level.

``Even though the vast majority of R.I.M.'s growth has been in a consumer market, they have
tremendous share in enterprise,'' said Charles S.~Golvin, a wireless analyst at Forrester Research.
``One big reason people stick with their BlackBerrys is because it's what their I.T. department
supports, despite personal preferences.

``If they want their work e-mail, that's their only choice,'' he said.

\section{For American Workers in China, a Culture Clash}

\lettrine{A}{s}\mycalendar{Dec.'09}{24} more Americans go to mainland China to take jobs, more
Chinese and Americans are working side by side. These cross-cultural partnerships, while beneficial
in many ways, are also highlighting tensions that expose differences in work experience, pay levels
and communication.

In the last few years, a growing number of Americans in their 20s and 30s have been heading to China
for employment, lured by its faster-growing economy and lower jobless rate. Their Chinese co-workers
are often around the same age.

``The tight collaboration of the two countries in business and science makes the Chinese-American
pairing one of the most common in the workplace in China,'' said Vas Taras, a management professor
at the University of North Carolina at Greensboro, a specialist in cross-cultural work group
management.

But the two groups were raised differently.

The Americans have had more exposure to free-market principles. ``Young Americans were brought up in
a commercial environment,'' said Neng Zhao, 28, a senior associate at Blue Oak Capital, a private
equity firm based in Beijing. ``We weren't. So the workplace is a unique learning process for my
generation.''

People in Ms.~Zhao's generation were born around or shortly after Deng Xiaoping opened up China to
the West, so China has evolved from a government-regulated economy to a more free-market system in
their lifetime. Therefore, they can face a steeper learning curve.

Sean Leow, 28, founder of Neocha, a social networking site based in Shanghai, says young Chinese
employees often enter jobs with less hands-on preparation. They may also have less understanding of
client services, he said.

In addition, he said, ``I know a lot of my Chinese colleagues did not do internships in college,''
in contrast to United States students.

Managers hiring workers in China appear to be paying a premium for Western experience. Foreigners
tend to earn 10 to 15 percent more than their Chinese counterparts in similar positions, said
Michael Norman, senior vice president at Sibson Consulting, an American firm.

That imbalance does not go unnoticed by Chinese workers. ``There is definitely the perception that
Americans get paid more for the same work,'' said Ting Wang, 25, an associate at WildChina, a travel
company based in Beijing.

The difference is a function of supply and demand, Mr.~Norman said. ``If you need the foreigner for
their specialized knowledge of the West, companies are willing to pay a little more.''

On the other hand, Chinese workers have a deeper understanding of the influences, like Confucianism
and Communism, that play a part in their country's culture and economy.

It is imperative for Americans working in China to adjust, said Mr.~Norman, who works on management
and work force issues for multinational companies operating in Asia.

``In the West, there is such a premium on getting things done quickly, but when you come to work in
China, you need to work on listening and being more patient and understanding of local ways of doing
business,'' he said.

Ming Alterman, 25, a senior account executive at Razorfish, a Shanghai-based digital media firm, is
the only American among 40 employees. He says Americans need to understand the importance of
building so-called guanxi (pronounced GWAN-she). The word means relationships, but has implications
beyond the obligatory happy hour, occasional lunches with the boss or networking.

``In China, it's really expected that you become friends with your boss and you go out and socialize
in a way that doesn't happen in the U.S.,'' Mr.~Alterman said.

The Chinese now rising in the work force were raised and educated in a system that tended to prize
obedience and rote learning. Their American counterparts may have had more leeway to question
authority and speak their minds. This can affect workplace communication.

When Corinne Dillon, 25, was working at a multinational company in Beijing, she noticed that her
Chinese colleagues were sometimes hesitant about expressing their opinions, which she thought was
rooted in views about hierarchy.

``Because foreigners are often in higher positions in companies, or even when they are not, there is
sometimes an implicit respect given to them that makes Chinese people not want to directly disagree
with them for fear of being perceived as impolite,'' said Ms.~Dillon, who is now director of sales
and marketing at That's Mandarin, a language school based in Beijing.

The difference cuts both ways. Ms.~Zhao, of Blue Oak Capital, recalled her first experience working
for an American at an American-run agency in Beijing. What her American boss perceived as directness
left her feeling humiliated, she said. ``I remember I was so embarrassed when my American boss told
me he didn't like something I was doing, right in front of me,'' she said. ``The Chinese way would
have been much more indirect.''

Communication styles, Professor Taras said, can create workplace challenges. ``Americans often
perceive the Chinese as indecisive, less confident and not tough enough, whereas the Chinese may see
Americans as rude or inconsiderate.''

This, he said, ``can lead to conflicts and misunderstandings, but also affect promotion and task
assignment choice, and ultimately performance.''

What is similar, though, is that both the Americans and the Chinese perceive a glass ceiling. ``Most
expats don't speak good enough Chinese, so their promotion prospects are limited, and on a social
and cultural level, young Chinese feel there are barriers that are hard to get past,'' said Ziyu
Wen, 28, who works with Americans in her job as a communications manager in Beijing.

Despite the tension, the Chinese-American pairing holds many economic and political benefits for
both countries.

``China needs workers who understand China and the West, so they can develop a business presence and
influence in overseas markets,'' Mr.~Norman said.

``Likewise, America needs people who truly understand the Chinese, in order to compete and
cooperate.'' Having Americans working alongside the Chinese in China, he said, ``is one of the best
ways to cultivate and internalize this understanding for the future.''

\section{Burnishing a Brand by Selecting an 'Idol'}

\lettrine{W}{hile}\mycalendar{Dec.'09}{24} companies focus most of their marketing on persuading
people to buy their products, internal marketing can also help their brands. These days, many
customers care about how companies treat employees.

That's the thought behind an ``American Idol''-style competition called Voice of McDonald's. Now in
its third iteration, the competition gives McDonald's employees from around the world a chance to
sing at a giant trade conference -- complete with vocal coaching and judging from stars like Ne-Yo
-- and win a \$25,000 top prize.

``We know that customer perceptions about your employment brand do have some impact in terms of
people's decisions to frequent certain restaurants or retail establishments,'' said Rich Floersch,
chief human resources officer for McDonald's.

Entrants submit audition videos. Professionals from the music industry select 30 semifinalists, and
the public is encouraged to vote on which 12 singers should make the finals.

``We wanted to be able to really leverage this concept that McDonald's is known for, which is
opportunity,'' Mr.~Floersch said. ``Our employees can become entrepreneurs and end up owning stores,
or come up through the ranks in the company.'' Twenty of the 50 top executives, including the chief
executive, James Skinner, began as restaurant workers, he said.

Each Voice of McDonald's finalist wins a trip, with a guest, to the McDonald's Worldwide Convention
in Orlando, Fla., in April. It's a giant trade show and networking event, with 15,000 franchisees,
employees and vendors who pay their own way to rub elbows.

The convention has 800,000 square feet of exhibit space, where the Ukrainians can check out what's
new on the Brazilian menu, and the Japanese can show off their cellphone-ordering technology to the
franchisee from Liechtenstein.

``It's all geared to, how do you run better restaurants, how do you contemporize your brand or your
operating system, what's new about the new menu items,'' Mr.~Floersch said.

The Voice competitors, once in Orlando, get full ``Idol''-like training and styling, including hair
and makeup, wardrobe, choreography and vocal coaching. Their final performances are before those
15,000 attendees, backed by a 50-piece orchestra. Judges include people in the music industry, and
usually a celebrity (Ne-Yo participated two years ago).

•

This year, there were 10,500 entries -- China had the most, followed by Japan -- which vocal coaches
narrowed to 30 semifinalists. The semifinalists had to pick from among a list of Sony Music songs
that McDonald's had licensed to use, leading to employees like Weijie Wu from China, Jesus Miguel
Molinares Espinoza from Peru and Franz Carl Guerrero from Canada all letting lose with versions of
``Stand By Me'' with a variety of accents.

Using Facebook, Twitter, YouTube and international social-media sites like Russia's VKontakte,
McDonald's asked the public to vote. There were 685,000 votes this year, versus the 46,000 from the
last competition two years ago.

This year, only two Americans made it to the 12-person final -- and they were both from towns about
a half-hour's drive from each other in Michigan.

Fatima Poggi, a 17-year-old senior at Pinckney Community High School in Pinckney, Mich., had been
working part time at McDonald's for a year and a half, handling the front counter and drive-through
window. Ms.~Poggi moved to the United States from Peru at age 10, and didn't speak English. Now, in
addition to the McDonald's job and school, she performs lead roles in community theater musicals
like ``Aida'' and ``Aladdin,'' and performs in events and contests around the country.

``The owner found out I was singing around the area,'' Ms.~Poggi said, ``and he talked to me about
the competition that they were going to be having.''

For Ms.~Poggi's semifinal song, she chose ``Alone'' by Heart. Her homemade video is dramatic,
showing Ms.~Poggi, in knee-high leather boots, strutting on a train track, interspersed with scenes
of her emoting on a dock.

Ms.~Poggi is the youngest finalist ever in the competition, and the oldest ever, 50-year-old Eddie
Davenport, is from Stockbridge, Mich., just west of Pinckney.

Mr.~Davenport moved to Michigan from North Carolina several years ago. After spending some time as a
youth minister, he got a job at McDonald's as a secondary maintenance man at two area restaurants.
He works 24 hours a week, doing jobs like filtering the frying oil and washing windows.
Mr.~Davenport has retinitis pigmentosa, a degenerative eye disease. ``Because of the vision, I
couldn't see how to do the registers and things like that. Mopping the floor is pretty easy,''
Mr.~Davenport said.

His semifinal submission is simple: showing him with a microphone, in front of a wall, singing
Babyface's ``When Can I See You Again.''

•

Last week, Mr.~Davenport arrived at work at 4:45 a.m. as usual, and was rearranging some stock items
when he was called out front. He found his wife, a video crew, and a McDonald's representative, who
told him he had made the finals.

He was thrilled, especially since he thought he was running out of time -- Voice of McDonald's
happens only every two years, and he may not be able to see, or work at McDonald's, long enough for
the next event.

``Honestly, I felt like this is my last shot, because the deterioration of my vision for the last
year has been so severe that if it continues the way that it has been going, I don't really feel
that I would be able to work here actively two years from now to try again,'' he said.

Mr.~Floersch, who will be in the audience at the April competition, said it was an important part of
positioning McDonald's as a good employer, especially since it needed 1.6 million workers worldwide.
``We want to make sure that we have an employment brand that is compelling, that draws people to the
organization,'' he said.

\section{Geely, Chinese Carmaker, Is Close to Buying Volvo}

\lettrine{T}{he}\mycalendar{Dec.'09}{24} Ford Motor Company and the Chinese automaker, Zhejiang
Geely Holding Group, said Wednesday that they had settled ``all substantive'' details on a sale of
Volvo, clearing the way for Geely to buy Ford's Swedish unit.

The automakers said that while final documentation, financing and government approvals remain to be
completed, ``Ford and Geely anticipate that a definitive sale agreement will be signed in the first
quarter of 2010, with closing of the sale likely to occur in the second quarter of 2010, subject to
appropriate regulatory approvals.''

The companies did not disclose a price. ``That kind of detail will come when we have a definitive
agreement,'' a Ford spokesman in London, John Gardiner, said.

Ford paid \$6 billion in 1999 to buy Volvo; Geely is expected to pay about \$2 billion. Ford
announced in October that it had selected Geely as the preferred bidder.

The joint announcement on Wednesday could ease Geely's efforts to obtain the necessary approvals in
Beijing. The Chinese government must give its approval for big overseas investments to go forward,
Mr.~Gardiner said.

Geely said in a statement that it ``expects to sign a definitive stock purchase agreement with Ford
in the first quarter of 2010.'' Zhejiang Geely is an industrial conglomerate with a controlling
stake in Geely Automobile, China's largest private automaker.

A Volvo deal would be one of the biggest moves yet by a Chinese car company in Europe or the United
States. Beijing Automotive Industry Holding last week agreed to pay \$200 million to acquire Saab
Automotive carmaking technology from General Motors.

Ford, the only Detroit carmaker to avoid bankruptcy this year, is seeking to raise money as it
refocuses on its core North American and European operations. It sold Aston Martin to a
British-Kuwaiti consortium in 2007, and sold Land Rover and Jaguar to the Indian automaker, Tata
Motors, in 2008. Ford last year also reduced its stake in Mazda Motor, the Japanese carmaker, to 13
percent from 33.4 percent.

Tata's purchase of Jaguar and Land Rover, however, has turned out to be something of a cautionary
tale. The companies continued to lose money, and the debt Tata took on to buy them has been a drag
on its earnings.

Tata has repaid the debt, but had to raise money by selling shares diluting the stake of existing
shareholders, including the Tata Group. Analysts say Tata Motors will have to nurse Jaguar and Land
Rover to health for some time.

Ford said it would continue to cooperate with Volvo, but it did not intend to retain a stake in the
Swedish company.

Geely has sought to assuage anxiety about the deal in Sweden, saying it intends to maintain Volvo
much as it is, including ``an independent management'' at its Gothenburg headquarters.

``Geely is committed to working with all stakeholders to complete the transaction in the best
interest of all parties,'' Li Shufu, the chairman of Geely, said in a statement. The company said it
had held ``constructive meetings'' in recent weeks with Volvo management, union leaders and
government officials in Sweden and Belgium.

Assuming the deal goes through, ``Volvo will retain its leadership in safety and environmental
technologies, and will be uniquely positioned as a world-leading premium brand to exploit
opportunities in the fast-growing China market,'' Geely said.

Meanwhile, General Motors is still working to dispose of its own Swedish carmaker, Saab. G.M. said
last week that it would shut down Saab, which is based in Trollhattan, after negotiations to sell
the company to Spyker Cars, a tiny Dutch automaker, fell through. G.M. said the sale of the older
Saab technology to Beijing Automotive would not inhibit a sale.

Although G.M. had already begun closing the company down, Spyker came back this week with a revised
offer, and G.M. said it was studying that offer as well as other potential bids. While talks were
continuing, analysts said Spyker's bid appeared to be a long shot.

\section{Hurry, Hurry, Just Hours to Shop}

\lettrine{I}{}\mycalendar{Dec.'09}{24} admit it. I am a last-minute shopper. And last-minute
shoppers pay the price literally and figuratively.

Literally, because when we rush to the nearest mall we often have to pay a premium to buy that
\ldots well, I don't know what -- I'm not done shopping.

Figuratively, we pay with anxiety, racing against the clock to find a suitable gift, not to mention
having to compete with other sharp-elbowed last-minute shoppers.

Fortunately, technology can make shopping go as smoothly as a sleigh over new snow. Software
add-ons, online ways to find gift ideas, mobile apps that hunt down a deal -- there is digital help
that even a Grinch could love.

Using these tools and following my strategy won't necessarily get you the perfect gift or the best
deal, but it should make sure no stockings go empty. Here's how to make your last-minute shopping
quick and easy.

PICKING OUT A GIFT Maybe the hardest part of last-minute shopping is thinking of a good present. You
could ask your social network for suggestions, but that tips off everyone that you've been
lollygagging and may tip them off to their gifts, too.

Instead, try doing the same thing anonymously through Aardvark, a query service that sends questions
to people who have designated themselves as experts in their fields. Ask for a gift for a cyclist,
for example, and the question will be sent to gift experts and cycling experts. It pings a few
people at a time until you get a suitable answer, so you're not bugging a lot of people. For best
results, ask a focused question, like ``what is a good gift for a 60-year-old opera lover?''

You can also anonymously see what people suggest on Twitter. I like to use Twitter Search. Enter
``Christmas gift list'' to see lists people have compiled for golfers, children and others. Be
advised that there is also a lot of promotional spam out there. A narrower search, like ``tennis
Christmas gift'' gets better suggestions. A site called the Christmas It List culls comments from 50
sites to determine the 25 hottest products. Or check its Twitter site to see day-by-day changes,
which show a larger selection.

FINDING THE RIGHT DEAL With some ideas, you can hit the Web, but rather than laboriously check each
product at several stores, let Web browser add-ons search out deals and discounts for you.

I like to use the Firefox browser with the add-on Invisible Hand. It works in the background on
Firefox and Google Chrome (Internet Explorer is coming, the company says), automatically checking
prices at about 50 online stores as you shop. A subtle alert in the toolbar tells you if it finds a
lower price than the one you have found yourself.

Another tool for smoothing online shopping is Billeo, which works on Firefox and Internet Explorer
browsers when using the search engines Yahoo, Google and Bing. Billeo shows an icon next to a search
entry if Visa, American Express, Bank of America or Chase offers a reward for the purchase, or if
the retailer is offering a discount. Billeo said it tracked offers from 720 retailers, like Target,
Macy's, Barnes \& Noble and Wal-Mart.

Billeo also streamlines one big online shopping headache, filling out forms at checkout. Fill out
Billeo's form once and it should auto-fill most checkouts you'll encounter, as long as the forms are
fairly standard and are not Flash-based.

Need something quick -- and cheap -- for the office Secret Santa exchange? I like the Web site
CouponMap.com. It shows available coupons in your immediate area (the results are shown on a map),
so you can print the coupon, run right out and get a discount gift. I've looked at dozens of coupon
sites, and find them to be pretty inaccurate, but CouponMap did well. It even found an extra hidden
discount on a gift certificate to a local crab house, reducing the price of a \$25 gift certificate
to \$7.50.

The Web site RetailMeNot.com also maps coupon deals for nearby businesses, a total of 40,000 stores
nationally, it says. The site then asks users to rate the effectiveness of the various offers,
helping to weed out expired deals.

COMPARING ON THE GO If inspiration has not struck and you find yourself milling around the mall,
there are still some excellent comparison-shopping tools to use on the spot.

RedLaser, a bar code-scanning app that recognizes the product you are looking at and then goes
online to find other prices, has proved to be a huge hit on the iPhone App Store.

Part of its brilliance is you don't have to take a picture of the bar code, you just frame the code
with the camera and it quickly and automatically scans and returns results. When you know the Web
price, you can decide if it's worth paying extra to have the gift right away -- think of it as an
optional late fee.

In the past I've relied mostly on the ShopSavvy app for Android, which is now also available for the
iPhone. Although the bar code reader can be a bit finicky, it returns not only online pricing but
also prices at nearby stores. It gives you an option of calling the store or getting directions.

Bar code scanners do have a weakness. They work best for products that are distributed nationally to
a variety of stores, items like books, CDs, DVDs and electronics. They don't work as well on items
exclusive to a single chain -- those products may have a proprietary code. They tend not to work
well on items like clothing and groceries, either.

Some of these weaknesses may be overcome by Google Goggles for Android phones. Using the camera, it
conducts a visual search. It can scan writing as well as bar codes, so it can search results based
on the writing on a toy's box, a logo or a label. Sometimes that returns a result when a bar code
doesn't work. But it's a last-ditch effort: the app is still a work in progress, and doesn't always
produce a useful result.

It may not be in the best holiday spirit, but when you're pressed for time in the gift-giving
season, you really can phone it in.

\section{Sorry, Vegans: Brussels Sprouts Like to Live, Too}

\lettrine{I}{}\mycalendar{Dec.'09}{24} stopped eating pork about eight years ago, after a scientist
happened to mention that the animal whose teeth most closely resemble our own is the pig. Unable to
shake the image of a perky little pig flashing me a brilliant George Clooney smile, I decided it was
easier to forgo the Christmas ham. A couple of years later, I gave up on all mammalian meat, period.
I still eat fish and poultry, however and pour eggnog in my coffee. My dietary decisions are
arbitrary and inconsistent, and when friends ask why I'm willing to try the duck but not the lamb, I
don't have a good answer. Food choices are often like that: difficult to articulate yet strongly
held. And lately, debates over food choices have flared with particular vehemence.

In his new book, ``Eating Animals,'' the novelist Jonathan Safran Foer describes his gradual
transformation from omnivorous, oblivious slacker who ``waffled among any number of diets'' to
``committed vegetarian.'' Last month, Gary Steiner, a philosopher at Bucknell University, argued on
the Op-Ed page of The New York Times that people should strive to be ``strict ethical vegans'' like
himself, avoiding all products derived from animals, including wool and silk. Killing animals for
human food and finery is nothing less than ``outright murder,'' he said, Isaac Bashevis Singer's
``eternal Treblinka.''

But before we cede the entire moral penthouse to ``committed vegetarians'' and ``strong ethical
vegans,'' we might consider that plants no more aspire to being stir-fried in a wok than a hog
aspires to being peppercorn-studded in my Christmas clay pot. This is not meant as a trite argument
or a chuckled aside. Plants are lively and seek to keep it that way. The more that scientists learn
about the complexity of plants -- their keen sensitivity to the environment, the speed with which
they react to changes in the environment, and the extraordinary number of tricks that plants will
rally to fight off attackers and solicit help from afar -- the more impressed researchers become,
and the less easily we can dismiss plants as so much fiberfill backdrop, passive sunlight collectors
on which deer, antelope and vegans can conveniently graze. It's time for a green revolution, a
reseeding of our stubborn animal minds.

When plant biologists speak of their subjects, they use active verbs and vivid images. Plants
``forage'' for resources like light and soil nutrients and ``anticipate'' rough spots and
opportunities. By analyzing the ratio of red light and far red light falling on their leaves, for
example, they can sense the presence of other chlorophyllated competitors nearby and try to grow the
other way. Their roots ride the underground ``rhizosphere'' and engage in cross-cultural and
microbial trade.

``Plants are not static or silly,'' said Monika Hilker of the Institute of Biology at the Free
University of Berlin. ``They respond to tactile cues, they recognize different wavelengths of light,
they listen to chemical signals, they can even talk'' through chemical signals. Touch, sight,
hearing, speech. ``These are sensory modalities and abilities we normally think of as only being in
animals,'' Dr.~Hilker said.

Plants can't run away from a threat but they can stand their ground. ``They are very good at
avoiding getting eaten,'' said Linda Walling of the University of California, Riverside. ``It's an
unusual situation where insects can overcome those defenses.'' At the smallest nip to its leaves,
specialized cells on the plant's surface release chemicals to irritate the predator or sticky goo to
entrap it. Genes in the plant's DNA are activated to wage systemwide chemical warfare, the plant's
version of an immune response. We need terpenes, alkaloids, phenolics -- let's move.

``I'm amazed at how fast some of these things happen,'' said Consuelo M.~De Moraes of Pennsylvania
State University. Dr.~De Moraes and her colleagues did labeling experiments to clock a plant's
systemic response time and found that, in less than 20 minutes from the moment the caterpillar had
begun feeding on its leaves, the plant had plucked carbon from the air and forged defensive
compounds from scratch.

Just because we humans can't hear them doesn't mean plants don't howl. Some of the compounds that
plants generate in response to insect mastication -- their feedback, you might say -- are volatile
chemicals that serve as cries for help. Such airborne alarm calls have been shown to attract both
large predatory insects like dragon flies, which delight in caterpillar meat, and tiny parasitic
insects, which can infect a caterpillar and destroy it from within.

Enemies of the plant's enemies are not the only ones to tune into the emergency broadcast. ``Some of
these cues, some of these volatiles that are released when a focal plant is damaged,'' said Richard
Karban of the University of California, Davis, ``cause other plants of the same species, or even of
another species, to likewise become more resistant to herbivores.''

Yes, it's best to nip trouble in the bud.

Dr.~Hilker and her colleagues, as well as other research teams, have found that certain plants can
sense when insect eggs have been deposited on their leaves and will act immediately to rid
themselves of the incubating menace. They may sprout carpets of tumorlike neoplasms to knock the
eggs off, or secrete ovicides to kill them, or sound the S O S.~Reporting in The Proceedings of the
National Academy of Sciences, Dr.~Hilker and her coworkers determined that when a female cabbage
butterfly lays her eggs on a brussels sprout plant and attaches her treasures to the leaves with
tiny dabs of glue, the vigilant vegetable detects the presence of a simple additive in the glue,
benzyl cyanide. Cued by the additive, the plant swiftly alters the chemistry of its leaf surface to
beckon female parasitic wasps. Spying the anchored bounty, the female wasps in turn inject their
eggs inside, the gestating wasps feed on the gestating butterflies, and the plant's problem is
solved.

Here's the lurid Edgar Allan Poetry of it: that benzyl cyanide tip-off had been donated to the
female butterfly by the male during mating. ``It's an anti-aphrodisiac pheromone, so that the female
wouldn't mate anymore,'' Dr.~Hilker said. ``The male is trying to ensure his paternity, but he ends
up endangering his own offspring.''

Plants eavesdrop on one another benignly and malignly. As they described in Science and other
journals, Dr.~De Moraes and her colleagues have discovered that seedlings of the dodder plant, a
parasitic weed related to morning glory, can detect volatile chemicals released by potential host
plants like the tomato. The young dodder then grows inexorably toward the host, until it can
encircle the victim's stem and begin sucking the life phloem right out of it. The parasite can even
distinguish between the scents of healthier and weaker tomato plants and then head for the hale one.

``Even if you have quite a bit of knowledge about plants,'' Dr.~De Moraes said, ``it's still
surprising to see how sophisticated they can be.''

It's a small daily tragedy that we animals must kill to stay alive. Plants are the ethical
autotrophs here, the ones that wrest their meals from the sun. Don't expect them to boast: they're
too busy fighting to survive.

\section{Teachers Defying Gravity to Gain Students' Interest}

\lettrine{B}{efore}\mycalendar{Dec.'09}{24} showing a video to the 11th and 12th graders in his
physics class, Glenn Coutoure, a teacher at Norwalk High School, warned them that his mouth would be
hanging open, in childlike wonderment, almost the whole time.

Mr.~Coutoure then started the DVD, showing him and other science teachers floating in an airplane
during a flight in September. By flying up and down like a giant roller coaster along parabolic
paths, the plane simulated the reduced gravity of the Moon and Mars and then weightlessness in
30-second chunks.

The teachers performed a series of experiments and playful stunts, like doing push-ups with others
sitting on their backs and catching in their mouths M \& M's that flew in straight lines, that they
hoped would help them better explain to their students the laws of motion that Sir Isaac Newton
deduced centuries ago.

``You see the ball just hangs there,'' Mr.~Coutoure said.

``That's hot,'' a student interjected.

The Northrop Grumman Foundation has sent science teachers on these flights of weightlessness in the
last four years to excite teachers and students about science and mathematics.

``All of a sudden, this teacher becomes a superhero,'' said Sandra Evers-Manly, the foundation's
president.

When the foundation polled the 205 teachers who flew in 2005 and 2006, nearly 92 percent reported an
increase in overall interest in science among their students. About 75 percent said more students
expressed a desire to continue studying math and science.

The flight on Sept.~29, aboard a modified 727 operated by Zero Gravity Corporation, which offers
commercial weightless flights at \$5,000 a ticket, took off from Stewart International Airport in
New Windsor, N.Y., and headed south. Off the New Jersey coast near Atlantic City, the jet began its
up-and-down gyrations -- one parabola at Martian gravity, two at the lesser gravity of the moon and
then a dozen weightless ones.

Anyone can feel weightless by falling. During the zero-gravity portion of the flight, the plane
flies toward the ground at the same free-falling speed and acceleration as the passengers. The plane
has few windows for viewing the approaching ground, and passengers feel as if they are floating in
level flight instead of a downward plummet. When the plane pulls out of the dive, the passengers,
who were told to lie down on the floor quickly, felt twice their usual weight.

For most of the teachers like Mr.~Coutoure, this mouth-open experience was the closest they would
ever get to space. (One, however, Rachael Manzer, a Suffield, Conn., middle school teacher, might
get there. She has been chosen as a candidate to fly on one of the suborbital spacecraft being
designed and built by private companies.)

As President Obama did last month when announcing \$260 million in public-private education
initiatives to generate interest in math and science, the teachers on the flight repeatedly spoke
about engaging and inspiring their students. They talked specifically about those in grades six
through eight, who typically seem to lose interest and proficiency.

Geoffrey Bergen of Whisconier Middle School in Brookfield, Conn., a teacher on the flight, said
traditional teaching methods struggled to hold students' interest in the age of video games and
fast-paced technology. ``It's really hard to teach out of a textbook when you consider the world we
live in,'' Mr.~Bergen said. The zero-gravity flight, he said, ``gives you a new tool for Newton's
laws.''

Seeing their teachers performing astronautlike feats, the students, he added, ``certainly are
engaged.''

After the flight, Adrienne Manzone, a teacher at Captain Nathan Hale School in Coventry, Conn.,
beamed. ``My coolness factor will go up 100 percent,'' she said.

Francis Q.~Eberle, executive director of the National Science Teachers Association, said programs
like Northrop Grumman's zero-gravity flights were helpful in improving science education. ``It helps
revitalize educators' passion for science,'' Dr.~Eberle said. ``The goal is to help them take that
excitement back and work with their students.''

But others, like Grover J.~Whitehurst, a senior fellow at the Brookings Institution and former head
of the federal Institute of Education Sciences, wonder if the emphasis on engagement and inspiration
for students can be somewhat misplaced.

``It certainly does no harm to do this,'' Dr.~Whitehurst said of the Northrop Grumman program. But,
he added, the \$5.1 million that the foundation spent on the program the last four years perhaps
could have been more effective elsewhere. ``If they had hired me as a consultant, I wouldn't have
imagined spending it on zero-gravity flights,'' he said. ``The question is, Does it really improve
their teaching? Does it have an impact? Those questions are not asked and answered with nearly
enough frequency.''

Dr.~Whitehurst said the efforts to make math and science more appealing to more students could even
be counterproductive if the message that hard work is needed to master the subjects was diluted. He
pointed to the tests called Trends in International Mathematics and Science Study, or Timss.
Countries in which a higher fraction of students said they liked math generally performed poorly,
while some countries where students disliked math, like South Korea, scored among the best. (There
are a few exceptions like Singapore, where students liked math and also did well.)

The notion that American science and math education is in crisis is overblown, Dr.~Whitehurst said.
In the most recent Timss tests, American students performed well above average and have improved
over the past decade, although some other international comparisons show American students in the
middle of the pack.

In Mr.~Coutoure's classroom, the video of the weightless flight was shown toward the end of class,
after he had lectured about the physics of motion, from Aristotle's mistaken belief that heavier
objects fall faster to Galileo's more rigorous observations and deductions.

He also performed simple demonstrations to illustrate the concepts, dropping a softball and a piece
of paper, then crumpling the paper into a ball and showing it then dropped as quickly as the
softball.

Jamie Heverin, one of the students, said, ``I'm curious about what he will pull out next time.''

\section{Weighing Medical Costs of End-of-Life Care}

\lettrine{T}{he}\mycalendar{Dec.'09}{24} Ronald Reagan U.C.L.A. Medical Center, one of the nation's
most highly regarded academic hospitals, has earned a reputation as a place where doctors will go to
virtually any length and expense to try to save a patient's life.

``If you come into this hospital, we're not going to let you die,'' said Dr.~David T.~Feinberg, the
hospital system's chief executive.

Yet that ethos has made the medical center a prime target for critics in the Obama administration
and elsewhere who talk about how much money the nation wastes on needless tests and futile
procedures. They like to note that U.C.L.A. is perennially near the top of widely cited data,
compiled by researchers at Dartmouth, ranking medical centers that spend the most on end-of-life
care but seem to have no better results than hospitals spending much less.

Listening to the critics, Dr.~J.~Thomas Rosenthal, the chief medical officer of the U.C.L.A. Health
System, says his hospital has started re-examining its high-intensity approach to medicine. But the
more U.C.L.A.'s doctors study the issue, the more they recognize a difficult truth: It can be hard,
sometimes impossible, to know which critically ill patients will benefit and which will not.

That distinction tends to get lost in the Dartmouth end-of-life analysis, which considers only the
costs of treating patients who have died. Remarkably, it pays no attention to the ones who survive.

Take the case of Salah Putrus, who at age 71 had a long history of heart failure.

After repeated visits to his local hospital near Burbank, Calif., Mr.~Putrus was referred to
U.C.L.A. this year to be evaluated for a heart transplant.

Some other medical centers might have considered Mr.~Putrus too old for the surgery. But U.C.L.A.'s
attitude was ``let's see what we can do for him,'' said his physician there, Dr.~Tamara Horwich.

Indeed, Mr.~Putrus recalled, Dr.~Horwich and her colleagues ``did every test.'' They changed his
medicines to reduce the amount of water he was retaining. They even removed some teeth that could be
a potential source of infection.

His condition improved so much that more than six months later, Mr.~Putrus has remained out of the
hospital and is no longer considered in active need of a transplant.

Because Dartmouth's analysis focuses solely on patients who have died, a case like Mr.~Putrus's
would not show up in its data. That is why critics say Dartmouth's approach takes an overly
pessimistic view of medicine: if you consider only the patients who die, there is really no way to
know whether it makes sense to spend more on one case than another.

According to Dartmouth, Medicare pays about \$50,000 during a patient's last six months of care by
U.C.L.A., where patients may be seen by dozens of different specialists and spend weeks in the
hospital before they die.

By contrast, the figure is about \$25,000 at the Mayo Clinic in Rochester, Minn., where doctors
closely coordinate care, are slow to bring in specialists and aim to avoid expensive treatments that
offer little or no benefit to a patient.

``One of them costs twice as much as the other, and I can tell you that we have no idea what we're
getting in exchange for the extra \$25,000 a year at U.C.L.A. Medical,'' Peter R.~Orszag, the White
House budget director and a disciple of the Dartmouth data, has noted. ``We can no longer afford an
overall health care system in which the thought is more is always better, because it's not.''

By some estimates, the country could save \$700 billion a year if hospitals like U.C.L.A. behaved
more like Mayo. High medical bills for Medicare patients' final year of life account for about a
quarter of the program's total spending.

Under the House health care legislation pending in Congress, the Institute of Medicine would conduct
a study of the regional variations in Medicare spending to try to determine how to reward hospitals
like Mayo for providing more cost-effective care. Hospitals identified as high-cost centers might
even be penalized, perhaps receiving lower payments from the government. The Senate bill calls only
for studies of Medicare spending variations, so it will be up to House-Senate negotiators to resolve
the matter in the final legislation.

That prospect worries Dr.~Rosenthal and his U.C.L.A. colleagues, who say that unless the distinction
can be clearly drawn between excellence and excess in medical care, efforts to cut wasteful spending
could be little more than blunt rationing.

``There's a real risk of doing harm here -- real harm,'' he said.

Indeed, U.C.L.A. and five other big California medical centers recently published their own research
results with a striking conclusion: for heart failure patients, the hospitals that spend the most
seem to save the most lives.

Testing the Thesis

Dr.~Rosenthal remembers a pivotal meeting back in 2005 when he and officials at the other California
hospitals met with Dartmouth researchers to discuss their findings.

``We were inspired,'' Dr.~Rosenthal recalled, saying he found himself agreeing with much of the
criticism aimed at his institution for its aggressive approach.

The Dartmouth analysis prompted Dr.~Rosenthal to seek further data. He collaborated with colleagues
at U.C.L.A. and four other medical centers affiliated with the University of California system, as
well as Cedars-Sinai Medical Center in Los Angeles, to design a study of why some hospitals spent so
much more on dying patients than others and what they got from their efforts.

To focus their analysis, the researchers chose to look only at a single category of patients:
elderly people with heart failure. The dead would be counted, as Dartmouth does, but so would the
living.

What they found seemed to contradict the Dartmouth thesis. The hospital that spent the most on heart
failure patients had one-third fewer deaths after six months of an initial hospital stay.

The researchers did not disclose which of the six hospitals had the best results. But for the
doctors involved, the implications were clear: spending more can sometimes save lives.

``It doesn't look like it is all waste,'' said Dr.~Michael K.~Ong, a U.C.L.A. internist and health
policy researcher who was one of the authors of the study, which was recently published in the
peer-reviewed medical journal Circulation: Cardiovascular Quality and Outcomes.

Another of the authors was Dr.~Michael A.~Gropper, a critical care specialist at the University of
California, San Francisco. The Dartmouth research has consistently portrayed his hospital as much
more cost-effective in end-of-life care than U.C.L.A. But the California study gave Dr.~Gropper a
new perspective. ``There's no doubt that additional investments may be worthwhile,'' he said.

Some other experts take a similar view.

``If you only look at the failures, you miss the benefit,'' said Dr.~Peter Bach, an epidemiologist
at Memorial Sloan-Kettering Cancer Center in Manhattan, who has examined the medical histories of
cancer patients who have died, including women with early-stage breast cancer.

``No one in their right mind would tell you not to treat these women,'' he said, ``even though some
of them will die.''

The California researchers say they also found much less variation among the six hospitals than the
Dartmouth data would indicate after they took account of significant differences among the patients
the hospitals treat, including the many patients who come to U.C.L.A. for a transplant, who are, by
definition, extraordinarily sick.

Over all, the California researchers found that the variation in spending among the six hospitals
was significantly less than the level reported by Dartmouth. When looking at all patients
hospitalized for heart failure, for example, the variation in use of resources was 27 to 44 percent
lower than when they looked at only the patients who died. And that corresponded with a separate,
informal analysis of Medicare spending by the Congressional Budget Office, which after adjusting for
the severity of illness in patients and differences in prices among regions, found less striking
variations in spending.

A report from the Medicare Payment Advisory Commission released this month also found less
variation.

Dr.~Rosenthal, who argues that there are also fundamental socioeconomic differences between patients
in the poorer sections of Los Angeles and those in the Mayo Clinic's small and solidly middle-class
hometown of Rochester, Minn., was co-author of an op-ed article in The Los Angeles Times last summer
making that case. ``Health care costs are significantly higher in areas of poverty,'' he wrote.

The Dartmouth researchers tend to dismiss such counterarguments, saying their conclusion -- higher
spending does not necessarily buy better patient outcomes -- is backed by decades of research. While
more spending may have yielded benefits among these six hospitals, a Dartmouth official said,
hospitals generally have not shown they deliver better results when they provide more care.

Still, Dr.~Elliott S.~Fisher, one of the lead investigators at the Dartmouth Atlas Project,
acknowledged that the California researchers' analysis might be better able than Dartmouth's to
identify cases in which more intensive care might prove beneficial. ``Sometimes more medical care is
better,'' he said, ``but the question is when.''

He says he believes that cost-effective hospitals with good medical outcomes should be financially
rewarded for their efforts and results. But he says that public policy aimed primarily at penalizing
high spenders would not be the solution. ``Simply reducing their prices,'' he said, ``won't fix
anything.''

'Hail Mary Pass'

Just how hard it is to determine who will most benefit from expensive care is clear in the case of
George Klidaras, a 49-year-old stay-at-home father of two who arrived by ambulance late at night in
mid-June in need of a heart transplant.

His age might have made him a good candidate for the procedure. And so might his overall state of
health. He was lucidly answering doctors' questions when he arrived. And although he had suffered a
stroke in his early 40s and had a chronic heart condition, as recently as March, Mr.~Klidaras had
been living a relatively normal life.

By the time Mr.~Klidaras arrived at U.C.L.A., though, his heart had weakened significantly. In the
preceding weeks he had already received a pacemaker and defibrillator, and his local cardiologist
decided it was time to ask U.C.L.A. -- a highly regarded transplant center and regional magnet for
cases beyond the skill of many other hospitals -- to tackle his case.

Coming to U.C.L.A. ``was our Hail Mary pass,'' said his wife, Andra, a postal worker.

Mr.~Klidaras's lungs were damaged from a pulmonary embolism, and he had a high white-blood-cell
count -- a possible sign of infection. So the first order of business was to try to determine the
source of infection so he would be well enough to undergo surgery.

The flurry of activity in the intensive care unit was ``overwhelming,'' Ms.~Klidaras recalled. Her
husband saw a dizzying array of specialists, including an infectious disease doctor and a
dermatologist after he developed a rash. ``They gave him every antibiotic and every test,'' she
said.

Mr.~Klidaras spent nearly five weeks in the intensive care unit, at a cost of about \$10,000 a day
and a total cost in the neighborhood of \$300,000. And the doctors never could stabilize his
condition enough for the transplant surgery.

After the doctors told Ms.~Klidaras there was nothing more they could do, she told them not to
resuscitate him if his heart stopped beating on its own. He died July 20.

``Until the last week, I believed he was going to make it,'' Mrs.~Klidaras said. ``I wanted them to
do everything they could to save him.''

Someone giving the strictest of reading to the Dartmouth doctrine might argue that, given the
outcome, the effort devoted to the Klidaras case was a futile expenditure of time and money.

Family Struggles

For U.C.L.A.'s doctors, deciding when enough is enough is not ultimately their call. Even when they
recommend against a patient's getting another procedure or test, it may be the patients and families
who cannot let go.

When doctors, patients and families have trouble agreeing when to stop medical treatment, the person
typically called in is Dr.~Neil S.~Wenger, a practicing physician who also serves as director of the
U.C.L.A. Health System Ethics Center.

``For someone to die who is in the clutches of medical care, decisions have to be made,'' Dr.~Wenger
said. ``Otherwise, you don't get to die a reasonable death.''

The decisions may include turning off life-prolonging technologies that were put in place when there
was still hope -- the dialysis machine to keep the kidneys functioning, the ventilator for the lungs
-- but now may be the only thing keeping the patient alive.

Dr.~Wenger often sees difficult family dynamics. He spent more than an hour recently counseling the
relatives of a woman with a traumatic head injury who was unlikely to ever regain consciousness.
When a family member suggested stopping treatment, a sibling protested, saying, ``You're killing my
sister.'' Such intense emotions are ``extremely common,'' Dr.~Wenger said.

Doctors, too, often have trouble letting go.

Many acknowledge that the current payment system encourages more care, because it rewards doctors
for providing additional tests and procedures -- not for spending the hours sometimes necessary to
guide patients and families through the long, difficult process of deciding when to stop.

``The more tubes you put in, the more you get paid,'' said Dr.~Patrick T.~Dowling, chairman of the
department of family medicine at U.C.L.A.

But the bigger challenge may be changing the ``we're not going to let you die'' culture at places
like U.C.L.A.

Doctors at other leading medical centers, like the one at the University of California, San
Francisco, say one big difference among institutions is how doctors and nurses talk with patients
and families about their choices.

``It's a cultural thing,'' said Dr.~Gropper. He says the doctors and nurses at the San Francisco
medical center take the time to keep talking with patients' families and even other doctors when
they seem reluctant to end treatments. ``You chip away at them,'' he said.

At U.C.L.A., such palliative care -- treatment devised to relieve pain and make patients more
comfortable, particularly at the end of life -- was essentially an afterthought until just a few
years ago, when an internal task force recommended that the hospital add it to its many other
specialties. The hospital now has a formal unit devoted to palliative care and is building up its
expertise in the field. Residents, in training to be doctors, are being taught how better to discuss
these issues with patients and their families.

Some doctors are resistant, particularly those with patients who had hoped for a transplant but were
removed from the list when it became clear they had no realistic chance of recovery.

Dr.~Bruce Ferrell, who helps lead the palliative care program, recalls a patient two years ago who
got a liver transplant but developed serious complications afterward and remained in the hospital
for a year. ``He had never, ever been told that he would have to live with a ventilator and
dialysis,'' Dr.~Ferrell said. ``He was never told that this is as good as it's going to get.''

Dr.~Ferrell talked with the patient about whether he might want to leave the intensive-care unit to
go home and receive hospice care. But when the surgeon overseeing the case found out, he was
furious.

``We do not use the h-word'' -- hospice -- ``on my patients,'' the surgeon told Dr.~Ferrell. ``Don't
ever come back.''

The patient chose to leave.

But lately, Dr.~Ferrell says, more of the transplant surgeons appreciate the value of what he is
trying to do.

``We're not the bad guys,'' he said. ``We offer options.''

\section{South Korea Stretches Standards for Success}

\lettrine{W}{ith}\mycalendar{Dec.'09}{24} acupuncture needles trembling from the corners of her
mouth like cat's whiskers, Moon Bo-in, 5, whined with fear. But the doctor, wearing a yellow gown
patterned with cartoon characters, poked more needles into her wrists and scalp.

``It's O.K., dear,'' said her mother, Seo Hye-kyong. ``It will help make you pretty and tall. It
will make you Cinderella.''

Swayed by the increasingly popular conviction that height is crucial to success, South Korean
parents are trying all manner of remedies to increase their children's stature, spawning hundreds of
growth clinics that offer hormone shots, traditional Eastern treatments and special exercises.

``In our society, it's all about looks,'' said Ms.~Seo, 35. ``I'm afraid my daughter is shorter than
her peers. I don't want her to be ridiculed and lose self-confidence because of her height.''

Ms.~Seo spends \$770 a month on treatments for her daughter and her 4-year-old son at one such
clinic, Hamsoa, which has 50 branches across the country and offers a mix of acupuncture,
aromatherapy and a twice-a-day tonic that contains deer antler, ginseng and other medicinal herbs.

``Parents would rather add 10 centimeters to their children's stature than bequeath them one billion
won,'' said Dr.~Shin Dong-gil, a Hamsoa doctor, invoking a figure in Korean currency equal to about
\$850,000. ``If you think of a child as a tree, what we try to do here is to provide it with the
right soil, the right wind, the right sunshine to help it grow. We help kids regain their appetite,
sleep well and stay fit so they can grow better.''

Koreans used to value what was perceived as a grittiness on the part of shorter people. ``A smaller
pepper is hotter,'' according to a saying here, and one need look no further for proof than to the
former South Korean strongman Park Chung-hee, or across the demilitarized zone to the North Korean
ruler Kim Jong-il, who claims to be 5-foot-5 (but adds inches with elevator shoes and a bouffant
hairstyle).

But smaller is no longer considered better, thanks in part to the proliferation of Western models of
beauty and success. ``Nowadays, children scoff if you mention Napoleon and Park Chung-hee,'' said
Park Ki-won, who runs the Seojung Growth Clinic. ``On TV, all young pop idols are tall. Given our
society's strong tendency to fit into the group and follow the trend, being short is a problem.
Short kids are ostracized.''

Concerns about the trend are growing, too, with some groups warning that growth clinics, while
operating within the limits of the law, promise far more than the evidence supports.

Yoon Myoung, a top researcher at Consumers Korea, a civic group that, with the help of scientists,
has been investigating the clinics, said parents should be more skeptical.

``There is no clinical proof or other evidence that these treatments really work,'' Ms.~Yoon said.
``They use exaggerated and deceptive ads to lure parents. But Korean families often have only one
child and want to do whatever they can for that child.''

Last month, the simmering discomfort over the trend exploded when a college student put it into
blunt words on national television.

``Being tall means being competitive,'' Lee Do-kyong, a student at Hongik University in Seoul, said
on a television talk show. ``I think short guys are losers.''

Bloggers vilified her, and lawmakers denounced the station, KBS-TV, for not editing her comments.
Viewers filed defamation lawsuits. Ms.~Lee was forced to apologize, and the Communications Standards
Commission ordered the show's producers to be reprimanded for ``violating human rights'' and
``stoking the looks-are-everything phenomenon.''

``She simply said what everyone thinks but doesn't dare say in public,'' said Dr.~Kim Yang-soo, who
runs a growth clinic called Kiness. ``Here, if you change your height, you can change your fate.''

At his clinic, Kim Se-hyun, a fifth grader, walked on a treadmill with her torso encased in a
harness suspended from an overhead steel bar. The contraption, the clinic maintains, will stretch
her spine and let her exercise with less pressure on her legs.

Nearby, sweat rolled off Lee Dong-hyun, 13, as he pedaled a recumbent bicycle while reading a comic
book. Behind him, his sister, Chae-won, the shortest girl in her first-grade class, stretched to
touch her toes on a blue yoga mat, squealing as an instructor pushed down against her back.

Two years ago, their mother, Yoon Ji-young, had tried giving Dong-hyun growth hormone shots, which
have also increased in popularity here. But many doctors will prescribe them only for exceptionally
small children with severe growth disorders. And parents have been discouraged by their high cost
and fears of side effects.

Ms.~Yoon said she was spending \$850 a month on the shots but stopped after eight months.

Now she drives her children to Kiness three times a week. ``Both my husband and I are short,'' said
Ms.~Yoon, 31, who is about 5 feet tall. ``I don't want my children to blame us for being short when
they grow up.''

Another mother at the clinic, Chang Young-hee, 54 and 4-foot-10, said her children had already
experienced height discrimination. Both her daughters are college graduates and have good jobs, but
when they reached marrying age, matchmakers regarded their short stature as a defect.

``It felt like a blow to the head,'' Ms.~Chang said. ``I learned a lesson. If you fall behind in
your studies, you can catch up later. But if you miss the time to grow, you miss it forever.''

Her daughters eventually married, but for the past four years, she has been taking her youngest
child, Seo Dong-joon, to Kiness. The boy, now 15, knows his goal.

``If I'm tall, I'll have an advantage selecting my future wife,'' he said, holding an English
vocabulary book, which he studies while exercising. ``Short guys are teased at school.''

South Koreans have been growing taller anyway, thanks to changes in their diet. Over the past 30
years the average height of high school senior boys in South Korea has increased 3.5 inches, to
5-foot-8, according to government data. Senior girls grew an average of 2 inches, to 5-foot-3.

Doctors at the growth clinics say that most children simply aspire to the new average height, but
with more tall teenagers, those who are not as tall seem even shorter. ``The gap between tall and
short has become more pronounced,'' said Dr.~Park of Seojung, who recently opened 36 joint-venture
growth clinics in China and said the quest to become taller was regionwide.

If so, one country that has been left behind is North Korea. Food shortages there have left children
stunted, according to the United Nations and private relief agencies. Dr.~Park cited the case of a
16-year-old who fled North Korea last July to join his mother, who had arrived in the South three
years earlier. The boy was 5 feet tall, almost four inches below the South Korean average.

``His height wasn't unusual for the North,'' Dr.~Park said. ``But when his mother saw him again, she
cried because the boy hadn't grown at all, and because she knew the disadvantages he'd face here.''

``My dream is to open growth clinics in North Korea,'' Dr.~Park said, ``so that, once we unify,
children from both sides will be able to stand shoulder to shoulder, not with one side a head taller
than the other.''

\section{Sawyer Sets Her Tone as ABC's Anchor}

\lettrine{T}{here}\mycalendar{Dec.'09}{24} were no high-tech frills or showy experiments in Diane
Sawyer's brisk first night as anchor of ``World News'' on ABC. Instead, on Monday, the network
draped its star in utter seriousness: Ms.~Sawyer gamely sparred with President Mahmoud Ahmadinejad
of Iran about his country's nuclear intentions and whether his government oppresses street
protesters in Tehran.

Mr.~Ahmadinejad didn't really answer her questions, but he did address her as ``respectable lady.''

Much fuss has been made about the fact that two of the three broadcast anchors are now women, but
what Ms.~Sawyer's first nightly news program really signaled is the return of the alpha anchors.
Behind a twinkly warmth, honed in countless ``Good Morning America'' appearances with pets, babies
and diet experts, Ms.~Sawyer appears to have the same high-octane drive and ambition as Brian
Williams and Katie Couric.

Seated across from Mr.~Ahmadinejad in a studio in Copenhagen this weekend, Ms.~Sawyer waved a copy
of a supposed Iranian document that was published in The Times of London this month and that
suggests Iran is working on a trigger for a nuclear bomb. ``Have you been testing a neutron
initiator?'' she asked, sounding a bit like a principal who has caught a student cheating. ``Would
you like to see this document?''

The Iranian leader declined, and expressed annoyance about Western suspicions: ``I think that some
of the claims about our nuclear issue have turned into a repetitive and tasteless joke.''

It was a timely interview and an expertly timed one. Iran has not been very good to the West of
late, but it has been handy for Ms.~Sawyer, who has reported there before and uses multiple clips
from past trips in Tehran in her ABC News promos, glossy ads that present the veteran co-host of
``Good Morning America'' as a globe-trotting reporter and a caring, experienced newshound.

Mr.~Ahmadinejad is an experienced interviewee, deft at filling time with evasive or coy answers. His
conversation with Ms.~Sawyer didn't shed new light on Iran's actions or intentions (in fairness,
interviews with heads of state rarely produce much news). Mostly, Ms.~Sawyer's trip to Copenhagen,
the site of the United Nations conference on climate change, provided a splashy tableau of
aggressive hard-news reporting.

Network anchors are a bit like European monarchs: they don't really matter much anymore, but people
still perk up for a royal wedding or a changing of the stars.

CBS, which weathered Dan Rather's scandal-tinted abdication and the splashy ascension of Katie
Couric, carries some of the melodrama and tabloid pizazz of Britain's royal family. Like NBC before
it, ABC News favors more orderly, tame transitions. ABC is closer to a Scandinavian monarchy: if
there is an internal scramble for the throne, nobody outside the realm really cares; once crowned,
the royals try to play down their privilege.

Ms.~Sawyer's was a studiously understated debut, but Monday's program did provide some clues of how
the tone and tempo of ``World News'' may change under her leadership. Charles Gibson, who stepped
down last week after fewer than four years, brought an affable, unpretentious, slightly rumpled
persona to the anchor job: he rarely seemed self-important or in a hurry, and he often let
correspondents take on stories that more insecure anchors would hoard for themselves. His farewell
on Friday was good-humored and not at all grandiose. He seemed genuinely happy to have had the job
and genuinely happy to let it go to someone else.

Ms.~Sawyer, who spent 10 years at ``Good Morning America'' waiting her turn for the top news job,
seems intent on exotic travel and big ``gets.''

It's a marked game change for her rivals: just as Ms.~Couric in 2006 became the first solo female
anchor of a network news show, Brian Williams has spent much of his tenure at ``NBC Nightly News''
as a young Turk competing against two older, avuncular and laid-back newsmen, Mr.~Gibson and Bob
Schieffer of CBS, who filled in after Mr.~Rather left his post ahead of schedule.

With Ms.~Sawyer's arrival, the three network anchors are evenly matched in looks and take-charge
personality.

Ms.~Sawyer waited longest of all for her turn in the evening news spotlight. It wasn't hard to
believe her when she smiled knowingly at the camera and said, ``Good evening, and it is so good to
be here with you tonight.''

\section{Death Toll Rises to 10 as Clashes in Iran Intensify}

\lettrine{P}{olice}\mycalendar{Dec.'09}{28} officers in Iran opened fire into crowds of protesters
on Sunday, killing at least 10 people, witnesses and opposition Web sites said, setting off a day of
chaotic street battles that threatened to deepen the country's civil unrest.

The protests, during the holiday commemorating the death of Imam Hussein, Shiite Islam's holiest
martyr, were the bloodiest and among the largest since the uprisings that followed the disputed
presidential election last June, witnesses said. Hundreds of people were reported wounded in cities
across the country, and the Tehran police said they had made 300 arrests.

One of the dead was Ali Moussavi, a 43-year-old nephew of the opposition leader Mir Hussein
Moussavi.

The decision by the authorities to use deadly force on the Ashura holiday infuriated many Iranians,
and some said the violence appeared to galvanize more traditional religious people who have not been
part of the protests so far. Historically, Iranian rulers have honored Ashura's prohibition of
violence, even during wartime.

In Tehran, thick crowds marched down a central avenue in midmorning, defying official warnings of a
harsh crackdown on protests as they chanted ``death to Khamenei,'' referring to Iran's supreme
leader, Ayatollah Ali Khamenei, who has expressed growing intolerance for political dissent in the
country.

They refused to retreat even as the police fired tear gas, charged them with batons and fired
warning shots. The police then opened fire directly into the crowd, opposition Web sites said,
citing witnesses. At least five people were killed in Tehran, four in the northwestern city of
Tabriz, and one in Shiraz in the south, the Web sites reported. Photographs of several victims were
circulated widely.

Unlike the other protesters reported killed on Sunday, Ali Moussavi appears to have been
assassinated in a political gesture aimed at his uncle, according to Mohsen Makhmalbaf, an
opposition figure based in Paris with close ties to the Moussavi family.

Mr.~Moussavi was first run over by a sport utility vehicle outside his home, Mr.~Makhmalbaf wrote on
his Web site. Five men then emerged from the car, and one of them shot Mr.~Moussavi. Government
officials took the body late Sunday and warned the family not to hold a funeral, Mr.~Makhmalbaf
wrote.

In some parts of Tehran, protesters pushed the police back, hurling rocks and capturing several
police cars and motorcycles, which they set on fire. Videos posted to the Internet showed scenes of
mayhem, with trash bins burning and groups of protesters attacking Basij militia volunteers amid a
din of screams.

One video showed a group of protesters setting an entire police station aflame in Tehran. Another
showed people carrying off the body of a dead protester, chanting, ``I'll kill, I'll kill the one
who killed my brother.''

By late afternoon, coils of black smoke rose over central Tehran from dozens of street fires, and
smaller groups of protesters continued to skirmish with police and Basij militia members. In the
evening, loudspeakers in Imam Hussein Square, where most of the clashes took place, announced that
gatherings of more than three people were banned, witnesses said.

There were scattered reports of police officers surrendering, or refusing to fight. Several videos
posted on the Internet show officers holding up their helmets and walking away from the melee, as
protesters pat them on the back in appreciation. In one photograph, a police officer can be seen
holding his arms up and wearing a bright green headband, the signature color of the opposition
movement.

The Tehran police denied firing on protesters and in an official statement late Sunday said five
people had been killed ``in suspicious ways.''

Ahmadreza Radan, deputy commander of state security forces in Tehran, said dozens of police officers
had been injured and ``some were killed,'' the semiofficial news agency ISNA reported.

Protests and clashes also broke out in the cities of Isfahan, Mashhad, Shiraz, Arak, Tabriz,
Najafabad, Babol, Ardebil and Orumieh, opposition Web sites said.

Foreign journalists have been banned from covering the protests, and the reports could not be
independently verified.

If the 10 deaths are confirmed, it would be the highest toll since the summer, when huge crowds took
to the streets to protest what they said was rampant fraud in the presidential election won by
Mahmoud Ahmadinejad.

The White House condemned what it called the ``unjust suppression'' of civilians by the Iranian
government on Sunday.

``Hope and history are on the side of those who peacefully seek their universal rights, and so is
the United States,'' said Mike Hammer, a spokesman for the National Security Council.

The turmoil revealed an opposition movement that is becoming bolder and more direct in its challenge
to Iran's governing authorities. Protesters deliberately blended their political message with the
day's religious one on Sunday, alternating antigovernment slogans with ancient cries of mourning for
Imam Hussein.

``This is the month of blood, Yazid will fall,'' the protesters shouted, equating Ayatollah Khamenei
with Yazid, the ruler who ordered Imam Hussein's killing.

The protests may have received a boost from the death last week of Grand Ayatollah Hossein Ali
Montazeri, a patriarch of Iran's Islamic Revolution who became a fierce critic of the country's
leaders, especially in recent months. His memorials have brought out not only the young activists
and students who have dominated the protests in recent months, but also older and more conservative
people, who revered him for reasons of faith as well as politics.

Sunday was the seventh day since his death, an important marker in Shiite mourning rituals. Late
Sunday, the authorities declared martial law in the city of Najafabad, Ayatollah Montazeri's
hometown, the Jaras Web site reported.

The government crackdowns on mourning ceremonies in the past week provoked many people in the more
traditional neighborhoods of south Tehran as earlier clashes did not, some residents said.

``People in my neighborhood have been going to the Ashura rituals every night with green fabric for
the first time,'' said Hamid, 33, a laborer who lives in the southern Tehran neighborhood of
Shahreh-Ray and declined to give his last name. ``They have been politicized recently, because of
the suppression this month.''

Yet few protesters expected the scale of the bloodshed that broke out on Sunday. The memory of Imam
Hussein is so potent among Shiites that killing for any reason is strictly forbidden on Ashura, and
Iranian leaders have always tried to avoid violence or even state executions during a two-month
period surrounding the holiday.

``Ashura is a very symbolic day in our culture, and it revives the notion that the innocents were
killed by a villain,'' said Fatemeh Haghighatjoo, a former member of the Iranian Parliament who is a
visiting scholar at the University of Massachusetts, Boston. ``Killing people on Ashura shows how
far Khamenei is willing to go to suppress the protests.''

In another sign of the breadth of the crackdown, security forces on Sunday raided the offices of a
clerical association in the holy city of Qum that has supported the opposition since the June
election, the Jaras Web site reported. Guards surrounded the house, and members of the association
and their families -- who had gathered inside the association's headquarters for an Ashura mourning
ceremony -- were not allowed to leave, the site reported.

Mr.~Radan, the police deputy commander, said that only one of the protesters killed in Tehran had
been shot. Two were run over by cars and one was thrown from a bridge, he said.

But a doctor working at Najmieh Hospital in Tehran said Sunday night that the hospital had performed
17 operations on people with gunshot wounds. They were treating 60 people with serious head
injuries, including three who were in critical condition, said the doctor, who spoke on the
condition of anonymity for fear of repercussions.

\vocabulary
\end{document}
